\chapter[Superficie di genere lineare uno]{Classificazione delle superficie di genere lineare \texorpdfstring{$p^{(1)}=1$}{p(1) = 1}}
\label{cp2-ch-vii}

Chapter VII (pp.~237--266) continues the alternatives established in Chapter VI, Section \ref{cp1-vi-2}. The absolute linear genus is computed on the smooth minimal model. Thus \(p^{(1)}=1\) means \(K_X^2=0\) in the nonruled cases here, while \(p=p_g=p_a(X)\) in the regular cases and \(P=P_2\). Enriques first treats the two numerically trivial canonical classes, then the surfaces whose pluricanonical systems yield a genus-one fibration. The genus of a divisor is always its arithmetic genus; a half-fibre may be reducible and need not be a smooth elliptic curve.

The double-plane and tetrahedral models are singular projective presentations. A double plane is normalized and resolved before its invariants are computed; the sextic double along the edges of a tetrahedron is nonnormal, and its invariants belong to its normalization and any necessary minimal resolution. These conventions are those of Part 1, especially Chapter V, Section \ref{cp1-v-7}, and are essential to the following comparisons.

\section[Superficie di genere zero e bigenere uno]{Superficie di genere zero e bigenere 1, con curva bicanonica d'ordine zero}\label{cp2-vii-1}

\subsection{The eight preliminary propositions (pp.~237--239)}\label{cp2-vii-1-1}

The minimal surfaces under discussion have \(q=p_g=0\), \(K_X\not\sim0\), and \(2K_X\sim0\). They are complex Enriques surfaces. Their plurigenus sequence is \(P_{2r}=1\), \(P_{2r+1}=0\). Indeed the even powers of the canonical line bundle are trivial, and a nontrivial numerically trivial line bundle has no nonzero section: a nonzero effective divisor numerically equivalent to zero would have zero intersection with an ample divisor. This proves the displayed values \(P_3=0\) and \(P_6=1\) once the classification criterion imported from Chapter VI is accepted. The converse criterion is assumed here, not proved afresh.

Adjunction and surface Riemann--Roch give, for a nonzero effective divisor \(D\),

\[
D^2=2p_a(D)-2,\qquad
\chi(\mathcal O_X(D))=1+\frac{D^2}{2}=p_a(D),\qquad
h^2(D)=h^0(K_X-D)=0.
\]

Consequently \(\dim\vert D\vert \ge p_a(D)-1\). If \(D\) is nef and big, vanishing applied to \(K_X+(D-K_X)\) gives \(h^1(D)=0\), so equality holds. An irreducible moving positive-square class has the required nefness: it has nonnegative intersection with every other irreducible curve, and positive self-intersection. This is a precise repair of the scope of propositions 2--3. ``Virtually without basepoints'' is not a proof that an actual linear system is generated by sections.

For these positive-square systems, adjunction sends \(D\) to \(D+K_X\); the two classes are distinct, have equal numerical invariants and equal dimensions, and have the same double; applying the adjunction operation twice returns the original class. Proposition 4 must retain this positive-square context. Taken for every effective divisor it contradicts proposition 7: a \((-2)\)-curve \(R\) has \(\vert R\vert \ne\varnothing\) but \(\vert R+K_X\vert =\varnothing\).

An irreducible arithmetic-genus-one curve has square zero by adjunction; it is either an isolated half-fibre or a moving fibre in the genus-one-pencil structure. For a primitive nef isotropic class \(E\), the standard Enriques theorem gives a genus-one pencil \(\vert 2E\vert \), with exactly two double fibres \(2E\) and \(2E'\), where \(E'=E+K_X\). Its adjoint fibre class is \(2E+K_X=E+E'\). The effective half-fibres can be reducible; they are not necessarily smooth elliptic curves. The terse argument on p.~238 identifies the expected vertical classes but relies on this structure theorem for effectivity, primitiveness, and completeness. It does not establish the entire fibre classification from the numerical equality alone. A horizontal curve meets \(2E\) in an even number of points, because \(E\) is an integral divisor. These standard facts are treated in Cossec et al. (2025, Generalities and Linear Systems) and Dolgachev (2016, \S\S~3--4).

Proposition 7 is correct. If \(R\) is an irreducible genus-zero divisor, then \(R^2=-2\). Were \(R+K_X\) effective, its negative intersection with \(R\) would force \(R\) to be a component, leaving an effective divisor in class \(K_X\), which is impossible. The argument needs this component-and-residual reasoning; negative intersection alone does not literally make two arbitrary effective divisors equal.

\textbf{Proposition 8 is false as printed.} It prohibits an irreducible rational bisection of any genus-one pencil on \(X\). A pencil \(\vert 2E\vert \) with a smooth rational curve \(R\) satisfying \(E R=1\) is precisely a special genus-one fibration. Such fibrations exist on nodal Enriques surfaces: Dolgachev and Kond\=o's 2024 author draft, Definition 6.3.2 and Theorem 6.3.3, printed p.~29, states that every nodal Enriques surface admits one. These theorem numbers refer to the 2024 author draft.

The internal failure can be pinpointed. Put \(C=2E\) and \(D=C+R\). Then

\[
D^2=0+4-2=2,\quad p_a(D)=2,\quad
D+K_X=R+E+E',\quad\chi(D+K_X)=2.
\]

The restrictions of members of \(\vert D+K_X\vert \) to the half-fibres have degree one, and therefore fixed intersection points when the half-fibres are smooth. Enriques infers that the remaining moving curves have no intersection with a general fibre. That inference does not follow: fixed intersections with two distinguished fibres do not make a horizontal moving part vertical. Moreover \(\vert C+R\vert \) need not be an irreducible pencil; fixed components must be tracked. Riemann--Roch is entirely compatible with the displayed effective adjoint. The contradiction claimed on p.~239 is therefore artificial. Restricting to an unnodal Enriques surface makes the prohibition vacuous, but does not repair its asserted general proof.

\subsection{The splitting argument and descent to low genus (pp.~239--243)}\label{cp2-vii-1-2}

For a nef big class \(C\) with \(p_a(C)=\pi\), the numerical calculations for \(2C\) are correct:

\[
(2C)^2=8\pi-8,\quad p_a(2C)=4\pi-3,\quad
\dim\vert 2C\vert =4\pi-4.
\]

The locus \(\Sigma\) of unordered sums of two members of \(\vert C\vert \) has dimension \(2\pi-2\), and the same holds for the locus \(\Sigma'\) from \(\vert C+K_X\vert \). Under the expected generic injectivity of factorization into members, their degree is

\[
i=\frac12\binom{2\pi-2}{\pi-1}.
\]

One modern derivation pulls the hyperplane class back to \(\mathbb P^{\pi-1}\times\mathbb P^{\pi-1}\), where it is \(h_1+h_2\), and divides the top intersection by the degree two of the ordered-to-unordered map. The map is generically two-to-one because a general member of either original system is irreducible. This justifies the degree computation in its intended situation.

In \(\mathbb P^{4\pi-4}\), two complementary-dimensional projective varieties of degree \(i\) have intersection degree \(i^2\) if the intersection is proper, and positive-dimensional intersection otherwise. This is a statement about cycle-theoretic length with multiplicities. It does not produce \(i^2\) distinct decompositions. A common divisor gives two different partitions of its components, yielding the relations written on pp.~239--240, including \(C_3\sim C_2+K_X\) and \(C_4\sim C_1+K_X\) after a compatible labelling. Excluding a rational irreducible component with an effective adjoint uses proposition 7 correctly. But arbitrary pieces in the partitions need not themselves be irreducible, so each later use of a genus or a moving system requires component control.

The proposed descent either finds a moving system of smaller positive genus or reduces to genus-one components. Its composed-system case is \(\vert C_1\vert=\theta+\vert sK\vert\), where \(K\) here denotes an elliptic fibre, not the canonical class. The argument that adding a horizontal rational component of \(\theta\) enlarges the moving system uses proposition 8 to replace a possible intersection two by an intersection at least four. That exclusion is invalid. A valid descent must retain fixed components and prove that its successive genus decreases terminate. Its intended conclusion is a genus-two system or half-fibres \(E_1,E_2\) with \(1\le E_1E_2=s\le\pi-1\). On p.~241 Enriques expressly allows a configuration \(C_1=\theta+\theta_1\), \(C_2=\theta+\theta_2\) with rational common component \(\theta\). Bilinearity gives \(C_1C_2=\theta^2+\theta\theta_1+\theta\theta_2+\theta_1\theta_2\); omitting the negative self-intersection changes the calculation. Thus the bound concerns virtual intersections. The text does not silently assume unnodalness.

The next step starts with half-fibres \(E_1,E_2\) satisfying \(E_1E_2=s>1\). Their sum has arithmetic genus \(s+1\). Since the degree calculation gives \(i^2\ge9\), the text seeks another decomposition beyond the displayed factorizations and thereby smaller \(s\). When the four half-fibres are integral, in particular in the unnodal case, regularity, infinitesimal rigidity, and uniqueness of factorization give the required local intersection calculation. These facts supply the omitted explanation in that setting.

\textbf{A cohomological completion, with its exact hypotheses.} For any effective isolated half-fibre \(E\), the exact sequence

\[
0\longrightarrow\mathcal O_X\longrightarrow\mathcal O_X(E)
\longrightarrow N_{E/X}\longrightarrow0
\]

and \(q=0\), \(h^0(E)=1\) give \(H^0(N_{E/X})=0\). Write the sections of the four half-fibres as \(a,a',b,b'\), where \(E_1\cap E'_1=E_2\cap E'_2=\varnothing\), and assume no two of the four divisors have a common component. At the common divisor \(D=E_1+E'_1+E_2+E'_2\), the affine tangent spaces of the two pair loci are

\[
 ab\,H^0(L)+a'b'\,H^0(L),\qquad
 ab'\,H^0(L+K_X)+a'b\,H^0(L+K_X),
 \quad L=E_1+E_2.
\]

If \(\delta\) lies in both, restriction to \(E_1\), followed by cancellation of the nowhere-zero \(a'\), gives \(b'u=bv\) for suitable restricted sections. Since \(b\) and \(b'\) never vanish simultaneously, the local quotients \(u/b=v/b'\) glue to a section of \(N_{E_1/X}\). It is zero, so \(\delta\) vanishes on \(E_1\). Repeating this for all four half-fibres shows that \(\delta\) is a multiple of the defining section of \(D\). Projectivizing leaves zero intersection between these two tangent branches. To conclude that these are the entire local pair loci, one must also exclude additional factorization preimages. This is automatic when all four half-fibres are integral: the only partition into two members of \(\vert L\vert \) is the first displayed pairing, and the only partition into two members of \(\vert L+K_X\vert \) is the cross-pairing. The finite addition maps are immersive at these unique preimages, so the image branches are smooth and the common point is transverse of length one. Hence this point cannot absorb the intersection degree \(i^2\ge9\). If the global intersection is positive dimensional, additional points exist anyway.

There is also a small parameter-space confusion in the printed phrase ``due punti'': the two displayed pairings have the same total divisor \(D\), so they are the two factorizations of one common point in the ambient \(\vert 2L\vert \), not two distinct common divisors. Counting one transverse point only strengthens the conclusion.

This completion applies in particular to the unnodal configuration. With reducible half-fibres, even when they share no components, one must check that further partitions of their rational components do not create extra local image branches. If half-fibres share rational components, as the text explicitly permits on p.~241, vanishing on each divisor does not imply divisibility by their product with the required multiplicities; the tangent calculation itself then stops at that step. Separate analysis of the local ideals and factorization branches is needed for these configurations, and the unrestricted reduction also still suffers from the false proposition 8. These are unresolved tasks in the particular historical descent proof, not open questions about existence of the modern models. Even the additional common divisor obtained above must be examined before asserting that its partitions decrease \(s\).

The \textbf{modern replacement} for this descent is the lattice-and-linear-system theory of Enriques surfaces: one uses nef isotropic sequences and the associated low-degree polarizations, allowing configurations of \((-2)\)-curves explicitly. In characteristic zero there are sufficiently long nef isotropic sequences to obtain the sextic and double-octic models; Dolgachev (2016, \S~5, pp.~15--16 of the author version) records the nondegeneracy result and its model consequence. This lattice-theoretic result establishes the model theorem independently of Enriques's descent.

Once suitable half-fibres with \(E_1E_2=1\) have been secured, the arithmetic of the proposed genus-three net is coherent. For

\[
L=E_1+E_2+E_2'=E_1'+2E_2,
\]

one has \(L^2=4\), \(p_a(L)=3\), and, when nef and big, \(h^0(L)=3\). The text locates the two simple basepoints at \(E'_1\cap E_2\) and \(E_1\cap E'_2\). Blowing them up lowers the moving self-intersection to two, producing a degree-two rational map to the plane. Their locations and simplicity require the relevant basepoint theorem or the generic configuration assumed by the construction; they are not consequences of the numerical equalities alone. Restricting the pencil \(\vert 2E_2\vert \) to a general member of \(\vert L\vert \) gives its degree-two map to \(\mathbb P^1\), explaining the hyperelliptic genus-three sections. The two basepoints are ramification points of this curve-theoretic map in the proposed generic model. The octic branch is \(D_8=D_6+p+q\): the contracted \(E'_1\) gives its quadruple point \(O=p\cap q\), which is a node of the residual sextic; the other two half-fibres give the two \([3,3]\)-clusters of the total branch, or tacnodes of \(D_6\) with tangents \(p,q\). The pencil of lines through \(O\) represents the genus-one pencil. This branch calculation is sketched and refers back to Chapter V, \S~7; it must be made on the resolved cover, rather than treating the singular plane branch as smooth. Its numerical invariants and the ten-parameter count do not prove its claimed exhaustive model theorem.

\subsection{Parameter counts, the tetrahedral sextic, and Comessatti's construction (pp.~244--247)}\label{cp2-vii-1-3}

The count on p.~244 contains a definite printing defect. Plane sextics form \(\mathbb P^{27}\). The written subtraction \(27-5\cdot3=12\) counts \textbf{five}, not four, assigned double points: the node \(O\), each of the two tacnodes, and one infinitely-near double point at each tacnode. The tacnodal directions are prescribed by the lines to \(O\). After allowing \(O,M,N\) to vary, six parameters are restored; subtracting \(\dim\operatorname{PGL}_3=8\) gives ten. Independence of these conditions and finite generic stabilizers are assumptions to verify on the appropriate open locus. The count is not, by itself, a proof that every Enriques surface has that model.

For the tetrahedral sextic, the space with a fixed coordinate tetrahedron has projective dimension thirteen. Its diagonal stabilizer has dimension three, giving ten parameters. The nonnormal model has degree six and ordinary double edges. A general transversal to two opposite edges has used four of its six intersection points and leaves two; this accounts for the double projection. A generic plane choice produces the stated degree-ten branch model, while the special choice through an adjacent edge yields the degree-eight version. Quadrics through three relevant edges induce genus-three hyperelliptic curves on the normalization. The inverse image of an ordinary double edge is a double cover of that line branched at four pinch points, hence is elliptic by Riemann--Hurwitz. Thus ``the neighbourhood of the double line'' must be interpreted through normalization; the singular line itself is rational.

A useful direct modern check of the general model is the conductor formula. If \(\nu:X\to Y\subset\mathbb P^3\) normalizes an ordinary tetrahedral sextic and \(\Delta\) is its conductor divisor, then

\[
\omega_X\simeq\nu^*\mathcal O_Y(2)(-\Delta).
\]

There is no quadric through all six edges, so \(p_g=0\); the standard conductor exact sequence gives \(q=0\); and the union of the four faces trivializes \(2K_X\). This verifies the Enriques invariants without confusing the hypersurface dualizing sheaf with the canonical bundle of its normalization. Dolgachev (2016, \S~2) gives this normalization calculation. Special degenerations require checking that the normalization or minimal resolution remains in the desired class.

The ten-dimensional conclusion also has an intrinsic proof. The canonical étale double cover of a complex Enriques surface is a K3 surface. Pullback shows \(H^0(T_X)=0\) and \(H^0(\Omega_X^1\otimes K_X)=0\). Therefore \(H^2(T_X)=0\) by duality. The tangent-bundle Riemann--Roch formula \(\chi(T_X)=2K_X^2-10\chi(\mathcal O_X)=-10\) gives \(h^1(T_X)=10\); the local deformation space is smooth of dimension ten. Global irreducibility requires the period/moduli theorem, not this local calculation. It is stated in Dolgachev (2016, \S~4). The discreteness of divisor classes when \(q=0\), invoked on p.~245, is relevant to comparing a fixed polarization family with abstract surfaces, but it does not alone prove every parameter-count quotient is faithful or finite.

The footnote about infinitely many birational automorphisms must remain generic. Translation along an elliptic pencil by the difference of multisection classes can give an infinite group if the resulting Jacobian section is nontorsion. Intersection of the original curves is not a general substitute for a nontorsion proof. Some special Enriques surfaces have finite automorphism group; the modern theory explicitly treats this distinction (Dolgachev \& Kond\=o, 2025). Since an Enriques surface is minimal with nonnegative Kodaira dimension, its birational self-maps are regular automorphisms.

The proposed Comessatti construction on pp.~245--246 has a stronger defect than an omitted dominance check: its stated cubic has a \textbf{connected}, rather than split, inverse image under the double cover with the full branch \(D_8=D_6+p+q\). The defect concerns the construction as stated by Enriques; the following divisor calculation uses precisely its printed incidences.

Let \(E=K_3\) be the smooth plane cubic and \(H=\mathcal O_E(1)\). The two distinct points \(P,Q\) have the same tangential point \(O\), so

\[
 p\vert _E=2P+O\sim H,\qquad q\vert _E=2Q+O\sim H.
\]

Put \(T=A_1+\cdots+A_4\). The stipulated second cubic \(K'_3\) gives

\[
3H\sim O+2P+2Q+T.
\]

The residual sextic therefore restricts with divisor \(2O+4P+4Q+2T\), the double of the divisor of \(K'_3\vert _E\). Up to a nonzero complex scalar its restriction is a square. But the actual branch section is the product of that sextic with \textbf{both} tangent lines. Its half-intersection divisor on \(E\) is

\[
D=2O+3P+3Q+T\sim3H+O+P+Q.
\]

Since the plane octic cover uses the line bundle \(\mathcal O_{\mathbb P^2}(4)\), normalization of its restriction to \(E\) gives the unramified quadratic algebra governed by

\[
\eta=\mathcal O_E(4H-D)
      =\mathcal O_E(H-O-P-Q)
      =\mathcal O_E(P-Q).
\]

This is nontrivial two-torsion: \(2(P-Q)\sim0\) follows from the two tangencies, while \(P-Q\not\sim0\) for distinct points on a smooth elliptic curve. Thus the normalized inverse image is the connected étale double cover associated with \(\eta\), itself of genus one. It is not two elliptic curves each mapping isomorphically to \(E\). Blow-ups and resolution cannot turn this nontrivial generic quadratic field extension into a split algebra. To obtain a split cubic, the residual half-contact divisor must be shifted by this nonzero two-torsion class; the printed condition puts it in the wrong class. This establishes a definite error in the reproduced construction, without disproving the independent modern tetrahedral-sextic model theorem.

The asserted \(\infty^5\) family for the \textbf{printed}, nonsplitting incidence is nevertheless correct and has a direct construction. Fix equations \(e,e'\) for \(E,K'_3\). Every compatible sextic has the form

\[
f=c(e')^2+e g,
\]

where \(g\) is a cubic through \(O\) and tangent to \(E\) at \(P,Q\), allowing higher order as a closed condition. These five linear conditions cut a five-dimensional vector space of cubics. More intrinsically, restriction to \(E\) has the one-dimensional kernel spanned by \(e\), while \(\mathcal O_E(3H-O-2P-2Q)\) has degree four and hence four sections. Adding the coefficient \(c\) gives six independent sextic sections, a projective space of dimension five. Conversely the prescribed contacts force \(f\vert _E=c(e'\vert _E)^2\), so the difference is divisible by \(e\), and the same singularity constraints force precisely these conditions on \(g\).

The singularities also have the intended generic type. In local coordinates at \(P\), write \(e=y+a x^2+\cdots\), \(e'=b y+c_2x^2+\cdots\), and \(g=d y+e_2x^2+\cdots\). Then \(f\) starts with a nonzero multiple of \(y^2\), a possible \(yx^2\) term, and an \(x^4\) term. A general choice gives two smooth tangent branches with intersection multiplicity two, an ordinary tacnode; the same calculation holds at \(Q\). The quadratic part at \(O\) is generically nondegenerate, giving a node, and the four \(A_i\) are ordinary contacts with \(E\). Bertini applies away from this finite base scheme. This repairs the dimension and nonemptiness of the specified incidence family while confirming, rather than removing, the global two-torsion obstruction.

Even after correcting that class, a proof that the corrected incidence construction dominates the branch family would require nonemptiness, independent tangency conditions, and a dominance calculation. Enriques's presentation does not establish these conditions for a corrected split cubic.

For a correctly chosen split cubic, the subsequent divisor construction would be numerically clear: add an elliptic bisection to the genus-two pencil induced by lines through one tacnode. The latter genus follows by Riemann--Hurwitz after resolving the fixed branch cluster. If its classes are \(A\) and \(B\), with \(A^2=2\), \(B^2=0\), \(AB=2\), then \((A+B)^2=6\), \(p_a(A+B)=4\), and \(h^0(A+B)=4\) under nefness. This yields the proposed sextic map. Its hyperplane sections are represented by plane sextics through the prescribed five-point cluster, with a variable double point and eight residual contacts with the branch. The printed dimension of this family is \(\infty^3\), not \(\infty^2\). Three suitable half-fibres meeting pairwise once account for one triangular set of double edges; their partners account for the complementary edges. This is a sketch of projective realization, not a complete verification of birationality and the entire conductor scheme. The resulting sextic theorem is attributed in the source to Enriques's 1906 \emph{Sopra le superficie algebriche di bigenere uno}; this historical reference does not repair the preceding split-cubic calculation. The Comessatti criterion invoked on p.~245 must be applied to the full branch, as in Part 1, Section \ref{cp1-v-8}.

The concurrent-face degeneration on p.~246 and the infinitely-near opposite-edge degeneration on p.~247 are also sketches. In the first, the four tetrahedral faces become concurrent and the proposed sextic acquires a quadruple vertex; projection gives a degree-ten branch curve said to reduce to the octic form by a quadratic transformation. In the second, the proposed surface has a triple line, an infinitely-near skew double line, and two further skew double lines; in the plane one tacnode of the residual sextic becomes infinitely near its node \(O\). These are separate geometric assertions. A rigorous treatment needs explicit equations or a flat family, the branch clusters of the indicated Cremona transformation, and normalization and resolution of the limiting nonnormal surface. Enriques does not construct such a family, and the correct invariants of the generic model do not establish these degeneration assertions.

\section{Le superficie con tutti i generi uguali ad 1}\label{cp2-vii-2}

\subsection{The exact dimension formula and its exceptions (p.~247)}\label{cp2-vii-2-1}

The second case of Chapter VI, Section \ref{cp1-vi-2}, has \(K_X\sim0\); its minimal regular complex surface is a K3 surface in the present complex projective setting. Every \(P_m=1\). For a divisor \(D\), adjunction gives \(D^2=2p_a(D)-2\); for a smooth member \(C\), its normal bundle is \(\mathcal O_C(C)=\omega_C\). This explains the canonical characteristic series without presupposing an embedding of \(X\).

The word \emph{connessa} in the displayed theorem on p.~247 must retain the numerical meaning used in Chapter IV, Section \ref{enriques-iv-17}, which the theorem explicitly cites. With this inherited convention the theorem is correct. Ordinary connected support is insufficient for a nonreduced divisor. The printed cross-reference to Chapter IV, \S~17 identifies numerical 1-connectedness as the intended hypothesis; the example below distinguishes that hypothesis from ordinary connected support. The exact result, for a nonzero effective \(D\), is

\[
h^1(\mathcal O_X(D))=h^0(\mathcal O_D)-1,
\qquad
\dim\vert D\vert =p_a(D)+h^0(\mathcal O_D)-1.
\]

To prove it, use \(0\to\mathcal O_X(-D)\to\mathcal O_X\to\mathcal O_D\to0\), the vanishings \(H^0(\mathcal O_X(-D))=H^1(\mathcal O_X)=0\), and Serre duality \(h^1(-D)=h^1(D)\). Surface Riemann--Roch then gives the second equality. A reduced connected proper curve has only constant global functions, so the desired equality holds in that case. It also holds for a 1-connected effective divisor. Huybrechts (2016, Chapter 2, Proposition 3.1 and Remark 3.2) explains the stronger connectedness/vanishing argument.

\textbf{A counterexample to the connected-support interpretation.} Take an elliptic K3 with fibre \(F\) and section \(R\). Then \(F^2=0\), \(R^2=-2\), and \(FR=1\). The connected nonreduced divisor \(D=4F+3R\) has \(D^2=6\), \(p_a(D)=4\), and \(DR=-2\). Thus \(R\) is fixed. The class \(D-R=4F+2R=2(2F+R)\) is nef and big: \((2F+R)R=0\), and every other irreducible curve has nonnegative intersection with both \(F\) and \(R\). Its square is eight, so vanishing gives \(h^0(D-R)=6\). Hence \(\dim\vert D\vert =5\), not four. Equivalently \(h^0(\mathcal O_D)=2\). It shows exactly why numerical connectedness cannot be weakened to connected support.

The example \(sF\), \(s>1\), is correct: \(p_a(sF)=1\), but \(h^0(sF)=s+1\), so \(\dim\vert sF\vert =s\). The claim that systems with a fixed rational component are always superabundant needs a convention. For \(D=sF+R\), \(s\ge1\), the restriction to the generic fibre has degree one and its unique effective divisor is the section point. Therefore every member contains \(R\), and \(h^0(D)=h^0(sF)=s+1\). But \(p_a(D)=s\), so \(h^1(D)=0\). For \(s>1\) it is superabundant only relative to the genus-one moving part, not relative to the arithmetic genus of the full divisor. Enriques's following exception about adding a rational unisection to a multiple fibre system is consistent with this same calculation. A general assertion about enlarging a linear system upon adding a component requires the restriction exact sequence and its connecting map, not just the number of intersection points.

The family \(aF+bR\) is useful for testing the scope of the assertion: the two fully computed cases above separate nontrivial functions on a nonreduced divisor from the mere presence of a fixed rational component. They provide no general rule that every fixed component causes nonzero \(H^1\).

\subsection{Low-genus projective models (pp.~248--249)}\label{cp2-vii-2-2}

For a nef big generated class \(H\), set \(H^2=2g-2\). Then \(h^0(H)=g+1\). The map has degree two onto \(\mathbb P^2\) for genus two, with branch sextic in the smooth generic case. The double-cover formula gives \(K_X=f^*(K_{\mathbb P^2}+\mathcal O(3))=0\) and \(q=0\). A nonhyperelliptic genus-three polarization gives a quartic in \(\mathbb P^3\); genus four gives a \((2,3)\) complete intersection in \(\mathbb P^4\); and the appropriate nontrigonal genus-five case gives an intersection of three quadrics in \(\mathbb P^5\). Smooth complete intersections of those degrees have trivial canonical bundle by adjunction and irregularity zero by their Koszul resolutions.

The genus-four calculation can be completed without assuming restriction-map surjectivity in advance. Once the image is known to be an integral nondegenerate degree-six surface in \(\mathbb P^4\), the inequality \(15>h^0(2H)=14\) gives at least one containing quadric. It is irreducible, since a reducible quadric would force the integral surface into a hyperplane. Two independent quadrics would intersect in a degree-four surface, which cannot contain this degree-six surface, so the quadric is unique. Likewise \(35-h^0(3H)=6\) gives at least six independent cubic equations. The linear multiples of the quadric span only five dimensions, so an independent cubic exists. Its complete intersection with the irreducible quadric has degree six and contains the image; Bézout and unmixedness identify the two surfaces. The Koszul resolution then verifies the exact dimensions and projective normality. Thus the genuine prerequisites are the birational, nondegenerate degree-six model and the cohomology values, not a separately assumed restriction-surjectivity theorem.

For genus five, \(h^0(2H)=18\), against twenty-one quadrics, gives at least three independent equations; projective normality makes this exactly a net. Having three equations does not alone prove a complete intersection: one must exclude a containing scroll or other exceptional base locus. The minimum-genus premise is intended to exclude smaller pencils, but its connection to these projective exceptions needs proof. For modern projective-model theory, see Saint-Donat (1974) and Huybrechts (2016, Chapter 2); the projective-normality statement used here is Huybrechts's Corollary 2.5.

One must keep the map and the embedded model distinct. Even an ample class may have a degree-two map in a special case. For example a genus-two class \(L\) has \(L^2=2\); \(2L\) has square eight and six sections, but maps two-to-one onto the Veronese surface of degree four in \(\mathbb P^5\). The same numerical self-intersection as a complete intersection of three quadrics does not force the same map.

\subsection{The conic induction and the count of nineteen (pp.~249--251)}\label{cp2-vii-2-3}

Enriques proposes an induction using a conic \(R\) on a model of degree \(2g-2\). Such a conic is rational, with \(R^2=-2\) and \(HR=2\). Then

\[
(H+R)^2=2g,\qquad p_a(H+R)=g+1,\qquad (H+R)R=0.
\]

A suitable generated \(H+R\) contracts \(R\) to an ordinary double point on a model of the next degree. Conversely projection from such a node recovers the conic model. These numerical identities are correct. Generation, the identification of the singularity as an ordinary double point, and smoothing inside a complete family need additional geometric arguments.

The count via hyperplanes with multiple tangencies predicts codimension one for the conic locus. Hyperplanes through its plane form \(\mathbb P^{g-3}\), and \(R(H-R)=2-(-2)=4\); substituting \(x=g-3\) into the printed expected codimension \(x-g+4\) does give one. These correct numerical equalities do not establish independence of the tangency conditions. In the period space, a realized additional conic class imposes a proper hyperplane, which supplies the local eighteen-dimensional interpretation. The text then uses triple-point rational models to argue that the full family has nineteen parameters and the nodal locus has eighteen. Its first triple-point count is wrong by one even in the quartic case: a triple point at a fixed point of \(\mathbb P^3\) imposes ten linear jet conditions, and a moving point restores three parameters, giving codimension seven. Quartic equations form \(\mathbb P^{34}\); modulo \(\operatorname{PGL}_4\), smooth quartics have nineteen parameters, and the ordinary triple-point locus has twelve. The later number six, for making an already nodal quartic triple, is correct as a relative codimension \(7-1\). The two occurrences must not both be corrected to seven.

The plane representation of a rational triple-point model uses twelve base points on a cubic subject to one divisor-class relation; its eleven independent points and the cubic's modulus motivate the twelve-parameter count. The source cites du Val's 1932 \emph{On rational surfaces whose prime sections are canonical curves}, and Enriques--Conforto, \emph{Le superficie razionali}, Book I, \S~41, p.~184, and \S~46, pp.~220, 223. These references identify the historical antecedents of the construction. The count does not prove that every surface in a proposed higher-genus family can be specialized to that rational model with independent conditions. On p.~251 Enriques explicitly admits that the reducible-section parameter argument, derived from Noether, is not rigorous. For a transverse union \(C_1+C_2\) meeting in \(m\) points, the correct arithmetic genus is \(g_1+g_2+m-1\). The subsequent ``number of parameters minus number of conditions'' calculation presumes independence and a dominating incidence map. Neither follows from arithmetic genus. This is an acknowledged sketch, not evidence that the true Noether--Lefschetz statement is false.

Here is a modern proof of the main dimension assertion. A nowhere-zero holomorphic two-form identifies \(T_X\) with \(\Omega_X^1\). Thus \(h^1(T_X)=h^{1,1}(X)=20\) and \(H^2(T_X)=H^0(\Omega_X^1)^*=0\). The Kuranishi space is smooth of dimension twenty. Requiring a fixed nonzero integral polarization class to stay of type \((1,1)\) cuts out one smooth equation: in the period domain it is one nontrivial orthogonality condition. Therefore the locally polarized deformation space has dimension nineteen. The global period theorem gives an irreducible nineteen-dimensional moduli space for each primitive positive degree. Huybrechts (2016, Chapter 6, \S~2 and Corollary 4.4) treats these period and moduli statements. The period argument establishes the dimension without assuming independence in the historical condition counts.

The source assigns dimension \(g_1+g_2\) to the two component systems and obtains an expected parameter deficit of \(-1\) for its reducible sections. That arithmetic uses the preceding genus formula, but the dimensions require the appropriate Riemann--Roch vanishing for both component systems, and the predicted deficit requires independent incidence conditions. It is no proof of generic nonexistence.

The period method proves the required Noether--Lefschetz assertion. In the polarized period domain, every additional independent integral \((1,1)\)-class imposes a proper hyperplane condition. There are countably many integral classes. Outside their union the Picard group is \(\mathbb Z H\). Accordingly ``general'' must mean \textbf{very general}, not all members of one Zariski-open set. With \(H\) primitive and \(H^2=2g-2\), every effective positive-square class is a positive multiple of \(H\), so the minimum genus of a moving positive-square system is \(g\). On a quartic, projective normality realizes each such multiple as an actual hypersurface section; the Picard statement alone is a statement of classes. Existence for every \(g\ge2\) follows from the existence and period theory of primitive degree-\(2g-2\) polarized K3 surfaces. A birational map between smooth minimal K3 surfaces is an isomorphism; their Picard lattices then distinguish the very general rank-one examples of different primitive degree.

\subsection{Franchetta's degeneration as reproduced by Enriques (pp.~251--254)}\label{cp2-vii-2-4}

The reproduced argument specializes a smooth quartic to a monoid \(\Phi_4\) with a triple point. Its twelve lines through that point correspond, under projection, to twelve points \(A_i\) on a plane cubic \(C_3\); the resolution is presented by blowing up these points. Let \(H\) be the pullback of a line and \(E_i\) the exceptional classes. The quartic hyperplane class is \(4H-\sum E_i\), with square four, while the cubic class \(3H-\sum E_i\) is contracted in the model. A curve in the limit is allowed to contain a moving plane part, the cubic with multiplicity \(s\), and exceptional components with multiplicities \(h_i\). Writing its resulting plane degree as \(n\) and multiplicities as \(a_i=r_i+s-h_i\), the two equalities printed on p.~252 are

\[
n^2-\sum a_i^2=2\pi-2,
\qquad n(n-3)-\sum a_i(a_i-1)=2\pi-2.
\]

Subtracting gives \(3n=\sum a_i\), equivalently zero intersection with the cubic class. These formal intersection calculations are correct provided the chosen total-transform classes, arithmetic genus, and degree are in fact preserved in the specialization.

The preceding opening argument on p.~252 also needs its exact conditional reading. If a specialized curve is a complete intersection with a degree-\(m\) surface, Enriques compares \(\vert C\vert \) with \(\vert mH\vert \), then invokes the impossibility of a \(g^{\pi}_{2\pi-2}\) on a smooth genus-\(\pi\) curve. The curve-theoretic statement is valid for \(\pi\ge2\): RR gives \(h^0(A)=\pi-1+h^0(K_C-A)\le\pi\) for a degree-\(2\pi-2\) line bundle \(A\). The unresolved part is the comparison of the original and limiting systems and their numerical intersections. If the degeneration has already established \(H C=mH^2\) and \(C^2=m^2H^2\) on the original quartic, the Hodge index theorem provides a clean modern route: \(C-mH\) has square zero and is orthogonal to the ample class \(H\), hence is numerically zero; torsion-freeness of the K3 Picard group then gives \(C\sim mH\). That is a conditional completion, not a substitute for proving the specialization hypotheses.

In the first case, with no added cubic or exceptional pieces, restricting the plane divisor to the cubic gives \(\sum r_iA_i\sim nG\), while the quartic base relation gives \(\sum A_i\sim4G\). Subtracting \(r_1\) times the latter gives
\[
\sum_{i=2}^{12}(r_i-r_1)A_i\sim(n-4r_1)G.
\]
For very general points subject to the one quartic relation, no other integral relation holds, so \(r_i=r\) and \(n=4r\). The class is then \(r(4H-\sum E_i)\), and its complete system is the corresponding multiple quartic system. The required genericity is stronger than arbitrary distinct points: special points on an elliptic curve can satisfy additional torsion or divisor relations.

The second case appends exceptional curves but no cubic: \(s=0\), some \(h_i>0\). If the cubic is fundamental for the moving system, the text reduces the non-enlarging possibility to one simple exceptional component and \(r=1\), then calls its impossibility as a limit ``obvious.'' If the cubic is not fundamental, it asserts regularity and enlargement upon adding an exceptional component. Neither alternative receives a flat-family or relative-Hilbert-scheme justification in Enriques's account. In the third case, \(s>0\), the allowed exceptional multiplicities satisfy \(h_i\le s\); if there are \(t\) of them and \(\ell=C_mC_3\), the printed inequality is \(\ell+ts-3s\ge0\). The argument that \(\ell+t\ge4\) enlarges \(C_m+C_3+\sum A_i\), whereas \(\ell+t<4\) forces \(s\le1\) and a loss of one dimension, relies on the specialized system remaining complete and on control of fixed parts. Upper semicontinuity alone does not preserve completeness of a selected limiting subspace of sections, nor does it prevent vertical components. The missing ingredients are a specified family, a relative line bundle or appropriate cycle, constancy of the relevant Hilbert polynomial, and a comparison of the actual limiting section spaces.

The gap concerns \textbf{the proof reproduced on these pages}. The period argument above establishes the modern very-general conclusion independently of the specialization argument attributed to Franchetta.

The concluding formula \(n=s\nu\), \(\pi-1=s(\gamma-1)\) is ambiguous because the text moves between multiples of divisor systems and multiple coverings of projective models. If \(H=mL\), then

\[
H^2=m^2L^2,\qquad p_a(H)-1=m^2(p_a(L)-1).
\]

If instead a map defined by \(H\) has degree \(s\) onto a surface of projective degree \(\nu\), then \(H^2=s\nu\). These are different integers. No blanket replacement of \(s\) by \(s^2\) can repair every interpretation. The following claim of two nineteen-parameter families, a multiple and a simple one, consequently needs a specified polarized moduli problem and a specified divisibility type. This distinction also governs the examples on p.~257.

\section{Notizia storica}\label{cp2-vii-3}

The historical notice occupies pp.~254--257. Enriques credits Castelnuovo's 1893 work, his own 1896 double-plane paper, his 1908 study of genus-one surfaces, Severi's 1909 work on a canonical curve of order zero, Franchetta's 1942 argument, and d'Orgeval's 1945 volume. He mentions Noether's 1883 method, Rohn's 1897 discussion, Fano's 1909 argument, Lefschetz's 1921 and 1923 treatments, and du Val's two 1932 notes. The transcendental input is referred to Enriques and Severi's \emph{M\'emoire sur les surfaces hyperelliptiques}, \emph{Acta Mathematica}, vol.~32, no.~27 in the source's locator; the abbreviated citation does not specify the meaning of ``no.~27''. The alternative projective constructions refer also to Enriques--Campedelli's \emph{Lezioni}, p.~311. These are \textbf{historical attributions made by Enriques}. The bibliography distinguishes the identifiable publications from citations whose bibliographical details remain uncertain. 

The reported 1896 characterization concerns the birational class described by its invariants; it does not say that every K3 surface has a genus-two polarization or that its given polarization defines a sextic double plane. The Kummer-surface argument starts with three moduli and sixteen nodes: principally polarized abelian surfaces have three moduli, and quotienting by inversion gives ordinary double points at the sixteen two-torsion points. It predicts nineteen after smoothing the nodes independently. Independence is the issue to prove, not something obtained by adding sixteen to three. Simultaneous resolution, lattice-polarized deformation theory, and the period description give the appropriate modern setting. The analogous argument using a quartic with four disjoint rational curves also requires independent deformation conditions. Enriques himself notes a small gap in Severi's proof; the passage gives neither its precise location nor a repair.

The elliptic examples on p.~255 are a useful construction of polarizations. If \(F\) is an elliptic fibre and \(H F=3\), then \((H+sF)^2=H^2+6s\), so the genus rises by \(3s\). Starting from genera two, three, and four supplies the three residue classes modulo three. This does not show that the general polarized K3 of those genera carries an elliptic pencil. It gives special examples and possible starting points for deformation. Generation and contraction of any \((-2)\)-classes must still be checked for the chosen divisor.

The projective examples on p.~257 distinguish two questions. A degree-eight polarized map may be a complete intersection of three quadrics, or may be the degree-two map associated with \(2L\), \(L^2=2\), onto the Veronese surface. The branch sextic has degree twelve after the Veronese embedding. Likewise a quartic with hyperplane class \(H\) has \((2H)^2=16\) and \(h^0(2H)=10\), giving a degree-sixteen embedding in \(\mathbb P^9\); this is not a four-sheeted cover of a quartic. The printed model is \(F_{16}\), not \(F_{18}\).

The final ``one family'' assertion is correct in a modern formulation with a fixed \textbf{primitive} polarization degree and the appropriate moduli space; Huybrechts's Chapter 6, Corollary 4.4 verifies irreducibility. It cannot be taken to say that primitive and imprimitive integral polarizations of the same square are automatically joined in a smooth family of polarized K3 surfaces. Divisibility in the integral cohomology lattice is locally constant under parallel transport. Nor do distinct maps on one K3 automatically give distinct unpolarized surfaces. For the two nodal degree-ten examples mentioned by du Val, a common smoothing would require a separate argument in their stated setting.

\section[Superficie di genere lineare uno con bigenere maggiore di uno]{Superficie di genere lineare \texorpdfstring{\(p^{(1)}=1\)}{p(1) = 1} con \texorpdfstring{\(P>1\)}{P > 1}}\label{cp2-vii-4}

\subsection{From linear genus one to a genus-one pencil (pp.~257--259)}\label{cp2-vii-4-1}

On a minimal surface of nonnegative Kodaira dimension with \(K_X^2=0\), there are two possibilities relevant here: \(\kappa(X)=0\) or \(\kappa(X)=1\). In the latter case abundance for surfaces makes a positive multiple of \(K_X\) semiample. Its Stein factorization is a morphism \(f:X\to B\) with connected fibres. A general fibre \(F\) has \(F^2=0\) and \(K_XF=0\), hence genus one by adjunction. The sketch in the text is in this direction, but needs semiampleness, removal of fixed components, and Stein factorization. Matsuki (2002, ``Birational Geometry of Surfaces,'' pp.~9--108) gives the general minimal-model framework. A regular surface has \(g(B)=0\), since pullback injects \(H^1(\mathcal O_B)\) into \(H^1(\mathcal O_X)\).

The regular \(\kappa=0\) cases are Enriques and K3. Enriques surfaces have genus-one pencils by the theorem used in \S~1. K3 surfaces may or may not have one: for instance a K3 with \(\operatorname{Pic}(X)=\mathbb Z H\), \(H^2>0\), has no nonzero isotropic divisor. Thus the exception in the italic conclusion on p.~258 means that an elliptic pencil is not automatic for K3 surfaces, not that no K3 surface can be elliptic.

For an effective pluricanonical divisor on a minimal \(K_X^2=0\) surface, \(p_a(mK_X)=1\), but this does not make its support a single smooth elliptic curve. It is usually a sum of vertical components. The distinction between arithmetic genus, reducibility, and support is indispensable.

Noether's formula gives \(e(X)=12\chi(\mathcal O_X)=12(p+1)\). For a genus-one fibration, the Euler number is the sum of the Euler contributions of singular fibres, since the smooth fibre has Euler number zero. Enriques's count \(S=12p+12\) can be read as the number of nodal singularities after an appropriate generic deformation, or as a weighted singular-fibre count. It is not the number of singular fibres for an arbitrary fibration: an \(I_n\)-fibre contributes \(n\), and the other Kodaira types have their own contributions. This is the necessary interpretation of the Zeuthen formula \(I=S-4=12p+8\) imported from Chapter V. Schütt and Shioda (2010, \S\S~4 and 6) supply the modern elliptic-surface context.

Enriques's determinant \(d\) is the minimum positive degree of a multisection, equivalently the index of the generic genus-one curve over \(\mathbb C(B)\). It belongs to the specified fibration: different genus-one fibrations on one surface can have different indices. It is not its discriminant or the determinant of the Néron--Severi lattice. A degree-one multisection gives a rational point of the generic fibre, and properness extends it to a section of the smooth model. A section excludes multiple fibres, since \(1=\theta\cdot(mF_0)=m(\theta F_0)\) forces \(m=1\); reducible fibres remain possible.

On the generic genus-one curve, a degree-\(d\) line bundle has \(d\) sections by Riemann--Roch. For \(d\ge3\) it separates every length-two subscheme: subtracting such a subscheme lowers the dimension by exactly two, since the residual degree is positive. It therefore embeds the curve in \(\mathbb P^{d-1}\). A genuine degree-two line bundle gives a degree-two map to \(\mathbb P^1\); when only a conic quotient has initially been constructed, its splitting over \(\mathbb C(t)\) follows from the vanishing of the conic obstruction over that field. The images globalize in a projective bundle using relative evaluation after a sufficiently positive twist on the base. Their union is a ruled family of projective spaces, not automatically a smooth scroll in an arbitrary projected image. In degree two the ruled quotient is rational, and birationally passing to the plane gives a double-plane presentation. Existence of that presentation or of a genus-one curve with a given index does not make the pair \((p,d)\) a complete birational invariant.

The Riemann--Roch step on p.~259 has a printed numerical error. If \(L F=d\), then

\[
(L+F)^2-L^2=2d,\quad
p_a(L+F)-p_a(L)=d,\quad
\chi(L+F)-\chi(L)=d.
\]

Only the first increment is \(2d\). For a rigorous eventual statement, let \(E=f_*\mathcal O_X(L)\), a torsion-free coherent sheaf of rank \(d\), hence a vector bundle on the smooth base curve. The projection formula gives \(H^0(X,L+sF)=H^0(\mathbb P^1,E(s))\); for large \(s\) the dimension grows by exactly \(d\) with each increment. Also \((L+sF)F=d\), so this twist preserves the required relative degree. This proves eventual section growth. It does not itself prove irreducibility; that follows by the generic-fibre map and an appropriate generated moving system of dimension two. In the degree-one case \(\vert L+sF\vert \) retains the section as a fixed horizontal part; one uses a higher multiple such as \(3L+sF\) to obtain the Weierstrass model.

For a section \(\theta\), \(\vert3\theta\vert\) on the generic fibre is its plane cubic embedding. Lines through the distinguished point cut the residual degree-two series \(\vert2\theta\vert\), giving the relative double cover. The printed spatial construction chooses a three-dimensional subsystem of \(\vert3\theta+sF\vert\): its surface of degree \(m\) has an \((m-3)\)-fold line, and the cubic fibres have a common flex at an \((m-2)\)-fold point on that line. Projection from the point is said to give a branch curve of degree \(2n\) with a \((2n-3)\)-fold point. These special projective incidences require the chosen subsystem and its base scheme to be checked; the generic-fibre argument proves the existence of some double plane, not the completeness of the subsequent two normal forms. The source footnote on p.~259 expressly acknowledges a calculation error in no.~3 of Enriques's 1912 \emph{Sulle superficie algebriche con un fascio di curve ellittiche}, and also refers to Enriques--Chisini, \emph{Lezioni}, vol.~III, p.~345. The footnote does not specify the erroneous calculation further.

\subsection{Types I and II do not exhaust determinant one (pp.~259--262)}\label{cp2-vii-4-2}

The two stated plane branch forms have the asserted invariants under their generic singularity assumptions. Type I has a branch curve of degree \(2n\) with an ordinary \((2n-3)\)-fold point. Blow up that point to obtain \(\mathbb F_1\), with negative section \(S\) and fibre \(F\), \(S^2=-1\), \(SF=1\). The total branch divisor for the normalized double cover is

\[
B=S+(3S+2nF)=4S+2nF=2(2S+nF).
\]

Since \(K_{\mathbb F_1}=-2S-3F\), the double-cover formula gives \(K_X\sim(n-3)F\) after resolving the ordinary branch crossings, and
\[
\chi(\mathcal O_X)=2+\tfrac12(2S+nF)((n-3)F)=n-1.
\]
Thus \(p_g=n-2\) and \(P_m=m(p-1)+1\) for \(p\ge1\). Here and below the displayed canonical fibre class on \(X\) is the pullback of the ruling fibre. There is a ramification section, so the index is one. Type II contains the branch line \(OB\) and an additional ordinary quadruple branch point \(B\). Removing that line leaves degree \(2n-1\), multiplicity \(2n-4\) at \(O\), and multiplicity three at \(B\). Resolving the quadruple point lowers \(\chi\) by one and produces \(p_g=n-3\), with the same minimal canonical and plurigenus formulas in terms of \(p\). The footnote's letters \(OP,P\) must be \(OB,B\). Varying \(n\) varies the birational invariant \(p_g\), so different resulting genera certainly give different birational classes.

The exceptional curves in type II can also be checked explicitly. \textbf{The resolution calculation for ordinary branch tangents.} Blow up \(O\), then the ordinary quadruple point \(B\), and resolve all the ordinary branch crossings of the exceptional section over \(O\) with the residual branch. One of those crossings is on the strict transform of \(OB\); call its exceptional curve \(E_1\). In total-transform notation, the relevant branch line is

\[
\ell=H-E_O-E_B-E_1,\qquad \ell^2=-2.
\]

It is a branch component, so \(\pi^*\ell=2R\) and \(R^2=-1\). The inverse image \(V\) of \(E_B\) is a smooth elliptic curve of square \(-2\), since it is a double cover branched at four distinct points. The inverse image \(W\) of \(E_1\) is rational of square \(-2\), branched at two points. Their fibre is

\[
2R+V+W,\qquad R V=R W=1,\qquad V W=0.
\]

Contracting \(R\) makes \(W\) a rational \((-1)\)-curve, gives \(V^2=-1\), and makes \(V W=1\). Contracting \(W\) then makes \(V^2=0\): it is the smooth fibre on the relatively minimal model. There are \textbf{two successive rational contractions}, not two disjoint lifts of the branch line, and not two rational curves both of square \(-1\) on the first resolved cover. This is the coherent interpretation of the footnote's compressed wording.

As a consistency check, after the first two base blow-ups,

\[
K_Y+L=(n-3)(H-E_O)-E_B,
\]

whose square is \(-1\). Ordinary branch-node resolutions do not change the cover's canonical square, so the resolved cover has \(K_X^2=-2\). The two contractions raise it to zero. The quadruple-point correction to \(\chi\) is \(\lfloor4/2\rfloor(\lfloor4/2\rfloor-1)/2=1\), while ordinary branch-node corrections are zero; hence \(\chi=n-2\) and \(p_g=n-3\). This resolves the geometric issue in the footnote under the stated genericity assumptions. Collisions or nonordinary branch tangents require a different local calculation. The compressed wording of the printed footnote does not justify counting two initial rational \((-1)\)-curves.

The claimed \textbf{completeness} of these two forms is the serious issue. There is a concrete test, not just a disagreement of formal counts.

\textbf{A fibre-theoretic test of the asserted completeness.} In generic type I, the residual branch curve \(3S+2nF\) meets \(S\) transversely in \(2n-3=2p+1\) points. At each crossing the double cover has an \(A_1\) singularity. In a local fibre coordinate \(u\) with section \(u=0\), the cover has the form \(w^2=u(u-t)\) up to units and the remaining nonvanishing branch factors. The special fibre's normalization is rational; resolving the surface node adds another rational component joining its two branches, yielding an \(I_2\)-fibre. These fibres persist on the relatively minimal model for \(p\ge1\). Type I therefore imposes \(2p+1\) reducible fibres.

In contrast a general Weierstrass surface over \(\mathbb P^1\), with \(\chi=p+1\), is given by

\[
y^2=x^3+A(t)x+B(t),\quad
A\in H^0(\mathcal O(4\chi)),\quad B\in H^0(\mathcal O(6\chi)).
\]

For general coefficients its discriminant has \(12\chi\) simple zeros and all singular fibres are \(I_1\). The coefficient dimension, minus one scaling and three base automorphisms, is

\[
(4\chi+1)+(6\chi+1)-1-3=10\chi-2=10p+8.
\]

Type I's printed count \(8p+7\) is lower by \(2p+1\), exactly the number of imposed \(I_2\)-fibres. This identifies a useful interpretation as a prescribed-fibre or lattice-polarized locus; matching the numbers alone is not a proof of a scheme-theoretic codimension or a complete stratification. For \(p>1\), the genus-one fibration is the Iitaka fibration and is determined birationally. A general Jacobian surface with only \(I_1\)-fibres cannot be birationally changed into this type-I fibration. Type II likewise retains branch-section crossings away from its distinguished quadruple-point fibre, and thus does not supply the omitted general surfaces. The claimed minimal-degree Cremona reduction must have omitted necessary infinitely-near or other singularity possibilities, or must be restricted to a smaller class. The text supplies no proof capable of deciding a complete corrected list.

The individual plane parameter calculations are still useful descriptions of their loci. Type I with fixed \(O\) has \(8n-3\) coefficients projectively; subtracting the six-dimensional stabilizer gives \(8n-9=8p+7\). Type II gives the printed \(8n-13\) and \(8n-19=8p+5\), conditional on the stated transformation group and condition independence. Those calculations cannot be promoted to the dimension of all index-one surfaces.

The Chapter V lower bound quoted here is \(M\ge9p-2p^{(1)}+12=9p+10\). Type I and II fall below it by \(p+3\) and \(p+5\), respectively. A bound for a full deformation family cannot be applied to an arbitrary special plane locus. The source responds that an index can drop on special members of a degree-two model family. This phenomenon is reasonable, but requires an actual family and does not cure the missing generic Jacobian models. Type \(\mathrm I'\) has a \((2n-4)\)-fold point and no forced branch section: on \(\mathbb F_1\) its branch is \(4S+2nF\). It has \(p=n-2\) and projective coefficient count \(10n-6\), hence quotient count \(10n-12=10p+8\). Type \(\mathrm{II}'\) adds an ordinary quadruple point, has \(p=n-3\), and gives \(10n-22=10p+8\). For \(p\ge2\), the difference \((10p+8)-(9p+10)=p-2\) is nonnegative. Their agreement with the natural genus-one parameter dimension does not prove dominance. With the generic branch conditions these are described as index-two families; a degree-two ruled-cover presentation proves only that the index divides two, since special members can have sections.

\subsection{Multiple fibres, the canonical bundle formula, and the printed constants (pp.~262--264)}\label{cp2-vii-4-3}

For a relatively minimal complex genus-one fibration \(f:X\to B\), write its multiple fibres as \(m_jF_j\), with primitive fibre divisors \(F_j\), and put \(\mathcal L=(R^1f_*\mathcal O_X)^{-1}\). Then

\[
K_X\sim f^*(K_B+\mathcal L)+\sum_j(m_j-1)F_j,
\qquad \deg\mathcal L=\chi(\mathcal O_X).
\]

This formula is for a relatively minimal smooth surface. Blow-ups add exceptional divisors to the canonical class. It is also specifically characteristic zero; wild fibres in positive characteristic require further terms. For the regular case \(B=\mathbb P^1\), \(\deg(K_B+\mathcal L)=p-1\), and taking pushforward of the pluricanonical sheaf yields

\[
P_m=\max\left\{0,\,1+m(p-1)+\sum_j
\left\lfloor\frac{m(m_j-1)}{m_j}\right\rfloor\right\}.
\]

The floor appears because \(m_jF_j\) is the pullback of a base point, whereas a residual coefficient strictly between zero and \(m_j\) remains fixed. This gives both the variable and fixed parts and proves all the required dimension formulas. The source's local-Jacobian argument recognizes the fixed multiplicity \(r-1\) along a primitive fibre of multiplicity \(r\); the global canonical bundle formula supplies the needed line bundle and rules out an inference based only on local vanishing orders. For the formula over an arbitrary base curve, see Bombieri and Mumford (1977, \S~1, Theorem~2, pp.~27--29); in characteristic zero there are no wild-fibre terms. Dolgachev (1981, Chapter II, \S~1, canonical-divisor theorem and first corollary, author-PDF pp.~32--35) treats the regular case over \(\mathbb P^1\); the corollary is on PDF p.~35. Schütt and Shioda's (2010) Theorem 6.8 is the section case and is not used alone to justify the multiple-fibre terms.

If there are \(s\) double fibres and no other multiples, the formula becomes

\[
P_{2i}=\max\{0,1+2i(p-1)+si\},\qquad
P_{2i+1}=\max\{0,1+(2i+1)(p-1)+si\}.
\]

Under the positivity assumptions of pp.~262--263 the maxima can be omitted. In particular \(P_2=2(p-1)+s+1\). The minus-one constant on p.~262 is a printed misprint. The missing \(+1\) in \(P=2(p-1)\) for \(\mathrm I',\mathrm{II}'\) on p.~263 is another. Finally \(p=n-2-s\) implies \(p-1=n-3-s\), not the printed \(n-3\). All three readings are visible in the print; their corrections follow from the canonical bundle formula.

Type III adds \(s\) branch lines through \(O\) and \(s\) \([3,3]\)-clusters. The residual curve has degree \(2n-s\), multiplicity \(2n-s-4\) at \(O\), and the stated tacnodes. The double-cover/canonical-resolution calculation gives \(p=n-2-s\), and these clusters correspond to double fibres. These calculations describe the intended generic configuration; they do not establish completeness of the normal-form list.

For degree \(a=2n-s\) and an assigned \((a-4)\)-fold point, the projective coefficient count is
\[
\binom{a+2}{2}-1-\frac{(a-4)(a-3)}2=5a-6.
\]
Counting each prescribed tacnodal incidence as six conditions gives \(5(2n-s)-6-6s\); this retains the source's generic incidence assumptions rather than asserting their independence in all degenerations. Allowing the \(s+1\) points and quotienting by \(\operatorname{PGL}_3\) gives the formal contribution \(2(s+1)-8=2(s-3)\), not \(2s-3\). Equivalently one varies the \(s\) points with \(O\) fixed and subtracts its six-dimensional stabilizer. The following printed total already uses the correct value:

\[
M=10n-9s-12=10p+s+8=8p+P_2+9.
\]

For \(s=1,2\), a negative formal contribution is not literally a negative-dimensional moduli space of configurations: the point configuration has a positive-dimensional stabilizer, which must act on the curve coefficients in the combined calculation. As always, independence and finite generic stabilizer of the full object are required. The specialization \(p=0,s=2\) gives \(P_2=1\) and \(M=10\), agreeing with Enriques surfaces and providing an internal check on the corrected constants.

For general determinant \(d\), every fibre multiplicity divides \(d\), because \(d=L\cdot(m_jF_j)=m_j(LF_j)\). This is a genuine necessary condition. A multiplicity-\(r\) fibre contributes \((r-1)F_0\), or \((r-1)/r\) of a full fibre, to the rational canonical class. If \(s_r\) counts such fibres, its rational base degree is
\[
p-1+\sum_{r\ge2}s_r\frac{r-1}{r}.
\]
The counts \(s_r\), together with \(p\), determine the displayed plurigenera via the floor formula. Moving the multiple-fibre locations and torsor data explains why the displayed locus dimensions can exceed \(10p+8\). These integers do not classify the surface: a Jacobian elliptic surface, a torsor class over its generic fibre, and global gluing, special-fibre positions and local monodromy remain necessary. Enriques acknowledges the difficulty of the arithmetic description for \(d>2\); the phrase that the classification is substantially solved should not be read as a theorem that these integers are complete birational invariants.

\subsection{The Jacobian and the irregular extension (pp.~265--266)}\label{cp2-vii-4-4}

Let \(C\) be the generic genus-one curve over \(K=\mathbb C(B)\), and \(J=\operatorname{Pic}^0(C)\). A degree-\(d\) divisor class identifies \(\operatorname{Pic}^d(C)\) with \(J\) by translation. This does not identify \(C\) with \(J\) over \(K\) unless the torsor has a \(K\)-rational point. They become isomorphic over an algebraic closure. The relatively minimal surface with generic fibre \(J\) is the Jacobian fibration and has a section.

The text's \([d,d]\) correspondence can be made precise without assuming a rationally varying choice of effective divisors over the whole Picard variety. \textbf{An explicit construction of the correspondence.} Fix a separable degree-\(d\) multisection divisor \(p_1+\cdots+p_d\) on \(C\). Over a splitting field take the union of the graphs of the maps

\[
x\longmapsto[x-p_i]\in J.
\]

The union is Galois invariant and descends to a correspondence \(C\dashrightarrow J\). Each projection has generic degree \(d\), and its closure gives the surface correspondence over the base. This is a multivalued correspondence, not a claimed single birational map. It avoids the descent and universal-divisor issue concealed in selecting representatives of all \(g_d^{d-1}\)'s at once.

A genus-one fibration and its Jacobian have corresponding reduced singular-fibre Euler contributions and the same \(\chi(\mathcal O)\). In particular their arithmetic genera agree. When they are regular, their geometric genera agree as well. Multiple smooth fibres contribute zero Euler number; multiple singular fibres must not simply be called smooth. In characteristic zero a multiple fibre of multiplicative type \(mI_b\) has underlying reduced Euler number \(b\), equal to the corresponding Jacobian contribution. Dolgachev (1981, Chapter II, \S~1, Propositions 1--2, author-PDF pp.~24--25) verifies the local comparison with the Jacobian; the Euler and \(\chi\) comparison is developed on author-PDF pp.~36--37. Its treatment also exhibits the torsor data omitted by numerical invariants alone.

For irregular surfaces with \(\chi>0\), the genus of the base equals \(q\). Indeed the fundamental line bundle has positive degree, so the \(H^0(B,R^1f_*\mathcal O_X)\) term in Leray vanishes, leaving \(h^1(\mathcal O_X)=g(B)\). In Enriques's notation this is \(g(B)=p_g-p_a\). The chapter refers this extension to continuous systems and to Chapter IX; it does not prove all of it here.

The \(p_a=-1\), equivalently \(\chi=0\), case needs care. An abelian surface has \(q=2,p_g=1\), and a bielliptic surface has \(q=1,p_g=0\). An abelian surface has an elliptic fibration only when it has an elliptic quotient; a simple abelian surface need not have one. The class also includes properly elliptic surfaces such as \(E\times B\) with \(g(B)\ge2\). For this product, \(K_X=\operatorname{pr}_B^*K_B\), \(K_X^2=0\), \(p_g=g(B)\), \(q=g(B)+1\), and translation in \(E\) gives a one-dimensional connected automorphism group. An abelian surface gives a two-dimensional translation group. The degree-zero higher direct image may now have sections, so the positive-\(\chi\) Leray argument does not imply \(q=g(B)\). Thus Enriques's historical words ``elliptic and hyperelliptic'' should not be equated with only the modern Kodaira-dimension-zero classes. The group statement is treated in Chapter XI.

The historical notice on p.~265 assigns these regular classification results chiefly to Enriques's notes of 1898 and 1914; the irregular extension on p.~266 refers to his Bologna note of December 1906. These are citations reported by the book. The explicitly acknowledged corrections in its notice do not by themselves certify the completeness of the remaining normal forms or identify every repair in the original papers.



\chapter{Superficie regolari canoniche e pluricanoniche}
\label{cp2-ch-viii}

Chapter VIII (pp.~267--324) studies canonical and pluricanonical models of regular surfaces, beginning with canonical images in three-dimensional projective space and ending with small values of the absolute linear genus. We retain the dictionary established in Part~1: on a smooth minimal complex surface of general type,
\[
p=p_g=p_a,\qquad q=0,\qquad p^{(1)}=K_X^2+1,\qquad
P=P_2=h^0(X,2K_X).
\]
An image in projective space must be normalized and resolved before its invariants are identified with those of \(X\). The degree of the image, the square of the moving divisor, and \(K_X^2\) are distinct unless the relevant map is birational and the fixed and base contributions have been removed. Here, as in Part~1, a simple system means a system defining a birational map, rather than a free or very ample system.

\section{Introduzione}
\label{cp2-viii-1}

\textbf{The scope of the reduction (pp.~267--268).}
Enriques initially supposes \(p>3\). The canonical map then has a surface as image in the cases he first constructs, but this is not a consequence of the dimension \(p-1\) alone: the system might be composed with a pencil, possess fixed components or base points, or define a map of degree greater than one. These possibilities are explicitly acknowledged, although not uniformly treated in the ensuing classification. The case \(p^{(1)}=1\), referred to the preceding classification, belongs to the elliptic discussion under its additional hypotheses; the equation \(K_X^2=0\) alone is not a characterization of an elliptic surface. See Part~1, Chapter VI, Sections~\ref{cp1-vi-1} and~\ref{cp1-vi-2}, for the distinction between numerical triviality, linear triviality, and properly elliptic behavior.

The assertion that sufficiently high pluricanonical systems give simple models is correct in this general-type setting. A precise uniform statement is that \(|mK_X|\) defines the canonical model for \(m\geq5\); it contracts the \(K_X\)-trivial exceptional configurations and need not embed the smooth minimal surface. Bombieri (1973, Main Theorem, p.~173) supplies this result. The more specific claims about third and fourth canonical multiples made at the end of the chapter require separate correction.

\textbf{The calculations used throughout the chapter.}
The Riemann--Roch, adjunction, and vanishing statements already discussed in Part~1 give
\begin{equation}
P_m=p+1+\frac{m(m-1)}2K_X^2\quad(m\geq2),
\qquad
p_a(mK_X)=1+\frac{m(m+1)}2K_X^2.
\label{cp2-viii-plurigenera}
\end{equation}
The latter is an arithmetic-genus identity, whether or not a member is smooth. If \(\sigma:\widetilde X\to X\) resolves a system and \(\sigma^*D=M+Z\), where \(|M|\) is free and \(Z\) is its total fixed divisor, then a generically finite induced map satisfies
\begin{equation}
M^2=\deg(\varphi)\deg\bigl(\varphi(\widetilde X)\bigr).
\label{cp2-viii-moving-degree}
\end{equation}
The contribution of a sequence of point blow-ups is the sum of the squares of the successive base multiplicities, not merely the number of visible points.

For the repeated double-plane checks, let \(\pi:W\to Y\) be a smooth double cover with branch divisor \(B\sim2L\). Then
\begin{equation}
\pi_*\mathcal O_W=\mathcal O_Y\oplus\mathcal O_Y(-L),\quad
K_W=\pi^*(K_Y+L),\quad
\chi(\mathcal O_W)=2\chi(\mathcal O_Y)+\frac{L(L+K_Y)}2.
\label{cp2-viii-double-cover}
\end{equation}
This is the cover calculation used in Part~1, Chapter V, Sections~\ref{cp1-v-7} and~\ref{cp1-v-8}. A branch point of multiplicity \(m\), blown up during canonical resolution, replaces \(L\) by its pullback minus \(kE\), where \(k=\lfloor m/2\rfloor\). It lowers \(\chi\) by \(k(k-1)/2\) and the square of the canonical divisor of the resolved cover by \(2(k-1)^2\). Odd exceptional components must be included in the new branch divisor. Subsequent contractions of \((-1)\)-curves raise \(K^2\) by one each. Consequently an ordinary quadruple branch point lowers \((\chi,K^2)\) by \((1,2)\), and an ordinary sextuple point by \((3,8)\). A standard \([3,3]\) point lowers \(\chi\) by one and, after the relevant exceptional \((-1)\)-curve is contracted, lowers the minimal \(K^2\) by one. These statements assume the indicated ordinary or infinitely near singularities and the absence of additional ones.

For a plane branch curve of degree \(2d\) without singularities the starting values are
\[
p_g=\frac{(d-1)(d-2)}2,\qquad
\chi=p_g+1,\qquad K^2=2(d-3)^2.
\]
Subtracting corrections to \(\chi\) gives a correction to \(p_g\) only after regularity is justified, usually through the corresponding adjoint system. Similarly, for a generic projection of a smooth surface of degree \(d\) with section genus \(g\), an ordinary double curve has degree
\[
\frac{(d-1)(d-2)}2-g.
\]
One cannot apply this formula to arbitrary nonordinary projections without verifying the singularity hypotheses.

\textbf{Independence of the invariants.}
Enriques invokes Noether's work of 1870 to emphasize that the superficial and linear genera are independent parameters, rather than universally linked by one equation. This does not say that every pair is realized. His local modifications of spatial models use an ordinary triple point, with corrections \((\Delta p_g,\Delta K^2)=(-1,-3)\), and the specified ordinary tacnodal configuration, with corrections \((-1,-2)\). The geometric genus calculations require independent adjoint conditions; the numbers cannot be transferred to a different singularity bearing the same informal name. The doubts on p.~268 about canonical pencils and fixed components are those of the book and its earlier Chapter IV discussion, not assertions that these questions remain unresolved today. The relevant qualifications appear in Part~1, Chapter IV, Sections~\ref{enriques-iv-18}--\ref{enriques-iv-20}.

\section{Superficie canoniche di genere \texorpdfstring{$p=4$}{p=4} e genere lineare \texorpdfstring{$p^{(1)}=6,5$}{p(1)=6,5}}
\label{cp2-viii-2}

\textbf{The quintic and its alternative spatial models (pp.~268--270).}
A smooth quintic \(X\subset\mathbb P^3\) has \(K_X\sim H\), \(p_g=4\), \(q=0\), and \(K_X^2=5\). Thus it realizes \(p^{(1)}=6\). Enriques's assertion that a simple canonical surface in \(\mathbb P^3\) first occurs in degree five concerns the complete canonical map with the indicated adjoint behavior; it must not be used to equate every numerical quintic with an embedded quintic.

The three septic models recalled from Noether are a septic with a triple conic and an additional triple point, a septic with a double sextic of genus three, and a septic with a double rational quintic together with its quadrisecant exceptional line. The first is reduced by a quadratic transformation of the first kind; the other two use cubic spatial Cremona transformations. Enriques describes the transforming systems and their residual curves. A complete proof must resolve their base schemes and verify which curves are contracted. The degrees of the systems alone do not prove that they are homaloidal.

The residual calculation in the second model can nevertheless be checked exactly. On the blow-up of the six points defining a smooth cubic surface, write \(H=3h-\sum_{i=1}^6E_i\). The displayed sextic has class
\[
A=6h-3E_1-2E_2-2E_3-2E_4-2E_5-E_6.
\]
Then \(HA=6\), \(A^2=10\), and \(K_YA=-6\), so \(p_a(A)=3\). The residual curve cut by another cubic has class
\[
B=3h-E_2-E_3-E_4-E_5-2E_6,
\]
with \(HB=3\), \(B^2=1\), \(K_YB=-3\), and \(p_a(B)=0\). These are precisely the numerical data of the residual twisted cubics used by Enriques. Their effectivity, the dimensions of the systems, and the Cremona property require the stated positions of the six points. In the third construction the quadrisecant line belongs to the base scheme of the cubic transformation; omitting it would change the transformation.

\textbf{A false exhaustion and its modern repair.}
The assertion on p.~269 that \(p_g=4\), \(K^2=5\), and an irreducible moving canonical system force a quintic canonical image is false. The divisibility argument that a canonical covering degree must divide five assumes a free canonical system; a base point changes the moving square to which that argument applies. Horikawa (1973, Theorem~1) distinguishes type I, with a free canonical system and a quintic image with rational double points, from type II, with one canonical base point and a canonical map of degree two to a quadric. After that point is blown up, the moving square is \(4=2\cdot2\), whereas \(K_X^2=5\). The target can be a smooth quadric or a quadric cone, represented on its resolution by a Hirzebruch surface. This is a missing class, not an error repaired by changing a projective degree.

The count \(55-15=40\) for smooth quintics is the dimension of their parameter space modulo projective transformations: quintic equations form \(\mathbb P^{55}\), and \(\operatorname{PGL}_4\) has dimension fifteen. Horikawa's Theorem~2 describes two local forty-dimensional branches meeting in dimension thirty-nine at an ample type IIb surface. That local assertion should not be promoted, on its own, into a proof of a global component count. His Theorem~3 proves that every minimal numerical quintic is a deformation of a smooth quintic. These distinctions replace the book's inference from its explicit models to an exhaustive family statement.

\textbf{The case \(p^{(1)}=5\).}
The lower possibility \(p_g=4\), \(K_X^2=4\) is realized by a double quadric. A clean check is to take \(Y=\mathbb P^1\times\mathbb P^1\), \(L=\mathcal O_Y(3,3)\), and a smooth branch divisor of bidegree \((6,6)\). Formula~\eqref{cp2-viii-double-cover} gives
\[
K_X=\pi^*\mathcal O_Y(1,1),\qquad K_X^2=4,\qquad
p_g=4,\qquad q=0.
\]
The canonical system is free and the canonical map is the degree-two map to the Segre quadric. The spatial branch curve has degree twelve. Thus the phrase on p.~270 asserting a canonical base point in this case is incorrect; no base point is forced.

The equivalent double-plane construction has a degree-twelve branch curve with two ordinary sextuple points. Starting from \((p_g,K^2)=(10,18)\), the resolution gives \((4,2)\). The joining line is unbranched on the resolved base and its inverse image splits into two \((-1)\)-curves; their contraction gives the minimal value \(K_X^2=4\). This is an explicit instance in which the resolved double cover is not yet the minimal surface.

Enriques also considers hypothetical degree-fourteen branch curves with prescribed \([8,6]\) or \([7,7]\) behavior in his attempted exclusion of alternative cases. Their genus and base-point computations are conditional on a fully resolved cluster and on the adjoint conditions; they do not prove the quintic exhaustion refuted by Horikawa. The spatial sextic with two skew double lines gives a further model through the congruence of lines meeting both lines. Its normalization and adjoints must be checked before its invariants are read off. The valid lower bound \(K^2\geq4\) for \(p_g=4\) is established independently in Section~\ref{cp2-viii-11}.

\section{Superficie con \texorpdfstring{$p=4$}{p=4} e \texorpdfstring{$p^{(1)}=7$}{p(1)=7}}
\label{cp2-viii-3}

\textbf{The simple sextic and the double cubic (pp.~271--273).}
A canonical sextic with double plane cubic has the latter plane as the fixed part of its adjoint quadrics; the residual planes give four canonical sections. With the stated ordinary singularities it has \(p_g=4\) and \(K^2=6\). The alternative sextic with six ordinary triple points has the same numbers: a smooth sextic has \((p_g,K^2)=(10,24)\), and six independent ordinary-triple-point corrections give \((4,6)\). Neither calculation by itself proves existence for arbitrarily assigned singular points.

For the double cubic, let \(Y\subset\mathbb P^3\) be a smooth cubic, \(K_Y=-H\), and branch along \(B\sim4H\). Taking \(L=2H\) gives
\[
K_X=\pi^*H,\qquad K_X^2=6,\qquad
\chi(\mathcal O_X)=5,\qquad p_g=4,\qquad q=0.
\]
The branch curve has degree twelve and is a complete intersection with a quartic. In the plane representation of the cubic this becomes a degree-twelve curve with six ordinary quadruple points; the same numbers follow from the branch corrections. The completeness assertion is supported in the source by intersecting with the conic pencils on the cubic: a different branch class would produce the unwanted genus-two pencils. A modern proof must make that divisor-class argument on the resolved cubic, rather than infer a class from its degree alone.

The pencil \(F_3^2+\lambda F_6=0\) is a flat family of sextic equations, with a double cubic at \(\lambda=0\). On the cubic, the sextic cuts degree eighteen; subtracting twice the double plane cubic leaves residual branch degree \(18-2\cdot3=12\). It suggests the asserted specialization of double cubics into simple sextics. Flatness of these equations does not establish flatness of their normalizations or simultaneous resolutions; these further steps are needed for the claimed statement about the moduli of smooth minimal surfaces.

The count on p.~272 concerns the \emph{class} of the dual of a general plane sextic. Its value is \(6\cdot5=30\). A doubled cubic contributes twice the class six of a smooth cubic, leaving eighteen pencils of lines through branch points; three ordinary nodes reduce the count by six, leaving twelve branch points on a plane section of the cubic. This agrees with \(B\sim4H\).

\textbf{Other canonical degrees and the limits of the family discussion.}
A degree-ten plane branch curve with two standard \([3,3]\) points gives
\[
(p_g,K^2)=(6,8)-2(1,1)=(4,6).
\]
Its canonical system has two base points, so its moving square is four and its canonical image can be a double quadric. The source counts thirty-nine parameters in one construction and thirty-eight in another. Such counts require independence of singularity conditions and control of stabilizers; different counts do not establish different components or an inclusion of one family in another.

Pompilj's triple quadric cone supplies another canonical behavior: the canonical curves are trigonal and the branch curve on the cone has degree eighteen. Enriques's proposed degeneration of the simple sextic through a triple conic uses twenty-four critical points, six pairwise collisions, and an alignment condition on three of the resulting points. This is a geometric sketch, not a verification of the total-space singularities or a proof that all the relevant smoothings exist.

The correct modern comparison is considerably more refined than the list of three spatial types. Bauer, Catanese, and Pignatelli (2006, Introduction and Theorem~0.1), recalling Horikawa's classification, describe four irreducible components and eleven nonempty strata for \(p_g=4\), \(K^2=6\), and prove that there are at most two connected components. They show that a surface of type IIIb can deform to type II; this is a specific adjacency. Neither the old models nor their parameter counts can be identified directly with these components without further work.

\section{Superficie con \texorpdfstring{$p=4$}{p=4} e \texorpdfstring{$p^{(1)}=8$}{p(1)=8}}
\label{cp2-viii-4}

\textbf{The double curve and the first parameter count (pp.~273--275).}
For a simple canonical septic, Enriques proposes a double curve \(C_7\) of degree seven and genus four on a quadric cone, passing triply through its vertex. The construction takes a half of the indicated degree-fourteen divisor class on the cone. Divisibility must hold in the appropriate divisor class group on the resolution; even degree alone does not establish such a half-class.

In the degree-seven plane representation of the curve on the cone, the putative multiplicity pattern with a quadruple point and an infinitely near triple point is excluded because a ruling would have eight intersections with a curve of degree seven. The proposed pattern is instead a quintuple point followed by a double point. The genus and intersection calculation is consistent with the required hyperelliptic curve, provided the infinitely near conditions and their proximity relations are respected.

The count is seventeen parameters for the curves on the cone, minus seven for the automorphisms of the cone, giving ten projectively distinct curves. For a fixed \(C_7\), removing the cone from a septic leaves a quintic through \(C_7\), counted in a projective system of dimension twenty-five; one additional septic gives a system of dimension twenty-six. Thus the claimed family has dimension \(10+26=36\). At this point the source presents the residual-system dimension inductively and assumes that an appropriate septic exists outside the reducible subsystem. The sextic with five ordinary triple points and the prescribed tacnode provides a separate numerical witness:
\[
(10,24)-5(1,3)-(1,2)=(4,7).
\]
It does not by itself establish the asserted incidence statement for every \(C_7\).

\textbf{Franchetta's auxiliary rational surface.}
The construction attributed to Franchetta (1939) is usefully expressed on a blown-up plane. With the proper and infinitely near points situated as in the construction, put
\[
D=11h-5E_A-3\sum_{i=1}^9E_{B_i}-\sum_{j=1}^8E_{D_j},
\qquad
T=7h-3E_A-2\sum_{i=1}^9E_{B_i}-\sum_{j=1}^8E_{D_j}.
\]
The numerical checks are
\[
D^2=7,\quad K_YD=7,\quad p_a(D)=8,\qquad
T^2=-4,\quad K_YT=8,\quad p_a(T)=3,\quad DT=0,
\]
and
\[
D\sim2T+K_Y.
\]
Thus, if \(T\) is contracted by the map defined by \(D\), the restriction of \(D\) is self-residual in the manner required for the auxiliary septic with an isolated quadruple point. The equality of divisor classes is exact; the construction of that map still requires effectivity, freedom of the moving system, birationality, and a local description of the contracted curve.

The source gives nineteen parameters for these rational septics. For a fixed double curve and fixed quadruple point it counts six parameters; allowing the point to vary raises this to nine, and hence suggests at least ten parameters for the double curves. This is not, by itself, a proof of dominance onto the entire proposed family of \(C_7\)'s. In particular, the conclusion that every assigned curve occurs needs an incidence argument and cannot be recovered simply by adding the dimensions of source and fiber.

The general parameter count for complete nonspecial \(g^3_7\)'s on hyperelliptic curves of genus four is
\[
\dim\mathcal H_4+\dim\operatorname{Pic}^7(C)+\dim\operatorname{PGL}_4
=7+4+15=26.
\]
A triple point corresponds to a neutral triple for the series; the source describes the condition as the containment of twice the hyperelliptic \(g^1_2\), leaving a degree-three divisor. This explains the expected one-condition count, but its use for dominance still requires the appropriate incidence calculation. The worn final digit in the superscript on p.~276 is consistent with \(26\), the reading required by this dimension calculation; it does not support an unqualified emendation to \(28\).

\textbf{The special triple-line model and the asserted degeneration (pp.~276--277).}
A septic with a triple line \(a\) and two double conics in planes through \(a\) is an explicit special configuration. Its limiting double-curve cycle is \(3a+C_2+C'_2\); the coefficient three records the contribution of the triple line to that cycle. Removing the two planes leaves quintics through the line and the conics, in a system of dimension thirty-one; restoring the planes gives thirty-three parameters, and the cross-ratio contributes one further parameter after the configuration is allowed to vary. The resulting count is thirty-four.

Enriques seeks to put these septics in the closure of the thirty-six-dimensional family. His construction projects a degree-eight curve in \(\mathbb P^4\) from a simple point and lets it split into two quartics, each double at the same point and sharing a tangent plane. The claimed nine conditions, the resulting triple line, and the asserted codimension two are consistent with the intended picture. They do not determine a flat family of the required embedded curve schemes or prove that the projected surfaces admit compatible normalizations. The proposed component inclusion therefore requires a flat family satisfying these incidence and normalization conditions.

\textbf{One canonical base point and a fixed exceptional line (pp.~277--279).}
The second simple model is a sextic with a double conic and a tacnode whose tangent plane is prescribed. For a minimal surface with \(K^2=7\), one simple canonical base point leaves moving square six. The source constructs sextic equations from a squared quadric times a squared plane and from the plane of the conic times a quintic with the indicated biplanar singularity.

In the intermediate nonminimal model, the fixed line \(a\) occurs twice in the canonical system. Write \(a^2=-\nu\). Adjunction gives
\[
p_a(2a)=-\nu-1,\qquad (2a)^2=-4\nu.
\]
The source obtains square \(10-4\nu\) and arithmetic genus \(8-\nu\) for the impure canonical divisor. The canonical identity relating them gives \(10-4\nu=(8-\nu)-1\), hence \(\nu=1\). Thus the relative linear genus is seven, whereas the absolute one is eight. The difference is an exceptional \((-1)\)-curve, not a change in the minimal invariant. The footnote's quadratic transformation gives the corresponding septic with triple conic and an infinitely near tacnodal condition; its exceptional genus-one contribution and its \((-1)\)-curve must be tracked on the resolved model.

The count of thirty-seven quintics, then thirty-eight sextics and a three-dimensional stabilizer, gives thirty-five parameters for this construction. These are parameters of the specified canonical sextic model; they are unrelated to the forty-dimensional moduli calculation for smooth canonical quintics in Section~\ref{cp2-viii-2}.

The proposed specialization from septics to this sextic plus a plane includes an infinitely near line in the limiting double curve. A projected degree-eight curve degenerates into a degree-six genus-three curve double at a point on the cone, together with a chord and a line through that point. Its projected limit is described by three times the conic, a line, and an embedded point at the projection center. Those scheme-theoretic features are essential to the arithmetic genus. Ignoring the embedded or infinitely near part would invalidate the proposed limit. The source does not supply a flat normalization or smoothing argument.

\textbf{The multiple canonical models and the explicit limitation (pp.~279--280).}
A double cubic with branch degree fourteen is represented by a double plane with degree-twelve branch curve, five ordinary quadruple points, and one \([3,3]\) point. The numerical check is
\[
(10,18)-5(1,2)-(1,1)=(4,7).
\]
Its canonical system has one base point and moving square six. In the pencil \(F_3^2+\lambda F_6=0\), the cubic section initially has degree eighteen; removing the contribution of the double conic leaves branch degree fourteen, including the line described in the source. As before, a pencil of equations is not yet a proof that the minimal surfaces belong to the same deformation family.

The double-quadric construction has a branch curve \(r+C_{11}\), an infinitely near pair of quadruple points \(A,A_1\) on \(r\), and \([3,3]\) points \(B,D\), with \(B\) on \(r\) and its first infinitely near point off \(r\). The canonical moving conics pass through \(A,D\), with \(r\) fixed in the indicated adjoint description. The printed collinearity \(A,A_1,D\) conflicts with the configuration just specified. Replacing the printed \(D\) by \(B\) is a proposed repair of this collinearity condition. The resulting double quadric has moving square four, leaving total canonical base contribution three, and the source counts thirty-three parameters. The independence of the cluster conditions and the proposed inclusion in a larger family remain unproved in this presentation.

Pompilj's exclusion of a triple quadric is invoked under its specific hypotheses on canonical curves and their trigonal series; it is not derived here. Enriques finally acknowledges that fixed canonical components have not been fully examined. That acknowledgement prevents the preceding list from being read as an exhaustive classification of all regular surfaces with \(p_g=4\), \(K^2=7\).

\section{Superficie con \texorpdfstring{$p=4$}{p=4} e \texorpdfstring{$p^{(1)}=9$}{p(1)=9}}
\label{cp2-viii-5}

\textbf{The octic construction (pp.~281--282).}
The proposed simple canonical octic has four triple points and a double curve \(C_{12}\) of degree twelve and genus ten on a cubic surface with four nodes. The divisor is obtained by halving the intersection of that cubic with an octic having multiplicity three at its four nodes. The text on p.~281 prints \(F_9\) at this step; \(F_8\) is required, since the complete intersection has degree \(3\cdot8=24=2\cdot12\). This is a printed degree slip.

In the plane representation used by Enriques, the half-system is a system of nonics with multiplicities \(4,1\) at each of the three pairs \(P,P_1\), \(Q,Q_1\), \(R,R_1\). The six proper and infinitely near points lie on the indicated conic; these conditions must be understood on the resolved cubic. A divisor of the correct degree is insufficient: the half-class must be integral in its Picard lattice and effective.

The book counts twenty-one parameters for \(C_{12}\). A cubic with four nodes is projectively unique in the specified class, and a fixed admissible \(C_{12}\) is expected to lie doubly on a thirteen-dimensional projective system of octics: removing the cubic leaves a twelve-dimensional system of quintics, and one additional octic completes the system. The resulting thirty-four parameters agree with
\[
10p-2p^{(1)}+12=40-18+12=34.
\]
The agreement is evidence for the intended construction, not proof that every proposed half-curve occurs as the double curve of an irreducible octic. The latter is precisely the missing incidence and existence assertion.

The second Franchetta construction has plane classes
\[
D=11h-3\sum_{i=1}^{12}E_{A_i}-\sum_{j=1}^{5}E_{B_j},
\qquad
T=7h-2\sum_{i=1}^{12}E_{A_i}-\sum_{j=1}^{5}E_{B_j}.
\]
They satisfy
\[
\begin{gathered}
D^2=8,\qquad K_YD=8,\qquad p_a(D)=9,\\
T^2=-4,\qquad K_YT=8,\qquad p_a(T)=3,\\
DT=0,\qquad D\sim2T+K_Y.
\end{gathered}
\]
These identities verify the numerical mechanism of the auxiliary rational octic with a quadruple point. The claimed seventeen parameters and the assertion that a quadruple point imposes seventeen conditions do not prove that the contraction has the required analytic type or that the resulting incidence map is dominant. Enriques himself marks the comparison as probable. A modern completion would require the linear system, its contracted divisor, and the relevant deformation space, rather than only the expected counts.

\textbf{The two double-cover families (pp.~282--283).}
A double quartic whose general section has genus two has spatial branch degree twelve and is represented by a degree-fourteen plane branch curve with one ordinary sextuple point and eight ordinary quadruple points. The corrected minimal invariants are
\[
(p_g,K^2)=(15,32)-(3,8)-8(1,2)=(4,8).
\]
The canonical adjoints are plane quartics double at the sextuple point and simple at the other eight points. The source's count
\[
119-19-64-8=28
\]
subtracts the imposed singularity conditions and projective equivalences under its independence assumptions. The comparison of twenty-eight with the expected thirty or more parameters in a larger family does not establish specialization or inclusion.

The double cubic has degree-sixteen branch curve on the cubic and a plane branch curve of degree twelve, with four ordinary quadruple points and two \([3,3]\) points:
\[
(10,18)-4(1,2)-2(1,1)=(4,8).
\]
There are two canonical base points, leaving moving square six, as required for a double cubic. The source assigns thirty-two parameters to this family. Its relation to the preceding double-quartic and simple-octic families is explicitly left unresolved; the numerical checks do not resolve it.

\section{Superficie con \texorpdfstring{$p=4$}{p=4} e \texorpdfstring{$p^{(1)}>9$}{p(1)>9}}
\label{cp2-viii-6}

For canonical degree \(n>8\), the construction on p.~283 calls for a half of a divisor on an adjoint surface of degree \(n-5\), with the appropriate multiplicities at its nodes. The first example is \(n=9\): the adjoint is a quartic with ten nodes, and on its resolution the candidate class is
\[
\frac{9H-3\sum_{i=1}^{10}E_i}{2}.
\]
Its numerical degree is compatible with the proposed double curve, but an arbitrary nodal quartic need not carry that integral half-class. Even when it exists, its effectivity, the dimension of its complete system, and the existence of an irreducible nonic with the prescribed double curve remain separate questions.

The source asks, rather than answers, whether the needed bisection exists on the nonrational adjoint, whether it is unique, and whether its curves are actually double curves of irreducible canonical surfaces. Different lattice or incidence types could lead to different families. The count \(50-2n\) is an expected moduli count for canonical degree \(n\), not an existence theorem. Here \(n=p^{(1)}-1\) is the canonical degree. The node count is a separate number, ten in the first example \(n=9\).

\section{Lacuna nella scansione}
\label{cp2-viii-7}

The scan breaks off after p.~283 and resumes at p.~286 with Section~8. Printed pp.~284--285, containing the end of Section~6 and Section~7, are missing from the only available scan. The edition marks this gap. The text and the heading of Section~7 remain unavailable; their contents cannot be determined from the adjacent pages.

\section{Superficie con \texorpdfstring{$p=5$}{p=5} e \texorpdfstring{$p^{(1)}=10$}{p(1)=10}}
\label{cp2-viii-8}

\textbf{A complete intersection of two cubics (pp.~286--287).}
The sextic with five ordinary triple points has \((p_g,K^2)=(5,9)\). Quadrics through its five triple points map the ambient three-dimensional space to the Segre cubic threefold; the ten joining lines are contracted to its ten nodes. The resulting canonical surface in \(\mathbb P^4\) has degree nine.

The assertion that it is the intersection of the Segre cubic with a second cubic has a short rigorous repair that does not assume projective normality. There are thirty-five ambient cubic forms, whereas
\[
h^0(X,3K_X)=\chi(\mathcal O_X)+3K_X^2=6+27=33.
\]
The kernel of restriction therefore has dimension at least two. One element is the Segre cubic, so there is an independent cubic through the image. The irreducible Segre cubic and this independent cubic have no common three-dimensional component. Their complete intersection has pure dimension two and degree nine. Since it contains the irreducible degree-nine canonical image, it is that image, with no residual component or extra generic multiplicity. This verifies the claimed complete intersection on the locus where the canonical map has the stated degree.

The Grassmannian of pencils in the thirty-five-dimensional space of cubics has dimension \(2(35-2)=66\). Subtracting \(\dim\operatorname{PGL}_5=24\) gives forty-two parameters for general \((3,3)\) complete intersections. This is a valid family calculation; it does not classify all surfaces with \(p_g=5\), \(K^2=9\).

\textbf{Projection and the even-node condition (pp.~287--288).}
A generic projection gives a nonic in \(\mathbb P^3\) with section genus ten and double curve of degree
\[
\binom{8}{2}-10=18.
\]
The six lines of the cubic threefold through the projection center give six triple points of the projected surface. Their tangent directions lie on the indicated quadric, so their projected points lie on a conic. The quartic adjoint consequently has six nodes lying on a trope conic. This special configuration carries divisor-class information that six arbitrary nodes on a quartic do not possess.

On the minimal resolution of the quartic, let \(H^2=4\) and \(E_i^2=-2\), with \(HE_i=0\). The class
\[
D_{16}=4H-\sum_{i=1}^{6}E_i
\]
has square fifty-two and arithmetic genus twenty-seven. If \(C\) is the trope conic, then \(C^2=-2\) and \(D_{16}C=2\). The resulting degree-eighteen class \(D_{18}=D_{16}+C\) has square fifty-four and genus twenty-eight. This checks the genus of the proposed half-system. Its projective dimension equals its genus when the required vanishing holds, for example for a nef and big divisor on the K3 resolution.

The more general assertion that every free linear system on a K3 surface has dimension equal to its arithmetic genus is false. If \(F\) is an elliptic pencil and \(m>1\), then \(p_a(mF)=1\) but \(\dim|mF|=m\). Thus the positivity hypothesis must be stated. Equally, the matching twenty-eight-dimensional counts for six-nodal quartics with and without the trope condition do not establish that an arbitrary such quartic admits the relevant even class.

Enriques writes \(x\leq28\) for the actual dimension of admissible degree-eighteen double curves, counts five surfaces for a fixed curve and twenty-eight quartics, and subtracts nineteen parameters for the ambient and projection choices:
\[
28+x+5-(15+4)=x+14.
\]
Comparison with forty-two gives \(x=28\) only if the incidence spaces have the asserted dimensions, dominate the relevant families, and have the stated generic fibers. This calculation does not prove that every abstract member of the complete half-system is the double curve of a canonical nonic.

\textbf{The multiple models and the next value of \(K^2\).}
The source asks whether a canonical system with one base point also occurs for \((p_g,K^2)=(5,9)\); in that case its moving square would be eight. It gives a double quartic from a degree-twelve branch curve with four ordinary quadruple points and one \([3,3]\) point:
\[
(10,18)-4(1,2)-(1,1)=(5,9).
\]
Its assigned forty-one parameters do not settle its deformation relation with the forty-two-dimensional complete-intersection family.

For \(p_g=5\), \(p^{(1)}=11\), the branch curve instead has three ordinary quadruple points and two \([3,3]\) points, giving \((5,10)\). Two canonical base points leave moving square eight, again appropriate to a double quartic. The branch degree on that quartic is sixteen. Indeed, a general canonical curve has genus eleven and maps doubly to a genus-two section, so Hurwitz gives
\[
2\cdot11-2=2(2\cdot2-2)+r,\qquad r=16.
\]
The \(C_{18}\) printed on p.~289 is therefore a slip, not a different branch calculation. The equality of the two proposed forty-dimensional counts does not prove that there are two distinct deformation families. The simple canonical decic is numerically supported by a sextic with four ordinary triple points and the specified tacnode:
\[
(10,24)-4(1,3)-(1,2)=(5,10).
\]
Existence of that model and completeness of an exhaustive classification are different assertions.

\section{Le superficie sopra una varietà a 3 o più dimensioni}
\label{cp2-viii-9}

\textbf{What is being extended (pp.~289--291).}
Enriques distinguishes general properties that extend formally from surfaces, properties that require new proofs, and genuinely new higher-dimensional phenomena. His references to Noether, Severi (1909), Rosenblatt (1928), and B.~Segre (1936) are historical sources for that program; their complete arguments are not reproduced here. In particular, this section does not license an unrestricted extension of surface regularity, rationality, or contraction criteria to threefolds.

For a smooth threefold \(V\), choose four sufficiently general sections of a divisor \(D\) defining a generically finite map to \(\mathbb P^3\), with no unresolved base contribution. The determinant of the differential gives the Jacobian divisor
\[
J(D)\sim K_V+4D.
\]
Consequently
\[
J(D+E)\sim J(D)+4E\sim J(E)+4D,\qquad
K_V\sim J(D)-4D,\qquad D'\sim D+K_V.
\]
These identities express the modern content of the Jacobian construction. If a chosen subsystem has a base locus or the model is singular, one must resolve it and include the discrepancy and ramification corrections before applying the displayed formulas.

A smooth hypersurface \(V_n\subset\mathbb P^4\) has \(K_V\sim(n-5)H\); the adjoint of a hyperplane section is therefore cut by hypersurfaces of degree \(n-4\). For a nonnormal threefold with an ordinary double locus, the adjoint condition is a conductor condition. Enriques lists the vanishing multiplicities \(i-1\), \(i-2\), and \(i-3\) along an ordinary multiplicity-\(i\) surface, curve, and point. These are statements about the indicated ordinary local models, not a formula for all singularities of the same multiplicity. On a resolution the precise replacement uses the conductor and the discrepancies.

\textbf{Adjunction, restriction, and the three genera.}
For a smooth divisor \(S\subset V\),
\[
K_S\sim(K_V+S)|_S.
\]
The converse assertion that a relation detected on a moving \(S\) determines a divisor on \(V\) needs injectivity of the relevant restriction on divisor classes. The assertion that all sections on \(S\) arise by restriction also needs the corresponding cohomological surjectivity. The absence of visible fundamental components alone proves neither. The exact sequence
\[
0\longrightarrow\mathcal O_V(D-S)\longrightarrow
\mathcal O_V(D)\longrightarrow\mathcal O_S(D)\longrightarrow0
\]
separates these issues.

The three historical genera singled out are the geometric genus \(P_g=h^0(V,K_V)\), the geometric genus of a canonical surface section, and the linear genus of that section. The text also introduces the numerical principal and superficial genera \(P_a\) and \(p_a\), referring to Severi for their definitions. These selected numerical and geometric characters are not a complete list of the modern Hodge or birational invariants of a threefold.

The numerical example on pp.~291--292 is correct under smoothness and effectivity assumptions. If \(S\in|K_V|\), put \(y=K_V^3\), and let \(C\) be the smooth intersection of two canonical divisors, of genus \(x\). Then
\[
K_S\sim2K_V|_S,\qquad K_C\sim3K_V|_C,\qquad
2x-2=3y.
\]
Thus
\[
K_S^2=4y,\qquad p^{(1)}(S)=4y+1=2x+y-1.
\]
In particular \(y\) is even whenever this construction has the indicated smooth complete-intersection curve. The arithmetic identities do not supply the existence of such a smooth \(S\) or \(C\).

For a smooth complete intersection of degrees \(m,n\) in \(\mathbb P^4\), adjunction gives \(K_S=(m+n-5)H\). In general, a smooth complete intersection of multidegree \((n_1,\ldots,n_c)\) in \(\mathbb P^r\) has
\[
K_S\sim\left(\sum_i n_i-r-1\right)H.
\]
For a residual intersection \(F+F^\alpha\), the source's adjoint condition is vanishing along the residual component. A rigorous formulation requires the scheme-theoretic residual, its locally Cohen--Macaulay or suitable duality hypotheses, and the corresponding linkage exact sequence. Merely subtracting degrees from an arbitrary reducible intersection does not prove the adjunction statement.

\section{Esempii di superficie canoniche iperspaziali}
\label{cp2-viii-10}

\textbf{The complete-intersection list (pp.~292--293).}
For a nondegenerate canonical complete intersection, \(K_X\sim H\) gives
\[
\sum_i n_i=r+2,\qquad n_i\geq2.
\]
There are \(r-2\) equations; hence \(2(r-2)\leq r+2\), or \(r\leq6\). With smoothness and completeness of the hyperplane system understood, the possibilities are
\[
\begin{array}{c|c|c|c}
\text{ambient space}&\text{multidegree}&p_g&K^2\\ \hline
\mathbb P^3&(5)&4&5\\
\mathbb P^4&(2,4)&5&8\\
\mathbb P^4&(3,3)&5&9\\
\mathbb P^5&(2,2,3)&6&12\\
\mathbb P^6&(2,2,2,2)&7&16 .
\end{array}
\]
This is a classification within the complete-intersection hypothesis. It is not a list of all canonical surfaces in these projective spaces.

\textbf{Sections of threefolds with prescribed anticanonical class (pp.~293--294).}
Let \(V\subset\mathbb P^{n+1}\) be a suitable threefold of degree \(n\) with \(K_V\sim-2H\). A cubic section \(X\in|3H|\) has \(K_X\sim H\). Assuming the stated complete hyperplane restriction,
\[
p_g=n+2,\qquad K_X^2=3n=3p_g-6,\qquad p^{(1)}=3p_g-5.
\]
The printed relation \(n=p-1\) is inconsistent with both the ambient dimension and these formulas; it must be \(n=p-2\). The source's rationality statement for the appropriate threefolds of degree greater than three must not be read as a claim that a smooth cubic threefold is rational.

Similarly, for a threefold of degree \(2r-4\) in \(\mathbb P^r\) with \(K_V\sim-H\), a quadric section is canonical and, with the same completeness hypotheses,
\[
p_g=r+1,\qquad K_X^2=4p_g-12,\qquad
p^{(1)}=4p_g-11,\qquad \deg V=2p_g-6.
\]
The existence and classification of the ambient threefolds are imported results, not consequences of these intersection computations. Enriques cites work of Fano, Godeaux, Bourniat, and Maxwell from 1936--1944 for the constructions. The book cites these constructions without proving their existence and classification results.

Cones over K3 surfaces require separate treatment at the vertex. The source itself warns that the preceding formulas need modification in that setting. Adjunction on a smooth threefold cannot be applied at a cone vertex without examining the relevant dualizing sheaf and the resolution.

\section{Minimo valore del genere lineare rispetto al genere superficiale}
\label{cp2-viii-11}

\textbf{Noether's inequality and the scope of the source proof (pp.~294--296).}
For a smooth minimal complex surface of general type,
\[
K_X^2\geq2p_g-4,
\qquad\text{equivalently}\qquad p^{(1)}\geq2p_g-3.
\]
Enriques proves the irreducible canonical case by applying Clifford's theorem to the characteristic series, and sketches the reduction when the moving part is reducible. The printed inequality \(p^{(1)}-1\leq2p-4\) on p.~295 has its sign reversed; the immediately preceding Clifford inequality and the subsequent argument require \(\geq\). The same page later prints \(2p+1\) where the preceding composed-pencil bound gives \(2p-1\), and changes the name of the base-point count from \(r\) to \(s\). These are errors in the printed formulas.

For a smooth canonical member \(C\), adjunction makes the characteristic bundle self-residual: \(2K_X|_C\sim K_C\). In the composed case the source obtains \(p^{(1)}\geq(p_g-1)(\pi-1)+1\geq2p_g-1\) under its hypothesis \(\pi\geq3\). Here is the complete regular-surface argument, including the cases not covered by that restriction. Resolve \(|K_X|\), writing
\[
\sigma^*K_X=M+Z
\]
with \(M\) free and \(Z\) effective. Suppose first that the image is a surface. For a smooth general \(C\in|M|\), the restriction sequence and \(q=0\) give \(h^0(C,\mathcal O_C(M))=p_g-1\). This series is special, since
\[
K_C-M|_C=K_{\widetilde X}|_C
\]
is effective. Clifford therefore gives \(M^2\geq2p_g-4\). Since \(M\) and \(\sigma^*K_X\) are nef,
\[
K_X^2-M^2=Z(M+\sigma^*K_X)\geq0.
\]
This proves the inequality without assuming that the original system was free or had no fixed curves.

If the canonical system is composed with a pencil, regularity makes the base of its Stein factorization rational. On \(X\), write its moving part as \(aF\), where \(a=p_g-1\) and \(F\) is a general irreducible member of the pencil. The divisor \(F\) is nef. If \(F^2=0\), then \(K_XF=2g(F)-2\), and general type forces \(g(F)\geq2\); hence
\[
K_X^2\geq aK_XF=2(p_g-1)(g(F)-1)\geq2p_g-2.
\]
If \(F^2\geq1\), nefness and the fixed-part intersections give
\[
K_X^2\geq a^2F^2\geq(p_g-1)^2\geq2p_g-4.
\]
This also treats pencils with base points. The cases \(p_g\leq1\) are immediate from \(K_X^2>0\). Thus genus-two pencils do not form an omitted exception to the theorem.

\textbf{Equality, freedom, and the double canonical image.}
Assume \(p_g\geq3\) and \(K_X^2=2p_g-4\). The pencil case above is excluded. In the surface-image argument all inequalities are equalities, so
\[
ZM=Z\sigma^*K_X=0,\qquad Z^2=0.
\]
The Hodge index theorem, since \(M^2>0\), implies that \(Z\) is numerically zero; a nonzero effective divisor cannot be numerically zero on a projective surface. Thus \(Z=0\): there are neither fixed components nor base points. This is the precise repair for the source's compressed passage from the Clifford equality to a free canonical system. Merely obtaining \(K_XZ=0\) would not suffice, because a minimal surface may contain \((-2)\)-curves.

For \(p_g\geq4\), the birational bound proved below excludes a simple canonical map on the Noether line. For \(p_g=3\), a birational canonical map to \(\mathbb P^2\) would make \(X\) rational and is likewise excluded. A nondegenerate surface in \(\mathbb P^{p_g-1}\) has degree at least \(p_g-2\); equation~\eqref{cp2-viii-moving-degree} now forces canonical degree two and image degree \(p_g-2\). The Clifford equality gives hyperelliptic general canonical curves. The \(S_{p-2}\) printed on p.~296 must therefore be \(S_{p-1}\), the canonical ambient space.

The image is a surface of minimal degree: a plane, the Veronese surface, or an appropriate rational scroll or cone. Its plane systems are respectively lines, conics, or curves of degree \(n\) with an \((n-1)\)-fold point and the specified additional simple points, including infinitely near ones. The passage to a complete list of double-plane branch models uses the classification of those rational systems and parity of the cover, not only Clifford's inequality.

Several numerical slips in the listed branch models can be repaired directly. A smooth octic branch gives \((p_g,K^2)=(3,2)\), and a smooth decic gives \((6,8)\), not the printed \(p=7\). For a degree-\(2n+6\) branch curve with one ordinary \(2n\)-fold point, the correct values are
\[
p_g=2n+1,\qquad K^2=4n-2,\qquad p^{(1)}=4n-1.
\]
The printed \(4n-2\) is the square \(K^2\), not the linear genus. Adding an ordinary sextuple point first gives \(p_g=2n-2\) and resolved \(K^2=4n-10\); the joining line lifts to two \((-1)\)-curves, whose contraction gives
\[
K_{\min}^2=4n-8,\qquad p^{(1)}=4n-7.
\]
Adding instead an ordinary quadruple point gives
\[
p_g=2n,\qquad K_{\min}^2=4n-4,\qquad p^{(1)}=4n-3.
\]
Thus the two further printed values \(4n-8\) and \(4n-4\) also omit the final \(+1\) in \(p^{(1)}\). The stated \([5,5]\) and \([3,3]\) configurations introducing canonical base points are excluded from this free equality case. All these branch calculations remain conditional on existence, the ordinary singularity type, independent adjoint conditions, and the stated contractions.

\textbf{Castelnuovo's inequality (pp.~296--298).}
If the canonical map is birational, then
\[
K_X^2\geq3p_g-7,
\qquad\text{equivalently}\qquad p^{(1)}\geq3p_g-6.
\]
The source invokes Bertini's lemma for a simple self-residual series \(g_n^r\) on a curve of genus \(\pi\). Two general divisors \(G,G'\) are chosen with \(r-1\) general points in common. The assertion that they have no further common points needs the general-position and simplicity hypotheses. Counting conditions on canonical divisors gives
\[
2n-3r+1\geq\pi-1.
\]
Since self-residuality gives \(n=\pi-1\), this yields \(\pi\geq3r\), and the characteristic series has \(r=p_g-2\), \(\pi=p^{(1)}\). The calculation is sound once those curve-theoretic hypotheses and completeness of the characteristic series are supplied.

The following repair handles fixed parts and canonical base points uniformly. Resolve the canonical system as above and take a smooth general \(C\in|M|\). Set
\[
n=M^2,\qquad r=p_g-2.
\]
The restriction series defines a birational map of \(C\) into \(\mathbb P^r\), while
\[
2g(C)-2=M(M+K_{\widetilde X})\geq2M^2
\]
because \(K_{\widetilde X}=M+Z+E\), with \(E\) effective and \(M\) nef. Thus \(g(C)\geq n+1\). Castelnuovo's genus bound for a nondegenerate birational curve of degree \(n\) in \(\mathbb P^r\) is strictly smaller than \(n+1\) when \(n\leq3r-2\). Hence \(n\geq3r-1\), and
\[
K_X^2\geq M^2\geq3p_g-7.
\]
Here \(r\geq2\); a birational series into \(\mathbb P^1\) would make \(C\) rational and contradict the genus estimate. This proves the form of the inequality actually used throughout the chapter, without silently assuming a free canonical embedding. The curve-theoretic input is the classical Castelnuovo bound; the surface steps are adjunction, nefness, and the restriction sequence.

\textbf{The equality examples and the scroll calculation (pp.~298--299).}
The quintic \((p_g,K^2)=(4,5)\) and the \((2,4)\) complete intersection \((5,8)\) attain equality. The second example is printed as \(F_8\), rather than \(F_6\). The third example lies on the degree-four \emph{threefold} cone over the Veronese surface in \(\mathbb P^6\): a quartic section containing a two-dimensional quadric cone leaves a residual surface of degree fourteen, with the asserted \((p_g,K^2)=(7,14)\). The degree computation is \(4\cdot4-2=14\). Enriques explicitly leaves the verification that the residual surface is canonical to the reader; its behavior at the vertex and the completeness of its canonical system require resolution.

For the general scroll construction, let \(V=\mathbb P(\mathcal E)\to\mathbb P^1\) have rank-three bundle \(\mathcal E\), tautological class \(H\), fiber-plane class \(F\), and \(\deg\mathcal E=p_g-3\). Then
\[
H^3=p_g-3,\qquad H^2F=1,\qquad
K_V\sim-3H+(p_g-5)F.
\]
For
\[
X\sim4H-(p_g-5)F,
\]
adjunction gives \(K_X\sim H|_X\), and
\[
K_X^2=H^2X=4(p_g-3)-(p_g-5)=3p_g-7.
\]
The splitting of \(\mathcal E\) matters: effectivity, smoothness, freedom, and completeness of the hyperplane system must be verified for the chosen scroll. The formula is not an existence assertion for an arbitrary rank-three bundle of that degree.

The spatial representation on p.~299 has degree \(2p_g-3\), a \((p_g-4)\)-fold line \(r\), and \(p_g-4\) triple lines meeting \(r\). Adjoint surfaces of degree \(2p_g-7\) split off the \(p_g-4\) planes through those pairs of lines, leaving degree \(p_g-3\), exactly the representing hyperplane system. This verifies the intended adjoint mechanism under the stated singularity assumptions. The printed ambient space \(S_{p-7}\) is \(S_{p-1}\).

\section{Limite superiore del genere lineare \texorpdfstring{$p^{(1)}$}{p(1)} rispetto al genere superficiale}
\label{cp2-viii-12}

On p.~299 Enriques quotes Rosenblatt's bound
\[
p^{(1)}\leq16p_a+27
\]
from a 1912 note, without reproducing its proof. The attribution identifies the historical source of the quoted bound. The subsequent Severi--Hodge bound has the transparent modern derivation
\[
p^{(1)}\leq2p_g+8p_a+11-\rho,
\]
where \(\rho\) is the Néron--Severi rank. Indeed, Noether's formula and Hodge decomposition give
\[
c_2=12\chi(\mathcal O_X)-K_X^2,\quad
b_2=c_2-2+4q,\quad
h^{1,1}=b_2-2p_g
=2p_g+8p_a+10-K_X^2.
\]
Since \(\rho\leq h^{1,1}\), the displayed inequality follows. This is a bound on the smooth minimal surface, with \(\rho\) interpreted as the rank of numerical divisor classes; the historical number of independent continuous systems is not itself a new topological definition.

The moduli expression
\[
M=10p-2p^{(1)}+12+\omega
=10\chi(\mathcal O_X)-2K_X^2+\omega
\]
continues Part~1, Chapter V, Sections~\ref{cp1-v-13} and~\ref{cp1-v-14}. Its first term is the deformation-theoretic expected dimension:
\[
\chi(T_X)=2K_X^2-10\chi(\mathcal O_X).
\]
For a minimal general-type surface \(h^0(T_X)=0\), and the local deformation space has dimension at least \(h^1(T_X)-h^2(T_X)=10\chi-2K^2\). Actual dimensions depend on obstruction equations; a single symbol \(\omega\) does not automatically encode a globally defined unobstructed excess or determine the number of components. Enriques's proposed interpretation through conditions imposed by cuspidal points on adjoints requires precisely the independence and deformation comparisons missing from several earlier counts. His statement that the geometric question remains open describes the state of his exposition, not the present state of geographical inequalities.

\section{Sistemi canonici appartenenti ad un'involuzione}
\label{cp2-viii-13}

\textbf{Double covers and the distinction between an obstruction and a converse (p.~300).}
If a surface is a double cover of a rational surface \(Y\), then
\[
\pi_*\omega_X=\omega_Y\oplus\omega_Y(L).
\]
Because \(h^0(Y,K_Y)=0\), all canonical sections come from the second summand, and ratios of canonical sections are invariant under the deck transformation. Thus the canonical map factors through the double cover. This proves the claim under the smooth-cover or resolved-cover hypotheses, including the correct branch corrections.

The inequality \(K^2<3p_g-7\) obstructs a birational canonical map, but its converse is false. A smooth double plane with branch degree \(2n\) has
\[
p_g=\frac{(n-1)(n-2)}2,\qquad K^2=2(n-3)^2,
\]
so \(K^2/p_g\) tends to four while the canonical map remains of degree two. A numerical region in which birationality is permitted is not a guarantee of birationality.

\textbf{The even twenty-node quintic (pp.~300--301).}
The proposed double quintic with twenty branch nodes must be interpreted on the resolution \(Y\) of a nodal quintic. If its exceptional curves form an even set,
\[
2L\sim\sum_{i=1}^{20}E_i,
\]
then \(L^2=-10\), \(K_YL=0\), \(K_Y^2=5\), and \(\chi(\mathcal O_Y)=5\). The double cover has \(\chi=5\) and resolved canonical square \(-10\); contracting the twenty ramification \((-1)\)-curves gives minimal \(K^2=10\). The conclusion \(p_g=4\) also requires \(q=0\). Twenty nodes alone do not provide the even divisor class or the vanishing needed for regularity. The proposed degeneration of a canonical decic to a doubled quintic is therefore an invitation to construct this particular even-node cover, not a proof that every twenty-nodal quintic supports it.

\textbf{The degree-four and degree-three canonical maps of B.~Segre.}
The sextic with two specified oscnodes is assigned \((p_g,K^2)=(4,8)\) and a canonical map of degree four to a quadric. Its local singularity corrections are imported from the cited construction, rather than proved from the word \emph{oscnodo}. An independent double-plane calculation uses a degree-eighteen branch curve with eight ordinary sextuple points:
\[
(28,72)-8(3,8)=(4,8).
\]
The canonical plane sextics double at those eight points belong to the Bertini involution, giving a canonical map of degree four to a quadric cone. The printed \(p^{(1)}=11\) for this supposedly identical case on p.~301 is wrong: \(K^2=8\) gives \(p^{(1)}=9\), in agreement with the first sextic. The assertion that the two constructions lie in the same family remains a question in the source.

The sextic with an oscnode and a triple point accompanied by an infinitely near double line gives \((p_g,K^2)=(4,6)\) and a triple quadric. The congruence of conics explains the three-to-one map, subject to the stated singularity calculation. This is the geometric realization of the quoted Castelnuovo observation about trigonal canonical curves.

\textbf{Pompilj's three cases (pp.~302--303).}
Under the specified irreducibility and trigonal hypotheses, the source imports Pompilj's list: a triple plane with \((p_g,K^2)=(3,3)\) and branch degree twelve; a triple quadric cone with \((4,6)\) and branch degree eighteen; and a triple rational cubic cone with \((5,9)\) and branch degree twenty-four. The classification is not proved in the book. Enriques cites the 1940--1941 papers as the historical sources for this classification.

Their displayed equations have the common form
\[
z^3+3p_{2d}z+q_{3d}=0,\qquad d=2,3,4,
\]
with discriminant
\[
-27\bigl(4p_{2d}^{\,3}+q_{3d}^{\,2}\bigr).
\]
Thus the reduced branch equation is \(4p_{2d}^{\,3}+q_{3d}^{\,2}=0\), with the appropriate cluster corrections. In the degree-eighteen case the specified multiplicities are \(p_6(O^3O_1^3)\) and \(q_9(O^5O_1^4)\); in the degree-twenty-four case they are \(p_8(O^5O_1^2O_2^2)\) and \(q_{12}(O^8O_1^4O_2^4)\). Their proximity data are part of the construction. In the third printed cubic equation on p.~303 the factor \(z\) after \(3p_8\) is missing; it is required by both the cubic equation and its displayed discriminant.

\textbf{The trigonal characteristic-series argument.}
Let a curve of genus \(g>3\) carry a \(g^1_3\) and a self-residual series. Choose a pencil of degree \(n=g-1\) in that series. If it is not composed with the trigonal pencil, Enriques proposes a plane model of degree \(m=n+3\), with an \(n\)-fold point \(A\) and a triple point \(B\). Lines through \(A\) cut the \(g^1_3\), while lines through \(B\) cut the \(g^1_n\). A pair of lines through \(B\), belonging to a canonical adjoint of degree \(m-3\), would leave a curve of degree \(m-5\) with multiplicity \(n-1=m-4\) at \(A\), an impossibility. This proves the intended contradiction once the birationality of the plane model and its adjoint interpretation are justified; they are assumptions in the compressed presentation.

The closing examples with further ordinary triple points use the already established correction \((-1,-3)\): they lower \(p_g\) to three in multiple-plane cases, or to two in pencil cases, while lowering \(K^2\) accordingly. They require independent adjoint conditions and do not classify all canonical pencils. Their role is to show that the construction survives beyond the initial \(p_g>3\) range, which motivates the pluricanonical models in the remaining sections.

\section{Superficie di genere lineare \texorpdfstring{$p^{(1)}=2$}{p(1)=2}: primo caso \texorpdfstring{$p=2$}{p=2}}
\label{cp2-viii-14}

\textbf{The numerical reduction and the double cone (pp.~303--305).}
When \(K_X^2=1\), Noether's inequality gives \(p_g\leq2\). For \(p_g=2,1,0\), respectively, equation~\eqref{cp2-viii-plurigenera} gives
\[
P_2=4,3,2,\qquad P_3=6,5,4.
\]
These are dimensions of spaces of sections; the corresponding projective spaces have dimensions one smaller. This distinction will matter in Sections~\ref{cp2-viii-20} and~\ref{cp2-viii-21}.

For \(p_g=2\), the canonical system is a genus-two pencil. In the standard canonical model it has one simple base point. If \(C\) is a smooth canonical member, then
\[
2K_X|_C\sim K_C.
\]
The degree-one divisor cut by another canonical member is therefore a Weierstrass point: its double is the hyperelliptic \(g^1_2\). The bicanonical system maps doubly to a quadric cone in \(\mathbb P^3\). Its branch curve has spatial degree ten, with additional special behavior at the vertex corresponding to the canonical base point. The vertex is not an ordinary branch point of a double cover of a smooth surface. A rigorous treatment either resolves the cone or uses the canonical ring.

\begin{example}[An explicit canonical model]
\label{cp2-viii-weighted-example}
Let
\[
X=\{w^2=f_{10}(x,y,z)\}\subset\mathbb P(1,1,2,5)
\]
be a general quasismooth hypersurface, with nonzero \(z^5\) and \(w^2\) coefficients. It avoids the isolated quotient singularities of the ambient weighted projective space and is smooth in its smooth locus. Weighted adjunction and the intersection calculation give
\[
K_X\sim\mathcal O_X(1),\qquad
K_X^2=\frac{10}{1\cdot1\cdot2\cdot5}=1,\qquad
H^0(X,K_X)=\langle x,y\rangle.
\]
The hypersurface cohomology calculation gives \(q=0\); hence \(p_g=2\). The bicanonical sections are
\[
x^2,\quad xy,\quad y^2,\quad z.
\]
Their first three coordinates satisfy the quadric-cone relation, and the omitted generator \(w\) gives the double cover.

The equation \(x=y=0\) defines one point on \(X\). On the chart \(z\ne0\), the two nonzero values of \(w\) are exchanged by the residual sign action; the stabilizer at either representative is trivial. The coordinates \(x,y\) are therefore local parameters. Thus the canonical pencil has exactly one simple base point. This calculation verifies both the claimed base point and its multiplicity, rather than inferring them only from \(K_X^2=1\).
\end{example}

The plane branch curve of degree ten with a \([5,5]\) point is a birational description of this geometry. Its parity resolution is not obtained by subtracting two independent ordinary quintuple-point corrections. The source's fixed-component argument says that a moving genus-two pencil and its adjoint lead back to the same canonical model. On a smooth resolution, any fixed components and \(K_X\)-trivial curves must be tracked; nefness alone does not make all such curves disappear. Bombieri (1973, Main Theorem, p.~173) gives the modern exceptional-case statement used below. The \(p=1\) at the end of the discussion on p.~305 is a printed slip for \(p=2\).

\section{Superficie con \texorpdfstring{$p^{(1)}=2$}{p(1)=2} e \texorpdfstring{$p=1$}{p=1} (\texorpdfstring{$P=3$}{P=3}, \texorpdfstring{$P_3=5$}{P3=5})}
\label{cp2-viii-15}

\textbf{The restriction argument and the projected sextic (pp.~305--307).}
Here \(K_X^2=1\), \(p_g=1\), \(P_2=3\), and \(P_3=5\). The unique canonical divisor has arithmetic genus two; Enriques first assumes it irreducible. A smooth member \(C\in|2K_X|\) has square four and genus four. Its complete characteristic series has projective dimension one, since regularity and restriction give
\[
h^0(C,\mathcal O_C(C))=h^0(X,2K_X)-1=2.
\]
The source uses this completeness to exclude the proposed hyperelliptic composite system. The restriction of \(|3K_X|\) to \(C\) is its complete canonical system: the restriction sequence has kernel \(H^0(K_X)\) and vanishing obstruction \(H^1(K_X)=0\). To deduce birationality on \(X\), one must also separate different bicanonical members and control fixed and base parts. Nonhyperellipticity of one smooth curve is not alone a global separation theorem. The true pluricanonical conclusion is covered by Bombieri's analysis.

When the indicated tricanonical map is free and birational, its image has degree nine in \(\mathbb P^4\), with section genus seven. The canonical curve maps to a triple line. Projection from a point of that line has degree \(9-3=6\), and the resulting sextic has three exceptional lines and three coplanar cuspidal double lines through the image point \(O\). A general plane section has genus seven, while a plane through \(O\) cuts a genus-four curve after normalization. The source excludes a single triple line because it would give the unwanted elliptic pencil.

The singularity in a plane through \(O\) is described on p.~306 as two branches with fifth-order contact, followed by the expression \emph{contatto sestipunto}. A modern analytic local equation has not been established from that phrase alone.

The adjoint quadric is the plane of the three double lines counted twice. The source counts a projective seven-dimensional system of adjoint cubics, with fourth-order contact along the two local branches, and concludes \(p_g=p_a=1\). To justify regularity, the cubic adjoints must cut the complete adjoint series on a section; this is a conductor and restriction statement, not a consequence of the number of equations written down. Once it is established, \(P_2=3=\chi+K^2=2+K^2\) indeed gives \(K^2=1\). The projected degree check \(6=9K^2-3\) agrees.

The constructive pencil begins with two cubic surfaces \(F_3,F'_3\) containing the three concurrent coplanar lines. The four-parameter contact system of second cubics has the stated contact at \(O\), is tangent along the three lines, and cuts an additional plane cubic on the first cubic. Their union is a sextic. It is combined with the sextic formed by three times the plane of the lines and a cubic cone with vertex \(O\). This is an actual construction sketch, but the preservation of the local singularity, the transition from tacnodal to cuspidal double lines, and the completeness of the adjoints still need local verification. The second occurrence is printed as \(F_3+F'_2\); it must be \(F_3+F'_3\). In equations, the union is represented by the product of the two cubic forms, not by their polynomial sum.

\textbf{The fixed-part calculation (pp.~307--308).}
Write \(|2K_X|=Z+|M|\), with \(M^2=n\) and \(p_a(M)=\pi\). The displayed identities become
\[
K_XM=2\pi-2-n\leq2,\qquad
MZ=4\pi-4-3n.
\]
The inequality \(MZ\geq2\) for nonzero \(Z\) uses two-connectedness, stronger than connectedness, from the earlier canonical-divisor theory. Together with the separate premise \(n>1\), these give
\[
(MZ)+n=2K_XM\leq4,\qquad n=2,\qquad \pi\leq3.
\]
This numerical step is valid under those hypotheses. The final exclusion, by representing the surface as a double plane branched in an octic with two elementary singularities and obtaining \(K^2=0\), requires the omitted analysis of the exact branch clusters and subsequent contractions. The inequalities alone are not an exhaustive double-plane classification. Accordingly, the source's declaration that the classification is now complete is stronger than the proof supplied on these pages.

\section{Superficie con \texorpdfstring{$p^{(1)}=2$}{p(1)=2} e \texorpdfstring{$p=0$}{p=0} (\texorpdfstring{$P=2$}{P=2}, \texorpdfstring{$P_3=4$}{P3=4})}
\label{cp2-viii-16}

\textbf{The false universal nonic and the correct tricanonical theorem (pp.~308--311).}
A regular surface with \(K^2=1\), \(p_g=0\) is a numerical Godeaux surface. Equation~\eqref{cp2-viii-plurigenera} gives \(P_2=2\), \(P_3=4\). If its tricanonical system is free and birational, its image in \(\mathbb P^3\) has degree nine, section genus seven, and ordinary double-curve degree
\[
\binom{8}{2}-7=21.
\]
The absence of quintic adjoints expresses \(p_g=0\), and the proposed five-dimensional vector space of degree-ten biadjoints cuts a pencil of bicanonical sextics of genus four after the double curve is removed. These are conditional model calculations; they depend on the normalization, conductor, and completeness of the biadjoint restriction.

The assertion that \(|3K_X|\) has no base points is false. Tricanonical birationality is true, as proved by Miyaoka (1976). Catanese (1981, Section V, Proposition~9, Proposition~11, Corollary~12, and Theorem~3, pp.~260--264) gives an accessible account of both statements. The base points are simple, are not infinitely near, and are governed by torsion classes \(t\ne-t\). Their number is half the number of such classes. For cyclic torsion of orders one or two there are none; for orders three or four there is one; for order five there are two. The resulting tricanonical image degrees are nine, eight, and seven, respectively.

The precise defect in the source's restriction argument is a change of quantifier. A \emph{general} bicanonical member may be smooth and irreducible, while the \emph{special} member through a tricanonical base point is reducible. In Catanese's description it is \(D_t+D_{-t}\), and the base point is their transverse intersection. Restricting to that special member as if it were a smooth genus-four curve invalidates the claimed base-point exclusion. The source's brief phrase that a similar analysis excludes base points supplies no missing proof. In the preceding fixed-component argument the printed \((K_2-|\theta|)'\) also misplaces the system bars: the residual adjoint meant is \(|K_2-\theta|'\). The following calculation verifies a counterexample satisfying the stated irreducibility hypothesis.

\begin{example}[An explicit septic tricanonical image]
\label{cp2-viii-godeaux-example}
Let
\[
Y=\{x_1^5+x_2^5+x_3^5+x_4^5=0\}\subset\mathbb P^3,
\]
and let \(G=\mu_5\) act with weights \(1,2,3,4\). Every nonidentity element has distinct eigenvalues, so its projective fixed points are the four coordinate points. None belongs to \(Y\). Thus \(X=Y/G\) is smooth and \(Y\to X\) is étale. Since \(K_Y^2=5\), one has \(K_X^2=1\). The canonical linearization has determinant weight \(1+2+3+4=10\equiv0\pmod5\), so the four canonical sections have the four nontrivial characters and \(p_g(X)=0\). Also \(q(X)=0\), because \(q(Y)=0\); equivalently, étale multiplicativity gives \(\chi(X)=\chi(Y)/5=1\).

The invariant quadrics are exactly
\[
x_1x_4,\qquad x_2x_3.
\]
A general nonzero linear combination defines a smooth quadric. Its intersection with \(Y\) is smooth by Bertini away from the base locus. At a base point, the differential of the quadric has a nonzero component transverse to the coordinate line containing that point, while the differential of the Fermat equation has its nonzero components along that line; hence the intersection is also smooth there. On the quadric the curve has class \((5,5)\), so it is ample and connected. A smooth connected curve is irreducible. Its étale quotient is therefore a smooth irreducible general bicanonical member on \(X\).

The invariant cubic sections are exactly
\[
s_0=x_1^2x_3,\quad s_1=x_1x_2^2,\quad
s_2=x_2x_4^2,\quad s_3=x_3^2x_4.
\]
Their common zero set on \(Y\) is the intersection with the two coordinate lines \(\langle x_1,x_4\rangle\) and \(\langle x_2,x_3\rangle\). Each line meets \(Y\) in five points forming one free \(G\)-orbit. Hence \(|3K_X|\) has two base points. On the first line, two displayed sections have independent linear terms in \(x_2,x_3\); on the second, they have independent linear terms in \(x_1,x_4\). Both base points are simple.

Birationality can also be checked directly. On \(x_1=1\), write \(u=x_2\), \(v=x_3\), \(w=x_4\). The image ratios satisfy
\[
r=\frac{u^2}{v},\qquad s=\frac{uw^2}{v},\qquad t=vw,
\]
and conversely
\[
u^5=\frac{t^2r^3}{s},\qquad
v=\frac{u^2}{r},\qquad w=\frac{tr}{u^2}.
\]
The five choices of \(u\) form exactly one \(G\)-orbit. The induced map is therefore generically one-to-one on \(X\). After blowing up the two simple base points its moving square is \(9-2=7\), so its tricanonical image is a septic. This example is not a construction attributed to Enriques.
\end{example}

\textbf{The Campedelli example: numerical consistency, historical failure, and later existence.}
The double plane invoked on p.~308 has branch degree ten, one ordinary quadruple point, and five \([3,3]\) points. The intended invariant calculation is
\[
(6,8)-(1,2)-5(1,1)=(0,1).
\]
The canonical forms are represented by conics through the six proper singular points, so \(p_g=0\) requires that those points not lie on a conic. It also requires the specified singularities and absence of additional ones. The bicanonical plane quartics are double at the quadruple point and pass through the five pairs of infinitely near triple points. Controlling that system and further exceptional curves is necessary to obtain the asserted minimal model. Werner (1994, Section~2, Proposition~1) makes these hypotheses explicit, including that the appropriate quartic system is exactly a pencil.

The original Campedelli existence input is defective. Bombieri (1973, Section~13, p.~218) expressly says that the analogous \(K^2=1\), \(p_g=0\) construction is incorrect. Dolgachev's account identifies the obstruction in the proposed three-conic-plus-quartic branch configuration: it would yield two independent order-two torsion classes, impossible for a numerical Godeaux surface. Catanese (1981, p.~263) confirms the exclusion of \((\mathbb Z/2\mathbb Z)^2\), citing Reid's lemma and Dolgachev's lectures. Enriques does not reproduce Campedelli's particular configuration, so the correction concerns the historical existence argument, rather than a proof that the numerical branch type is impossible.

Werner (1994, Sections~2--4) supplies a valid later octic-plus-conic construction with the required branch and adjoint conditions and free tricanonical system. Together with Miyaoka's birationality theorem it gives a tricanonical nonic: \(P_3=4\) and \((3K)^2=9\). Thus such nonics do exist. This existence statement does not rescue their asserted universality. The book's four-parameter count
\[
65-8-5\cdot9-8=4
\]
still requires independence of the conditions and a valid quotient count, and lies below the deformation baseline eight; Werner's construction does not establish that count. The \(F_8\) printed in the comparison on p.~308 is a local slip for the nonic \(F_9\) being discussed, although genuine tricanonical octics exist for the separate torsion cases above. Modern terminology reserves \emph{numerical Campedelli surface} for \(K^2=2\), \(p_g=q=0\), not this \(K^2=1\) case.

\textbf{The conditional exclusion of a triple cubic (pp.~309--311).}
Assuming the false preliminary freeness assertion, a nonbirational tricanonical map has degree three onto a cubic, since its moving square would be nine and the nondegenerate image has degree at least three. A general tricanonical curve has genus seven and maps triply to an elliptic cubic; Hurwitz gives total ramification degree twelve. Calling it the degree of a reduced branch curve requires the corresponding simple-branch hypothesis.

For a smooth cubic target, the source chooses complementary twisted-cubic systems \(C_3,L_3\) with \(C_3+L_3\sim2H\). Their pullbacks have self-intersection three. The argument that both have genus four uses the total branch intersection twenty-four and excludes a lower genus: genus below three would force vanishing plurigenera, while genus three would force a hyperelliptic pencil incompatible with a free trigonal pencil. Those assertions need the completeness and fixed-part hypotheses of the earlier adjoint theory.

The characteristic \(g^1_3\) on the genus-four pullback curves is then asserted to be self-residual. Otherwise the projected adjoint curve has genus seven and adjunction gives
\[
7=4+p^{(1)}+3-1,
\]
forcing \(p^{(1)}=1\), contrary to \(p^{(1)}=2\). The printed reference to the genus of \(C\) at this point must be to the adjoint \(C'\). The calculation is valid once the proposed high-degree spatial model and the genus-seven projected adjoint have been established.

Next, the complete \(|2C|\) is assigned projective dimension six while the minimal system of products has dimension five. The inference that exactly one bicanonical member belongs to that subsystem, and also to the complementary subsystem \(|2L|\), requires a cohomological intersection argument; the difference of dimensions alone does not prove that particular uniqueness statement. In the plane representation, one subsystem is represented by conics and the other by decics quadruple at six points not on a conic. The claimed impossible common conic is a further incidence and cluster assertion, not a consequence of the total degrees alone.

For a singular cubic target, the plane sections through a double point are again assigned genus-four pullbacks. The argument uses their intersection with bicanonical curves to produce a fundamental curve \(\theta\), and then identifies the residual tricanonical system with the curves in question. It obtains the contradictory genus \(7-1=6\), instead of four. This step needs a resolution of the target, its branch, and the adjoint restriction across \(\theta\); those data are not supplied for every possible singular double point. The conditional cubic exclusions therefore do not repair the false base-point argument. Miyaoka's theorem supplies the true birationality conclusion independently. The final word \emph{bicanoniche} on p.~311 is a printed slip for \emph{tricanoniche}.

\section{Superficie di genere lineare \texorpdfstring{$p^{(1)}=3$}{p(1)=3}: primo caso \texorpdfstring{$p=3$}{p=3}}
\label{cp2-viii-17}

Noether's inequality gives \(p_g\leq3\) when \(K^2=2\), leading to the four cases \(p_g=3,2,1,0\) on pp.~311--324. The historical classification for the positive-genus cases is attributed to Enriques's note of May 1897, with the simple \(p_g=2\) model of the next section added here.

For \(p_g=3\), a smooth double plane branched in an octic has
\[
K_X=\pi^*\mathcal O_{\mathbb P^2}(1),\qquad
K_X^2=2,\qquad p_g=3,\qquad q=0,\qquad P_2=6.
\]
Its general canonical curve has genus three. The free canonical system and double-plane map also follow from the Noether-equality argument in Section~\ref{cp2-viii-11}.

The source's direct fixed-part calculation is more conclusive than an isolated reading of \(K_XZ=0\) suggests. Write \(K_X=M+Z\), \(M^2=n\), \(p_a(M)=\pi\). Nefness and adjunction give
\[
K_XM=2\pi-2-n,\qquad MZ=2\pi-2-2n\geq0,\qquad K_XZ\geq0.
\]
For the net under discussion \(n\geq2\). The displayed inequalities force
\[
n=2,\quad \pi=3,\quad K_XM=2,\quad MZ=0,\quad K_XZ=0.
\]
Since \(K_X^2=M^2=2\), they also give \(Z^2=0\). The Hodge index theorem applied to \(M\) now implies \(Z\equiv0\), and effectivity gives \(Z=0\). Moreover, after resolving any base points the nonrational surface must still map with degree at least two onto \(\mathbb P^2\); the moving square is at least two and already \(M^2=2\), so there is no base contribution. This completes the freedom argument. This excludes the fixed divisor, without excluding other \(K_X\)-trivial curves on the surface.

\section{Superficie con \texorpdfstring{$p^{(1)}=3$}{p(1)=3} e \texorpdfstring{$p=2$}{p=2} (\texorpdfstring{$P=5$}{P=5})}
\label{cp2-viii-18}

\textbf{Canonical restrictions and the two bicanonical behaviors (pp.~312--315).}
For \(p_g=2\), \(K^2=2\), one has \(P_2=5\). Assuming the moving canonical pencil has a smooth genus-three member \(C\), the sequence
\[
0\longrightarrow\mathcal O_X(K_X)\longrightarrow
\mathcal O_X(2K_X)\longrightarrow\omega_C\longrightarrow0
\]
is surjective on sections because \(H^1(K_X)=0\). Thus the bicanonical system cuts the complete \(g^2_4\) on \(C\). Products of the two canonical sections distinguish the parameters of the canonical pencil. Consequently, once fixed and base parts have been controlled, the nonhyperelliptic general member gives a birational bicanonical map, while the hyperelliptic general member gives a degree-two map. The completeness of the curve restriction explains this distinction; it does not by itself establish the global freedom asserted at the start of the source argument.

In the nonhyperelliptic case the image is an octic in \(\mathbb P^4\), with normalized section genus seven. The canonical curves are plane quartics with two common points. Their planes form the ruling of a rank-three quadric threefold with vertex line \(a\); the line itself is not contained in the surface. In coordinates the quadratic relation among the three products \(x^2,xy,y^2\) of canonical sections explains this quadric and its vertex line.

Projection from a suitable point of the quadric gives an octic in \(\mathbb P^3\) with two infinitely near quadruple lines and a double conic in the stated plane. Although the degree-eight image is a quadric--quartic complete intersection, its normalization is not a smooth \((2,4)\) complete intersection: the latter has \(p_g=5\), \(K^2=8\), and section genus nine. The conductor singularities must account for the drop by two in the section genus.

The proposed equation combines a square of a quartic with the product of the fourth power of the plane \(\omega\), the square of the conic plane \(\gamma\), and a quadric. The adjoints split off the prescribed powers of \(\omega\) and then \(\gamma\), leaving planes through the quadruple line that cut the canonical pencil. The biadjoint \(\omega^5\gamma^2\) of degree seven is intended to meet the surface only along the singular curve scheme, making the residual plane sections bicanonical. These are precise degree and factor calculations, but the conductor multiplicities, normalization, and completeness of restriction require proof. The phrase that one immediately verifies the model does not supply those checks.

\textbf{Two explicit special models inside the simple case.}
A degree-twelve plane branch curve with eight ordinary quadruple points gives
\[
(p_g,K^2)=(10,18)-8(1,2)=(2,2).
\]
The canonical pencil is represented by cubics through the eight points. The full bicanonical system has five sections. The invariant subsystem represented by plane sextics double at the eight points has only four sections and is not complete. The double-cover splitting supplies the additional section; equivalently, the full system contains the degree-twelve contact curves described in the text, while the four-section subsystem is composed with a map of degree four. Thus projection of the full octic from an external point can have a quadruple quadric-cone image without making the complete bicanonical map degree four.

For the K3 construction on p.~314, let \(Y\) be the resolution of a quartic with a trope conic and six nodes on it. Require
\[
2L\sim H+\sum_{i=1}^6E_i.
\]
Then \(L^2=-2\), and a double cover has
\[
\chi=3,\qquad K_{\mathrm{resolved}}^2=-4.
\]
The six ramification exceptional curves are \((-1)\)-curves, whose contraction gives \(K_{\min}^2=2\). The trope makes \(L\) effective with \(h^0(L)=1\), yielding \(p_g=2\), \(q=0\). Integrality of \(L\) and this section count are real hypotheses; six arbitrary nodes do not supply them.

The source writes the cover by adjoining \(\sqrt{\alpha\beta}\), with \(\alpha\) the trope plane and \(\beta\) the plane of the branch quartic. The conic is an apparent critical curve, while the six exceptional node curves carry the additional odd branch contribution. The involution pairs the general canonical members and fixes two: one above the branch quartic, and a hyperelliptic one above the conic, with its two intersections with the quartic and the six nodes giving eight branch points. These details explain why the cover lies in the simple bicanonical case even though it admits a second degree-two projection.

\textbf{The hyperelliptic case and the branch degrees.}
A degree-ten plane branch curve with two ordinary quadruple points and two \([3,3]\) points gives
\[
(6,8)-2(1,2)-2(1,1)=(2,2).
\]
The canonical pencil is represented by conics, each lifted as a double cover with eight branch points and hence genus three. The bicanonical plane system consists of quartics with two double base points and the two indicated tangency conditions. The source also gives a reducible branch presentation as a line plus a nonic with three triple points on the line and a pair of infinitely near triple points away from it; its parity and cluster data are part of that alternative description.

The bicanonical target is the appropriate degree-four del Pezzo surface in \(\mathbb P^4\), an intersection of two quadrics with the stated two double points at the canonical-pencil base points. Its general hyperplane section has genus one, and its double pullback has genus seven. Hurwitz gives branch degree twelve \emph{on that quartic}. This is compatible with the degree-ten branch in its plane representation; these are different polarizations on different models. Their confusion produces a printed error in the summary of Section~\ref{cp2-viii-21}.

\textbf{The moduli claim and the fixed canonical part.}
The claimed dimension \(24=9p-2p^{(1)}+12\) on p.~315 cannot be the dimension of a complete moduli component for these invariants. The deformation lower bound is
\[
10\chi-2K^2=10\cdot3-4=26.
\]
The earlier formula in the book used \(10p\), not \(9p\). Moreover, equality of two construction dimensions would not by itself establish two distinct components. The projected degeneration from an octic to a doubled quartic with two infinitely near double lines describes a candidate common special locus; a deformation analysis is needed to decide the actual relation of the loci. The parameter counts do not settle these component questions.

For a reducible canonical system, the moving part is still an irreducible pencil: composition with more than one fiber over the rational base would give at least three independent canonical sections. Write \(K_X=M+Z\), \(M^2=n\), \(p_a(M)=\pi\). If \(Z=\sum_j a_j\theta_j\), \(\rho_j=p_a(\theta_j)\), and \(\nu_j=\theta_j^2\), adjunction gives the exact identity
\[
(2\pi-2-n)+\sum_j a_j(2\rho_j-2-\nu_j)=K_X^2=2.
\]
Each summand \(2\rho_j-2-\nu_j=K_X\theta_j\) is nonnegative by nefness. Thus, with multiplicities included and self-intersections signed, the source uses
\[
K_XM=2\pi-2-n\leq2,\qquad
MZ=2\pi-2-2n\geq0,\qquad \pi\geq2,
\]
and obtains
\[
(\pi,n)=(2,0),\quad(2,1),\quad(3,2).
\]
For \((2,0)\), the adjoint \(K_X+M\) has square six and a four-dimensional space of sections. The proposed degree-two map to a cubic scroll in \(\mathbb P^3\) would contradict completeness if the polarization pulls back from the linearly normal cubic scroll, which has at least five independent sections. Establishing that resolved scroll map and identifying its pullback line bundle is the missing input in the source's compressed normality argument.

For \((2,1)\), the printed genus four and degree seven of the adjoint on p.~316 are incorrect. Adjunction gives \(K_XM=1\), so
\[
(K_X+M)^2=2+2+1=5,\qquad
p_a(K_X+M)=1+\frac{(K_X+M)(2K_X+M)}2=5.
\]
There is a direct repair: \(MZ=0\), \(M^2=1\), and \(Z^2=(K_X-M)^2=1\), contradicting Hodge index. Thus the case is excluded, but not by the erroneous numerical argument printed in the book.

For \((3,2)\), one has \(K_XM=2\), \(MZ=0\), and \(Z^2=0\). Hodge index again gives \(Z=0\), so this returns to the already discussed genus-three canonical pencil. The further assertion that this case necessarily has the double-quartic bicanonical model does not follow from these equalities: the distinction between hyperelliptic and nonhyperelliptic canonical members remains necessary.

\section{Superficie con \texorpdfstring{$p^{(1)}=3$}{p(1)=3} e \texorpdfstring{$p=1$}{p=1} (\texorpdfstring{$P=4$}{P=4})}
\label{cp2-viii-19}

\textbf{The birational octic and its contact surfaces (pp.~316--318).}
For \(p_g=1\), \(q=0\), \(K^2=2\), the bicanonical system has four sections. In the birational case its free map has degree-eight image in \(\mathbb P^3\), section genus seven, and ordinary double-curve degree
\[
\binom{7}{2}-7=14.
\]
The source identifies a degree-seven biadjoint \(\Phi_7\) and a quartic adjoint \(\Phi_4\), with the latter cutting the unique canonical curve of degree \(4\cdot8-2\cdot14=4\) away from the double curve. Its equation is of the form
\[
F_8=\Phi_4^2+\alpha\Phi_7.
\]
The degree-fourteen curve is described as the double curve of the rational degree-seven surface \(\Phi_7\), whose general hyperplane section is elliptic. This imports the earlier theory of rational surfaces with elliptic sections and its adjoint completeness hypotheses.

The phrase that \(C_{14}\) is the complete intersection of \(F_8\) and \(\Phi_7\) must be interpreted with its multiplicity. Their intersection has degree \(8\cdot7=56\), so the cycle is \(4C_{14}\), rather than a reduced complete intersection of degree fourteen. The book's own division by four records this fact. Scheme structure along the double curve remains relevant to the normalization and the canonical ring.

The source counts twelve invariants for \(C_{14}\) or \(\Phi_7\) and four parameters for octics with that curve, obtaining sixteen. This agrees with the dimension of the component constructed rigorously later, but the old parameter count alone does not prove completeness, smoothness, or irreducibility of a moduli component.

\textbf{The omitted degree-four maps.}
Enriques asserts that every multiple bicanonical image in this case is a double quartic, even when the unique canonical curve is assumed irreducible. This is false. Catanese and Debarre (1989, pp.~1--3) explicitly identify this error and list five cases, all of which occur:
\[
\begin{array}{c|c}
\text{degree of the bicanonical map}&\text{image}\\ \hline
4&\text{quadric cone}\\
4&\text{smooth quadric}\\
2&\text{rational quartic}\\
2&\text{K3 quartic}\\
1&\text{octic}.
\end{array}
\]
The torsion-\(\mathbb Z/2\mathbb Z\) case forms an irreducible rational sixteen-dimensional family and has degree-four cone image. The smooth-quadric degree-four locus is an eleven-dimensional unirational family of simply connected surfaces. It is therefore not enough to append a reducible-canonical-divisor qualification to Enriques's list.

\begin{example}[A degree-four bicanonical map with irreducible canonical curve]
\label{cp2-viii-bidouble-example}
On \(Y=\mathbb P^1\times\mathbb P^1\), choose smooth transverse divisors, with no triple intersection, of types
\[
D_1=(1,1),\qquad D_2=(3,1),\qquad D_3=(1,3).
\]
Form the smooth bidouble cover \(\pi:X\to Y\), assigning the three distinct nonzero inertia elements of \((\mathbb Z/2\mathbb Z)^2\) to the three branch divisors. The character line bundles are
\[
L_1=(2,2),\quad L_2=(1,2),\quad L_3=(2,1),
\qquad 2L_i\sim D_j+D_k.
\]
Its algebra is
\[
\pi_*\mathcal O_X=
\mathcal O_Y\oplus L_1^{-1}\oplus L_2^{-1}\oplus L_3^{-1}.
\]
The local covers at a transverse intersection of two branch curves have independent square-root coordinates, so the stated branch hypotheses give a smooth surface. The character decomposition and Riemann--Roch give
\[
p_g(X)=\sum_{i=1}^3h^0(Y,K_Y+L_i)=1,
\]
\[
\chi(\mathcal O_X)=4+\frac12\sum_{i=1}^3 L_i(L_i+K_Y)=2,
\qquad q(X)=0.
\]
The ramification formula gives
\[
2K_X\sim\pi^*\mathcal O_Y(1,1),\qquad K_X^2=2.
\]
In particular \(K_X\) is ample and this surface is minimal.

The three nontrivial summands in \(H^0(X,2K_X)\) are the spaces of sections of \((-1,-1)\), \((0,-1)\), and \((-1,0)\), all zero. Thus
\[
H^0(X,2K_X)=H^0(Y,\mathcal O_Y(1,1)),
\]
and the complete bicanonical map is the degree-four cover followed by the Segre embedding.

Let \(R_1\) be the reduced ramification divisor above \(D_1\). Since \(K_Y+L_1\sim0\), the canonical section is represented by \(R_1\); hence \(K_X\sim R_1\). The map \(R_1\to D_1\simeq\mathbb P^1\) has degree two and is branched at the eight transverse intersections with \(D_2+D_3\):
\[
D_1(D_2+D_3)=(1,1)(4,4)=8.
\]
It is connected, smooth, and of genus three. The unique canonical curve is therefore irreducible. This directly verifies a counterexample satisfying the source's own hypothesis, using the branch types listed by Catanese and Debarre.
\end{example}

\textbf{The double K3 quartic and its even nodes (pp.~318--320).}
For the first double-quartic case, the target has surface geometric genus one and smooth plane sections of genus three. The printed assertion on p.~318 that it is a quartic of genus two is a slip for surface genus one; it must not be repaired by assigning genus one to its plane sections.

On the K3 resolution \(Y\) of the quartic, the needed branch data are
\[
2L\sim H+\sum_{i=1}^{10}E_i,\qquad L^2=-4.
\]
The double cover has \(\chi=2\) and resolved \(K^2=-8\). Contracting its ten ramification \((-1)\)-curves gives \(K_{\min}^2=2\). The asserted \(p_g=1\), \(q=0\) additionally requires \(h^0(Y,L)=0\). The ten nodes must have the indicated parity with \(H\); they are not interchangeable with any ten nodes of an arbitrary quartic. This cover calculation also explains the source's number \(t=10\) of the relevant triple points in the octic construction.

The pencil \(\lambda F_8+F_4^2=0\) gives a doubled quartic as a limiting equation. The apparent contact curve \(C_{14}\) contributes even multiplicity, while the plane branch quartic and the exceptional curves over the nodes contribute odd multiplicity. The square-root expression \(\sqrt{\alpha\Phi_7}\), and its replacement by \(\sqrt{\alpha F_3}\), must be interpreted in the function field using an actual equality of contact divisors modulo twice a divisor. Merely replacing a polynomial by another of the same degree does not preserve the cover.

The numerical contact calculation can be made precise. With \(H^2=4\), \(E_i^2=-2\), and \(HE_i=0\), put
\[
D_{14}=\frac{7H-\sum_{i=1}^{10}E_i}{2}.
\]
Then
\[
\begin{array}{c|c|c}
\text{class}&\text{square}&\text{arithmetic genus}\\ \hline
D_{14}&44&23\\
D_{14}-H&20&11\\
D_{14}-2H&4&3 .
\end{array}
\]
The degree-fourteen and degree-ten restrictions to a smooth plane quartic have respectively twelve and eight sections by curve Riemann--Roch. To subtract these numbers from the dimensions of surface systems, as Enriques does, requires surjectivity of restriction; this follows from an appropriate vanishing statement or direct calculation for the stated nef and big classes, not from arithmetic genus alone. With these hypotheses the successive projective dimensions are \(23,11,3\), yielding the genus-three contact sextics of the cubic surfaces.

The complementary-sextic calculation also needs the distinction between the singular quartic and its resolution. Write
\[
D_6=D_{14}-2H=\frac{3H-\sum E_i}{2}.
\]
The total residual divisor in a cubic section is \(3H-D_6=D_6+\sum E_i\). Its intersection with \(D_6\) is fourteen, of which ten arise from the exceptional node contributions; the intersection with the proper residual sextic has degree \(D_6^2=4\). On a smooth member \(C_6\), adjunction identifies those four residual points with its canonical series \(g^2_4\). The actual spatial contact sextic must be distinguished from a plane section of the quartic. Its nonhyperelliptic canonical model is a plane quartic, but that is a different presentation. The count fourteen includes the exceptional contributions, while the proper residual curves meet in four points.

The source assigns one condition to the existence of a genus-three sextic and nine further conditions to the ten fixed points with canonical residual series, for ten conditions in total. It then obtains
\[
34-10+3-15=12
\]
parameters after choosing the branch plane and taking projective equivalence. Catanese and Debarre confirm dimension twelve for this degree-two K3-quartic locus. The old count by itself does not establish independence of the conditions or its claimed inclusion in a specified family.

\textbf{The double rational quartic and the actual inclusion result (pp.~320--321).}
The second degree-two case has a rational quartic target whose normalized plane sections have genus two; the spatial model has a double line. Its degree-twelve branch on the quartic corresponds on a plane model to a decic with one ordinary quadruple point and four \([3,3]\) points. The numerical check is
\[
(p_g,K^2)=(6,8)-(1,2)-4(1,1)=(1,2).
\]
The unique canonical curve is represented by the conic through the five proper singular points and is hyperelliptic. The bicanonical system is represented by quartics double at the quadruple point and satisfying the four infinitely near tangency conditions. The source counts
\[
65-8-36-8=13
\]
parameters, in agreement with the thirteen-dimensional unirational locus described by Catanese and Debarre.

The inequalities \(12<15\) and \(13<15\) printed in the book do not prove inclusion in the birational-octic family. The relevant deformation lower bound is sixteen, and a family of smaller dimension might lie in several different components. Catanese and Debarre prove a precise replacement: the degree-two K3-quartic and rational-quartic loci, together with the degree-four smooth-quadric locus, belong to a sixteen-dimensional irreducible component \(\mathcal M_E\) with a birational-bicanonical open locus. In their 1989 paper, irreducibility of the \emph{entire} torsion-zero moduli space is expressly a conjecture; their theorem concerns \(\mathcal M_E\).

Enriques's universal inclusion claim also misses the torsion-\(\mathbb Z/2\mathbb Z\) family with degree-four cone image, and he expressly declines to examine all reducible canonical divisors. The explicit bidouble example shows that the error already occurs with an irreducible canonical curve. The constructed double-quartic loci and their later inclusion in \(\mathcal M_E\) are valid, but they do not exhaust the possibilities.

\section{Superficie con \texorpdfstring{$p=0$}{p=0} e \texorpdfstring{$p^{(1)}=3$}{p(1)=3}}
\label{cp2-viii-20}

\textbf{Numerical Campedelli surfaces and three printed errors (p.~321).}
For \(p_g=q=0\), \(K^2=2\), the plurigenera are \(P_2=3\), \(P_3=7\). If the bicanonical net is free, its map has degree eight to \(\mathbb P^2\). The complete tricanonical system takes values in \(\mathbb P^6\), not in the \(S_7\) printed in the book: seven sections give projective dimension six. The arithmetic genus of a tricanonical member is
\[
p_a(3K_X)=1+6K_X^2=13,
\]
not the printed \(6p^{(1)}-2=16\). A general ordinary spatial projection of the degree-eighteen birational image therefore has double-curve degree
\[
\binom{17}{2}-13=123,
\]
not the printed \(120\). All three erroneous values occur in the printed text.

The proposed double plane has degree-ten branch curve and six \([3,3]\) points whose proper points lie on no conic. Under the stated branch and minimality hypotheses,
\[
(p_g,K^2)=(6,8)-6(1,1)=(0,2).
\]
The no-conic condition is the canonical-adjoint condition needed to obtain \(p_g=0\). Existence is invoked from Campedelli, rather than proved on this page. The count \(65-54-8=3\) describes a proposed special construction and is smaller than the local deformation lower bound \(10\chi-2K^2=6\); it does not compute the dimension of a whole component.

\textbf{The characteristic series and the modern tricanonical repair.}
For a smooth \(C\in|2K_X|\), regularity gives
\[
h^0(C,\mathcal O_C(C))=P_2-1=2.
\]
The complete characteristic series is therefore \(g^1_8\), not the printed \(g^2_8\). Since \(g(C)=7\), curve Riemann--Roch gives
\[
h^0(C,\mathcal O_C(2K_X))-h^1(C,\mathcal O_C(2K_X))
=8-7+1=2,
\]
so it is \emph{nonspecial}. Hyperellipticity of \(C\), even together with freedom of this complete series, does not by itself force the series to be composed with the hyperelliptic pencil. If that additional composition were established, freedom would identify the line bundle with four times the hyperelliptic pencil, giving five sections and contradicting the dimension two. The missing step is therefore a global argument establishing the asserted composition, together with the necessary fixed-part control. Freeness and irreducibility alone do not establish this composition.

Catanese proves that the tricanonical map of a numerical Campedelli surface is birational, following the earlier work with Bombieri (Catanese, 1981, Section IV, Theorem~2, pp.~258--259). His proof explicitly uses the absence of tricanonical base points and then excludes degrees two and three. Thus the degree-eighteen tricanonical model is valid. This modern theorem repairs the conclusion; it does not retroactively validate the incomplete restriction argument printed on p.~321.

\section{Superficie pluricanoniche semplici e multiple}
\label{cp2-viii-21}

\textbf{The three recalled bicanonical types and the double-plane reduction (pp.~321--323).}
The summary recalls:
\[
\begin{array}{c|c|c}
\text{type}&(p_g,K^2)&\text{plane branch description}\\ \hline
\mathrm I&(3,2)&\text{octic}\\
\mathrm{II}&(2,1)&\text{decic with a }[5,5]\text{ point}\\
\mathrm{III}&(2,2)&\text{decic with two quadruple and two }[3,3]\text{ points}.
\end{array}
\]
The \(C_{12}\) printed for the plane branch in type III on p.~322 is a slip for \(C_{10}\). The degree twelve previously obtained on its bicanonical quartic target is correct, but refers to a different model and polarization.

For types I and II, the bicanonical members are hyperelliptic and the tricanonical map is composite, as the explicit rings below verify. In general, canonical restriction to a nonhyperelliptic bicanonical member separates its points, but concluding birationality on the surface additionally requires separation of distinct members and control of fixed and base parts. The source then attempts to prove that types I and II are the only regular surfaces with \(p_g>0\), \(K^2>0\) having a multiple tricanonical map. Its strategy reduces the images of hyperelliptic bicanonical curves on a rational double-plane quotient to systems of lines, conics, or curves of degree \(n\) with an \((n-1)\)-fold point and \(s\) further simple points. The dimension of the last plane system is
\[
\binom{n+2}{2}-1-\frac{n(n-1)}2-s=2n-s,
\]
provided the conditions are independent and the infinitely near cluster is admissible. The permitted configurations of distinct and infinitely near simple points belong to the classical classification of rational plane systems; this is an imported ingredient.

The line case would give \(P_2=3\), \(p_g=1\), \(K^2=1\), excluded by the preceding analysis of the tricanonical map in Section~\ref{cp2-viii-15}. The conic case gives the double octic. In the third case, the \(n=2\) model with an octic branch and a \([3,3]\) point is quadratically related to a special decic with a \([5,5]\) point and an additional node. This is a relation between the specified special branch models, not a statement that an arbitrary decic is quadratically transformed to an octic.

For the general calculation, write the fixed part as
\[
\sum_i h_i\,OO_i,\qquad d=\sum_i h_i.
\]
The second-adjoint degree condition gives \(2m=n+d+6\). The branch curve must have multiplicity \(2m-5\) at \(O\), while the naive second-adjoint multiplicity is \(2m-8\), one less than required. The source therefore introduces a further infinitely near point of multiplicity \(2m-5\). This inference needs the full proximity and adjoint calculation for the cluster. If the joining line is not itself a branch component, Bézout gives
\[
2(2m-5)\leq2m,\qquad m\leq5.
\]
The positivity assumptions in this case leave \(m=4,5\), recovering the two branch degrees. If the line is a branch component, it must first be accounted for; the displayed Bézout inequality cannot simply be applied unchanged. This is a useful outline of the exceptional cases, but fixed divisors, branch parity, and completeness prevent it from replacing the modern global theorem.

\textbf{The two tricanonical exceptions and the errors in their stated models.}
For type I, \(K_X=\pi^*\mathcal O_{\mathbb P^2}(1)\), and all tricanonical sections are pullbacks of plane cubics. Thus \(P_3=10\), the map has degree two, and its image is the degree-nine Veronese plane in \(\mathbb P^9\). A general hyperplane section of that image is elliptic. This part of the summary is correct.

For type II, use Example~\ref{cp2-viii-weighted-example}. The complete tricanonical basis is
\[
x^3,\quad x^2y,\quad xy^2,\quad y^3,\quad xz,\quad yz.
\]
The target is \(\mathbb P^5\). The sole base point \(x=y=0\) is simple, since the last two sections have independent linear terms there. The moving square is therefore \(9-1=8\). The map has degree two, and its image is the rational normal scroll \(S(1,3)\), of degree four, isomorphic to the Hirzebruch surface \(\mathbb F_2\). Its general hyperplane section is rational. The source's sextic in \(S_6\) with elliptic sections is thus wrong in degree, ambient dimension, and sectional genus. Bombieri (1973, Main Theorem, p.~173(iii)(a)) explicitly gives the degree-four scroll.

A general tricanonical member in this model is smooth at the simple base point and has genus seven. Its degree-two map to a rational hyperplane section makes it hyperelliptic. This directly refutes the premise on p.~324 that the tricanonical curves under discussion are never hyperelliptic.

\textbf{The tetracanonical counterexample.}
The same canonical ring is generated in degrees \(1,1,2,5\), with its relation in degree ten. Its degree-four part has the basis consisting of five binary quartics in \(x,y\), the three products of \(z\) with binary quadratics, and \(z^2\). Hence \(P_4=9\). None of those sections involves \(w\), so the tetracanonical map still factors through the involution \(w\mapsto-w\). It is free, has degree two, and has image degree
\[
\frac{(4K_X)^2}{2}=\frac{16}{2}=8
\]
in \(\mathbb P^8\). Bombieri describes the underlying quadric cone in this higher embedding; it is not a degree-two surface in \(\mathbb P^8\).

At degree five the new generator \(w\) appears. On \(x\ne0\), the sections \(x^5,x^4y,x^3z,w\) recover
\[
y/x,\qquad z/x^2,\qquad w/x^5,
\]
so \(|5K_X|\) is birational for this explicit model. Thus both the nonhyperellipticity premise and the universal tetracanonical conclusion for \(p_g>0\) fail. They cannot be repaired by changing only the printed degree of the tricanonical image.

\textbf{The correct uniform theorem and an adjoint proof of birationality.}
For every smooth minimal complex surface of general type, \(|mK_X|\) is free and birational for \(m\geq5\); on the canonical model it gives an embedding. Bombieri (1973, Main Theorem) is the primary source. The following birationality argument uses Reider (1988), in the precise formulation of Peters (n.d., Theorem~7.2, p.~18).

Set \(L=(m-1)K_X\). It is nef, with \(L^2\geq16\). Reider's base-point criterion would produce an effective divisor \(D\) with
\[
(LD,D^2)=(0,-1)\quad\text{or}\quad(1,0).
\]
The second possibility is excluded because \(LD\) is divisible by \(m-1\geq4\). In the first, \(K_XD=0\), and \(D(D+K_X)\) would be odd, contrary to adjunction. Thus \(|K_X+L|=|mK_X|\) has no base points.

If two general points were not separated, Reider's separation criterion, valid for \(L^2\geq10\), would give an effective \(D\) through them with
\[
\begin{array}{c|c}
LD&D^2\\ \hline
0&-1\text{ or }-2\\
1&-1\text{ or }0\\
2&0 .
\end{array}
\]
Divisibility excludes \(LD=1,2\). If \(LD=0\), every irreducible component of \(D\) is \(K_X\)-trivial and lies in the exceptional locus of the canonical contraction. General pairs of points avoid that finite union of curves. Hence the map separates general points and is birational. This argument uses smoothness, characteristic zero, and nefness and bigness of \(K_X\); it does not assert very ampleness on a minimal resolution containing \((-2)\)-curves. The stronger canonical-model embedding statement is supplied by Bombieri's theorem.

\textbf{The separate \(p_g=0\) conjecture.}
At the end of p.~324 Enriques explicitly marks extension to \(p_g=0\) as probable and leaves it unproved. The tetracanonical assertion in that case is true: if \(p_g=q=0\), then \(|4K_X|\) is birational. For \(K^2\geq2\), this follows from Bombieri's Main Theorem. For the remaining numerical Godeaux case \(K^2=1\), Catanese (1981, Section III, Theorem~1, pp.~254--255) proves injectivity on a nonempty Zariski open subset, hence birationality in characteristic zero. This resolves the conjectural statement positively through a separate argument. It does not validate the preceding false assertion for \(p_g>0\).

\begin{remark}[The retained examples and the limits of the reconstruction]
\label{cp2-viii-limits}
The three worked examples above isolate different geometric phenomena. The Fermat quotient separates irreducibility of a general member from freedom of its linear system. The bidouble cover separates the degree of the complete bicanonical map from a proposed classification based only on restrictions to one canonical curve. The weighted hypersurface separates the self-intersection of a divisor from the moving square after base points are resolved, and shows exactly which graded generator first makes the pluricanonical map birational.

A way to test the older half-curve dimension arguments is to identify the precise lattice-polarized K3 locus on which the half-class is integral, nef, and effective, then determine the image, fibers, and boundary of the incidence map to surface moduli. This specifies the additional lattice and deformation problems required by those particular classical arguments, without asserting that they are open problems in current research.

The intersection, cohomology, and graded-ring calculations establish the examples and numerical comparisons just given. The projected sextic and septic constructions still require local analytic and conductor calculations; the Franchetta incidence constructions require dominance and compatible deformations. An abbreviated construction must be assessed through those specific requirements, independently of the truth of the modern classification theorem.
\end{remark}


\chapter[Superficie irregolari e sistemi continui]{Superficie irregolari e sistemi continui di curve disequivalenti}\label{cp2-ix}

% Original text: pp. 325--369. All source sections are retained in order.



\section{Introduzione}\label{cp2-ix-1}

The modern setting is a smooth, connected, projective surface \(X\) over \(\mathbb C\). A singular projective model in Enriques's examples is replaced by a smooth model, and assigned multiple points are incorporated by blowing up when intersection numbers or characteristic systems require this. Neither the analytic arguments with periods nor the equality between the numerical irregularity and the dimension of the Picard variety are statements in arbitrary characteristic. Page references follow the original printed pagination.

The chapter's principal valid conclusion is
\[
q=h^1(X,\mathcal O_X)=p_g-p_a=\dim\operatorname{Pic}^0(X).
\]
There is a complete algebraic family of sufficiently positive effective divisors whose linear-equivalence classes fill a \(q\)-dimensional Picard component. This theorem does \textbf{not} assert that every effective divisor belongs to such a dominant family with complete characteristic system. The chapter itself eventually exposes precisely that stronger assertion's failure, especially in \S\S~7--8.

Three distinct smoothness/completeness questions must therefore be kept separate.

\begin{enumerate}
\item The Picard scheme parametrizes line bundles. Its tangent space at the identity is \(H^1(\mathcal O_X)\); in characteristic zero its identity component is smooth.
\item The divisor scheme parametrizes effective Cartier divisors together with their sections up to scalar. At \(C\), its tangent space is \(H^0(C,\mathcal O_C(C))\), and it can be singular or nonreduced even when the Picard variety is smooth.
\item A chosen classical family has a characteristic linear series given by the image of its tangent space in \(H^0(\mathcal O_C(C))\). Completeness requires surjectivity of this characteristic map. A tangent vector is not automatically the derivative of an actual curve in that chosen reduced family.

\end{enumerate}
This is the modern formulation of the gap recognized on pp.~331--333. The revisers' footnote on p.~333 expressly says that the following argument was left incomplete and contains many lacunae. The passage therefore presents an incomplete argument for the algebraic theorem. Kleiman's historical and scheme-theoretic accounts distinguish the original theorem, overextended completeness assertions, and later rigorous approaches (Kleiman, 2005, \S~5, especially Remark 5.18; Kleiman, 2014, author's manuscript dated 2013).

\subsection{Notation dictionary and exact sequence}\label{cp2-ix-1-1}

For a smooth effective irreducible curve \(C\subset X\), put \(L=\mathcal O_X(C)\), \(N=\mathcal O_C(C)\), \(n=C^2\), and \(\pi=p_a(C)\). Enriques's index of speciality is \(i=h^2(X,L)=h^0(X,K_X-C)\), his superabundance is \(\omega=h^1(X,L)\), and the complete linear-system dimension is \(r=h^0(X,L)-1\). Surface Riemann--Roch and adjunction give
\[
\chi(L)=\chi(\mathcal O_X)+\frac{C(C-K_X)}2
       =p_a+n-\pi+2,\qquad
r=p_a+n-\pi+1-i+\omega.
\]
Here \(p_a=\chi(\mathcal O_X)-1\), and is allowed to be negative. It is not the genus of a curve.

The sequence
\begin{equation*}\label{cp2-ix-eq1}
0\longrightarrow\mathcal O_X\longrightarrow L\longrightarrow N\longrightarrow0
\tag{IX.1}
\end{equation*}
gives the exact tangent sequence
\begin{equation*}\label{cp2-ix-eq2}
0\longrightarrow H^0(L)/\mathbb C
\longrightarrow H^0(N)
\overset{\partial}{\longrightarrow}H^1(\mathcal O_X)
\overset{s}{\longrightarrow}H^1(L).
\tag{IX.2}
\end{equation*}
Thus the deficiency of the characteristic series cut by the complete linear system is
\begin{equation*}\label{cp2-ix-eq3}
\delta=h^0(N)-r
       =\dim\ker s
       =q-\operatorname{rank}s,\qquad 0\leq\delta\leq q.
\tag{IX.3}
\end{equation*}
The projective characteristic series cut by \(\vert C\vert \) has dimension \(r-1\), whereas the complete normal series has dimension \(h^0(N)-1\). The subtraction of one is essential throughout the dimension counts. For a family \(T\) of dimension \(R\), an injective characteristic map at a general point gives a series of dimension \(R-1\). In characteristic zero one may take a reduced, generically immersive parameter space; inseparability makes this automatic identification unavailable in characteristic \(p\).

\subsection{The numerical bound and the reduction to positive systems}\label{cp2-ix-1-2}

On a regular surface \(q=0\), (IX.2) says that the characteristic series of a complete linear system is complete. More generally, a family containing complete linear systems of general dimension \(r\), with \(d\) parameters of distinct linear-equivalence classes, has dimension \(r+d\). Its characteristic series has dimension \(r+d-1\), provided the tangent map has the expected rank. Castelnuovo's deficiency bound then gives \(d\leq q\). The argument is sound for the stated general smooth members and separable families; it is not a proof that an arbitrary divisor scheme is smooth.

The passage from reducible curves to irreducible residual curves uses a sufficiently positive auxiliary system. Modernly, if the original divisors form a bounded algebraic family, one chooses a high tensor power of an ample line bundle uniformly over its parameter space. Relative Serre vanishing and relative generation provide residual systems, and Bertini gives smooth irreducible general residual members. This is the missing explicit quantifier behind ``sufficiently high order.'' It is a useful reduction, not a claim that all individual residual members are irreducible.

Enriques distinguishes the numerical quantity \(p_g-p_a\) from a geometric maximum of the number of parameters of inequivalent complete systems. Over \(\mathbb C\), the latter is \(\dim\operatorname{Pic}^0\). The fundamental theorem identifies the two. The birational invariance used in changing projective models is correct for smooth surfaces: \(h^1(\mathcal O)\), \(h^2(\mathcal O)\), and hence \(\chi(\mathcal O)\), are unchanged by point blowups.

\subsection{Characteristic and hypothesis boundary}\label{cp2-ix-1-3}

The identities from Riemann--Roch, adjunction, and the cohomology of (IX.1) remain valid over an algebraically closed field for a smooth projective surface. The identification
\[
\dim\operatorname{Pic}^0_{\mathrm{red}}=h^1(\mathcal O_X)
\]
does not hold in arbitrary positive characteristic: the Picard scheme may be nonreduced. Its tangent dimension is still \(h^1(\mathcal O_X)\), while its reduced dimension measures actual class parameters. Characteristic-zero smoothness, Hodge symmetry, \(b_1=2q\), and integration of complex periods must not be transported unchanged to characteristic \(p\). Kleiman (2005, Corollaries 5.13--5.14 and Remark 5.15) explicitly separates these facts and examples.

Even in characteristic zero, smoothness of the Picard variety does not imply smoothness of the divisor scheme, and semiregularity of one divisor does not establish completeness for every family on the surface. Nor does generic vanishing apply to Albanese dimension one without additional hypotheses. These are distinct limitations on the three theorems.



\section{Condizione aritmetica perchè una curva, sopra una superficie d'irregolarità geometrica \texorpdfstring{\(q\)}{q}, appartenga ad una serie continua \texorpdfstring{\(\infty^q\)}{infinity to q} di curve disequivalenti}\label{cp2-ix-2}

The displayed inequality is
\begin{equation*}\label{cp2-ix-eq4}
p_a+n-\pi+1-i\geq0,
\quad\text{equivalently}\quad
\chi(L)-h^2(L)\geq1.
\tag{IX.4}
\end{equation*}
It has a clean proof \textbf{at the level of effective line-bundle classes}. For \(\alpha\in\operatorname{Pic}^0(X)\), the Euler characteristic \(\chi(L\otimes\alpha)\) is constant. Upper semicontinuity gives \(h^2(L\otimes\alpha)\leq h^2(L)\) in a neighborhood of the origin. Consequently
\[
h^0(L\otimes\alpha)
=\chi(L)+h^1(L\otimes\alpha)-h^2(L\otimes\alpha)
\geq\chi(L)-h^2(L)\geq1.
\]
Thus a neighborhood of the original class consists of effective classes. The image of the relevant proper divisor parameter space in the Picard component is closed; an image containing an open subset of the irreducible Picard variety contains the whole component. Equivalently, effective twists exist in every class after taking the closure of the dominant family.

There is no defect in interpreting the theorem as a series of complete linear systems or effective classes. A separate question is whether a prescribed section \(s\in H^0(L)\), and hence its particular divisor, extends along such a family. If \(h^0\) jumps at the original class, many sections need not extend. Section \S~8 recognizes this explicitly for canonical divisors. The stronger assertion about extending every prescribed section is not part of the theorem in \S~2.



\section{La varietà di Picard corrispondente ad una superficie irregolare}\label{cp2-ix-3}

Choose an ample \(H\) sufficiently positive that
\[
H^j(X,H\otimes\alpha)=0\quad(j>0,\ \alpha\in\operatorname{Pic}^0(X)).
\]
Uniformity follows by relative Serre vanishing for the Poincaré line bundle on \(X\times\operatorname{Pic}^0(X)\). Its direct image \(E\) after tensoring with \(H\) is locally free, of rank \(\chi(H)\). The projective bundle \(\mathbb P(E)\) parametrizes all effective divisors in these classes and has dimension \(\chi(H)-1+q\). Its characteristic map at a smooth divisor is the full tangent space \(H^0(N)\). This constructs the desired complete continuous system and proves the numerical/geometric equality once the Picard scheme and characteristic-zero smoothness have been established. It does not use the unproved all-order construction of \S~5.

On p.~328, the quotient by linear equivalence is a Picard component \(\operatorname{Pic}^L\), a torsor under \(\operatorname{Pic}^0\). Choosing \(L\) as origin identifies the torsor with the group. Taking the residual with respect to a fixed large divisor corresponds to \(M\mapsto H\otimes M^{-1}\), hence to inversion followed by translation. The operation is algebraic on the Picard scheme. The original discussion anticipates this structure, but existence, representability, and properness are substantive theorems, not consequences of set-theoretic quotienting alone.

Nor is the image of every continuous curve family an abelian variety. For example, on \(X=B\times\mathbb P^1\) with \(g(B)>1\), the fiber family is parametrized by \(B\), whose image in \(\operatorname{Pic}^1(B)\) is a curve rather than a group. Sums of independently varying members may generate a larger abelian subvariety. Merely multiplying every class by a fixed integer does not increase image dimension. This distinction clarifies the language about multiples on p.~328.



\section{Proprietà caratteristica delle superficie algebriche irregolari: teorema fondamentale per \texorpdfstring{\(p_g=0\)}{pg = 0}}\label{cp2-ix-4}

For \(p_g=0\), write \(p_a=-p\). Take regular nonspecial systems \(\vert C\vert \), \(\vert D\vert \), with invariants \((n,\pi,r)\), \((\nu,\rho,r')\), and transverse intersection number \(m\). Then
\[
r=-p+n-\pi+1,\qquad r'=-p+\nu-\rho+1.
\]
For the sum \(E=C+D\),
\[
N=n+\nu+2m,\qquad
\Pi=\pi+\rho+m-1,
\]
and hence
\begin{equation*}\label{cp2-ix-eq5}
R=\dim\vert E\vert =-p+N-\Pi+1=r+r'+p+m.
\tag{IX.5}
\end{equation*}
Near an ordinary transverse union, requiring \(m\) unspecified nodes imposes at most \(m\) conditions. The connectedness argument rules out an irreducible \(m\)-nodal curve specializing, with exactly those \(m\) persistent nodes, to the disjoint normalization \(C\sqcup D\). A modern proof simultaneously normalizes the persistent nodes after appropriate base change: proper connected fibers cannot acquire a disconnected special normalization in such an equinodal family. Thus the relevant equinodal component stays split. Its dimension is at least \(R-m=r+r'+p\), giving at least \(p\) parameters of linear-equivalence classes after allowing the linear motions of the two summands. Combined with the upper bound \(d\leq p_g-p_a=p\), this gives equality.

The argument needs the transverse nodal setup, an actual flat family, and a normalization/connectedness theorem. It is a geometric sketch with a valid modern completion; it is not the unrestricted assertion that imposing nodes always gives independent conditions. On p.~330 the printed inequality is \(q\geq-p\), although \(p=-p_a>0\). The argument proves \(q\geq p\). This is an original sign misprint, not a questionable transcription.

There is an especially short independent modern proof in this case. If \(C\) is sufficiently ample, \(H^1(L)=H^2(L)=0\), and (IX.1) gives \(H^1(N)=0\) because \(H^2(\mathcal O_X)=0\). The divisor scheme is therefore smooth at \(C\); (IX.2) gives tangent dimension \(r+q\). In addition, \(H^2(\mathcal O_X)=0\) kills the obstruction space for the Picard functor, so the Picard scheme is smooth even in positive characteristic. This special argument is stronger in characteristic than the general complex-period proof (Kleiman, 2005, Proposition 5.19).

The last paragraph's lower bound for arbitrary \(p_g\) with \(p_a<0\) is only \(q_{\mathrm{geo}}\geq-p_a\). It is not the full equality \(q_{\mathrm{geo}}=p_g-p_a\).



\section{Teorema fondamentale per \texorpdfstring{\(p_g>0\)}{pg > 0}}\label{cp2-ix-5}

\subsection{The first-order calculation and its precise limitation}\label{cp2-ix-5-1}

For general \(p_g\), the sum-system count becomes
\[
R=r+r'+m-p_a.
\]
To produce \(q=p_g-p_a\) additional class parameters by splitting, the desired codimension is \(m-p_g\), rather than \(m\). The curves through the intersection group \(G=C_0\cap D_0\) supply infinitesimal splitting directions. In modern terms this is information in a tangent or normal space. Enriques correctly stops on p.~331 to distinguish differential existence from integral existence.

A local calculation shows why. At a node \(xy=0\), a first-order deformation
\[
xy+t(ax+by)
\]
is first-order equivalent to a split product. But
\[
(x+tb)(y+ta)=xy+t(ax+by)+t^2ab,
\]
so the original deformation has a second-order smoothing term unless an additional correction is made. First-order factorization does not imply persistence of a node to all orders. Likewise, the scheme \(y^2=0\) has a larger tangent space than the dimension of its reduced support. Tangent excess is not evidence of a missing reduced family.

The desired tower \(C_1,C_2,\ldots\) on p.~332 must consist of \textbf{compatible} lifts through all infinitesimal thickenings. Constructing a second-order neighbor for some choice at each stage does not by itself establish compatibility, unobstructedness, convergence, or algebraization. A convergent analytic Lie flow does integrate a tangent vector on an already constructed smooth abelian variety. It cannot establish the smoothness and existence of that variety by assuming the group on which it is to flow.

\subsection{The branched-cover construction, pp.~333--335}\label{cp2-ix-5-2}

Projecting a suitably embedded curve to a line gives the familiar branch count
\[
b=2N+2\Pi-2.
\]
For a nodal union \(C+D\), the counts from \(C\) and \(D\), together with two branch points for each of the \(m\) nodes after smoothing, add correctly:
\[
(2n+2\pi-2)+(2\nu+2\rho-2)+2m
=2(n+\nu)+2(\pi+\rho+m-1)-2.
\]
An ordinary branch point has transposition monodromy; the pair coalescing at an apparent node has product the identity. The existence theorem for branched covers allows one to separate such branch points and construct nearby abstract covers. With a sufficiently positive embedding and the necessary deformation theory, this also supplies an ambient spatial picture.

It does \textbf{not} show that the deformed curve remains on the fixed surface \(X\). Nor does the branch count prove that every first-order section of its normal sheaf on \(X\) integrates. The distinction between ambient deformations and deformations constrained to \(X\) is the principal missing bridge. The claimed dimensions \(R\) and \(R-m\) require unobstructed ambient deformations and a surjective map to the chosen node-smoothing parameters, hypotheses that the Riemann-existence discussion alone does not establish. The printed \(r=p+n-\pi+1\) on p.~335 also drops the subscript \(a\), an independent typographical error.

\subsection{A concrete degree error in the auxiliary-curve construction}\label{cp2-ix-5-3}

The difficulty on pp.~335--337 is not only a vague complaint about rigor. Let the auxiliary curve \(L\) have genus \(P\), let \(M=C\cdot L\), and let the restricted complete series \(\vert C\vert _L\) have dimension \(r>0\), as in the text. Choose an imposed divisor \(A\) of degree \(r\) that makes \(C_0\) unique.

If one adds \(S\) of degree \(P-M\) to the \textbf{whole} divisor \(C_0\cap L\), the result \(G\) has degree \(P\). However,
\[
h^0(L,\mathcal O_L(G))
\geq h^0(L,\mathcal O_L(C_0))=r+1>1.
\]
Riemann--Roch on \(L\) says that a nonspecial degree-\(P\) divisor has exactly one section. Therefore this \(G\) cannot be nonspecial as claimed. If one first removes the \(r\) imposed points \(A\), the residual divisor has degree \(M-r\), and a degree-\(P\) completion requires
\[
\deg S=P-M+r,
\]
not \(P-M\). The text alternates between referring to the removed group and counting the full intersection, and later changes the added group's letter. Both notational changes occur in the printed text. A corrected bookkeeping convention can remove this particular contradiction, but does not prove the subsequent transversality or integrability claims.

\subsection{Independence, nodal auxiliary curves, and descent}\label{cp2-ix-5-4}

The assertion that speciality of \(G\) and of its infinitesimally adjacent \(G_1\) imposes two independent conditions requires a transversality statement. Two nearby equations may have the same differential; first-order adjacency is precisely a situation in which their dependence needs analysis. Dimension subtraction alone is insufficient. If the auxiliary curve degenerates to a nodal curve, one must work with its arithmetic genus and dualizing sheaf, or normalize and impose the appropriate gluing condition at the two preimages of the node. These requirements cannot be replaced by counting free points on the normalization.

The pencils \(C_0+D_0+\lambda E\) appearing on p.~337 usually have irreducible general members. Their infinitesimal splitting does not provide an actual factorization of the pencil. Finally, \S~5 acknowledges possible exchange of the two components under monodromy. To obtain a well-defined algebraic locus \(G''\), one must establish descent through that monodromy, independence from the auxiliary \(D\), constant degree, and compatibility with the proposed iteration. These are not supplied.

The displayed addition rule in the p.~338 footnote is valid as \textbf{linear equivalence of divisors}. On a Picard torsor, the proposed parallelogram construction becomes \([G'']=2[G']-[G]\). This clarifies the geometric idea but does not establish the existence of every effective representative invoked in the argument.

\subsection{Rigorous replacement and the role of semiregularity}\label{cp2-ix-5-5}

For \(H^1(L)=H^2(L)=0\), the projective-bundle construction in IX.1 supplies a full replacement of the chapter's fundamental theorem. More generally, the semiregularity homomorphism is
\[
\sigma:H^1(C,N)\longrightarrow H^2(X,\mathcal O_X).
\]
A divisor is semiregular when this map is injective. By (IX.1), this is equivalent to the surjectivity of \(H^1(\mathcal O_X)\to H^1(L)\). The semiregularity theorem gives smoothness of the divisor space at the expected dimension
\[
R=h^0(L)-1+h^1(\mathcal O_X)-h^1(L)
  =r+q-\omega.
\]
It does not say that every divisor is semiregular, nor that semiregularity is necessary for every accidental smooth point. Kodaira and Spencer's 1959 completeness theorem and the later algebraic treatments belong here, with their hypotheses retained (Kodaira \& Spencer, 1959; Kleiman, 2005, Remark 5.18). A different modern integration criterion is the transversality theorem of Mendes Lopes, Pardini, and Pirola (2013, Theorem 1.1), which states explicit cup-product exactness conditions; it is not an unconditional validation of first-order splitting.

One instructive classical counterexample can be expressed without any period calculation. Let \(E\) be an elliptic curve and \(V\) a nonsplit extension \(0\to\mathcal O_E\to V\to\mathcal O_E\to0\). In the ruled surface \(\mathbb P(V)\), the section \(D\) defined by the quotient has normal bundle \(\mathcal O_E\). Hence \(h^0(N_D)=1\), but the section is isolated among the corresponding algebraically equivalent sections: a distinct such section would be disjoint from it and split the extension. The nonzero first-order motion does not integrate. Kleiman (2005, Remark 5.18) identifies this as the Severi--Zappa example and points to Mumford's treatment. It directly refutes any assertion that smoothness of \(\operatorname{Pic}^0\) alone forces all complete characteristic systems.



\section{Storia della teoria dei sistemi continui}\label{cp2-ix-6}

The historical section gives Enriques's retrospective account of the competing priority claims. Its references document a chain of ideas: rationally parametrized systems; the relation between inequivalent curves and irregularity; integrals of the first and second kind; the 1904 characteristic-system announcement; criticisms and attempted repairs; and Poincaré's analytic proof.

A rationally parametrized family of divisor classes is constant because every morphism \(\mathbb P^1\to A\) to an abelian variety is constant. Thus such a family is contained in one complete linear system. This is the precise modern statement behind the opening observation. A family may be a proper subfamily of that linear system.

The trace construction of first-kind integrals from an incidence correspondence is a valid route to holomorphic one-forms. Over \(\mathbb C\),
\[
h^0(X,\Omega_X^1)=h^1(X,\mathcal O_X)=q,
\qquad b_1(X)=2q.
\]
Second-kind integrals modulo exact ones describe the relevant de Rham cohomology; one must not count arbitrary meromorphic primitives without quotienting by exact differentials. The equality of the two Hodge summands and the period-lattice statement are characteristic-zero ingredients.

The genus-two example on p.~341 has two distinct models. A theta divisor on the abelian surface \(J(B)\), \(g(B)=2\), has genus two and self-intersection two. Its translates provide a two-dimensional family with complete normal series \(K_B=g^1_2\), although \(\vert \Theta\vert \) has dimension zero. The symmetric square \(B^{(2)}\), however, is the blowup of \(J(B)\) at the canonical class; its coordinate curves \(B+b\) have self-intersection one. The example must be placed on the abelian model when the characteristic degree is two.

The intrinsic construction \((C+D)\vert _C-D\vert _C=\mathcal O_C(C)\) correctly explains why the complete normal series is independent of the chosen residual divisor \(D\). That independence does not imply its integrability. Pages 342--343 acknowledge that the original completeness argument exceeded what had been demonstrated. The later statement on p.~347 that the construction completes the geometric proof cannot override the revisers' p.~333 warning or the explicit missing steps described above.

The sentence on p.~344 that surfaces with \(p_g=1,p_a=0\) always have a ``pencil of elliptic curves'' needs careful interpretation. Such a surface has \(q=1\), hence has a pencil \textbf{over an elliptic base}, as \S~12 correctly says, with fibers of arbitrary genus. If the sentence means fibers of genus one, it is false without an additional classification hypothesis: general-type surfaces with \(p_g=q=1\) and genus-two Albanese fibers exist. Lewis and Lyons (2018, \S~2 and Theorem 2.1) explicitly describe examples with \(K_X^2=2\) as double covers of \(E^{(2)}\) and identify their Albanese fibers as genus two. The source's two phrasings are not interchangeable. The claim that a surface with trivial canonical bundle and \(p_g=1\) cannot have \(q=1\) is valid and is examined separately on p.~357.

The historical references preserve Enriques's account of his and Castelnuovo's communications. The modern Picard-scheme and completeness theorems cited below establish the mathematical statements with their stated hypotheses; they do not by themselves settle the priority questions raised in this historical notice.

\section{Su diversi tentativi di dimostrazione e d'estensione del teorema fondamentale}\label{cp2-ix-7}

\subsection{The dimension paradox and what it actually proves}\label{cp2-ix-7-1}

Segre's attempted argument retains \(m-p_g\) nodes while smoothing a union \(C+D\). The normalization's characteristic series is assigned dimension
\[
2p_g+(n-\pi)+(\nu-\rho)+1.
\]
The naïve product of the two component characteristic series has dimension one less:
\[
2p_g+(n-\pi)+(\nu-\rho).
\]
That discrepancy is a contradiction only if the \textbf{entire} limiting linear series has been identified. It is not enough to compute the limits of a preferred set of cutting curves that remain proper intersections.

The plane-octic example is decisive. A degree-eight curve with fifteen nodes has normalization genus \(21-15=6\); its canonical series is \(g^5_{10}\), obtained from quintics through the nodes. Degenerate it to two smooth plane quartics meeting in sixteen points, retaining fifteen nodes and leaving one connecting point \(O\). Each quartic has genus three, and the connected nodal normalization has genus six. A quintic through fifteen points of the \((4,4)\) complete intersection passes through the sixteenth as well, by Cayley--Bacharach. Restricting only proper quintic intersections therefore seems to lose a parameter.

The missing limits arise when the cutting quintic acquires one quartic as a component. In a modern limit-linear-series description, the aspects may be represented by \(K_C+2O\) on one component and \(K_D\) on the other, or the reverse, with the requisite gluing and scaling. The residual base point at \(O\) is not present in every limiting aspect. The source's hyperbola example makes the same point in elementary form: for \(xy=t\), a line tending an asymptote can produce residual points on that asymptote, which are absent if one substitutes \(t=0\) before taking the closure of the intersections.

The edition's apparatus proposes changing ``curves \(C_8\)'' to \(C_5\) in the p.~349 sentence about curves through fifteen points not being forced through \(O\). That emendation is unsupported and would make the sentence false: quintics through fifteen of the sixteen points are forced through \(O\), whereas degree-eight forms with those simple vanishing conditions are not. The printed \(C_8\) naturally concerns degree-eight deformations of the quartic union, distinct from the quintic canonical cutters. The printed \(C_8\) should be retained. For example, restriction of degree-eight forms to \(C_4\), after the fifteen imposed points, has degree \(32-15=17\); it is not forced to vanish at the sixteenth point. This distinction is independent of the need to include all limiting canonical divisors.

On p.~350, the disputed differential lemma is asked to force all \(p_g\) connecting points to remain fixed in every limit, including limits with a common component. The passage does not prove that property. Enriques explicitly withdraws his own proposed solution using irreducible cutters. His criticism is a substantive warning about nonproper intersection and specialization, not a proof of a substitute lemma.

\subsection{A verifiable modern splitting statement}\label{cp2-ix-7-2}

Felisetti and Fontanari (2022, Theorem 1.1) give a precise version. The family is flat over a normal base, all curves are nodal, the general member is irreducible, and the special member is \(C+D\) with irreducible components. Write \(\Gamma=C\cap D=\Gamma_1\sqcup\Gamma_2\), with \(\Gamma_1\) the persistent nodes. Assume the complete system \(\vert \mathcal O_D(D)\vert \) is nonempty and the stated generality of \(C\) relative to \(D\), in particular that \(\vert C(-\Gamma_1)\vert \) has no base point on \(D\). Then either \(\vert \Gamma_2\vert \geq p_g+1\), or \(\Gamma_2\) imposes dependent conditions on \(H^0(K_X)\).

The conclusion concerns the actual connecting nodes and canonical dependence. It does not identify every classical limiting characteristic series and does not validate every hypothesis-free version historically called the splitting principle. The article explicitly notes the missing nodality and generality assumptions in the original formulation.

\subsection{Cohomological meaning of the three cases, pp.~351--352}\label{cp2-ix-7-3}

Set \(\epsilon=\operatorname{rank}[H^1(\mathcal O_X)\to H^1(L)]\). Then
\[
\delta=q-\epsilon,\qquad 0\leq\epsilon\leq\min(q,\omega).
\]
Under the \textbf{express hypothesis} that the chosen complete continuous family has complete characteristic system, its dimension is \(r+q-\epsilon\). The cases are:

\begin{enumerate}
\item \(i>0,\omega=\epsilon=0\): the family dimension is \(r+q\), with \(q\) parameters of complete linear systems. Generic constancy of the fiber dimension gives the same speciality index.
\item \(i=0,\omega>0,\epsilon=0\): the family dimension is again \(r+q\). The image in the Picard variety has dimension at most \(q\); semicontinuity gives generic fiber dimension at most \(r\). Equality of total dimensions forces both bounds to be equalities, so the generic superabundance remains \(\omega\).
\item \(i=0,\omega=\epsilon>0\): the family dimension is \(r+q-\epsilon\). Its image dimension can lie between \(q-\epsilon\) and \(q\), while generic linear-system dimension lies between \(r-\epsilon\) and \(r\). The extreme alternatives described by Enriques are compatible with this formula.

\end{enumerate}
The family-completeness premise on p.~351 is indispensable. This passage does not assert it for all families; it explicitly retreats from the earlier universal claim.

\subsection{The ruled quintic and the two components}\label{cp2-ix-7-4}

On the normalization of the ruled quintic \(F_5\), the base has genus two, so \(p_g=0,p_a=-2,q=2\). The hyperplane divisor \(C\) has \(C^2=5\), genus two, \(K_XC=-3\), and \(\chi(\mathcal O_X(C))=3\). The displayed \(r=3\) gives \(h^0=4\) and \(\omega=1\). Its normal degree is five, so \(h^1(N)=0\) and \(h^0(N)=4\); thus \(\delta=1\), exactly as stated. The divisor family has dimension four, although its general members lie in two-dimensional complete linear systems varying over a two-dimensional Picard component. A three-dimensional special linear system is a jumping fiber, not evidence for a one-dimensional class parameter.

Adding two ruling fibers produces the degree-nine model \(F_9\) in \(\mathbb P^6\). The special pair of incident generators produces the node used in the projection back to \(F_5\). A second four-dimensional family arises from hyperplanes through this node and one ruling generator: it has one parameter of classes and three-dimensional linear systems. The pair of preimages \(A,B\) of the node imposes only one condition on the special system and two on a general twist. This is the mechanism of ``neutral'' conditions and the meeting of distinct divisor components.

There is a visible omission in the p.~353 residual formula. The prose says to subtract a ruling generator after imposing \(A,B\), but the displayed expression \(\vert C+\bar a+\bar b-A-B\vert \) omits that generator. On the blowup at \(A,B\), its displayed square is \(9-2=7\); subtraction of a ruling fiber is needed to get square five and the stated \(g^3_5\). Thus the displayed residual system lacks the ruling generator required by the stated dimension and degree.

\subsection{The cone examples and the precise meaning of exorbitance}\label{cp2-ix-7-5}

Let \(B\subset\mathbb P^2\) be a smooth plane curve of degree \(n\geq3\), genus \((n-1)(n-2)/2\), and \(L=\mathcal O_B(1)\). On the standard ruled resolution of its cone, let \(E\) be the negative section with \(E^2=-n\). The hyperplane class is \(H=E+f^*L\), with \(H^2=n\) and \(h^0(H)=h^0(L)+h^0(\mathcal O_B)=4\). For a nontrivial degree-zero twist \(\alpha\), the section contributed by \(H^0(\alpha)\) disappears; effective divisors are the vertex section plus the fibers over a divisor in \(\vert L\otimes\alpha\vert \).

For \(n=3\), the generic fiber linear system has dimension two and the class locus dimension one. For \(n=4\), these dimensions are one and three. For \(n\geq5\), the effective class locus \(W_n(B)\) has dimension \(n\) and its general system is zero-dimensional. These are exactly the counts on p.~353. ``Contained'' is correct for the class of \(H\) in the closure of effective classes; it does not say that every hyperplane section belongs to the closure of the generic effective divisors. This is the geometric content of an exorbitant linear system.



\section{Il sistema paracanonico}\label{cp2-ix-8}

\subsection{The general formula and the invalid extension of a section}\label{cp2-ix-8-1}

The paracanonical component is \(\operatorname{Pic}^{K_X}(X)\). For \(\alpha\ne\mathcal O_X\) in \(\operatorname{Pic}^0(X)\),
\[
h^2(K_X\otimes\alpha)=h^0(\alpha^{-1})=0.
\]
Indeed, an effective nonzero divisor with numerically trivial class would have positive intersection with an ample divisor, a contradiction. Hence
\begin{equation*}\label{cp2-ix-eq6}
h^0(K_X\otimes\alpha)=\chi(\mathcal O_X)+h^1(K_X\otimes\alpha),\qquad
\dim\vert K_X\otimes\alpha\vert =p_a+h^1(K_X\otimes\alpha).
\tag{IX.6}
\end{equation*}
If \(p_a\geq0\), every nontrivial twist is effective. At the canonical class itself, \(h^2(K_X)=1\), so one must not apply the \S~2 inequality with \(i=0\) there. Formula (IX.6), rather than that substitution, establishes the p.~354 effectivity assertion.

For the sufficiently positive auxiliary \(C\), the adjoint system has dimension \(p_a+\pi-1\), and its image on \(C\) is a canonical subseries of dimension \(\pi-1-q\). The group \(G=C\cap K_X\) has degree \(2\pi-2-n\). The flawed old proof assumes that every canonical divisor can be reached as a limit from any nearby paracanonical linear system, so that arbitrary normal directions are represented by nearby residual canonical curves. Pages 354--355 correctly reject this: a smaller nearby linear system can specialize to a proper subsystem of \(\vert K_X\vert \). The predicted universal dimension \(2p_g-p_a-1\) is therefore not established.

\subsection{What generic vanishing proves, and what the source conjectures}\label{cp2-ix-8-2}

The source calls the alternative dimension \(p_g\) probable under \(q>1,p_a\geq0\), and smooth irreducible canonical curves. It is not stated as a proved theorem. Under the additional hypothesis of \textbf{maximal Albanese dimension}, generic vanishing gives
\[
h^i(K_X\otimes\alpha)=0\quad(i>0)
\]
for general \(\alpha\). The main paracanonical component then has dimension
\[
\chi(\mathcal O_X)-1+q=p_g.
\]
This is an appropriate modern theorem explaining the symmetric-square examples, with the Albanese hypothesis retained (Green \& Lazarsfeld, 1987; Mendes Lopes et al., 2013). Generic vanishing is not a theorem for all irregular surfaces merely because their canonical system has a smooth irreducible member.

\textbf{A counterexample to the conjecture as literally phrased.} Let \(B\) have genus \(b>1\), put \(Y=B\times\mathbb P^1\), and choose
\[
L=A\boxtimes\mathcal O_{\mathbb P^1}(d),\qquad
a=\deg A\gg0,\quad d\geq3.
\]
Let \(\pi:X\to Y\) be the smooth double cover branched over a general member of \(\vert 2L\vert \). The standard double-cover algebra is
\[
\pi_*\mathcal O_X=\mathcal O_Y\oplus L^{-1},\qquad
K_X=\pi^*(K_Y+L).
\]
Künneth gives \(H^1(Y,L^{-1})=0\), hence \(q(X)=b\). Also
\[
p_g(X)=h^0(Y,K_Y+L)=(a+b-1)(d-1).
\]
The bundle \(K_Y+L=(K_B+A)\boxtimes\mathcal O(d-2)\) is ample and globally generated; a general member of its pullback canonical system is smooth, and connectedness of an ample divisor makes it irreducible. The pullback identifies \(\operatorname{Pic}^0(B)\) with \(\operatorname{Pic}^0(X)\): it is injective, and dimensions are equal. For every \(\alpha\in\operatorname{Pic}^0(B)\),
\[
h^0(X,K_X\otimes f^*\alpha)=(a+b-1)(d-1)=p_g(X).
\]
The dimensions are constant because \(\deg(K_B+A)\) is above the special range. Consequently the paracanonical family has dimension \(p_g+b-1>p_g\). These examples have \(p_a=p_g-b\geq0\) and the source's stated irreducibility properties. The missing maximal-Albanese-dimension hypothesis is therefore substantive, not cosmetic.

\subsection{Symmetric squares: the cases \texorpdfstring{$q=3,4,5$}{q=3,4,5} are computable}\label{cp2-ix-8-3}

Let \(X=B^{(2)}\), where \(B\) is a smooth curve of genus \(g\geq3\). Then
\[
q=g,\quad p_g=\binom g2,\quad
p_a=\frac{g(g-3)}2,\quad
\chi(\mathcal O_X)=\binom{g-1}{2},\quad
K_X^2=(g-1)(4g-9).
\]
For a degree-zero line bundle \(\alpha\) on \(B\), the determinant-tautological description gives
\begin{equation*}\label{cp2-ix-eq7}
H^0(X,K_X\otimes\alpha^{(2)})
\cong\bigwedge\nolimits^2 H^0(B,K_B\otimes\alpha).
\tag{IX.7}
\end{equation*}
For \(\alpha\ne\mathcal O_B\), curve Riemann--Roch gives \(h^0(K_B\otimes\alpha)=g-1\) and \(h^1=0\). Thus the paracanonical linear systems have dimension \(p_a\), and their total main component has dimension \(p_g\). This directly answers the source's questions about completeness of the trace constructions. Sheridan (2022; author manuscript, \S~1 and Example 1.3) supplies the determinant construction and a modern description of these divisor components.

For \(g=3\), the general paracanonical system consists of one divisor, with a three-dimensional family of classes; \(\vert K_X\vert \) has dimension two. The source's \(K_X^2=6\), canonical-curve genus seven, and Clifford obstruction to a nonhyperelliptic \(g^3_6\) are consistent. For \(g=4\), the general system is \(\mathbb P^2\), with four parameters of classes, giving dimension six. The canonical \(\mathbb P^5\) is an extra component; its intersection with the main component is the Plücker quadric \(\operatorname{Gr}(2,4)\). Nondecomposable two-forms explain why a general canonical divisor does not deform out of \(\vert K_X\vert \). For \(g=5\), the general paracanonical system is \(\mathbb P^5\), with five class parameters, giving dimension ten. Pencils in a four-dimensional space of sections form \(\operatorname{Gr}(2,4)\), a quadric that spans \(\mathbb P(\bigwedge^2\mathbb C^4)=\mathbb P^5\); the complete system is larger than just its decomposable trace divisors.

There is a separate slip in the explanation on p.~356: the printed text says that the chord surface belonging to a ``special linear complex'' cannot occur as a paracanonical limit. With the standard classical identification of a linear complex with an alternating form on the four-dimensional canonical section space, a special complex has rank two, hence is decomposable. These are exactly the points of \(\operatorname{Gr}(2,4)\), the limit locus just described. It is the general \textbf{nonspecial}, rank-four complex that lies outside the main component. Thus the asserted exorbitance of the whole canonical system is correct, but the special/nonspecial description of the excluded members is reversed. The correction concerns that explanation, not the dimension six or the existence of the extra canonical component.

One can also see the parity phenomenon directly. A canonical two-form \(s\) induces, by cup product and duality, an alternating map \(H^1(\mathcal O_X)\to H^1(K_X)\). Its kernel determines the normal-series deficiency. An odd-dimensional alternating matrix has nontrivial kernel; an even-dimensional one may be invertible. This explains why deficiency one is plausible in the odd-\(q\) examples but cannot be asserted for every canonical divisor with even \(q\).

\subsection{Elliptic surfaces: two original numerical errors}\label{cp2-ix-8-4}

On p.~356 the printed text states that a canonical curve is composed of \(p_g-1\) elliptic fibers; this confuses projective dimension with degree on the base. A clean subclass already disproves the printed count. Let \(f:X\to B\) be a relatively minimal elliptic surface with section, \(\chi(\mathcal O_X)>0\), and fundamental line bundle \(\mathcal L\) of degree \(\chi(\mathcal O_X)\). Then
\begin{equation*}\label{cp2-ix-eq8}
K_X=f^*(K_B+\mathcal L),\qquad
q=g(B),\qquad
\deg(K_B+\mathcal L)=2q-2+\chi(\mathcal O_X)=p_g+q-1.
\tag{IX.8}
\end{equation*}
For every \(\eta\in\operatorname{Pic}^0(B)\), this base line bundle has degree greater than \(2q-2\), so
\[
h^0(K_X\otimes f^*\eta)=p_g,\qquad
\dim\vert K_X\otimes f^*\eta\vert =p_g-1.
\]
The total dimension is therefore \(p_g-1+q=2p_g-p_a-1\), including when \(p_a=0\). Formula (IX.8) is the canonical-bundle formula, not an inference from the desired dimension (Schütt \& Shioda, 2010, Theorems 6.8 and 6.12).

Page 357 instead puts \(2q-2=2p_g-2\) fibers in the \(p_a=0\) canonical divisor and lowers the generic linear-system dimension to \(p_g-2\). That corresponds to forgetting the positive degree-one fundamental line bundle \(\mathcal L\). For example, take an elliptic base \(B\), \(\deg\mathcal L=1\), and a general Weierstrass equation with coefficients in \(\mathcal L^4,\mathcal L^6\). Its minimal smooth surface has \(p_g=q=1,p_a=0\), canonical divisor one fiber, and a one-dimensional family of unique paracanonical divisors. The printed prediction gives total dimension zero. This counterexample requires no multiple-fiber subtleties. Multiple fibers in a more general elliptic fibration add canonical-bundle terms and must be treated separately; they cannot rescue a statement already false for elliptic surfaces with section.

The final observation on p.~357, that \(K_X\cong\mathcal O_X\) and \(p_g=1\) cannot coexist with \(q=1\), is correct. Choose an ample smooth \(C\). Then \(N_C=K_C\), and for a nontrivial \(\alpha\in\operatorname{Pic}^0(X)\), vanishing gives \(h^0(C+\alpha)=\pi\) if \(q=1\). Restriction is injective because \(h^0(\alpha)=0\). For sufficiently ample \(C\), restriction of \(\alpha\) remains nontrivial, so \(K_C\otimes\alpha\vert _C\) is a noncanonical degree-\(2\pi-2\) line bundle with only \(\pi-1\) sections. This contradiction makes the source's short curve-theoretic proof rigorous.



\section{Digressione sulle varietà abeliane}\label{cp2-ix-9}

A modern abelian variety is a connected projective algebraic group. Analytically it is a complex torus \(\mathbb C^q/\Lambda\), with \(\Lambda\) a lattice of rank \(2q\), admitting a positive integral Riemann form. The polarization/algebraicity condition matters: a general complex torus is not projective. Enriques includes algebraicity and nondegeneracy in his definition, but the phrase ``\(2q\)-periodic functions'' alone would not be a sufficient modern definition.

Translations form a simply transitive group: exactly one translation sends a specified point to another. On a Picard torsor, residual involutions have the form \(x\mapsto a-x\), and their even products are translations. The \textbf{mixed} group containing translations and these inversions is not simply transitive in the literal modern sense. Both a translation and an inversion send any given point to a prescribed target; inversions have fixed points satisfying \(2x=a\). Thus ``simply transitive'' on p.~357 is accurate for the connected translation subgroup, not for the whole mixed group.

The variety of effective degree-\(g\) divisors \(B^{(g)}\) is birational to \(\operatorname{Pic}^g(B)\); it is not generally itself an abelian variety. The Abel map contracts special complete linear systems. Enriques's ``exceptional variety'' belongs to the symmetric-product model, whereas the nonsingular Jacobian has no such exceptional divisor and has a regular translation action. Keeping these models separate resolves much of the apparent tension between ``without exception'' and later contractions.

The reducible-period discussion expresses the existence of an abelian subvariety of a Jacobian and its complementary factor \textbf{up to isogeny}. For a polarized abelian variety \(A\) and subvariety \(B\), the connected kernel of the dual map associated to the polarization supplies a complementary \(B'\) such that \(B\times B'\to A\) is finite and surjective. It need not be a product isomorphism (Milne, 2008, I, Proposition 10.1).

Restriction from \(X\) to a sufficiently ample smooth curve embeds \(\operatorname{Pic}^0(X)\) in its Jacobian. Every complex abelian variety also arises as a Picard variety of a smooth projective surface: for dimension at least two take a smooth complete-intersection surface in its dual and apply weak Lefschetz; in dimension one take the dual elliptic curve times \(\mathbb P^1\). Cutting further to a sufficiently ample curve and dualizing the resulting Jacobian quotient gives an embedding into a Jacobian. This supplies the content of the ``Torelli theorem'' named here; it is not the modern Torelli assertion that a polarized Jacobian determines its curve.

The proof on p.~359 is a geometric precursor of the polarization map
\[
\phi_L:A\longrightarrow\widehat A,\qquad
a\longmapsto t_a^*L\otimes L^{-1}.
\]
For ample \(L\), this is an isogeny. Thus the translation orbit of a sufficiently general divisor maps finitely to its linear-equivalence classes, and the number \(s\) in the text is the order of the finite kernel. A positive-dimensional connected subgroup cannot act nontrivially through projective transformations on a fixed complete linear system: a homomorphism from a proper connected group to the affine group \(\operatorname{PGL}\) is constant. The source's fixed-point argument gestures toward this fact but presupposes algebraicity and regularity of the action. More explicitly, \(\deg\phi_L=(L^q/q!)^2\); for an elliptic curve and a degree-\(d\) line bundle it is \(d^2\).

Pages 360--361 prove that a connected finite unramified cover of an abelian variety is again abelian. In the complex setting, this has a direct complete proof: if \(A=\mathbb C^q/\Lambda\), a connected finite covering corresponds to a finite-index sublattice \(\Lambda'\subset\Lambda\). Its total space is \(\mathbb C^q/\Lambda'\); pullback of a polarization makes it projective. Choosing an origin makes the covering an isogeny. All translations lift, and the deck group is finite abelian.

The source's monodromy argument is sound in outline once this covering-space setup is supplied. Its final commutator calculation is also correct: if \(T\Pi T^{-1}=\Pi\Omega\) and \(\Omega\) lies in a finite deck group, then \(\Omega\) is constant as \(T\) varies in a connected group. At the identity it is \(1\); therefore the lifted translations commute. This does not establish connectedness or a group law from a bare multivalued rational correspondence without the earlier covering assumptions.



\section{Segue: il genere delle varietà abeliane}\label{cp2-ix-10}

Let \(B\) be a smooth curve of genus \(p\), and let \(B^{(q)}\) denote its \(q\)-th symmetric product. The claim concerns \(1\leq q\leq p\), and the notation \(J_q\) in the source here denotes this symmetric product, not an abelian variety of dimension \(q\).

For a \(q\)-dimensional subspace \(W\subset H^0(B,K_B)\), evaluate sections on a divisor \(\xi\) of degree \(q\). The determinant of
\[
W\longrightarrow H^0(\xi,K_B\vert _\xi)
\]
vanishes exactly when \(\xi\) lies in a member of the associated \(g^{q-1}_{2p-2}\). Its divisor is the determinant-tautological divisor \(N_{K_B}\), and
\begin{equation*}\label{cp2-ix-eq9}
N_{K_B}=K_{B^{(q)}},\qquad
H^0(B^{(q)},K_{B^{(q)}})\cong\bigwedge\nolimits^qH^0(B,K_B).
\tag{IX.9}
\end{equation*}
One can verify the canonical identification locally off the diagonals by writing the determinant of the differentials at the \(q\) points; the alternating determinant cancels the quotient ramification along the diagonals and extends across collisions. This provides a precise construction of the canonical divisor described in the text.

For \(q=2\), the coordinate curve \(M_A=A+B\) is isomorphic to \(B\), has normal bundle \(\mathcal O_B(A)\), and self-intersection one. Adjunction gives
\[
K_{B^{(2)}}\vert _{M_A}=K_B-A.
\]
The divisor cut by the chosen canonical pencil consists of the \(2p-3\) points other than \(A\) in its member through \(A\). The missing point \((A,A)\) is the characteristic point of the moving coordinate curves. Thus the source's intersection calculation correctly identifies the adjunction residual.

For the induction, \(M_{q-1}=A+B^{(q-1)}\) has normal divisor \(M_{q-2}\), and
\[
K_{B^{(q)}}\vert _{M_{q-1}}=K_{M_{q-1}}-M_{q-2}.
\]
This is the clean line-bundle statement behind the generalization of the ``third equivalence criterion.'' The source gives a sketch of that generalization, not a proof of the criterion in every dimension. The printed text has two index/letter slips: on p.~362 the variety called the image of \(g^{q-1}_{2p-2}\) should be \(K_{q-1}\), although it prints \(M_{q-1}\); on the first line of p.~363 the restriction is to \(M_{q-1}\), not the printed \(M_{q-2}\).

For \(q=p\), (IX.9) gives \(h^0(K_{B^{(p)}})=1\). Its canonical divisor is the exceptional divisor of the Abel map \(B^{(p)}\to J(B)\). The canonical bundle on the Jacobian itself is trivial. These are compatible statements on different birational models.

The extension to all abelian varieties has an immediate proof independent of the symmetric-product model. Translation identifies every tangent space with \(T_0A\), so
\[
T_A\cong\mathcal O_A^{\oplus q},\qquad
\Omega_A^1\cong\mathcal O_A^{\oplus q},\qquad
K_A\cong\mathcal O_A.
\]
Because \(A\) is connected and projective, \(h^0(\mathcal O_A)=1\), hence its geometric genus \(h^0(K_A)\) is one. An étale isogeny preserves this canonical triviality. This validates the conclusion on p.~363 without treating contractions on a symmetric product as if they occurred on an abelian variety.

The alternative genus-two argument on p.~364 is geometrically informative but abbreviated. A base-point-free degree-four pencil on a genus-two curve has ten simple branch points. Its unordered-pair trace curve is a degree-six cover of \(\mathbb P^1\); a transposition of four letters acts on their six two-element subsets as two transpositions, so it gives two simple ramification points. Riemann--Hurwitz yields
\[
2g(L)-2=-12+20=8,\qquad g(L)=5.
\]
On the minimal abelian surface, its divisor class is a translate of \(2\Theta\), with square eight. The net in the construction has two simple base points, leaving effective intersection degree six. The complete normal line bundle on a smooth \(L\) is \(K_L\). These statements justify the source's genus and degree counts. They should not be read as saying that a two-dimensional linear net by itself cuts the entire four-dimensional canonical series: its tangent image and the complete normal series are distinct. The auxiliary \(g^2_6\) completed in \(A+B+g^2_6\subset g^4_8\) likewise must be distinguished from the pencil cut by the net's other members. The source invokes Zeuthen's correspondence formula rather than writing this ramification calculation.



\section{Fascio irrazionale sopra le superficie di genere geometrico nullo}\label{cp2-ix-11}

\subsection{Surfaces in abelian varieties and the missing hypothesis in a generalization}\label{cp2-ix-11-1}

If a smooth surface maps generically finitely onto a two-dimensional subvariety of an abelian variety, it has a nonzero holomorphic two-form. At a point where the differential has rank two, choose two invariant one-forms whose restrictions are independent. Their wedge pulls back nontrivially. After resolving a singular subvariety, the same argument applies. This proves the p.~365 conclusion \(p_g>0\) directly.

The translation orbit of a surface in a \(q\)-dimensional abelian variety has dimension at least \(q-2\), since its stabilizer has dimension at most two. This dimension estimate is correct. The broader assertion that a family of surfaces of dimension at least \(q-2\) in \textbf{any} \(q\)-fold of positive geometric genus must have members of positive geometric genus requires an additional coverage/characteristic-rank hypothesis.

\textbf{A counterexample to the unrestricted generalization.} Let \(V=S\times T\), where \(S,T\) are projective K3 surfaces and \(S\) contains a smooth rational curve \(R\). Let \(D\) vary in a sufficiently large complete linear system on \(T\). The surfaces \(R\times D\) form a family of dimension at least two, while \(p_g(R\times D)=0\). Yet \(h^0(K_V)=1\). The family fills only \(R\times T\), not \(V\). This violates the general statement as literally phrased.

A sufficient corrected hypothesis is that the evaluation map of a family of smooth surfaces dominates the ambient smooth \(q\)-fold and has normal characteristic rank \(q-2\) at a general point. Contract a nonzero ambient \(q\)-form with \(q-2\) independent characteristic sections. Adjunction gives a nonzero two-form on a general member. Translation families satisfy this condition; an arbitrary large family need not. The source's induction silently uses this rank when it assumes the one-parameter subfamilies sweep genuine threefolds and the successive families sweep the ambient variety.

\subsection{A complete modern proof of the \texorpdfstring{$p_g=0$}{pg=0} pencil theorem}\label{cp2-ix-11-2}

Suppose \(p_g(X)=0\) and \(q>0\). Let \(a:X\to A=\operatorname{Alb}(X)\). Its image cannot be two-dimensional, by the preceding wedge-of-one-forms argument. It cannot be a point, because pullback identifies invariant one-forms on \(A\) with all holomorphic one-forms on \(X\). Thus its image is a curve. Stein factorization gives a morphism
\[
f:X\longrightarrow B
\]
with connected fibers and smooth projective base \(B\).

Pullback of one-forms gives \(g(B)\leq q\). Conversely, the Albanese image generates \(A\), so the induced \(J(B)\to A\) is surjective, giving \(q\leq g(B)\). Therefore
\begin{equation*}\label{cp2-ix-eq10}
g(B)=q=-p_a.
\tag{IX.10}
\end{equation*}
This proves the required irrational pencil. The genus in the theorem is the genus of the base, not the genus of its moving curves. If there were two independent connected pencils over positive-genus curves, the product map would be generically finite onto the product of the bases. Pullback of the tensor products of their one-forms would give nonzero two-forms on \(X\), contrary to \(p_g=0\). Thus the Albanese pencil is unique up to an isomorphism of its base.

The converse sentence on p.~366 prints \(p_a=0\); it must read \(p_g=0\). With \(p_a=0\) and \(p>0\), the stated conclusion \(p_a=-p\) is self-contradictory. Under \(p_g=0\), a connected irrational pencil is the Albanese pencil and has genus \(q\), so the converse is correct.

\subsection{The original incidence proof, pp.~366--368}\label{cp2-ix-11-3}

For a one-parameter family of curves of index \(i>2\), let \(C_i(x)\) be the sum of the \(i\) members through \(x\). The source distinguishes the dimension of the image of \(x\mapsto[C_i(x)]\) in the appropriate Picard/Jacobian variety:

\begin{itemize}
\item Dimension zero: the sums are linearly equivalent. The cited Severi--Castelnuovo correspondence criterion implies that the original family has constant linear class, so it is contained in a linear system.
\item Dimension one: fibers of the class map give an irrational pencil. Restrictions of the original curves to a general fiber have constant class; their differences are vertical, and the earlier equivalence criterion expresses them as differences of fibers.
\item Dimension two: the map is generically finite onto a surface in an abelian variety, forcing \(p_g>0\).

\end{itemize}
For \(p_g=0\), an inequivalent family therefore yields the middle case. The rank-zero and vertical-equivalence steps are applications of the earlier chapter's equivalence criteria, not proved afresh here. In particular, reducible singular fibers require care: their components form a discrete lattice of divisor classes. A connected algebraic family cannot vary its discrete component class, but this does not by itself make two named effective components literally identical. The \(K_1=\bar K_1\) line on p.~367 abbreviates this cancellation in a class calculation.

The displayed equalities in this argument are statements of linear equivalence, even where the book uses an ordinary equals sign. Modern notation \(C+\sum K\sim C_0+\sum K'\) prevents confusion with equality of effective divisors.

The first two paragraphs of p.~368 are obscured by overlapping leaves in the available copy. Their wording in the edition is an editorial restoration; the obscured words remain uncertain. The mathematical content is consistent with (IX.10) and the independent-two-pencil argument. In the exposed historical notice, the transcendental argument is the same wedge obstruction: functionally independent first-kind integrals would produce a nonzero two-form; otherwise their level sets yield the Albanese pencil.

\section{Nota sulle superficie di irregolarità uno}\label{cp2-ix-12}

For every smooth projective complex surface with \(q=1\), the Albanese variety is an elliptic curve and its image is the whole curve. The Albanese morphism has connected fibers. To see this, factor through a smooth curve \(B\). Since \(g(B)\leq q=1\) and \(B\to\operatorname{Alb}(X)\) is nonconstant, \(B\) has genus one. The universal property yields a reverse homomorphism from \(\operatorname{Alb}(X)\) to \(B\) splitting the finite map; this forces that finite map to be an isomorphism. Thus the pencil is over an elliptic base, with fiber genus \(\pi\geq0\) unrestricted.

Moreover \(f^*:\operatorname{Pic}^0(B)\to\operatorname{Pic}^0(X)\) is an isomorphism. If \(C,C_0\) are algebraically equivalent divisors, their difference is
\begin{equation*}\label{cp2-ix-eq11}
\mathcal O_X(C-C_0)=f^*\mathcal O_B(b-b_0),
\qquad C\sim C_0+F_b-F_{b_0}.
\tag{IX.11}
\end{equation*}
Fibers here mean the full scheme-theoretic fibers, including multiplicities. This is the exact content of the final assertion about continuous systems. It does not say that every twist is effective or that all effective representatives belong to one irreducible divisor component. The formula concerns line-bundle classes and is compatible with exorbitance.

Printed p.~370 is blank apart from marks and a handwritten folio. It contributes no mathematical claim.

\chapter{Le superficie di genere geometrico nullo}
\label{cp2-x}

\section{Introduzione}
\label{cp2-x-1}

The programme on pp.~371--372 has two logically different outcomes. The criteria for ruledness, and the description of the seven bielliptic families at the end of the chapter, are valid with the hypotheses and repairs explained below. The intervening assertion that every nonruled surface with $p_g=0$ and $p_a=-1$ has a one-dimensional elliptic group of automorphisms is false. Section~11 misses elliptic fibrations with nontranslation monodromy. A free quotient constructed in Section~\ref{cp2-x-11} has $p_g=0$, $q=1$, $K_X^2=0$, and Kodaira dimension one, but no nonzero vector field. It disproves the claimed exhaustiveness of Sections~11--13; it does not disprove the plurigenus criterion in Section~14. Keeping these conclusions separate is essential to reading the chapter correctly.

As in Part~1, $X$ denotes a smooth connected projective surface over $\mathbb C$, $p_g=h^0(X,K_X)$, $q=h^1(X,\mathcal O_X)$, $p_a=p_g-q$, and $P_m=h^0(X,mK_X)$. We use $\sim$ for linear equivalence and $\equiv$ for numerical equivalence. Enriques's $p^{(1)}$ is $K_{X_{\min}}^2+1$ when it is the absolute linear genus; his relative linear genus belongs to the particular model. His fibre letter $K$ will generally be replaced by $H$ or $F$, so that it cannot be mistaken for $K_X$. A ``fascio ellittico'' means a pencil with an elliptic \emph{parameter curve}; its members can have any genus. Conversely, a ``fascio lineare di curve ellittiche'' has base $\mathbb P^1$ and genus-one general fibre. These two uses of ``elliptic'' govern nearly every distinction below.

The input from Chapter~IX is precise: if $p_g=0$ and $q>0$, the Albanese map has one-dimensional image, its connected-fibre factorization has base genus $q$, and two independent irrational pencils would force $p_g>0$. These are the Albanese and independence-of-pencils results of Chapter~IX. Enriques proposes to use them first to show that $q>1$ forces ruledness, and then to handle $q=1$ by constructing a second pencil. The first conclusion is valid. The second is valid when the Albanese fibres have genus greater than one, but fails in the genus-one case as printed.

\section{Le superficie contenenti un sistema di curve di genere \texorpdfstring{$\pi$}{π} e di grado \texorpdfstring{$n>2\pi-2$}{n>2π−2}}
\label{cp2-x-2}

\subsection{The criterion and the successive adjoints, pp.~372--374}

Let $C$ be a general integral member of a moving system, with $C^2=n$ and arithmetic genus $p_a(C)=\pi$. Adjunction gives
\[
 K_X\cdot C=2\pi-2-n.
\]
Thus the numerical hypothesis is $K_X\cdot C<0$. The intersections of $C$ with its successive adjoint classes are
\[
 (C+jK_X)\cdot C=n+j(2\pi-2-n),\qquad j\geq1,
\]
which yields exactly the printed sequence $2\pi-2$, $4\pi-4-n$, $6\pi-6-2n$, and so on. Since $C$ is nef, an effective class cannot have negative intersection with it. The adjoint sequence must therefore terminate. For the same reason $mK_X$ cannot be effective for any $m>0$, so all plurigenera vanish. This last point supplies the justification for the passage to $p_g=0$ that is compressed in the text.

The regular case is Castelnuovo's rationality criterion, already treated in Part~1, Chapter~VI, ``Il criterio di razionalit\`a''. Enriques's citation on p.~372 points to Chapter~V, \S~4, although the developed criterion belongs to the later discussion as well. In the irregular case one has an Albanese pencil $f:X\to B$ with $g(B)=q$. The text studies the last nonempty adjoint and distinguishes $q>1$ from $q=1$.

For $q>1$, an integral horizontal component $D$ mapping to $B$ with degree $d\geq2$ satisfies Riemann--Hurwitz:
\[
 g(\widetilde D)\geq d(q-1)+1>q.
\]
This is a correct obstruction to the proposed terminal adjoint. However, the preceding assertions that $C\cdot H>1$, and that a reducible terminal adjoint has a component meeting $H$ at least twice, are not proved by merely exhibiting the twists $C+f^*\alpha$, $\alpha\in\operatorname{Pic}^0(B)$. One must control effectivity, dimensions, and the horizontal components throughout the adjunction process. The text gives the classical continuous-system argument in compressed form; it does not give those controls.

For $q=1$, the residual possibility is a family of elliptic curves. Enriques says that a family of positive self-intersection gives rational curves by taking intersections of equivalent pairs, and reduces self-intersection $\nu>1$ to $\nu=1$ by imposing $\nu-1$ points. The intended parameter curve for the pairs is a complete $g^1_2$ on an elliptic curve, hence $\mathbb P^1$. It is not the whole symmetric square of the elliptic parameter curve. Existence and irreducibility after the point conditions, and the assertion that the intersection construction sweeps the surface, require justification. The recollection on p.~376 identifies the intended $g^1_2$ argument but does not fill all of these family-theoretic steps. The proof should consequently be classified as a sketch of a correct theorem.

\begin{theorem}
\label{cp2-x-negative-canonical}
If a smooth projective complex surface contains an integral curve $C$ which is not a $(-1)$-curve and satisfies $K_X\cdot C<0$, then it is birational to a ruled surface. In particular, the moving-system criterion of pp.~372--374 is valid.
\end{theorem}

Here the genus in the equivalent inequality $C^2>2p_a(C)-2$ is arithmetic genus. Replacing it by the geometric genus of a singular curve is not legitimate. The hypothesis excluding an exceptional curve is also essential: blowing up a nonruled surface supplies a $(-1)$-curve with negative canonical intersection.

A short modern proof uses the minimal-model result established in Part~1, Chapter~VI, \S~1, rather than trying to complete every moving-system construction. Suppose that $X$ is nonruled, and contract it to a minimal surface $\mu:X\to Y$. Then $K_Y$ is nef and
\[
 K_X=\mu^*K_Y+E
\]
with $E$ an effective exceptional divisor. Any integral curve outside $\operatorname{Supp}E$ has nonnegative intersection with both terms. A curve in the exceptional locus is smooth rational and has negative self-intersection. By adjunction, its canonical intersection is $-2-C^2$, which is negative only when $C^2=-1$. This proves the assertion. The minimal-model input, including the nefness of $K_Y$ in the nonruled case, is the substantive modern theorem; the argument is not an independent proof of it. See the treatment of minimal models in Part~1 and Beauville (1978).

\subsection{A single curve, exceptional curves, and the historical outline, pp.~374--376}

The extension from a system of dimension at least two to a single curve is explicitly made in the book. For $sC$ the arithmetic quantities are
\[
 p_a(sC)=s\pi+\frac{s(s-1)}2n-s+1,\qquad (sC)^2=s^2n,
\]
and
\[
 (sC)^2-2p_a(sC)+2=s(n-2\pi+2)>0.
\]
The Riemann--Roch estimate printed on p.~374 follows once $h^2(\mathcal O_X(sC))=0$. When $C^2>0$, the curve is nef and big, and this vanishing holds for sufficiently large $s$: an ample intersection excludes the effectivity of $K_X-sC$. The growing spaces of sections then produce moving divisors. Riemann--Roch alone does not prove that every $|sC|$ is irreducible or has no fixed components. The modern proof of Theorem~\ref{cp2-x-negative-canonical} avoids that unnecessary strengthening.

If $\pi=0$, the inequality gives $n\geq-1$. Excluding a $(-1)$-curve leaves $n\geq0$. The borderline case $n=0$ consists of a smooth rational curve of self-intersection zero; its deformation or linear-system argument gives a rational fibration. This is the exact reason that the exceptional-curve exclusion matters. If there are infinitely many exceptional curves, the surface is ruled. Indeed, a nonruled surface has only finitely many $(-1)$-curves, since they belong to the finite exceptional locus over its minimal model. It does not say that every ruled surface has infinitely many of them.

The historical outline on p.~375 is expressly advertised as incomplete. For a terminal adjoint of genus $\pi$, dimension $r$, degree $n$, and nodal-fibre number $\delta$, the displayed relation and Riemann--Roch give
\[
 \delta-n-4\pi=-12q-\dot p^{(1)}+9,
 \qquad r\geq n-\pi+1-q,
\]
whence
\[
 5\pi+r\geq11q+\dot p^{(1)}-8+\delta.
\]
The dot appears in the first printed formula and its footnote, but is absent from the next formula. Unless a passage to the absolute model is justified separately, the same relative invariant must be used throughout this calculation. This is a printed notational slip, not a transcription error.

The subsequent bound $r\geq3\pi-5$ is part of an omitted case analysis, not a freestanding rationality criterion for arbitrary systems. Taken without that context it is false: a genus-two polarization on a K3 surface has $r=2\geq1=3\pi-5$. The book itself warns that the cases $q<3$ need special treatment. Consequently one should retain the argument as a historical sketch with its stated restrictions, and use the proved negative-canonical-intersection criterion for the conclusion.

Enriques attributes the original adjunction criterion to Castelnuovo--Enriques (1901), the simpler pencil-based presentation to Enriques--Campedelli (1934), and the single-curve extension to Campedelli (1933), with the titles given on pp.~374--376. The terminal conclusion, $p^{(1)}\geq1$ for the minimal nonruled model, means precisely $K_{X_{\min}}^2\geq0$.

\section{Lemma sulle curve riducibili di un fascio}
\label{cp2-x-3}

\subsection{The local contribution, pp.~376--378}

Let $f:X\to B$ be a fibration with connected fibres and smooth general fibre of genus $g$. For a possibly nonreduced fibre $F$, define
\[
 \delta(F)=e(F_{\mathrm{red}})-(2-2g).
\]
This is the local contribution to the topological Euler characteristic, hence to the Zeuthen--Segre invariant. On a minimal nonruled surface, Enriques asserts $\delta(F)\geq0$, with equality for a degenerate fibre only when it is a multiple of a smooth elliptic curve. The assertion is correct. The restriction to a minimal nonruled model, or an equivalent relative-minimality hypothesis with the appropriate positivity argument, is doing work here.

For a simple normal crossings fibre $F=\sum m_iD_i$ with smooth components of genus $g_i$, the expression underlying the calculation on p.~377 is
\[
 \delta(F)=\sum_i(m_i-1)(2g_i-2)
 +\sum_{i<j}(m_i+m_j-1)D_i\cdot D_j.
\]
This is the Godeaux calculation recalled from Chapter~V. A node on a component can either be counted separately, using its geometric genus, or absorbed into arithmetic genus with the corresponding local correction. These conventions must not be mixed. In particular, the inequality obtained from adjunction is
\[
 D_i^2\leq2p_a(D_i)-2,
\]
not the corresponding inequality with the genus of the normalization unless singularity corrections are retained.

The book mentions higher intersections and singular components as possible complications and replaces an equality by an inequality. It does not give a complete local computation for arbitrary fibre singularities. A uniform repair is available directly on the smooth surface. Write $D=F_{\mathrm{red}}$. Since $F^2=0$, adjunction gives $2g-2=K_X\cdot F$. For a reduced complex curve with isolated singularities,
\[
 e(D)=2-2p_a(D)+\sum_{x\in\operatorname{Sing}D}\mu_x(D),
\]
where $\mu_x$ is the plane-curve Milnor number. Therefore
\begin{equation}
\label{cp2-x-defect}
 \delta(F)=K_X\cdot(F-D)-D^2+
 \sum_x\mu_x(D).
\end{equation}
Every term on the right is nonnegative: $K_X$ is nef, $F-D$ is effective, $D^2\leq0$ by the intersection form on vertical divisors, and $\mu_x\geq0$. This proves the lemma, including non-normal-crossings fibres.

For the equality case, all Milnor numbers vanish, so $D$ is smooth. Since fibres are connected, $D$ is a single smooth component. Its square is zero. If $F=mD$ with $m>1$, equality also gives $(m-1)K_X\cdot D=0$, and adjunction then gives $g(D)=1$. A reduced smooth fibre, of course, also has contribution zero; the book's exception concerns a \emph{degenerate} fibre. This distinction prevents the erroneous reading that every zero-contribution fibre must have genus one.

The alternative proof attributed to Franchetta uses the negative semidefiniteness of the intersection matrix. If a vertical effective divisor had positive square, Riemann--Roch would make its multiples grow quadratically, whereas divisors supported in fibres come from a linear series on the base and have at most linear growth. One must include the eventual vanishing of $h^2$ when using that argument. There are two errors in the printed dimension estimate on p.~378. If $D^2=i>0$ and $p_a(D)=\rho$, surface Riemann--Roch gives, for sufficiently large $r$,
\[
 \dim|rD|\geq p_a(X)+r^2i-r\rho-\frac{r(r-1)}2i+r
 =p_a(X)+\frac{r^2i}{2}-\frac r2(2\rho-2-i).
\]
Thus the printed $-r$ must be $+r$, and the final $+1$ must be removed from both expressions: it belongs to the bound for $h^0$, not to the projective dimension. Both defects occur in the printed formulas. Neither affects the quadratic leading term needed for the contradiction. The intersection-form proof of \eqref{cp2-x-defect} avoids the growth argument entirely.

The same footnote prints $\bar\delta=(t+1)(2\varrho-2)$ for a multiple smooth component. The correct coefficient is $t-1$, as follows from the preceding general formula or directly from adjunction. At $\varrho=1$ both expressions vanish, which hides the error in the special equality case under discussion.

The historical footnote cites Castelnuovo--Enriques (1901), Godeaux (1920), and Campedelli (1934). The journal name is printed ``Hainart''; the correct bibliographic place-name is Hainaut. The place-name is a misprint in the book. The Milnor-number argument above is a separate modern proof.

\section{Superficie di genere \texorpdfstring{$p_g=0$}{pg=0} con \texorpdfstring{$p_a<-1$}{pa<−1}}
\label{cp2-x-4}

The argument on pp.~378--379 is complete using the pencil theorem and the preceding local lemma. Set $q=-p_a>1$. The Albanese fibration has base genus $q$. If its general fibre is rational, $X$ is ruled over that base. Suppose instead that its genus is $g\geq1$ and that $X$ is minimal nonruled. Noether's formula and $p_g=0$ give
\[
 e(X)=12(1-q)-K_X^2<0.
\]
On the other hand, the fibration formula is
\[
 e(X)=4(q-1)(g-1)+\sum_{b\in B}\delta(F_b)\geq0,
\]
by Section~\ref{cp2-x-3}. This contradiction proves ruledness. The identities are Noether's formula and the Zeuthen--Segre fibration formula from Chapter~V. The ruling has base genus $q$, because irregularity is a birational invariant and equals the base genus for a ruled surface.

In the book's notation the comparison forces the absolute $p^{(1)}$ into an impossible range; the Euler-characteristic formulation shows the sign directly. No assumption of general singular fibres is necessary. Multiplicities and nonreduced fibres have already been accounted for by $\delta(F_b)$. Nor does the proof assert that all surfaces with $p_g=0$ are ruled: the case $q=0$ and the nonruled part of $q=1$ remain.

\section{Superficie con \texorpdfstring{$p_g=0$}{pg=0}, \texorpdfstring{$p_a=-1$}{pa=−1} e un fascio ellittico di curve di genere \texorpdfstring{$\pi>1$}{π>1}: lemma I}
\label{cp2-x-5}

\subsection{What the numerical calculation establishes, pp.~379--381}

Let $f:X\to B$ be the Albanese fibration, where $B$ is elliptic and the general fibre $H$ has genus $\pi>1$. On a minimal nonruled model,
\[
 0=12\chi(\mathcal O_X)=K_X^2+e(X),\qquad
 e(X)=\sum_b\delta(F_b).
\]
Both summands are nonnegative. Hence $K_X^2=e(X)=0$, equivalently $p^{(1)}=1$ and $\Delta=0$ in the book. Section~\ref{cp2-x-3} then shows that \emph{every} fibre is smooth and reduced. Indeed, for a multiple fibre $F=mD$ with smooth reduced support $D$, adjunction gives $2\pi-2=m(2g(D)-2)$ and its local contribution is $(m-1)(2g(D)-2)$. If $m>1$ and $\pi>1$, this contribution is positive. Thus $f$ is a smooth family, not merely a fibration without nodal fibres. This calculation is complete.

Lemma~I asserts the existence of an effective parabicanonical divisor, and more generally an effective para-$m$-canonical divisor for $m\geq2$, whose components are disjoint elliptic curves. In modern notation it asks for
\[
 0\ne D_m\in|mK_X+\alpha_m|,
 \qquad\alpha_m\in\operatorname{Pic}^0(X),
\]
with $D_m\cdot H=m(2\pi-2)$, $D_m^2=0$, and $p_a(D_m)=1$. Here ``paracanonical'' has the Picard-variety meaning of Chapter~IX: the twisting class belongs to $\operatorname{Pic}^0(X)$, so linear equivalence to $mK_X$ is not asserted.

For $A_m=H+mK_X$, direct calculation gives
\[
 A_m^2=4m(\pi-1),\qquad
 p_a(A_m)=1+(2m+1)(\pi-1),
\]
and Riemann--Roch gives
\[
 h^0(X,A_m)\geq(2m-1)(\pi-1),\qquad
 \dim|A_m|\geq(2m-1)\pi-2m.
\]
The vanishing $h^2(A_m)=0$ follows by intersecting $K_X-A_m$ with the nef fibre $H$. For $m=2$ this is $A_2^2=8\pi-8$, $p_a(A_2)=5\pi-4$, and $\dim|A_2|\geq3\pi-4$. The genus $3\pi-2$ printed on p.~380 belongs to the \emph{first} adjoint $H+K_X$, not to $A_2$; the text uses it to calculate the dimension of the next adjoint and is consistent on that point.

The restriction of $A_m$ to $H$ is $mK_H$, whose complete series has dimension $(2m-1)(\pi-1)-1$ for $m\geq2$. If the restriction map has a kernel, an effective divisor in $|mK_X|$ follows. Otherwise Enriques varies $A_m$ through its twists by $\operatorname{Pic}^0(B)$ and claims a limiting member containing $H$. This needs a proper incidence family and an argument that its positive-dimensional fibre actually meets the boundary in the required way. A dimension count alone does not imply that a member acquires $H$ as a component. Furthermore, the number on p.~381 is the lower bound for each projective linear system; the continuous family has one additional parameter when its distinct twists are all effective. Confusing those two dimensions obscures the intended argument.

The effective divisor, once obtained, has the stated component structure. Every rational curve maps constantly to $B$, and a smooth fibre of genus greater than one has no rational component. Thus $X$ has no rational curves. Write $D_m=\sum a_iC_i$. Since $K_X$ is nef and $K_X\cdot D_m=0$, one has $K_X\cdot C_i=0$. The Hodge index theorem gives $C_i^2\leq0$, whereas adjunction and the absence of rational curves give $C_i^2=2p_a(C_i)-2\geq0$. Hence each $C_i$ is a smooth elliptic curve of square zero, and the equality $D_m^2=0$ makes distinct components disjoint. Arithmetic genus one alone would not justify this conclusion without these additional arguments.

\subsection{A modern existence proof and its scope}

The smooth family $f$ is isotrivial; a complete proof is given in Section~\ref{cp2-x-10}. After a finite \emph{étale} cover $E'\to B$ it becomes $H\times E'$, with a finite abelian group $G$ acting on $H$ and by translations on $E'$. Decompose $H^0(H,mK_H)$ into characters of $G$. A nonzero eigensection becomes invariant after twisting by the corresponding torsion line bundle on $B$, and descends to a nonzero section of $mK_X+f^*\alpha_m$. Its divisor upstairs is the product of a zero-divisor on $H$ with $E'$, so downstairs it has exactly the disjoint elliptic components required by Lemma~I. This proves the lemma for every $m\geq2$, with no assumption that the untwisted $P_m$ is positive.

The warning at the end of p.~381 is therefore justified: different such divisors can have the same reduced support. Lemma~I by itself does not provide the two different supports demanded in Lemma~II. The latter requires further work.

\section{Lemma II: disegno della dimostrazione}
\label{cp2-x-6}

The exact statement on pp.~381--384 is that there are two elliptic divisors with different reduced supports, suitable multiples of which are para-pluricanonical and restrict to equivalent divisors on every $H$. Enriques tries to obtain the second divisor from a nontrivial torsion twist on the fibres. The distinction between a torsion \emph{class}, a line bundle representing it over the function field, and an effective divisor in that line bundle is correctly recognized in his two simplifying hypotheses.

For even $\pi$ he chooses $0\ne\eta\in\operatorname{Pic}^0(H)[2]$ and sets $L_H=K_H+\eta$. Then
\[
 2L_H\sim2K_H,\qquad
 h^0(H,L_H)=\pi-1,
\]
whereas $h^0(H,K_H)=\pi$. The equality for $L_H$ follows from Riemann--Roch and $h^0(H,-\eta)=0$. A divisor restricting to $L_H$ has an adjoint restricting to $2K_H+\eta$, which is half of $4K_H$ without being bicanonical. For odd $\pi$ the same construction uses a prime $s\mid\pi$ and a nonzero $s$-torsion class. The point of this choice is
\[
 \gcd(s,\pi-1)=1.
\]

Hypothesis~I asks for a rational choice of the torsion class along the family. The finite étale group scheme $\operatorname{Pic}^0(X/B)[s]$ may have monodromy, so it need not have a nonzero global section. Passing to a connected component of its cover kills that obstruction. Hypothesis~II asks for a divisor representing the chosen series over $\mathbb C(B)$. This is a genuine descent issue over a general field, and it should not be dismissed as a failure to choose coordinates. Over the present function field it is resolved by Tsen's theorem and the vanishing of its Brauer group; see Section~\ref{cp2-x-7}.

Once the fibre restrictions are equivalent, their difference is vertical by the second equivalence criterion of Chapter~III. Because $f$ is smooth with irreducible fibres, vertical divisors are pullbacks from $B$. The parity or $s$-divisibility calculation in Sections~7--8 is then needed to replace a vertical divisor by $s$ times a divisor from the base. Equality of numerical characters yields a Riemann--Roch lower bound; it does not automatically give equality of actual dimensions. The text's ``same genus, degree, and hence dimension'' should be read with that qualification.

Several printed formulas lose information in the outline. On p.~382, a divisor pushed down through a degree-$m$ étale cover has intersection $m(4\pi-4)$ with $H$; it cannot be compared with the original parabicanonical divisor of intersection $4\pi-4$ without also multiplying that divisor. On p.~383 the expression $sK$ in the comparison with the $\pi$th adjoint must be $sK^\pi$, and ``para-$s$-canonical'' must be ``para-$s\pi$-canonical''. The scans confirm the missing indices. On p.~384 the fibres $K^*$ are called ``ellittiche'', although their genus remains $\pi>1$; it is their parameter curve that is elliptic. These are source slips, not alternative mathematical conventions.

More substantially, a nontrivial torsion class on a covering family can have trivial norm downstairs. Pushing down the constructed divisor therefore does not preserve the claimed non-paraequivalence without an additional norm calculation. Nor does non-paraequivalence by itself force different reduced supports. These are gaps in the proposed proof of Lemma~II. They do not invalidate the lemma in the present $\pi>1$ setting: the quotient description in Section~\ref{cp2-x-10} supplies an elliptic pencil, and sufficiently divisible pluricanonical systems contain two divisors supported on distinct fibres of that pencil. This is a replacement proof, not a claim that the printed torsion construction already establishes descent.

\section{Primo caso: il genere \texorpdfstring{$\pi$}{π} delle curve \texorpdfstring{$K$}{K} sia pari}
\label{cp2-x-7}

\subsection{The torsion cover and the parity argument, pp.~384--387}

The $2^{2\pi}$ square roots of $2K_H$ are exactly $K_H+\eta$ for $\eta\in\operatorname{Pic}^0(H)[2]$. Only $\eta=0$ is special: it has projective dimension $\pi-1$, whereas the remaining roots have dimension $\pi-2$. The roots form a finite étale local system because the family is smooth and the characteristic is zero. A connected monodromy-orbit cover $B'\to B$ is étale, hence $B'$ is again elliptic. Its base change $X'\to X$ is also étale.

The invariants which are automatically preserved in this situation are $K_{X'}^2=0$ and $\chi(\mathcal O_{X'})=0$, since both multiply by the degree. The individual values $p_g=0$ and $q=1$ need not be preserved. For example, the product cover in Section~\ref{cp2-x-10} has $q=\pi+1$ and $p_g=\pi$. Thus ``the same characters'' must not be expanded into a claim that all the original invariants survive. Where the later proof uses $p_a=-1$ and $p^{(1)}=1$ only, this issue causes no error; where it imports a construction with stronger hypotheses, it requires rechecking.

The parity argument itself can be made exact. After subtracting complete fibres and using the divisibility of $\operatorname{Pic}^0(B)$, suppose a divisor $A$ satisfies $2A\sim tH$ with $t\in\{0,1\}$. Then $A^2=0$ and adjunction gives
\[
 2p_a(A)-2=K_X\cdot A=t(\pi-1).
\]
If $t=1$, the right side is odd because $\pi$ is even, a contradiction. Thus $t=0$. More generally, if $sA\sim tH$, $0\leq t<s$, then
\[
 s\bigl(p_a(A)-1\bigr)=t(\pi-1),
\]
and $\gcd(s,\pi-1)=1$ forces $t=0$. This is the numerical core of the construction; it remains valid even when the effective-divisor and descent arguments surrounding it need repair.

The formula on p.~386 that compares the doubled adjoint systems contains an extra unindexed fibre term; the corrected comparison must have the same degree of the vertical divisor on the base on both sides. An extra whole fibre cannot disappear by passing to paracanonical equivalence. The printed comparison is $|2L'-2H''-2(H+H_1+\cdots+H_r)|=|H_0|$, whereas the half-divisor in the next line is $L'-H''-(H_1+\cdots+H_r)$. The unindexed $H$ must be deleted from the first expression.

\subsection{Representing a rationally given linear series}

Enriques's plane-quartic example on p.~385 is meant to show that a rationally specified family of divisor classes need not come with a rational choice of one divisor. In modern language its complete linear series can be a Brauer--Severi variety over the field of definition. Over an arbitrary field the obstruction is real. Over $k=\mathbb C(B)$, however,
\[
 \operatorname{Br}(k)=0
\]
by Tsen's theorem. The Stacks Project, Theorem~59.67.10 (Tag~03RD) and Lemma~59.67.11 (Tag~03RF), state precisely the needed facts. A Galois-invariant Picard class on the generic curve therefore lifts to a line bundle. A positive-dimensional space of sections over $k$ then contains a nonzero $k$-section, whose divisor closes up on the surface. This settles Hypothesis~II after Hypothesis~I has been achieved.

The book instead constructs a divisor in an odd multiple of the given series using quadrics. For $i=\pi-1$, the class
\[
 i(K_H+\eta)=iK_H+\eta
\]
still differs from $iK_H$, while its double is $2iK_H$. The parity of $i$ is exactly what preserves the nontrivial $2$-torsion twist. This algebra is correct. The geometric construction on pp.~387--388 invokes a pencil of quadrics and its conjugate $(\pi-1)$-gon. It needs the appropriate restriction map for quadrics, the claimed codimension of the subsystem, and control of degenerate pencils. These conditions are not proved in the text. In particular, the asserted codimension two does not follow without the required surjectivity of the multiplication or restriction map. The reference to the harmonizing invariant of Battaglini and Rosanes explains the construction for a suitable pencil; it does not verify that the curve supplies such a pencil in all cases.

On p.~387 the phrase ``semibicanonica'' for $\overline g+g$ should be ``semiquadricanonica'': its class is $2K_H+\eta$, of degree $4\pi-4$. For the later odd multiple $i\overline g$, the divisor $D$ restricts to $iK_H+\eta$, so $2D$ and $iC$ do have the same degree and class on the fibre when $C$ is parabicanonical. That doubled comparison is arithmetically consistent. What does not follow is the assertion that $D$ has a component absent from $iC$; a change in coefficients on the same support can account for nonequivalence. The next example isolates this issue.

\subsection{Different classes and different supports, pp.~388--389}

Even linear nonequivalence need not change the reduced support. On a product $E\times B$ of elliptic curves, choose fibres $F_p,F_q$ with $p-q$ of order four in $\operatorname{Pic}^0(B)$. The divisors $3F_p+F_q$ and $F_p+3F_q$ have the same reduced support, square zero, and arithmetic genus one; their difference $2(F_p-F_q)$ is nonzero two-torsion. They remain paraequivalent, since that difference belongs to $\operatorname{Pic}^0(E\times B)$.

The implication from nonequivalence to different components fails even in the numerical setting of the lemma. Let $H$ be a genus-two curve with hyperelliptic involution $\iota$, let $E'$ be elliptic, and choose a nonzero point $a\in E'[2]$. The action generated by $(\iota,T_a)$ on $H\times E'$ is free. Its quotient $X$ has $p_g=0$, $q=1$, and a smooth genus-two Albanese fibration over $E'/\langle T_a\rangle$. The other projection gives an elliptic fibration $h:X\to H/\langle\iota\rangle=\mathbb P^1$ with six double fibres $2E_i\sim F$, one for each Weierstrass point $P_i$ of $H$. The canonical bundle formula gives
\[
 K_X\sim-2F+\sum_{i=1}^6E_i,\qquad 2K_X\sim2F.
\]
Now put
\[
 C=2E_1+2E_2\in|2K_X|,\qquad D=3E_1+E_2.
\]
These effective divisors have identical reduced support, square zero, arithmetic genus one, and degree four on an Albanese fibre. Yet $2(D-C)\sim0$, while $(D-C)|_H=P_1-P_2$ is a nonzero two-torsion class. Every class in $\operatorname{Pic}^0(X)$ comes from the Albanese base and restricts trivially to $H$. Hence $D$ and $C$ are not even paraequivalent. Their different coefficients supply the nonequivalence without supplying a new irreducible component. This disproves that step of the argument, while the elliptic pencil itself confirms the conclusion of Lemma~II in this example.

The last step, pushing an elliptic divisor from the torsion cover down to $X$, also needs attention. Its image is elliptic if it is a nonconstant étale multisection of $f$, because the image maps to the elliptic base and Riemann--Hurwitz excludes a rational image. But its torsion class can become trivial under norm, and coefficients on common components can change. Neither issue is settled by the phrase that it ``responds'' to the desired divisor downstairs.

For $\pi=2$ the six Weierstrass points of each fibre provide a concrete divisor without a Brauer--Severi construction. Their relative locus is finite étale over $B$, because the family is smooth; its connected components are elliptic curves. On a genus-two curve the Weierstrass divisor is linearly equivalent to $3K_H$, hence its double is $6K_H$. This verifies the special ``semiparasesticanonical'' construction on p.~389, with the possible splitting into several elliptic components understood. It does not by itself settle the general distinct-support descent problem.

\section{Secondo caso: il genere \texorpdfstring{$\pi$}{π} delle curve \texorpdfstring{$K$}{K} sia dispari}
\label{cp2-x-8}

On pp.~389--392 take an odd prime $s\mid\pi$ and a nonzero class $\eta\in\operatorname{Pic}^0(H)[s]$. The noncanonical root $K_H+\eta$ has the same nonspecial dimension $\pi-2$ as in the even case. The étale torsion cover again handles the rational choice of the class, and the divisibility computation of Section~\ref{cp2-x-7} applies because $s$ is coprime to $\pi-1$.

The distinguished combination of series is
\begin{align*}
 \overline{\overline g}
  &=(2\pi-2)\overline g-(\pi-2)g\\
  &=\pi K_H+(2\pi-2)\eta.
\end{align*}
It has degree $\pi(2\pi-2)$, and its $s$-fold multiple is $s\pi K_H$. Since $s\mid\pi$ is odd, $s$ does not divide $2\pi-2$; hence the twist $(2\pi-2)\eta$ is nonzero. This is the exact algebraic reason that the constructed series is not pluricanonical. The arithmetic is sound.

Enriques tries to extract a rational divisor in this class by making a member of $\overline g$ have contact of order $\pi-2$ at a moving point and summing the residual intersections. The class computation gives the displayed combination. What is missing is a proof that the required jet conditions are independent and yield a rationally varying residual divisor in all cases, including special or ramified configurations. The construction is a sketch of a covariant, not a complete existence proof. Over $\mathbb C(B)$ the Tsen argument already supplies a divisor in $K_H+\eta$ itself, so this more elaborate multiple is unnecessary for the modern proof.

The system comparisons on pp.~390--391 are governed by the same formula $sA\sim tH$ and the same arithmetic-genus obstruction. The argument that equal numerical characters ensure sufficiently many sections uses Riemann--Roch as a lower bound, not equality of actual dimensions. The final cover on p.~392 makes \emph{Hypothesis~I} hold; the printed reference to Hypothesis~II is a wrong label. After that cover, the norm and support issues noted in the preceding section remain. Consequently this section establishes the torsion and divisibility calculations, while its global proof of Lemma~II remains incomplete. Section~\ref{cp2-x-10} supplies a complete replacement for the lemma in the stated setting.

\section{Lemma III}
\label{cp2-x-9}

Lemma~III, pp.~392--393, starts with a para-$i$-canonical divisor $C$ whose components are disjoint elliptic curves, and an effective divisor $D$ with equivalent restrictions to every fibre $H$. Apart from common components, $D$ is assumed disjoint from $C$. It concludes that $D$ is also para-$i$-canonical and has the same elliptic component structure. The conclusion that $D^2=0$ is not an additional hypothesis. The horizontal intersection $C\cdot H>0$ and the smooth irreducible fibres are essential. The second equivalence criterion of Chapter~III gives
\[
 D-C\sim f^*L
\]
for a line bundle $L$ on $B$. One way to see this is to trivialize the difference on the generic fibre, take the divisor of a rational trivializing section, and observe that the remaining divisor is vertical. Irreducibility and reducedness of the fibres make it a pullback from $B$.

The degree of $L$ is not automatically zero from fibrewise equivalence alone. It is zero here because
\[
 0=C\cdot(D-C)=(C\cdot H)\deg L,
\]
because every component of $C$ has square zero and distinct components of $C$ are disjoint. Thus $L\in\operatorname{Pic}^0(B)$, so $C$ and $D$ are paraequivalent. It follows that $D^2=0$ and $p_a(D)=1$; the component argument in Section~\ref{cp2-x-5} gives the asserted elliptic structure. On an elliptic curve every degree-zero line bundle is represented by a difference of two points, explaining the book's expression $K_1-K_2$.

The alternative proof in the text describes the same vertical difference by moving fibres. Its numerical conclusion is correct, but the intermediate degree-zero calculation should be made explicit. If reducible or multiple fibres were allowed, vertical components could appear which are not whole pullbacks, and the statement would require modification. In the present $\pi>1$ case their absence was proved in Section~\ref{cp2-x-5}; it is not a newly imposed genericity assumption.

\section{Conclusione: un fascio lineare di curve ellittiche}
\label{cp2-x-10}

\subsection{The printed construction, pp.~393--394}

For Albanese fibre genus $\pi>1$, the theorem is correct: $X$ has a base-point-free pencil over $\mathbb P^1$ whose general fibre is elliptic. Enriques uses the two divisors of Lemma~II to define a pencil of degree $n=i(2\pi-2)$ on each $H$. After removing common components and base points, Riemann--Hurwitz bounds the number of branch values by $2n+2\pi-2$. The union of the corresponding divisors on the fibres is then compared with multiples of $C$ and $D$ using Lemma~III.

Several steps need more than the numerical bound. The relative branch divisor must be constructed with its multiplicities; its claimed avoidance of the two chosen divisors uses the unramifiedness of a nonconstant map between smooth elliptic curves. The resulting ruled model must carry the two sections with the claimed complementarity. Finally, scaling the two sections by a $\mathbb G_m$-action on that ruled model does not by itself prove that its inverse images on the branched cover stay elliptic. One needs control of the branch divisor under the action. These are genuine missing family and lifting arguments. The earlier support and norm gaps in Lemma~II also enter here. Accordingly the displayed geometric construction is a sketch, although its conclusion in this case is valid.

If a moving elliptic system has disconnected members, one must take its connected-fibre factorization. Its base cannot be an additional irrational curve because $p_g=0$, by the independence-of-pencils result of Chapter~IX. Thus the irreducible components form a pencil with rational base. The Euler calculation then gives total fibre defect zero. The constancy of the elliptic modulus follows because the $j$-map has no poles and extends across the multiple smooth fibres; a regular function on the complete base is constant. Absence of nodes by itself would not be a definition of constant modulus, but the extension supplies the missing justification.

\subsection{A complete modern proof}

The smooth genus-$\pi$ family $f:X\to B$ from Section~\ref{cp2-x-5} is isotrivial. Here is the analytic reason, valid over $\mathbb C$. Pull it back to the universal cover $\mathbb C\to B$ and trivialize the integral first-cohomology local system. Its period map lands in Siegel upper half-space, which has a bounded realization. Liouville's theorem makes this period map constant. Torelli then makes the isomorphism class of the curve constant. Since a curve of genus greater than one has a finite automorphism group, the relative isomorphism torsor is finite étale; after a finite étale base change the family is a product. See Del Padrone and Mistretta (2006, Theorem~1.3) for this argument.

Choose the connected Galois cover $E'\to B$ which kills the monodromy. Since $B$ is elliptic and the cover is étale, $E'$ is elliptic and its deck group $G$ consists of translations. In particular $G$ is abelian. We obtain
\begin{equation}
\label{cp2-x-translation-quotient}
 X\cong(H\times E')/G,
\end{equation}
where the diagonal action is free. Kernels can be divided out so that the action on $H$ is faithful. Taking invariant holomorphic one-forms gives
\[
 q(X)=g(H/G)+g(E'/G)=g(H/G)+1.
\]
Since $q(X)=1$, the quotient $H/G$ is $\mathbb P^1$. The first projection therefore descends to
\[
 h:X\longrightarrow H/G\cong\mathbb P^1
\]
with elliptic general fibre $E'$. Its exceptional fibres are smooth elliptic curves with multiplicity equal to the stabilizer order on $H$. Translations of $E'$ commute with the diagonal action and descend to a one-dimensional elliptic group of automorphisms of $X$.

This proves the theorem and all the existence conclusions required from Lemmas~I--III in this case. For example, choose a sufficiently divisible $m$ which kills every inertia denominator. Then $mK_X$ is pulled back from a positive-degree line bundle on $\mathbb P^1$. Two different effective divisors of that line bundle give para-pluricanonical divisors with different reduced supports, as required by Lemma~II. This proof imports isotriviality and Torelli explicitly; it does not disguise them as consequences of the elementary dimension counts in the book.

\section{Il caso di un fascio ellittico di curve di genere uno}
\label{cp2-x-11}

\subsection{What remains true and where the proof fails, pp.~394--397}

For a minimal nonruled surface with $p_g=0$, $q=1$ and genus-one Albanese fibres, the calculation still gives $K_X^2=e(X)=0$. Every fibre is therefore either smooth and reduced or a multiple of a smooth elliptic curve. The fibration is isotrivial. None of these facts excludes multiple fibres, nor do they imply that the monodromy acts by translations on the fibre.

Enriques embeds the general fibres as elliptic normal curves of degree $m$. The number of projective isomorphisms between two such polarized curves is $2m^2$, $4m^2$, or $6m^2$, according as the elliptic curve has general, harmonic, or equianharmonic modulus. This count is correct: $m^2$ translations preserve a degree-$m$ polarization, and the automorphism group fixing the origin has order $2$, $4$, or $6$. An orbit of a point under these identifications defines a multisection. It is elliptic only if its cover of the elliptic base is unramified. The very degeneration to a multiple fibre that the text discusses can introduce ramification into that multisection. The proof cannot assume its genus before excluding that degeneration.

On p.~395 the printed $n=6$ must be $m=6$, as the edition's apparatus observes. The subsequent counts $36$ and $18$ also do not match the preceding total $2m^2=72$ for a general degree-six polarized elliptic curve. They might count one chosen species of isomorphisms, but no such restriction is stated. This internal ambiguity has no bearing on the more serious failure: the proof attempts to exclude an actually existing class of multiple fibres.

On p.~396 it represents the surface as a cover of a cubic cylinder with branch divisors in two transverse directions. If their inertia orders are $r,s>1$, Enriques adds the numbers
\[
 (s-1)\frac ns,\qquad(r-1)\frac nr
\]
of transpositions in minimal factorizations of the two monodromies, and infers that their product needs at least $n$ transpositions. Minimal transposition length is not additive under multiplication. Even if each factor has a specified cycle structure, their product can have much smaller length; for example a permutation multiplied by its inverse is the identity. Thus the stated lower bound does not follow.

There is a second, independent local problem. A singular inverse image of an oblique curve does not force the covering surface to be singular. The finite map
\[
 (u,v)\longmapsto(u^s,v^r)
\]
has a smooth total space and branch along two transverse axes. The inverse image of a general line through their crossing has equation $u^s+\lambda v^r=0$ and is singular at the origin for $r,s>1$. This is exactly the kind of singular curve that the argument mistakes for a singularity of the surface. Neither the local inference nor the theorem it is used to prove can be retained.

\subsection{An explicit counterexample within the stated hypotheses}

\begin{example}
\label{cp2-x-counterexample}
Let $B$ be an elliptic curve. Choose a connected double cover $D\to B$ branched at two distinct points $b_1,b_2$, and a connected étale double cover $B'\to B$. Such a branched cover exists: choose $L\in\operatorname{Pic}^1(B)$ with $L^2\cong\mathcal O_B(b_1+b_2)$ and take the usual double-cover algebra. Set
\[
 C=D\times_B B'.
\]
The two quadratic extensions are distinct, one ramified and one unramified, so their fibre product is connected. It is smooth of genus three: it is an étale double cover of the genus-two curve $D$. Write $G=\langle t,c\rangle\cong(\mathbb Z/2)^2$, where $t$ is the branched-cover involution and $c$ the étale-cover involution. Only $t$ has fixed points on $C$; both $c$ and $tc$ act freely because they move the $B'$ coordinate.

Let $E$ be another elliptic curve, and choose $0\ne a\in E[2]$. Let
\[
 t(x)=x+a,\qquad c(x)=-x
\]
on $E$. These involutions commute. The diagonal action of $G$ on $C\times E$ is free: the only nonidentity element with fixed points on $C$ is $t$, and $t$ has no fixed points on $E$. Therefore
\[
 X=(C\times E)/G
\]
is a smooth projective surface with an étale product cover of degree four.

Since $K_{C\times E}$ is the pullback of $K_C$, it is nef and has square zero. Its Kodaira dimension is one. These properties descend through the finite étale cover, so $X$ is minimal, $K_X^2=0$, and $\kappa(X)=1$. Furthermore,
\[
 \chi(\mathcal O_X)=\tfrac14\chi(\mathcal O_{C\times E})=0,
 \qquad
 q(X)=g(C/G)+g(E/G)=1+0=1.
\]
Thus $p_g(X)=0$ and $p_a(X)=-1$, exactly as required in the statement on p.~396.

Nevertheless,
\[
 H^0(X,T_X)
 =H^0(C\times E,T_{C\times E})^G
 =H^0(E,T_E)^G=0,
\]
because $g(C)=3$ gives $H^0(C,T_C)=0$ and $c$ acts as $-1$ on the one-dimensional space $H^0(E,T_E)$. Hence $\operatorname{Aut}^0(X)$ is trivial.

The Albanese fibration $f:X\to B$ has elliptic general fibre and exactly two double smooth elliptic fibres, over $b_1,b_2$. The other quotient map $X\to E/G\cong\mathbb P^1$ has general fibre of genus three, not one. There cannot be an additional transverse genus-one fibration: two independent fibre classes of square zero have a sum of positive square, and a nef $K_X$ orthogonal to both would be numerically trivial by the Hodge index theorem. That would contradict $\kappa(X)=1$ here. Thus the second elliptic pencil promised by the book does not exist.
\end{example}

This quotient belongs to the nontranslation case in Catanese, Frapporti, Glei\ss ner, Liu, and Sch\"utt (2025, Lemma~1.1(2)). Their Appendix~B includes exactly the group $\mathbb Z/2\times\mu_2$ with elliptic-base monodromy $(\epsilon,0;t,t)$.

The example also makes the missing line-bundle datum visible. Let $\eta\in\operatorname{Pic}^0(B)[2]$ be the nontrivial line bundle defining $B'\to B$, and write the two double fibres as $2F_1,2F_2$. The canonical formula is
\[
 K_X\sim f^*\eta+F_1+F_2.
\]
The Hodge line $\eta$ is nontrivial because the involution $c$ acts by $-1$ on the elliptic differential. Pushing forward its multiples gives
\[
 P_m=h^0\!\left(B,\eta^m\otimes
 \mathcal O_B\bigl(\lfloor m/2\rfloor(b_1+b_2)\bigr)\right).
\]
Consequently $P_1=0$ and
\[
 P_m=2\lfloor m/2\rfloor\quad(m\geq2).
\]
In particular $(P_2,P_3,P_4,P_6,P_{12})=(2,2,4,6,12)$. The nonzero higher plurigenera are consistent with the ruledness criterion proved later. The counterexample targets the asserted group action and second elliptic pencil, not that criterion.

\subsection{The two paraelliptic examples, p.~397}

The first example is the double cover
\[
 u^2=(z-a)(z-b)(z-c)(z-d)\psi(x,y),\qquad f(x,y)=0,
\]
where $f=0$ is an elliptic plane cubic and the quadric $\psi=0$ cuts it in six simple points. Compactify the cylinder as $B\times\mathbb P^1$. The branch divisor has six vertical components and four horizontal components. If $M$ denotes its half class, then $M$ has degrees $(3,2)$ on the two factors. The double-cover formula, already treated in Part~1, Chapter~V, \S~7, gives
\[
 \chi(\mathcal O_X)=2\chi(\mathcal O_{B\times\mathbb P^1})
 +\tfrac12 M\cdot(K_{B\times\mathbb P^1}+M)=3.
\]
The $24$ crossings of branch components produce ordinary double points, whose resolutions do not change $p_g$ or $q$. One has $q=1$, and hence
\[
 p_g=3,\qquad p_a=2.
\]
The printed values $p_a=3$, $p_g=4$ are therefore wrong. Both values occur in the printed text. Equivalently, the surface is birational to a diagonal quotient of a genus-four double cover of $B$ and an elliptic double cover of $\mathbb P^1$; the invariant space of two-forms has dimension three.

In the second example
\[
 u^3=(z-a)(z-b)(z-c)\psi(x,y),\qquad f(x,y)=0,
\]
the stated cubic $\psi$ has three triple contacts, one double contact, and one simple intersection with the cubic $f$. Their total intersection multiplicity is $3\cdot3+2+1=12$, contradicting B\'ezout's value nine. One possible repair is to replace ``three'' triple contacts by ``two''; another is to make $\psi$ a quartic. The printed text says three and cubic; the two proposed changes are alternative mathematical repairs.

With either repair the residual cyclic-cover branching on $B$ has weights two and one at two points. The corresponding cyclic triple cover has genus three. Its two nontrivial holomorphic-differential eigenspaces each have dimension one; the cyclic triple cover of $\mathbb P^1$ branched at $a,b,c$ is elliptic. Taking the diagonal quotient therefore gives $q=1$, $p_g=1$, and $p_a=0$, agreeing with the stated invariants after the contact-count correction. The quotient singularities are rational, so resolving them does not alter these Hodge numbers.

The proposed classification problem for ``paraelliptic'' surfaces on pp.~397--398 should be read historically. The counterexample above shows that $p_a=-1$ does not separate the two monodromy behaviours in the way suggested by the terminology. Modern isotrivial quotient data provide the appropriate distinction.

\section{Superficie ellittiche}
\label{cp2-x-12}

\subsection{The valid translation class and the invalid exhaustive summary, pp.~398--400}

The opening summary on p.~398 repeats the false conclusion of Section~11. It must be restricted to the translation class constructed in Section~\ref{cp2-x-10}, together with the bielliptic case discussed below. It is not an exhaustive description of all irregular surfaces with $p_g=0$.

Within that class the broader calculation is correct. If
\[
 X=(C\times E)/G
\]
with $G$ acting freely diagonally and by translations on the elliptic factor $E$, and if $\rho=g(C/G)$, then invariant differentials give
\[
 q(X)=\rho+1,\qquad
 p_g(X)=\rho,\qquad
 \chi(\mathcal O_X)=0.
\]
The second equality follows because the holomorphic differential on $E$ is invariant, so invariant two-forms are precisely invariant one-forms on $C$ multiplied by that differential. This proves the equality $p_g=\rho$ more directly than the canonical-series comparison in the book. It also makes clear why the equality can fail after introducing nontranslation monodromy: that differential then carries a nontrivial character.

The book calls these surfaces ``ellittiche'', calls the nonruled ones ``proprie'', and includes elliptic ruled surfaces as ``improprie''. Its ``proper'' includes the bielliptic surfaces of Kodaira dimension zero. Modern ``properly elliptic'' means Kodaira dimension one; these expressions must not be equated. The determinant $n=C\cdot H$ is the degree of the finite map to the product of the two pencil bases. It belongs to the chosen pair of pencils or a normalized quotient presentation. A claim that it is independent of every possible presentation needs an additional minimality convention.

Translations of $E$ descend and give the elliptic group whose orbits are the elliptic fibres. The pointwise proof on p.~399 chooses an isomorphism of fibres respecting the intersection groups. The assertion of uniqueness requires a marking: special elliptic automorphisms and the stabilizer of the isogeny data otherwise leave choices. The quotient construction avoids this ambiguity and works uniformly for general, harmonic, and equianharmonic modulus.

Conversely, a nontrivial elliptic group acting birationally on a projective surface leads, after regularizing the action, to this product-quotient description up to birational equivalence. A useful justification uses an ample line bundle. Its restrictions to the elliptic orbits define a morphism to a Picard variety of the acting elliptic curve, and translation on an orbit induces multiplication by the positive degree of that line bundle. Taking a fibre of this morphism produces a transverse curve; the finite kernel of the multiplication isogeny gives a presentation as a finite quotient of the product of that curve with an elliptic curve. This explains both the second pencil and the rank-at-most-two translation group. Projectivity is essential to this argument, through the ample line bundle. On a minimal nonruled surface the birational action is already regular.

The printed converse proof invokes the transverse pencil before fully constructing it and compresses this equivariant geometry into a continuous-system argument. Its invariant calculation also contains the sign error
\[
 \Delta=13-12p_a-p^{(1)}
\]
on p.~400. The correct expression is $13+12p_a-p^{(1)}$, as follows immediately from Noether's formula. The printed sign is minus. The subsequent inequality for $p_a$ uses the corrected sign. In the ruled case $p^{(1)}\geq1$ is not available; that case must be handled separately, as the product-quotient argument does.

\subsection{The determinant-one model, p.~401}

The simultaneous intersections of the two pencils form the fibres of a finite abelian cover of their product. Its group is a finite subgroup of an elliptic translation group, hence has at most two generators. A determinant-one model is birational to the product of an elliptic curve with the other pencil base; in the $p_g=0$ case it is $E\times\mathbb P^1$, the cubic-cylinder model. A cyclic group leads to one Kummer radical, and a two-generator group to two independent radicals. This is a statement about the translation subgroup, not about every finite subgroup of the full automorphism group of an elliptic curve: the latter can be nonabelian and is crucial in the omitted family.

The construction in the next section is valid for this class. It cannot restore the false exhaustion claim merely by allowing more choices of the same abelian radicals.

\section{Costruzione delle superficie ellittiche di genere \texorpdfstring{$p_g=0$}{pg=0}}
\label{cp2-x-13}

\subsection{Connectedness in both directions, pp.~401--403}

The two conditions on p.~402 are necessary and distinct. The cover of a ruling line must be connected, and the cover of a transverse elliptic section must also be connected. Connectedness of the total surface does not imply connectedness of both sets of fibres. Enriques's example illustrates this correctly:
\[
 u^2=(z-a)(z-b)(z-d)(z-e)\psi(x,y)^2,
 \qquad f(x,y)=0.
\]
After replacing $u$ by $u/\psi$, this is birational to the product of the elliptic curve $f=0$ with the elliptic double cover of $\mathbb P^1$ branched at $a,b,d,e$. It is irreducible and has $p_g=1$, $q=2$. The fibres over a fixed $z$ split into two elliptic components, and the connected pencil has elliptic base. Its determinant is one, despite the degree-two map to the original cylinder. This example is sound.

For an unramified cyclic cover of an elliptic curve $B$ one should formulate the Kummer construction with a rational function $\varphi$ and
\[
 \operatorname{div}(\varphi)=nD,
 \qquad\eta=\mathcal O_B(D)\in\operatorname{Pic}^0(B)[n].
\]
Here $D$ is a degree-zero divisor; in the plane model it is the difference between a degree-three contact divisor and a line section. Divisibility of \emph{all} valuations, including poles and points at infinity, is what makes the normalization unramified. Counting only affine contact points is insufficient.

The cover $w^n=\varphi$ is connected of degree $n$ if and only if $\eta$ has exact order $n$. The text's criterion that the contact triple not belong to the plane-line series says only $\eta\ne0$. It is sufficient when $n$ is prime, but not for composite $n$. For instance, let $n=4$ and let $\eta$ have order two. One can arrange $\varphi=h^2$ with $\operatorname{div}(h)=2D$. Then
\[
 w^4-\varphi=(w^2-h)(w^2+h)
\]
is reducible although the degree-three contact divisor is not a line section. Exact order repairs the asserted irreducibility criterion on pp.~403--405.

For two radicals of orders $p,q$, connectedness of degree $pq$ means that the corresponding character data are independent: the homomorphism from $\mathbb Z/p\oplus\mathbb Z/q$ to $\operatorname{Pic}^0(B)$ determined by the two torsion classes is injective. Requiring both classes to be nonzero, or merely different, is not enough. This is the precise meaning needed for the source's unspecified ``dissomiglianza''. In the later order-nine case, for example, $\eta$ and $-\eta$ are distinct nonzero $3$-torsion classes but generate only a cyclic group of order three.

\subsection{Ramification and the abelian constraint, pp.~403--405}

In the cyclic case write
\[
 u^n=\varphi(x,y)\prod_{i=1}^t(z-a_i)^{r_i}.
\]
The assertion that infinity is unramified is equivalent to $\sum r_i\equiv0\pmod n$. Put $h_i=\gcd(n,r_i)$ and $s_i=n/h_i$. Above $a_i$ there are $h_i$ points, each of ramification index $s_i$. Therefore
\begin{equation}
\label{cp2-x-rh}
 2g(C)-2=n\left(t-2-\sum_{i=1}^t\frac1{s_i}\right),
\end{equation}
where $C$ is the connected curve over the ruling line. This is the correct equality on p.~404. The repeated equality on p.~405 prints $2n-2\pi-2$ instead of $2n+2\pi-2$ on its right-hand side; the minus sign is a source misprint.

Connectedness of this cyclic curve also requires that its local monodromies generate the full cyclic group, equivalently $\gcd(n,r_1,\ldots,r_t)=1$. It is not supplied by the sum-zero condition. Together with exact order of the elliptic torsion datum, these conditions ensure the intended two connected pencils.

For an abelian cover of $\mathbb P^1$, let $S_i$ be its inertia elements. They satisfy
\[
 S_1\cdots S_t=1.
\]
If $M=\operatorname{lcm}(s_j:j\ne i)$, commutativity gives
\[
 S_i^M=\left(\prod_{j\ne i}S_j^{-1}\right)^M=1,
\]
so $s_i\mid M$. This proves the least-common-multiple condition on p.~404. It is only a necessary condition: one must still choose inertia elements of the prescribed orders, with product one, generating a group which embeds into the torsion of an elliptic curve. The printed calculation with the product of the orders does not alone yield the sharper least-common-multiple conclusion; the argument above does.

Finally, the degree normalization of $\varphi$ on p.~405 is valid at the level of divisor classes and function fields. For a fixed torsion class $\eta$, choose an effective divisor of degree three in the class of a line section plus $\eta$. Its $n$-fold multiple is cut by a degree-$n$ plane curve. Changing the chosen divisor within the same class changes the Kummer function by an $n$th power, hence gives a birationally identical cover. This justifies the admissible continuous variation within a fixed torsion class. It does not say that arbitrary variation of branch points or of torsion data preserves the surface. Nor must a polynomial of higher degree literally factor as $\varphi_1^n\varphi$ in the ambient polynomial ring: the corresponding relation is a rational-function identity on the cubic, after the necessary change of coordinate $u$. This is the precise repair of the ``si riconosce facilmente'' reduction.

\section{Superficie coi plurigeneri nulli: caratterizzazione delle rigate}
\label{cp2-x-14}

\subsection{The canonical divisor and the exact floor formula, pp.~405--408}

For the translation class with $p_g=0$, let $h:X\to\mathbb P^1$ be the elliptic pencil and write its multiple fibres as $s_iF_i$, $i=1,\ldots,t$. Let $F$ be a general fibre, so $s_iF_i\sim F$. The canonical formula is
\begin{equation}
\label{cp2-x-canonical}
 K_X\sim-2F+\sum_{i=1}^t(s_i-1)F_i.
\end{equation}
It follows either by ramification from the product cylinder or by descending differentials from the quotient model. It is a linear-equivalence formula, including the reduced multiple fibres. Replacing every $F_i$ by $F/s_i$ yields only a rational numerical relation and loses the information needed for the plurigenera.

The book removes exceptional components from a singular cylinder model and calls the result the ``pure'' canonical curve. The modern interpretation is the canonical class on the smooth minimal model, not the arbitrary deletion of fixed components from a canonical system. On that model \eqref{cp2-x-canonical} is exact. The assertion that $\sum(s_i-1)F_i$ has only one section follows rigorously from connected fibres and the pushforward calculation below; the degeneration argument invoking connectedness of fibres is an intuitive version of it. In the preliminary subtraction on p.~406, the bound obtained for the coefficients is $0\leq k_i\leq s_i-1<s_i$, rather than the unnecessarily strict printed $k_i<s_i-1$. The weaker, justified bound is exactly what the pushforward argument needs.

For any integers $c_i\geq0$,
\[
 h_*\mathcal O_X\left(\sum_i c_iF_i\right)
 \cong\mathcal O_{\mathbb P^1}\left(\sum_i\left\lfloor\frac{c_i}{s_i}\right\rfloor b_i\right),
\]
where $b_i\in\mathbb P^1$ is the point under $F_i$. Equivalently, a rational function pulled back from the base has valuation divisible by $s_i$ along $F_i$, so a pole of order at most $c_i$ allows a base pole of order at most $\lfloor c_i/s_i\rfloor$. Thus
\begin{equation}
\label{cp2-x-plurigenus}
 P_m=\max\left\{0,\ 1-2m+
 \sum_i\left\lfloor m\left(1-\frac1{s_i}\right)\right\rfloor\right\}
 =\max\left\{0,\ 1+m(t-2)-\sum_i\left\lceil\frac m{s_i}\right\rceil\right\}.
\end{equation}
This is the printed formula with its necessary zero truncation. The residual coefficient of $F_i$ is the unique $0\leq h_i<s_i$ for which $m+h_i$ is divisible by $s_i$, so $(m+h_i)/s_i=\lceil m/s_i\rceil$. The ceiling calculation, rather than treating the $F_i$ as fractional copies in the Picard group, explains the formula.

For a nonruled member of this class, Riemann--Hurwitz gives
\[
 t-2-\sum_i\frac1{s_i}\geq0,
\]
so $t\geq3$. Consequently $P_2=t-3$. If $t>3$, a nonzero bicanonical section has nonzero powers, and both $P_4$ and $P_6$ are positive. If $P_2=0$, then $t=3$.

The display on p.~408 omits the factors $1/s_i$ from the ramification sum. The next inequalities give the weaker bound $\sum(1-1/s_i)\geq1$, and hence $\sum1/s_i\leq t-1$. Those weaker inequalities are true but do not imply the following bound $\sum1/s_i\leq1$ when $t=3$. The required stronger bound is $\sum(1-1/s_i)\geq2$, yielding $t-2$. The missing factor is an error; the weaker inequalities themselves are true, but are insufficient for the subsequent deduction.

\subsection{The small plurigenera: correcting the case table, pp.~409--410}

For $t=3$, let $a_j$ denote the number of multiplicities equal to $j$. The exact formulas are
\begin{align*}
 P_3&=\max(0,1-a_2),\\
 P_4&=\max(0,2-a_2-a_3),\\
 P_6&=\max(0,4-2a_2-a_3-a_4-a_5).
\end{align*}
They follow immediately from \eqref{cp2-x-plurigenus} and are safer than the coarse four-case table in the book. The partition on p.~409 should have $s_1=3$ in its third case, since the fourth treats $s_1\geq4$. Even after that correction the asserted constant values of $P_4$ and $P_6$ over whole cases are false. These constant values occur in the printed case table.

For a valid concrete example take a cyclic degree-twelve cover of $\mathbb P^1$ with exponents $4,3,5$ at three branch values. Their sum is twelve and they generate $\mathbb Z/12$. The inertia orders are $(3,4,12)$, and Riemann--Hurwitz gives genus three. Pair it with an étale cyclic degree-twelve cover of an elliptic curve, using the diagonal free action. The resulting surface is in Enriques's translation class and has
\[
 P_3=1,\qquad P_4=1,\qquad P_6=2.
\]
It belongs to the source's third case, where $P_4=0$ and $P_6=1$ are printed. Thus these are errors on actual surfaces, not merely on inadmissible integer triples. Similarly, $(2,5,10)$ gives a cyclic example with $P_4=P_6=1$, contradicting the printed $P_6=0$ in the second case.

The conclusion needed for ruledness nevertheless follows from the corrected formulas. The inequality $\sum1/s_i\leq1$ allows at most one multiplicity two. If $P_4=0$, either the triple is $(2,3,s)$ with $s\geq6$, or it is $(3,3,s)$ with $s\geq3$. In the former case $P_6=1$; in the latter $P_6\geq1$. Therefore every nonruled surface in the translation class satisfies
\[
 P_4+P_6\geq1.
\]
This argument does not even require the abelian least-common-multiple condition. That condition becomes indispensable for the finer $P_{12}>1$ distinction.

\subsection{The twelfth plurigenus and the existence condition, pp.~410--412}

For $t$ multiple fibres,
\[
 P_{12}=\max\left(0,12t-23-\sum_i\left\lceil\frac{12}{s_i}\right\rceil\right).
\]
The floor brackets are missing in the printed first formula on p.~411. They cannot be omitted: $\lceil12/s_i\rceil$ is generally not $(s_i+11)/s_i$.

If $t\geq5$, each ceiling is at most six, so
\[
 P_{12}\geq6t-23\geq7.
\]
The printed bound for the sum is $24$, independent of $t$, and the following bound $12t-47$ is consequently wrong. For example six multiplicities equal to two already give a sum of $36$. If $t=4$, the minimum is one, attained exactly for $(2,2,2,2)$; if one multiplicity is at least three, then $P_{12}\geq3$. This portion of the source's argument is correct after the earlier floor correction.

For $t=3$, suppose the orbifold degree in \eqref{cp2-x-rh} is positive and impose the abelian least-common-multiple condition. Order $s_1\leq s_2\leq s_3$. If $(s_1,s_2)=(2,3)$, divisibility forces $s_3=6$, which has degree zero and is excluded. If $(s_1,s_2)=(2,4)$, it forces $s_3=4$, also excluded. In particular the printed triple $(2,4,8)$ is impossible: $8$ does not divide $\operatorname{lcm}(2,4)=4$. Its inclusion is an error in the enumeration.

The remaining cases give the required bound transparently. For $(2,5,s_3)$, divisibility and positive orbifold degree force $s_3=10$, and the ceiling sum is $6+3+2=11$. For $s_1=2$, $s_2\geq6$, the sum is at most ten. For $s_1=3$, the pair $s_2=3$ would force $s_3=3$ and degree zero, so $s_2\geq4$ and the sum is at most ten. If $s_1\geq4$, it is at most nine. Therefore $P_{12}>1$ in every positive-degree case. The last line of the book's enumeration also prints $s_1>4$, omitting equality; it should read $s_1\geq4$.

The footnote about Geppert is mathematically useful. The formal triple $(2,5,5)$ has positive orbifold degree but $P_{12}=1$. It violates the abelian condition because two does not divide five. Thus the strict inequality cannot be deduced from the floor formula and positive degree alone. Enriques explicitly says that his 1914 presentation did not make this condition clear, and attributes the criticism to Geppert (1931) and the explicit correction to Enriques--Campedelli (1934, \S~15).

When the orbifold degree is zero, the only multiplicity lists are $(2,2,2,2)$, $(2,3,6)$, $(2,4,4)$, and $(3,3,3)$. Their canonical indices divide twelve, and $P_{12}=1$. In the translation class this is exactly the bielliptic case. Positive degree gives Kodaira dimension one and $P_{12}>1$.

\subsection{A complete ruledness criterion including the omitted family}

\begin{theorem}
\label{cp2-x-ruled-criterion}
For a smooth projective complex surface $X$, the following are equivalent:
\[
 X\text{ is birational to a ruled surface},\qquad
 P_4(X)=P_6(X)=0,\qquad P_{12}(X)=0.
\]
\end{theorem}

The book's proof is incomplete because it applies the translation-class calculation to every nonruled irregular $p_g=0$ surface. The missing case can be supplied without altering the theorem. If $P_4=P_6=0$, powers of sections give $p_g=P_2=0$. When $q=0$, Castelnuovo's criterion applies; when $q>1$, Section~\ref{cp2-x-4} applies. Suppose $q=1$ and $X$ is minimal nonruled. Then $K_X^2=e(X)=0$. If the Albanese fibres have genus greater than one, Sections~\ref{cp2-x-10} and \ref{cp2-x-14} give the contradiction $P_4+P_6\geq1$.

If the Albanese fibres have genus one, write $f:X\to B$ with $B$ elliptic. The canonical bundle formula is
\[
 K_X\sim f^*L+\sum_i(m_i-1)F_i,
 \qquad \deg L=\chi(\mathcal O_X)=0.
\]
Here $L$ is the Hodge line and the multiple fibres are $m_iF_i$. For the general canonical bundle formula, see Bombieri--Mumford (1977, \S~1, Theorem~2, pp.~27--29). In the present characteristic-zero setting its multiple-fibre term follows locally from the ramification coordinate, and its base term from the relative dualizing sheaf. The monodromy character on the elliptic differential is retained in $L$; this is precisely the datum omitted in the failed translation argument.

If there is any multiple fibre, then
\[
 P_4=h^0\!\left(B,L^4\otimes
 \mathcal O_B\left(\sum_i\left\lfloor4\left(1-\frac1{m_i}\right)\right\rfloor b_i\right)\right)>0,
\]
because the line bundle has positive degree, at least two. If there are no multiple fibres, the fibration is a smooth isotrivial elliptic bundle. Its differential monodromy has order $r\in\{1,2,3,4,6\}$. Thus $L^r\cong\mathcal O_B$. The case $r=1$ gives $p_g=1$ and is excluded; for each remaining order, either $r\mid4$ or $r\mid6$, again giving $P_4+P_6>0$. This covers nontranslation monodromy as well as the bielliptic examples.

Ruled surfaces have all positive plurigenera zero, as established in Part~1. Finally, a nonzero section of $4K_X$ or $6K_X$ has a nonzero power in $12K_X$. Hence $P_{12}=0$ implies $P_4=P_6=0$, and the reverse implication follows through ruledness. This completes the proof. For the same reason, among minimal surfaces with $p_g=0,q=1$, the condition $P_{12}=1$ still singles out the bielliptic case: a multiple fibre over the elliptic Albanese base contributes at least six to the degree defining $P_{12}$, while the translation case was handled above.

\section{Nota storica}
\label{cp2-x-15}

The historical discussion on pp.~412--414 is unusually candid about gaps. It identifies the original 1905 argument for Lemma~II as incomplete: the special members through a fixed bicanonical group could all coincide, counted with multiplicity. It then attributes the torsion-root strategy to Enriques's 1930 note and its exposition to Enriques--Campedelli (1934). It also acknowledges the rational-divisor selection problem called Hypothesis~II. The book itself therefore recognizes both defects.

The proposed covariant constructions do not, as printed, resolve every restriction-map, degeneration, and descent issue identified in Sections~\ref{cp2-x-6}--\ref{cp2-x-8}. Tsen's theorem repairs the function-field representability issue; isotriviality and the product quotient give a complete proof of the geometric conclusion when the Albanese fibres have genus greater than one.

The historical paragraph also presents the exclusion of multiple Albanese fibres in the genus-one case as a refinement based on Chisini's analysis of crossing branch curves, and attributes it to Enriques--Campedelli (1934, \S~12). Example~\ref{cp2-x-counterexample} shows that this refinement is false in the stated generality. A historical claim that a proof had been completed cannot override the counterexample. The valid ruledness criterion and the valid translation-surface constructions must be separated from that false exhaustion assertion.

Finally, the book attributes the noncyclic examples to Bagnera--De Franchis (1907) and the general two-radical construction to Chisini (1921). For a modern presentation of the seven bielliptic families, see Nuer (2025, \S~2.1, Table~2 and Theorem~2.2).

\section{Superficie con curve pluricanoniche d'ordine zero: il modulo generale}
\label{cp2-x-16}

\subsection{Numerical triviality and the two ordinary types, pp.~414--417}

The final classification concerns the minimal surfaces with $p_g=0$, $q=1$, and $K_X\equiv0$, now called bielliptic surfaces. It does not classify all minimal surfaces with $p_g=0,q=1,K_X^2=0$: Example~\ref{cp2-x-counterexample} has these latter invariants and $\kappa=1$. In the translation class, the distinction is the genus of the transverse curve in the product quotient. The phrase ``canonical curve of order zero'' must be read as numerical triviality on the minimal surface, with a nontrivial torsion canonical line bundle, not as $K_X\sim0$.

Adjunction gives $D^2=2p_a(D)-2$ for every divisor when $K_X\equiv0$. In the Kodaira-dimension-one class it gives $D^2\leq2p_a(D)-2$, with strict inequality only when $K_X\cdot D>0$. The blanket strict inequality in the parenthesis on p.~415 is too broad if it includes vertical elliptic fibres, for which both sides are zero. It is correct for a moving divisor transverse to the canonical elliptic fibration. Likewise, ``order $2\pi-2$'' on p.~414 means intersection with the selected genus-$\pi$ fibre, not a degree independent of the chosen polarization.

Write a bielliptic surface as
\[
 X=(E_0\times E_1)/G,
\]
with $G$ acting on $E_0$ by translations and $E_1/G\cong\mathbb P^1$. The group is abelian of rank at most two because it embeds into the torsion of $E_0$. Every automorphism of $E_1$ has the affine form $x\mapsto\lambda x+a$. A translation by $b$ commutes with it precisely when $(\lambda-1)b=0$.

For general modulus the only nontrivial multiplier is $\lambda=-1$, and commuting translations lie in $E_1[2]$. A subgroup with rational quotient must contain such an involution. The rank-at-most-two condition leaves
\[
 G\cong\mathbb Z/2\quad\hbox{or}\quad(\mathbb Z/2)^2.
\]
Without that rank condition the automorphism group of an elliptic curve does contain $(\mathbb Z/2)^3$, generated by all two-torsion translations and inversion. Thus the restriction comes from the surface's translation factor, not from the elliptic curve alone. This makes explicit a hypothesis suppressed in the short group argument on p.~415.

Type $I_a$ has a degree-two map $E_1\to\mathbb P^1$ with four branch points and is represented by
\[
 u^2=(z-a_1)(z-a_2)(z-a_3)(z-a_4)\varphi(x,y),
 \qquad f(x,y)=0.
\]
The printed radical includes $\varphi$. The contact divisor of the quadric $\varphi$ must define nontrivial two-torsion on the base cubic. Since the order is prime, the book's noncollinearity condition is sufficient here. The diagonal quotient construction makes the total surface smooth after normalization and gives determinant two.

For type $I_b$, let
\[
 A=\varphi(x,y)(z-a)(z-b),\qquad
 B=\psi(x,y)(z-d)(z-e),\qquad u=\sqrt A+\sqrt B.
\]
The four branch values have inertia order two, and the degree-four cover has total ramification eight. Riemann--Hurwitz gives genus one. The two quadratic torsion classes on the base cubic must be nonzero and distinct; over $\mathbb F_2$ this is exactly independence. Hence the stated contact condition is sufficient in this case. The printed $(x-b)$ on p.~416 must be $(z-b)$, as the scan and the subsequent formulas show.

The two involutions are correctly expressed by
\[
 u\longmapsto\frac{A-B}{u},\qquad
 u\longmapsto\frac{B-A}{u}.
\]
Their product is $u\mapsto-u$, the fixed-point-free translation involution. The displayed identity on p.~417 verifies the first map directly.

The asserted plane-quartic singularity description, however, does not follow from the general displayed equation. Eliminating radicals gives
\[
 u^4-2(A+B)u^2+(A-B)^2=0.
\]
Its coefficient of $z^4$ is $(\varphi-\psi)^2$, so it is generally of degree four in $z$, not degree two. For general coefficients the two nodes occur at $u=0$ and the two roots of $A-B$, rather than as a double point and an infinitely near double point at infinity. Special coordinate and coefficient choices can change that plane model, but they are not imposed in the source. The normalization still has genus one by the already correct ramification computation. This is an error in the auxiliary plane-model description, not in the existence of type $I_b$.

\section{Segue: caso armonico}
\label{cp2-x-17}

\subsection{The groups and the cyclic equation, pp.~417--419}

For harmonic modulus, $j(E_1)=1728$ and the multipliers are $1,-1,i,-i$. The composition table on p.~417 is a table of these four multiplier species, with subscripts recording changing translations. It is verified uniformly by
\[
 (a,\lambda)(b,\mu)=(a+\lambda b,\lambda\mu).
\]
A group involving multiplier $i$ contains a cyclic order-four subgroup. The translations commuting with it form $\ker(i-1)$, which has order two. Thus the two new groups are
\[
 \mathbb Z/4,\qquad\mathbb Z/4\oplus\mathbb Z/2.
\]
They yield determinants four and eight, respectively. The earlier ordinary groups remain possible on a harmonic elliptic curve; the cases here are the additional ones.

The order-four quotient has inertia orders $(4,4,2)$ and the cyclic equation
\[
 u^4=(z-a)(z-b)(z-d)^2\varphi(x,y)
\]
is correct. The exponents sum to four, so there is no additional branching at infinity. Riemann--Hurwitz gives genus one. The torsion class attached to $\varphi$ must have exact order four; noncollinearity alone does not suffice for this composite order.

The inverse-monodromy equation with exponents $(3,3,2)$ also gives a harmonic elliptic curve. The statement on p.~418 that it is birationally distinct merely because the substitutions are inverted is not justified. For constant coefficients, both curves are elliptic with $j=1728$ and are isomorphic over $\mathbb C$. More explicitly, replacing $u$ by $(z-a)(z-b)(z-d)/u$ changes the first exponent vector to the second and replaces the coefficient by its inverse. In a family, this also changes the torsion marking; distinctions between marked covers or twists over a function field must be stated separately from birational distinctions between unmarked complex surfaces. Inversion of the branch generators alone does not establish the latter.

\subsection{The correct degree-eight equation and the two incorrect variants, pp.~419--421}

The first formula on p.~419 is correct:
\begin{equation}
\label{cp2-x-correct-iib}
 u=\sqrt{(z-b)(z-d)\varphi(x,y)}
   +\sqrt[4]{(z-a)(z-b)^3\psi(x,y)},
 \qquad f(x,y)=0.
\end{equation}
In $\mathbb Z/2\oplus\mathbb Z/4$, its inertia elements at $a,b,d$ are
\[
 (0,1),\quad(1,3),\quad(1,0).
\]
They generate the group, have orders $4,4,2$, and sum to zero. These are exactly the three permutations written out on p.~419. Their product is the identity, as stated on p.~420, and the cover is unramified at infinity. Its genus is one.

The ``similar'' expression printed later on p.~420 is
\[
 u=\sqrt{(z-a)(z-b)^3\varphi}
   +\sqrt[4]{(z-a)(z-d)\psi}.
\]
It is not another genus-one representative. The inertia orders at $a,b,d,\infty$ are $(4,2,4,2)$. The fourth-root radicand has degree two, and therefore creates the omitted order-two inertia at infinity. The degree-eight normalization has
\[
 2g-2=8\left(-2+\tfrac34+\tfrac12+\tfrac34+\tfrac12\right)=4,
\]
so $g=3$. The projective tangent and tacnode description in the accompanying footnote cannot serve as a description of the asserted three-branch elliptic cover.

The final displayed type $II_b$ on p.~421 is likewise wrong as printed:
\[
 u=\sqrt{(z-a)(z-b)^3\varphi(x,y)}
   +\sqrt[4]{(z-b)(z-d)\psi(x,y)}.
\]
Its inertia orders are $(2,4,4,2)$, again including infinity, and again its normalized curve has genus three. The printed radical indices are two and four, with the displayed exponents; those values cause the discrepancy. Equation~\eqref{cp2-x-correct-iib}, already present two pages earlier, is a valid replacement. With independent elliptic torsion data, the erroneous final equation would instead construct a surface of Kodaira dimension one in the translation class, with four multiple fibres and $P_{12}=7$, rather than a bielliptic surface with $P_{12}=1$.

For the corrected equation, the base-cubic data are a class $\eta$ of exact order two and a class $\theta$ of exact order four, with
\[
 \eta\ne2\theta.
\]
This ensures that the two Kummer extensions are independent. The source's final comparison of the contact sextuples expresses this last inequality correctly. Its preceding noncollinearity test for the order-four root still needs the exact-order strengthening.

The rationality of the deck transformations also needs a different proof from the one printed on p.~420. Two polynomial relations and their resultant show that $u'$ is algebraic over $\mathbb C(z,u)$, not that it is rational in $z,u$. For the independent radicals in \eqref{cp2-x-correct-iib}, the compositum is a Galois field of degree eight. The eight sums $\pm A+i^jB$ are distinct for general $z$, so $u=A+B$ is a primitive element. Hence every deck conjugate belongs to $\mathbb C(z,u)$ and is a rational function of $z,u$. This supplies exactly the conclusion needed, with its field-theoretic reason.

\section{Caso equianarmonico}
\label{cp2-x-18}

\subsection{Classification and multiplier table, pp.~421--423}

For equianharmonic modulus, $j(E_1)=0$ and the multiplier group is $\mu_6$. If a group contains a multiplier of order three, commuting translations lie in $\ker(\zeta_3-1)$, a group of order three. If it contains a primitive sixth root, that kernel is trivial. The three new possibilities are therefore
\[
 G\cong\mathbb Z/3,\qquad\mathbb Z/6,
 \qquad(\mathbb Z/3)^2.
\]
This proves the classification sketched by lifting projectivities of the three branch points on p.~422. The rank restriction from the other elliptic factor remains understood.

The five-row composition array on p.~422 is again verified by multiplying the six linear multipliers. One entry is a genuine misprint: the fourth row repeats $\pi\omega=\pi$ after already giving $\pi\omega=\omega_1$. It must be $\tau\omega=\pi$. All other displayed species products agree with the affine composition rule. The small stop after the third-row final $\tau$ is punctuation, not another group-theoretic symbol.

Type $III_a$ is represented by
\[
 u^3=(z-a)(z-b)(z-d)\varphi(x,y),\qquad f(x,y)=0.
\]
Its inertia orders are $(3,3,3)$, and the base-cubic torsion datum must be nonzero of order three. Type $III_b$ is
\[
 u^6=(z-a)(z-b)^2(z-d)^3\varphi(x,y),\qquad f(x,y)=0,
\]
with inertia orders $(6,3,2)$ and a base-cubic torsion class of exact order six. Both exponent sums are divisible by the radical index, so infinity is unramified, and both normalized transverse curves have genus one. For type $III_b$, the source's condition of noncollinearity must again be strengthened: a class of order two or three is nonzero but does not produce a connected degree-six elliptic cover.

\subsection{The two cubic radicals, pp.~423--424}

The degree-nine equation is correct:
\[
 u=\sqrt[3]{(z-a)^2(z-d)\varphi(x,y)}
   +\sqrt[3]{(z-b)(z-d)^2\psi(x,y)},
 \qquad f(x,y)=0.
\]
The inertia elements in $(\mathbb Z/3)^2$ are $(2,0)$, $(0,1)$, and $(1,2)$. They generate the group and sum to zero. The three products of three-cycles displayed on p.~424 give exactly these translations on the two indices. Consequently there is no inertia at infinity, the total ramification is eighteen, and Riemann--Hurwitz gives genus one. The existence of order-three automorphisms fixing points makes the elliptic curve equianharmonic.

On the base cubic, the two nontrivial $3$-torsion classes must be linearly independent over $\mathbb F_3$. The printed requirement that their contact triples be nonequivalent excludes equality but does not exclude negatives. If the second class is the negative of the first, the compositum has degree three rather than nine. This is a real missing connectedness condition in the stated normal form. With independence imposed, the primitive-element argument from Section~\ref{cp2-x-17} applies to the nine sums and verifies that the branch permutations extend to birational deck transformations. No resultant argument is required.

\section{Riassunto}
\label{cp2-x-19}

\subsection{The seven families and the corrected plurigenus table, pp.~424--425}

The seven-family classification is correct after repairing the normal forms and connectedness conditions. Every type exists: embed the relevant group $G$ into the torsion translations of an elliptic curve $E_0$, let it act on $E_1$ with the specified multiplier and translation subgroup, and take the free diagonal quotient. Conversely, the commuting-translation calculation in Sections~\ref{cp2-x-16}--\ref{cp2-x-18} gives exactly these seven groups. Nuer (2025, \S~2.1, Table~2 and Theorem~2.2/Table~3) provides a modern presentation of the families, fibre multiplicities, and canonical orders.

\begin{center}
\begin{tabular}{c c c c c}
Type & $G$ & $n$ & Multiple-fibre orders & $\operatorname{ord}(K_X)$\\
\hline
$I_a$ & $\mathbb Z/2$ & $2$ & $(2,2,2,2)$ & $2$\\
$I_b$ & $(\mathbb Z/2)^2$ & $4$ & $(2,2,2,2)$ & $2$\\
$II_a$ & $\mathbb Z/4$ & $4$ & $(2,4,4)$ & $4$\\
$II_b$ & $\mathbb Z/2\oplus\mathbb Z/4$ & $8$ & $(2,4,4)$ & $4$\\
$III_a$ & $\mathbb Z/3$ & $3$ & $(3,3,3)$ & $3$\\
$III_b$ & $\mathbb Z/6$ & $6$ & $(2,3,6)$ & $6$\\
$III_c$ & $(\mathbb Z/3)^2$ & $9$ & $(3,3,3)$ & $3$
\end{tabular}
\end{center}

The modulus restrictions concern $E_1$, namely the book's curves $K$ on which the nontranslation automorphisms act. The summary on pp.~424--425 instead writes $C$ when calling the curves general, harmonic, or equianharmonic. That is a swapped fibre label. The orbit elliptic curve $C$ can have general modulus even in the harmonic and equianharmonic cases.

Let $r=\operatorname{ord}(K_X)$. A torsion line bundle on a projective variety has a nonzero section only if it is trivial: an effective divisor numerically equivalent to zero has zero intersection with an ample divisor, hence is zero. Therefore
\begin{equation}
\label{cp2-x-periodic-plurigenera}
 P_m=\begin{cases}1,&r\mid m,\\0,&r\nmid m.\end{cases}
\end{equation}
This gives the complete table, rather than only selected values:
\[
\begin{array}{c|ccccc}
\text{type}&P_2&P_3&P_4&P_6&P_{12}\\\hline
I_a,I_b&1&0&1&1&1\\
II_a,II_b&0&0&1&0&1\\
III_a&0&1&0&1&1\\
III_b&0&0&0&1&1\\
III_c&0&1&0&1&1
\end{array}
\]
The first two printed rows are correct. The equianharmonic rows are not: $P_4=1$ should be $P_4=0$ in all three equianharmonic types, and $III_c$ must have $P_3=1$, not the $P_3=0$ printed when it is grouped with $III_b$. The scans confirm all these entries. They also disagree with the book's own floor formula, so no change of modern terminology can reconcile them.

The equality $P_{12}=1$ and the numerical condition $K_X\equiv0$ are valid characterizations within the minimal $p_g=0,q=1$ setting explained in Section~\ref{cp2-x-14}. The bare existence of an elliptic Albanese fibration is not such a characterization, because multiple fibres with nontranslation monodromy give the counterexample of Section~\ref{cp2-x-11}.

\subsection{The canonical cover and the final adjoint equation, pp.~425--426}

The historical ``hyperelliptic'' or ``abelian proper'' covering surface is an abelian surface. Its modern construction is the canonical cyclic cover
\[
 \widetilde X=\operatorname{Spec}_X
 \left(\bigoplus_{j=0}^{r-1}\omega_X^{-j}\right),
\]
after choosing a trivialization of $\omega_X^r$. Exact order $r$ makes it connected, and characteristic zero makes it étale. In the product presentation, it is the quotient of $E_0\times E_1$ by the kernel of the multiplier character on the holomorphic two-form. That kernel acts by translations on both factors, so its quotient is an abelian surface. Hence
\[
 p_g(\widetilde X)=1,\qquad q(\widetilde X)=2,
 \qquad\chi(\mathcal O_{\widetilde X})=0,
 \qquad P_m(\widetilde X)=1\quad(m\geq1).
\]
This proves the assertion about the covering surface directly and clarifies that its canonical triviality is linear, whereas the original surface has a nontrivial torsion canonical class.

For type $I_a$, $r=2$. The book describes the cover by extracting a square root of a biadjoint expression on a model $F_m\subset\mathbb P^3$. The biadjoint degree should be $2m-8$, twice the adjoint degree $m-4$, with the usual adjoint conditions along the singular locus of the model. The text prints $\varphi_{2n-8}$ twice although the surface degree is denoted $m$; this is an index slip in the book. One must take the normalization and interpret the equation on the smooth model: the apparent branch divisor has even multiplicities and the resulting canonical cover is étale. Merely adjoining a square root of an arbitrary polynomial of that degree would not establish this.

The closing historical notes attribute the seven families to Bagnera--De Franchis and independently to Enriques--Severi, and the algebraic presentation here to Enriques's 1934 note using Chisini's work. The chapter's lasting results are the ruledness criterion and the bielliptic classification; their modern proofs require the explicit corrections above, and the asserted exhaustion of all $p_g=0,q=1$ surfaces by the translation class must be discarded.

\chapter{Classificazione generale delle superficie}
\label{cp2-ch-xi}

Chapter XI (pp.~427--464) assembles the numerical classification, the geometry of continuous transformation groups, and the distinction between canonical divisors of positive and zero order. Its four principal classes survive as the four Kodaira dimensions. A substantial subsidiary identification does not: a surface with \(p_a=-1\), equivalently \(\chi(\mathcal O_X)=0\), need not have a positive-dimensional group of birational transformations. The missing case is an isotrivial elliptic surface whose finite monodromy acts nontrivially on the elliptic differential. This distinction, already necessary in Chapter X, must be retained throughout the present chapter. Here ``abelian surface'' translates Enriques's \emph{superficie iperellittica}; his \emph{superficie ellittica}, when defined by an elliptic transformation group, is a narrower notion than a modern surface admitting a genus-one fibration.

\section{Introduzione}
\label{cp2-xi-1}

The opening quotation from Castelnuovo's Bologna lecture of 1928 describes a method of constructing examples, testing conjectures, and only then seeking proofs (pp.~427--428). It supplies historical evidence about that method, not a proof of the classifications subsequently announced. Enriques's comparison between ordering surfaces by degree and ordering them by sectional genus is likewise an observation about the choice of parameter problem. There is no single parameter space in which the two informal uses of ``more general'' are being compared.

For a smooth member \(H\) of a base-point-free system on a specified smooth model, adjunction, already established in Part~1, gives
\[
2p_a(H)-2=H^2+K_XH.
\]
Thus \(n>2\pi-2\), \(n=2\pi-2\), and \(n<2\pi-2\) express respectively \(K_XH<0\), \(K_XH=0\), and \(K_XH>0\), when \(n=H^2\) and \(\pi=p_a(H)\). The distinction concerns the moving system after its base points have been resolved. For ample \(H\) on a minimal surface with \(K_X\) nef, equality forces \(K_X\equiv0\) by the Hodge index theorem. Equality for an arbitrary fibre of a pencil does not force this: an elliptic fibre has \(H^2=K_XH=0\) on many surfaces of Kodaira dimension one. Consequently the sectional descriptions in the classification must concern sufficiently positive systems, not every base-point-free divisor indiscriminately.

The four families announced on p.~429 become
\[
\kappa(X)=-\infty,\quad0,\quad1,\quad2.
\]
The first consists of rational and ruled surfaces; the second has torsion canonical bundle on its minimal model; the third has a pluricanonical genus-one fibration; the fourth has a canonical model with ample canonical divisor and at worst rational double points. These are announcements at this point, justified by the preceding chapters together with the qualifications below. The phrase about infinitely many exceptional curves must not be interpreted as saying that every minimal ruled surface contains infinitely many \((-1)\)-curves: \(\mathbb P^2\) and \(\mathbb P^1\times\mathbb P^1\) contain none. The older language also concerns rational moving curves and their behaviour under birational changes. Part~1's distinction between a contractible exceptional curve and a fibre of square zero remains essential.

\section[Trasformazioni e rigate]{Superficie che ammettono una serie continua di trasformazioni birazionali in se stesse: casi che conducono alle rigate}
\label{cp2-xi-2}

\subsection{The four criteria and the meaning of a continuous group}

On pp.~429--430 Enriques states four ruledness criteria: an effective rational one-dimensional transformation group; a connected group acting nontrivially on the members of a positive-dimensional linear system; a group of dimension at least two with one-dimensional general orbits, or of dimension at least three with two-dimensional general orbits; and families whose iterated products have unbounded dimension. The footnote to the second criterion explicitly interprets ``transitively'' as permuting the curves in infinitely many ways. It does not require a transitive action on the entire projective parameter space.

These criteria have clean modern proofs. Work birationally on the minimal model when \(X\) is not ruled. Birational self-maps of this model are regular. Its connected automorphism group is an abelian variety: by Chevalley's theorem a nonproper connected automorphism group would contain \(\mathbb G_a\) or \(\mathbb G_m\), whose general orbit closures give a uniruling. This is Brion's Proposition~7.1.4 in \emph{Some structure theorems for algebraic groups}; the representability and tangent-space statements are in Theorem~7.1.1. For complex surfaces, uniruledness implies ruledness. A connected affine one-dimensional group is therefore impossible on a nonruled surface, proving the first criterion.

For the second criterion a connected group fixing a line-bundle class acts projectively on its sections. An abelian variety has no nonconstant algebraic homomorphism to the affine group \(\operatorname{PGL}\), so this action is trivial. Passing to the graph of the linear system gives the same conclusion for a rationally defined system. The finite stabilizer of an ample line bundle on an abelian variety is a particularly transparent instance; it must not be confused with the connected group preserving its numerical class. The distinction is the one already made in Chapter IX between linear equivalence and numerical equivalence.

For the dimension criterion, an effective action of an abelian variety on a smooth projective variety has finite stabilizers after passing to the induced Albanese action. Equivalently, the homomorphism from the acting abelian variety to \(\operatorname{Alb}(X)\) has finite kernel; a positive-dimensional kernel fixing a point would act trivially. Therefore the dimension of an effective connected group is at most the dimension of its general orbit. On a nonruled surface the possibilities are dimension one with elliptic orbits, and dimension two with an open orbit. This proves exactly the exclusions in criterion~3.

Finally a connected algebraic family through the identity lies in the finite-dimensional group \(\operatorname{Aut}^0(X)\); its product families have bounded dimension. This proves criterion~4 in the sense actually used on p.~433. One must retain that meaning. An automorphism group can have infinitely many discrete components, and translating a connected family into a nonidentity component may generate infinitely many components without unbounded dimension. Such a phenomenon does not meet the dimension-growth hypothesis of Enriques's proof.

\subsection{What the original arguments supply}

The first proof imports the Noether--Enriques ruledness theorem, citing Conforto, Book~II, Chapter~I, \S\S4--6, and attributes an earlier group-action argument to Painlev\'e. The second imports the finiteness theorem for automorphisms of a curve of genus at least two, called Schwarz's theorem, and the fact that an elliptic curve has only finitely many automorphisms preserving a positive-degree line bundle. Both facts are adequate for the one-dimensional orbit argument. The book's further reduction of a general linear system to an invariant pencil (p.~432) is a sketch: a centralizer need not fix each point fixed by a chosen element, and the existence of the required nontrivially acted-on invariant line is not supplied by the sentence about centralizers. The projective-representation argument above replaces that reduction.

The stabilizer and infinitely-near-point arguments on pp.~432--433 similarly require independence of imposed point conditions, algebraicity of the stabilizers, and eventual detection of a nontrivial transformation by a finite jet. A formal automorphism acting identically at every order is indeed the identity, but the claimed one-dimensional stabilizer does not follow merely by counting the chosen points. These are gaps in the local proof of a correct criterion, not counterexamples to it. The preceding abelian-group argument supplies the missing structure.

Picard's observation on p.~431 is stated with an insufficient justification: the transformations of one fixed degree are parametrized by algebraic conditions, but this set is generally not a group, since composition can increase degree. On a fixed projective minimal nonruled model the automorphism scheme and its identity component give the required algebraic group. The argument cannot be transplanted unchanged to arbitrary bounded-degree pieces of the Cremona group.

The numerical calculation in the fourth proof is valid after resolving the base cluster. If the birational transforms of a general curve have fixed moving square \(n\), fixed geometric genus \(\pi\), and multiplicities \(h_i\), then
\[
N_r=n+\sum_i h_i^2,\qquad
\Pi_r=\pi+\sum_i\frac{h_i(h_i-1)}2,
\]
and hence
\[
N_r-2\Pi_r=n-2\pi+\sum_i h_i.
\]
If the image curves form a genuinely \(r\)-dimensional family, the characteristic-system estimate \(h^0(N_{C/X})\ge r\), together with \(h^0(L)\le\deg L+1\) for an effective line bundle on an integral curve, gives \(N_r\ge r-1\). The source does not justify that the chosen curve has finite stabilizer in every product family; this is the missing premise of its dimension growth. Conditional on that premise, its inequalities
\[
\sum h_i^2>(2\pi-n)^2,\qquad
\sum h_i\ge\sqrt{\sum h_i^2}>2\pi-n
\]
do imply \(N_r>2\Pi_r-2\), even with more strength than needed. Chapter X's negative-canonical-intersection criterion then applies. The numerical identities themselves are not errors.

\section{Superficie ellittiche e iperellittiche}
\label{cp2-xi-3}

\subsection{Surfaces which actually possess the group}

The forward classification on pp.~433--438 starts with an existing positive-dimensional group. Subject to this hypothesis, its central conclusion is correct. A two-dimensional connected group acting effectively on a nonruled minimal surface is an abelian surface acting by translations, and the surface itself is an abelian surface after choosing an origin. Its open orbit is already closed because the group is proper. A stabilizer is finite; quotienting the group by that stabilizer yields the effective translation group. In an effective transitive commutative action a stabilizer fixes every point and is therefore trivial. This proves simple transitivity more directly than the auxiliary parameter surface of p.~436.

The calculation of invariants is immediate from translation-invariant differential forms:
\[
\Omega_A^1\simeq\mathcal O_A^{\oplus2},\quad
K_A\sim0,\quad q(A)=2,\quad p_g(A)=1,
\quad p_a(A)=-1,\quad P_m(A)=1.
\]
These are geometric facts about the minimal abelian surface. Blowing up a point introduces a positive exceptional canonical divisor, without changing \(p_g,q\), or any plurigenus. Thus ``canonical curve of order zero'' always requires the pure, minimal-model convention of Part~1.

For a one-dimensional elliptic group acting on a nonruled surface of Kodaira dimension one, the modern description is a free diagonal quotient
\[
X=(C\times E)/G,
\]
where \(g(C)\ge2\), \(E\) is elliptic, and the finite group \(G\) acts on \(E\) by translations. Let \(B=C/G\) have genus \(b\). The two projections give the elliptic-orbit fibration over \(B\) and a fibration over the elliptic curve \(E/G\). Cohomology of the finite \'etale quotient gives
\[
q=b+1,\qquad p_g=b,\qquad\chi(\mathcal O_X)=0,
\qquad K_X^2=0.
\]
The cases of Kodaira dimension zero with a one-dimensional group are the bielliptic surfaces, already treated in Chapter X; abelian surfaces can also contain one-dimensional elliptic translation subgroups. This proves the forward assertions of p.~434 without using their later false converse. The quotient description and its translation/nontranslation distinction are explicitly recorded in Catanese, Frapporti, Glei\ss ner, Liu, and Sch\"utt (2025, Lemma~1.1).

The source's computation \(\Delta=0\), \(12p_a-p^{(1)}+13=0\) is consistent with these invariants. To conclude the two separate equalities \(p_a=-1\) and \(p^{(1)}=1\), however, one also needs the prior nonruled inequalities or the quotient structure; a single equation alone does not determine two integers. If a transverse curve \(T\) maps to \(B\) with degree \(n\) and the orbit fibration has multiple fibres of orders \(s_i\), Riemann--Hurwitz gives
\[
g(T)=1+n(b-1)+\frac n2\sum_i\left(1-\frac1{s_i}\right).
\]
This is the formula on p.~435. It presupposes the described ramification and no additional branch points. The letters \(C\) and \(K\), used consistently for the two fibrations on p.~434, are interchanged in the explanatory list on p.~435; the invariant interpretation removes that ambiguity. The determinant and the multiplicities do not determine the complex structure or the monodromy representation by themselves. The strict comparison \(\pi>\varrho\) later printed on p.~439 also needs qualification: Riemann--Hurwitz gives only \(\pi\ge\varrho\) in the stated range. The determinant-one product \(B\times E\) has transverse fibres isomorphic to \(B\), so equality occurs.

\subsection{Irregularity of Jacobians and abelian varieties}

The assertion \(q(A)=\dim A\) on pp.~437--438 is correct. The successive restriction argument, extending the third equivalence criterion to arbitrary dimension, is only sketched. It needs an actual injectivity theorem for restriction on \(\operatorname{Pic}^0\), rather than the unsupported claim that every pair of inequivalent divisors stays inequivalent on each special coordinate subvariety. Part~1, ``Criteri d'equivalenza'', and Chapter IX's treatment of Picard varieties and equivalence give the earlier dependencies. For a complex torus \(A=\mathbb C^g/\Lambda\), invariant \((0,1)\)-forms give \(h^1(\mathcal O_A)=g\); equivalently its dual abelian variety has dimension \(g\). This supplies a complete modern calculation without the restriction induction.

The intermediate claim about the symmetric square is also correct: \(q(C^{(2)})=g(C)\), as follows by taking the symmetric invariant part of the one-forms on \(C\times C\). For a divisor on an abelian subvariety, the source's lower bound from its translation orbit requires a sufficiently ample divisor with finite stabilizer. An arbitrary divisor can have a positive-dimensional translation stabilizer. The claimed increase from \(r\) to \(r+p-q\) independent divisor classes is not proved merely by sweeping a divisor with the complementary group; the dual product isogeny supplies the missing Picard-variety statement.

Poincar\'e reducibility supplies a complement to an abelian subvariety only up to isogeny. After a finite \'etale cover the product calculation is valid, and \(\chi(\mathcal O)\) multiplies by the degree. Thus the printed formula
\[
24(p'_a+1)=24n(p_a+1)
\]
has the expected consequence \(p'_a=-1\) when \(p_a=-1\). The factor 24 is harmless. The references on p.~438 are not harmless locators: ``Cap.~X, \S9'' should refer to Chapter IX's abelian-variety discussion, and ``Cap.~V; \S28'' to Chapter V, \S8. The p.~437 dimensions of subgroup and ambient group are also reversed in print; they must be \(q\) and \(p\), respectively. These are source misprints, not transcription mistakes.

\section[Caratterizzazione mediante i generi]{Caratterizzazione delle superficie ellittiche e iperellittiche mediante i valori dei generi}
\label{cp2-xi-4}

\subsection{The missing monodromy and an explicit counterexample}

The attempted converse on pp.~438--445 is false as stated. It does not follow from \(p_a=-1\) and \(p^{(1)}=1\) that the surface admits an elliptic transformation group. The proof on p.~439 considers two possible genera for the base of the canonical elliptic fibration, namely \(b=q=p_g+1\) and \(b=q-1=p_g\), and excludes the first by omitting a degree-zero factor in the canonical bundle. That first case really occurs.

\begin{example}[A smooth isotrivial fibration with \(p_g=1\), \(q=2\), and no continuous automorphisms]
\label{cp2-xi-monodromy-example}
Let \(B\) be a smooth curve of genus two and choose a nontrivial two-torsion line bundle \(\eta\in\operatorname{Pic}^0(B)\). The associated connected \'etale double cover \(C\to B\) has genus three and a fixed-point-free deck involution \(\sigma\). Let \(E\) be any elliptic curve, and set
\[
X=(C\times E)/\langle(\sigma,[-1])\rangle.
\]
The diagonal action is free because \(\sigma\) is free, so \(X\) is smooth and the quotient map is \'etale of degree two. Since \(K_{C\times E}=\operatorname{pr}_C^*K_C\), one has \(K_X^2=0\), \(K_X\) nef, and \(\kappa(X)=1\). Invariant one-forms come only from \(B\), so \(q=2\). The anti-invariant part of \(H^0(C,K_C)\) is \(H^0(B,K_B\otimes\eta)\), of dimension one; tensoring with the anti-invariant differential on \(E\) gives \(p_g=1\). Consequently \(\chi=0\) and \(p_a=-1\).

The fibration \(f:X\to B\) is smooth, with every fibre isomorphic to \(E\), but
\[
\omega_X\simeq f^*(\omega_B\otimes\eta),\qquad
P_m=h^0(B,\omega_B^{\otimes m}\otimes\eta^{\otimes m}).
\]
Thus \(P_1=1\) and \(P_m=2m-1\) for \(m\ge2\); in particular \(P_4=7\). Finally,
\[
H^0(X,T_X)=H^0(C\times E,T_{C\times E})^{\langle(\sigma,[-1])\rangle}=0,
\]
because \(H^0(C,T_C)=0\) and \([-1]\) negates the one-dimensional space of vector fields on \(E\). In characteristic zero the identity component of the automorphism group is therefore trivial. Since \(X\) is minimal and nonruled, this also excludes a continuous birational transformation group.
\end{example}

This example directly contradicts the p.~442 assertion that a surface with \(q=2\) and a pencil of base genus two must have \(p_g\ge2\). It also contradicts the claimed global equivalences on pp.~443, 460--461, and the \(p_a=-1\) row on p.~464. Taking \(B\) of any genus \(b\ge2\) gives \(q=b\), \(p_g=b-1\), and the same absence of continuous automorphisms, so the issue includes the \(p_g>1\) case of p.~438. The general distinction between translation and nontranslation monodromy is the one in Catanese, Frapporti, Glei\ss ner, Liu, and Sch\"utt (2025, Lemma~1.1).

For a relatively minimal genus-one fibration in characteristic zero, put \(\mathcal L=(R^1f_*\mathcal O_X)^\vee\). The missing term is visible in the canonical bundle formula
\[
\omega_X\simeq f^*(\omega_B\otimes\mathcal L)
 \otimes\mathcal O_X\left(\sum_i(m_i-1)F_i\right),
\qquad f^{-1}(b_i)=m_iF_i,
\]
where \(\deg\mathcal L=\chi(\mathcal O_X)\). Bombieri and Mumford (1977, \S1, Theorem~2, pp.~27--29) give this formula for an arbitrary nonsingular base. Their line bundle denoted \(L\) is the dual of \(\mathcal L\) used here; in characteristic zero their torsion term vanishes and all multiple fibres have coefficient \(m_i-1\). For a quasi-bundle \(\mathcal L\) is torsion. Trivial \(\mathcal L\) gives \(q=b+1\), \(p_g=b\); nontrivial \(\mathcal L\) gives \(q=b\), \(p_g=b-1\). The vanishing of Euler defects forces constant \(j\), but does not trivialize \(\mathcal L\). This is exactly the logical failure on p.~439, not a failure of the isotriviality conclusion.

\subsection{The Albanese lemma and the subsequent branch calculation}

The lemma on pp.~440--441 remains correct with its stated hypothesis: if \(q=2\) and there is no pencil with base genus two, the Albanese morphism has two-dimensional image, hence is generically finite onto an abelian surface. Indeed an Albanese image of dimension one has, after Stein factorization, genus exactly \(q\); this is Chapter IX's Picard--Albanese result in modern form. It is preferable to Enriques's incidence construction, whose restriction and equivalence steps require the precise earlier criteria. The source's \(L\)-curves lie on \(\Phi\), not on \(F\), and its \(L_{I'}\) on p.~441 is a printed slip for \(L_{P'}\).

The next step, asserting that \(p_a=-1,p_g=1\) excludes that genus-two pencil, is the false one. In the displayed Euler formula on p.~442, the coefficient \(8(\pi-1)\) is also incorrect: a fibration over a base of genus two contributes \(4(\pi-1)\). The corrected equality is
\[
e(X)=\Delta+4(\pi-1)=12\chi(\mathcal O_X)-K_X^2.
\]
For \(\chi=0\) it still forces \(K_X^2=\Delta=0\) and \(\pi=1\), using nefness on the nonruled minimal model. It does not force trivial elliptic monodromy, and hence gives no lower bound \(p_g\ge2\).

When the Albanese map really is generically finite, the ramification identity \(K_X=a^*K_A+R=R\) explains why absence of ramification leads to an \'etale cover of an abelian surface. Such a cover is again abelian, by the finite-cover result of Chapter IX. If the branch locus is a union of elliptic curves, lifting a connected translation group through a branched cover requires preservation of the covering data, not just of the reduced branch set. The source omits that lifting justification. A finite cover of the acting elliptic group can supply the lift in the appropriate equivariant situation; it does not remedy the missing Albanese-dimension-one surfaces.

For the translation case with base genus one and at least one multiple fibre, the plurigenus calculations on pp.~442--443 are sound. The canonical coefficients are \((m_i-1)F_i\). Squaring makes one whole fibre move for each multiple fibre, with residue \((m_i-2)F_i\). If there are \(r\ge2\) such fibres, the induced base line bundle has degree \(r\) and hence \(r\) sections. If there is just one, the fourth power has base degree at least two and therefore at least two sections. Thus \(P_4>1\). An abelian surface instead has \(P_4=1\). This distinction survives, although \(P_4>1\) does not imply a continuous group: the preceding example has \(P_4=7\).

\subsection{What remains true of \texorpdfstring{$\chi=0$}{χ=0}}

The numerical conclusion \(\chi=0\Rightarrow K_X^2=0\) for a minimal nonruled complex projective surface is correct. The source's p.~445 proof establishes it only in a restricted branch situation and acknowledges the missing case of arbitrary multiple branch components in a footnote. For an ordinary simply branched cover with reduced nodal--cuspidal branch curve of normalization genus \(\varrho>1\), its Severi inequality and characteristic-series estimate read
\[
24\chi(\mathcal O_X)\ge6(\varrho-1)-2\tau,
\qquad
2\delta+3\tau\le2(\varrho+\delta+\tau-1).
\]
The second gives \(\tau\le2(\varrho-1)\), so the first is strictly positive. This numerical deduction is valid under the branch and moving-family hypotheses. It is not a proof for an arbitrary generically finite map with nonreduced discriminant.

A modern replacement uses the classification and Noether's formula: minimal general-type surfaces have \(\chi>0\); the remaining nonruled cases with \(\chi=0\) are abelian surfaces, bielliptic surfaces, and properly elliptic quasi-bundles. The last class has both translation and nontranslation monodromy. The correct automorphism criterion within that class is that the finite group act by translations on the elliptic factor. Simple abelian surfaces have no elliptic subgroup; those containing an elliptic curve are isogenous to a product of two elliptic curves. The product decomposition is up to isogeny, not generally an isomorphism. These facts explain the valid discussion on p.~444 while correcting the surrounding overstatement. The forward summary on p.~443 also needs the bielliptic case: these surfaces have \(p_g=0,q=1\) and an elliptic automorphism group, so the printed restriction \(p_g\ge1\) omits cases already present in Chapter X.

The historical note on p.~446 attributes the analytic group classification to Picard and Painlev\'e and the geometric formulations to Enriques and Severi. It specifically acknowledges that Picard's omission of the canonical-divisor condition had been encouraged by an error of Noether. These are historical attributions, distinct from the modern monodromy correction required by the numerical converse.

\section[Superficie con genere geometrico e lineare uno]{Superficie coi generi \texorpdfstring{\(p_g=1\)}{pg = 1} e \texorpdfstring{\(p^{(1)}=1\)}{p(1) = 1}}
\label{cp2-xi-5}

\subsection{The trichotomy, with its overlaps made explicit}

On pp.~447--450 the three families are K3 surfaces, abelian surfaces, and surfaces with a genus-one fibration. They are not disjoint as written: some K3 and abelian surfaces have such fibrations. A disjoint birational classification under \(p_g=1\), \(K_X^2=0\) is obtained by separating \(\kappa=0\) from \(\kappa=1\). In the former case the minimal surface is K3, with \(q=0\), or abelian, with \(q=2\). In the latter case it is properly elliptic and \(q\in\{0,1,2\}\). For \(q=2\), both monodromy cases of the preceding section must be allowed.

The assertion \(p_a\ge-1\), equivalently \(q\le2\), is true. The lengthy construction in an ambient Picard variety on pp.~448--449 does not prove it: the chosen intermediate varieties need not be smooth, their incidence dimensions are assumed, and restriction of the claimed characteristic system need not have the asserted rank. In particular, the displayed calculation of a dimension at least three is conditional on these unproved incidence statements. A precise replacement is the Castelnuovo--de Franchis argument in Section~\ref{cp2-xi-8}. If \(q>2\) and \(p_g=1\), it produces a fibration over a base of genus at least two. Its fibres cannot be rational, since then \(p_g=0\). For a relatively minimal such fibration both \(K_X^2\) and its topological Euler number are nonnegative, contradicting
\[
12\chi(\mathcal O_X)=12(2-q)<0.
\]
The same argument validates the weaker assertion on p.~449 that a surface with \(p_g<3\) and \(q>2\) has an irrational pencil.

The printed three-line inequality at the bottom of p.~449 cannot be repaired by changing just one sign. For a fibration of fibre genus \(g\) over base genus \(b\), the correct formula is
\[
e(X)=\Delta+4(g-1)(b-1),\qquad
12(p_a+1)-p^{(1)}+1=e(X)\ge0
\]
when \(g,b\ge1\). Thus \(12(p_a+1)\ge p^{(1)}-1\ge0\). The source reverses the first comparison; in the middle expression it both reverses the inequality and omits the final \(+1\); and in the equality case it prints \(p_a=1\) instead of \(-1\). Thus all three printed lines require correction. The intended conclusions \(p_a\ge-1\) and \(p_a=-1\Rightarrow p^{(1)}=1\) do follow from the corrected formula.

The regular assertion imported from Chapter VII is that \(p_g=q+1=1\), \(P_2=1\), and \(K_X^2=0\) characterize the K3 case; positive canonical order instead gives a genus-one fibration and \(P_2>1\). The abelian criterion \(p_g=1,q=2,P_4=1\) also survives. Proposition~c), as printed across pp.~447--448, first assumes canonical order zero and then allows an elliptic alternative with \(P_4>1\). Those two conditions are incompatible: a trivial canonical bundle has every \(P_m=1\). The second alternative belongs to the positive-canonical-order case. The intended wording admits more than one possible emendation. In the remaining case \(p_g=q=1\), one has \(\chi=1\), so the canonical fibration has elliptic base and \(\deg\mathcal L=1\). The canonical formula gives
\[
P_2=h^0\left(B,\mathcal L^{\otimes2}\otimes
 \mathcal O_B\left(\sum_i\left\lfloor2(1-1/m_i)\right\rfloor b_i\right)\right)\ge2.
\]
The nontrivial paracanonical classes have one moving fibre, together with the fixed fractional-fibre contributions. This proves the conclusion on p.~450. It also explains why a special adjoint system cannot be handled merely by its virtual dimension: in the source's argument the auxiliary curve must be sufficiently positive for the displayed adjoint dimension and deficiency to have the claimed values.

For the \(p_g=q=1\) example referred to on p.~450, one must use the corrected branch description in Chapter X, \S11. The text's invitation to classify all such surfaces is a historical research invitation; it is not an assertion that their classification remains unchanged or entirely open today.

\subsection{Canonical torsion and the twelfth plurigenus}

The last paragraph of p.~450 supplies the four classes with torsion canonical bundle: Enriques, bielliptic, K3, and abelian surfaces. Their invariants on the minimal model are
\[
\begin{array}{c|c|c|c|c}
\text{type}&p_g&q&\chi&\operatorname{ord}(K_X)\\\hline
\text{Enriques}&0&0&1&2\\
\text{bielliptic}&0&1&0&2,3,4,6\\
\text{K3}&1&0&2&1\\
\text{abelian}&1&2&0&1
\end{array}
\]
For a torsion canonical bundle of order \(r\),
\[
P_m=\begin{cases}1&r\mid m,\\0&r\nmid m.\end{cases}
\]
Indeed a nontrivial numerically trivial line bundle cannot have a nonzero section: its effective zero divisor would have positive intersection with an ample divisor. Every listed order divides 12, so \(P_{12}=1\). The converse belongs to the effective plurigenus classification, not just to this elementary torsion calculation. The corrected arguments in Chapter X and the canonical formula supply the Kodaira-dimension-one part; general type gives \(P_{12}=\chi+66K_X^2>1\), and ruled surfaces give zero. The last bielliptic label on p.~450 is printed \(\mathrm{III}_e\); the established list and p.~457 require \(\mathrm{III}_c\).

\section[Teoria geometrica delle superficie iperellittiche]{Nota sulla teoria geometrica delle superficie iperellittiche}
\label{cp2-xi-6}

\subsection{Polarizations, periods, and moduli}

The moduli counts on pp.~450--451 refer to specified algebraic families: 10 for Enriques surfaces, two for the two general bielliptic types and one for each special elliptic-modulus type, 19 for polarized K3 surfaces, and three for polarized abelian surfaces. Chapter VII already supplies the required qualifications to the K3 count. None of these numbers asserts that an unpolarized surface has a unique projective model, or that discrete polarization data classify its complex structure.

An ample line bundle on an abelian surface has polarization type \((d_1,d_2)\), with \(d_1\mid d_2\), and
\[
L^2=2d_1d_2,\qquad h^0(A,L)=d_1d_2,\qquad
p_a(L)=1+d_1d_2.
\]
The primitive type \((1,\delta)\) is the setting of Enriques's divisor \(\delta\) and of his normal-period display. The period matrix additionally satisfies the Riemann bilinear relations, including positive-definite imaginary part; three complex period parameters remain. The integer \(\delta\) belongs to the chosen polarization, and is not an absolute invariant of an unpolarized abelian surface admitting several polarizations. For an ample curve of genus \(\pi\), Riemann--Roch and vanishing give \(\dim|L|=\pi-2\). Very ampleness is an additional condition: small polarizations can define a multiple map or no embedding at all.

A maximal isotropic subgroup of the finite polarization kernel gives an isogeny to a principally polarized abelian surface. In type \((1,\delta)\) its order, and hence this isogeny's degree, is \(\delta\). A principally polarized abelian surface is either the Jacobian of a smooth genus-two curve or a product of elliptic curves with its product polarization. The latter is the compact-type genus-two degeneration which the source itself uses later. These statements replace the insufficient inference on p.~452 that equal moduli counts by themselves establish domination by Jacobians. Isogenies have dual isogenies; composing with the dual gives multiplication by an integer, explaining the source's passage between quotients by finite translation groups and unramified multiple covers. This argument uses the actual polarization and isogeny structure, not a comparison of parameter counts. Polarization descent is stated precisely by Milne (2008), Part~I, Proposition~13.8, Remark~13.9, and Corollary~13.10. The Jacobian/product dichotomy is Voight (2021), Theorem~43.5.1.

\subsection{The Jacobian and Kummer constructions}

For a genus-two curve \(C\), the Abel map \(C^{(2)}\to J(C)\) is birational and contracts its unique canonical pencil. It does not make the symmetric square identical to the minimal abelian model at every point. A residual involution associated with a degree-four divisor acts on \(\operatorname{Pic}^2(C)\) as \(x\mapsto a-x\). Its fixed points solve \(2x=a\), so there are \(2^4=16\). On the principally polarized Jacobian, the quotient by such an involution is the quartic Kummer surface; the 16 fixed points become ordinary double points. Products of two residual involutions are translations. These are the geometric mechanisms described on pp.~451--452 and developed in Part~1, Section~\ref{cp1-v-8}, and Chapter IX's discussion of Kummer surfaces and translation groups.

The geometric genus-two-curve argument on pp.~452--453 is only a sketch. A linear system of dimension \(\pi-2\) has expected dimension zero after imposing \(\pi-2\) unspecified nodes, but an expected dimension does not prove existence, finiteness, irreducibility, or nodality. The product principal polarization already illustrates the danger: its only theta divisor is reducible. The later argument also needs a curve whose actual translation stabilizer is trivial; otherwise its orbit is an isogenous quotient of \(A\), not automatically \(A\) itself. For a smooth genus-two curve, a nonidentity translation preserving it would induce a free automorphism; Riemann--Hurwitz excludes a nontrivial free finite action on a genus-two curve. With that hypothesis the orbit argument works. The same argument applies to an integral singular curve of geometric genus two: translations preserving it lift to its normalization, and a fixed point upstairs would give a fixed point on the abelian surface. Thus, once irreducibility and geometric genus two are proved, the stabilizer assertion needs no extra exception for nodes. The isogeny argument in the preceding subsection supplies the classification without presupposing an unproved nodal-curve existence theorem.

An abelian variety of dimension \(g\) has polarized moduli dimension \(g(g+1)/2\), whereas smooth genus-\(g\) curves have dimension \(3g-3\) for \(g\ge2\). The strict inequality used on p.~452 holds for \(g\ge4\), not for every \(g\). It proves that a general abelian variety of dimension at least four cannot be a Jacobian; allowing isogenies adds countably many correspondences and still does not fill the moduli space. This is a statement about general abelian varieties, not about the existence of a Jacobian containing a particular abelian variety up to isogeny. The latter is a different theorem, already separated in Chapter IX.

\subsection{The genus-three illustration and the last quotient remarks}

For a genus-three curve mapping with degree \(n\) to an elliptic curve, four simple branch points satisfy Riemann--Hurwitz. The branch choices must also satisfy transitivity and the global monodromy relation; arbitrary substitutions do not necessarily define a connected cover. For fixed admissible discrete monodromy, the expected Hurwitz-space dimension is
\[
1+4-1=4:
\]
one elliptic modulus, four branch positions, and subtraction of simultaneous translation. The elliptic factor of \(J(C)\) has a complementary abelian surface up to isogeny. The addition map from the product of the two factors has finite kernel equal to their intersection. The corresponding finite quotient of the ambient Jacobian is a product of isogenous factors. Identifying its abelian-surface factor with a genus-two Jacobian requires a principal polarization, which is not supplied merely by taking the intersection subgroup. After a suitable additional isogeny to a principally polarized surface, one obtains either a genus-two Jacobian or a product of elliptic curves. The former product with the elliptic factor is the Jacobian of a compact-type degeneration with genus-two and elliptic components. The latter requires a further decomposed compact-type interpretation. This distinguishes the valid isogeny construction from the source's stronger identification of its particular quotient.

Subtracting the elliptic modulus from four does not, by itself, prove that the complementary surfaces vary in three independent moduli or fill a prescribed polarization space. One would need a dominance statement for the Prym map. The source does not establish it. Its valid contribution is the construction of a complementary abelian surface and the isogeny relation, not a complete proof of a moduli theorem.

On p.~455 ``rank'' means the generic number of points in a period fundamental domain representing one point of the quotient; it is the degree of the finite quotient map. A nontranslation free quotient in the special product situation can give a bielliptic surface. The quotient by inversion gives a singular Kummer K3 surface, whose resolution has 16 disjoint exceptional \((-2)\)-curves. Changing the polarization changes the available projective models, not the number of two-torsion fixed points. The displayed degrees \(2\pi-2\) and dimensions \(\pi\) belong to suitable polarized K3 models and need the same base-point-free and birationality conditions as Chapter VII.

The claimed complete classifications of regular exceptional quotients by Enriques--Severi and Bagnera--De Franchis are imported from the cited 1908 memoirs; no proof or complete group list is supplied on this page. Likewise the reference to Godeaux's school reports earlier work, rather than proving that every finite correspondence on a K3 surface is a quotient by automorphisms. A quotient must be Galois before that conclusion is warranted. The assertion depends on the precise family hypotheses and group lists of those memoirs; the abbreviated citation on this page does not supply them. The printed ``Cap.~XII, \S2'' on p.~451 should be ``Cap.~VII, \S2''; the cited book has no Chapter XII. The other p.~451 locator, to Enriques--Chisini's \emph{Lezioni}, should read Book~VI, Chapter~III, \S40, Volume~IV, p.~246; the relevant theorem is on p.~250. The printed Book~V and p.~24 are both erroneous.

\section{Classificazione generale delle superficie algebriche}
\label{cp2-xi-7}

\subsection{The four classes and the twelfth plurigenus}

The classification on pp.~455--458 is a synthesis, not a new proof of its constituent theorems. With the absolute linear genus interpreted on the minimal model, its four alternatives become
\[
\begin{array}{c|c|c|l}
\text{class}&P_{12}&K_X^2&\text{birational type}\\\hline
A&0&\text{not prescribed}&\kappa=-\infty:\ \text{rational or ruled}\\
B&1&0&\kappa=0:\ \text{torsion canonical bundle}\\
C&>1&0&\kappa=1:\ \text{genus-one fibration}\\
D&>1&>0&\kappa=2:\ \text{general type}
\end{array}
\]
In the first row a minimal ruled surface can have negative canonical square, and the square of a chosen nonminimal model is not a birational invariant. In the remaining rows the square is that of the minimal surface. The effective assertion that the single index 12 distinguishes these classes requires the characteristic-zero plurigenus results of Chapter X; it is stronger than the definition of Kodaira dimension by eventual growth. It is not inferred from \(K_X^2=0\) alone.

For class \(A\), the rationality criterion \(p_a=P_2=0\) is the Castelnuovo criterion of Part~1, Section~\ref{cp1-vi-4}, and Chapter X. An irrational ruled surface over a curve of genus \(b\) has
\[
p_g=0,\quad q=b,\quad p_a=-b,\quad P_m=0\quad(m>0).
\]
The extra criterion \(p_a<-1\) is established in the next section. When \(p_a=-1\), the conditions \(P_4=P_6=0\), equivalently \(P_{12}=0\) in this classification, distinguish the ruled case from the nonruled cases. The source's extra closing parenthesis after the last equality on p.~456 is a printing error. Its statements about infinitely many exceptional curves and continuous transformation families carry the interpretation and qualifications of Section~\ref{cp2-xi-2}.

For class \(B\), numerical triviality of \(K_X\) gives \(H^2=2p_a(H)-2\) by adjunction. For an ample \(H\), vanishing and Riemann--Roch give
\[
h^0(X,H)=\frac{H^2}{2}+\chi(\mathcal O_X),\qquad
\dim|H|=\pi+\chi(\mathcal O_X)-2.
\]
Thus the dimensions in the book are \(\pi\), \(\pi-1\), and \(\pi-2\) for K3, Enriques, and abelian or bielliptic surfaces respectively. These are complete ample systems; a low-degree polarization can fail to embed the surface, as the examples in Chapters VII and X already show. The canonical-section condition also has a precise meaning. Adjunction gives \(K_H=(K_X+H)|_H\), so if \(K_X\) is trivial the hyperplane bundle restricts to \(K_H\). If \(K_X\) has order \(s>1\), only
\[
\mathcal O_H(sH)\simeq K_H^{\otimes s}
\]
is automatic. This explains the distinction \(B'\), consisting of K3 and abelian surfaces, and \(B''\), consisting of Enriques and the seven bielliptic types. It does not imply that each torsion restriction has exactly the same order as \(K_X\).

The two last paragraphs concerning \(B'\) and \(B''\) on p.~457 interchange their labels. The eight families are the Enriques family and the seven bielliptic deformation types, hence belong to \(B''\). These surfaces possess genus-one pencils. A general polarized K3 or abelian surface in \(B'\) need not possess such a pencil. The dimensions 10, two, and one are the previously stated dimensions of the appropriate families; the infinitely many polarized K3 and abelian families are distinguished by discrete polarization data and have dimensions 19 and three. A single underlying surface can admit several polarizations, so this is not a partition of unpolarized surfaces by an absolute degree invariant. The stronger assertion that all the indicated projective models have been constructed by the displayed equations is not proved by these moduli counts.

\subsection{The genus-one class and the data it requires}

For class \(C\), the pluricanonical map of a sufficiently divisible order is a fibration \(f:X\to B\) with connected genus-one fibres. The existence of this fibration follows from the arguments in Chapters VII and X, interpreted on a relatively minimal model. The relevant formula is
\[
\omega_X\simeq f^*(\omega_B\otimes L)\otimes
\mathcal O_X\left(\sum_i(m_i-1)F_i\right),\qquad
\deg L=\chi(\mathcal O_X),
\]
where \(m_iF_i\) are the multiple fibres. It retains the line bundle \(L\), not only its degree. For \(\chi>0\), Leray gives \(q=g(B)\), which is the source's \(\varrho=p_g-p_a\). For \(\chi=0\),
\[
q=g(B)+h^0(B,L^{-1}),
\]
so \(g(B)=q-1=p_g\) when \(L\simeq\mathcal O_B\), and \(g(B)=q=p_g+1\) when \(L\) is nontrivial. The latter case is omitted on p.~458 and in the final table. The degree-zero line bundle is torsion in the isotrivial projective situation, but torsion does not mean trivial. Section~\ref{cp2-xi-monodromy-example} gives a smooth example with no multiple fibres in which the second alternative is indispensable.

Enriques's determinant \(n\) is the least degree of a rationally defined effective point group on the generic genus-one fibre, equivalently the least degree of a multisection. Over \(\mathbb C(B)\) it is the index of the genus-one curve; the standard index interpretation uses that the index is realized by a rational effective divisor in this setting. It is not the degree of an arbitrary projective model and does not assert the existence of a section. The number and multiplicities of multiple fibres, \(\chi\), the base genus, and this index are useful discrete organizing data. They do not determine a complex surface up to isomorphism. Branch positions, the Jacobian fibration, its monodromy, the Hodge line, and the torsor class may vary. Likewise equal plurigenera are not a complete isomorphism invariant. The source calls extension of its regular-case classification to the irregular case an object of research; this is an explicitly stated programme, not an established theorem.

\subsection{General type and bounded families}

For class \(D\), the finite-family assertion on p.~458 is correct after replacing the unspecified pluricanonical model by a uniformly controlled one. Enriques's argument does not supply a uniform exponent or a fixed Hilbert polynomial; merely observing that each surface has some projective model does not prove boundedness. Bombieri's Main Theorem (1973, p.~173) supplies a uniform exponent: \(5K_X\) defines the canonical model, contracting precisely the canonical exceptional configurations, and the resulting model has only rational double points. Its Hilbert polynomial with respect to \(5K\) is
\[
\chi\bigl(\mathcal O_{X_{\rm can}}(5tK)\bigr)
=\chi(\mathcal O_X)+\frac{25t^2-5t}{2}K_X^2.
\]
Rational singularities preserve this polynomial and \(\chi\). Fixing \(K_X^2\) and \(\chi\) therefore places the canonical models in a Hilbert scheme of finite type, with only finitely many irreducible components. This proves finite bounded algebraic families, not finitely many isomorphism classes. Within a component moduli may vary in positive dimension.

The stronger wording on p.~464, finitely many types for fixed \(p^{(1)}\), can be interpreted in the same bounded-family sense. Noether's inequality gives \(p_g\le(K_X^2+4)/2\), and general type gives \(\chi\ge1\). Since \(\chi=1-q+p_g\le1+p_g\), there are only finitely many possible positive \(\chi\) for a fixed \(K_X^2\). The preceding Hilbert-scheme argument applies to each value. This repairs the conclusion without pretending that a full classification of those components or their points is available.

The observations that regular surfaces predominate among the examples constructed in Chapter VIII, and that some small-square constructions appear necessarily regular, are observations about those constructions. They are not proofs of regularity for every surface with the same numerical invariants. The two inequalities in the footnote on p.~459 are more serious: as printed, they are false. Their hypothesis is merely that the canonical system is not composed with a pencil. They assert
\[
K_X^2\ge6\chi(\mathcal O_X)-3\quad(q\ge3),
\qquad
K_X^2\ge3\chi(\mathcal O_X)\quad(q=2).
\]
Both fail even for smooth double covers with a two-dimensional canonical image. The following constructions are included to repair this specific claim, rather than to extend the classification programme.

\begin{example}[A double cover with \(q=2\)]
\label{cp2-xi-false-bound-q2}
Let \(B\) be a curve of genus two, let \(A\) be a line bundle of degree \(a\ge4\), and put
\[
Y=B\times\mathbb P^1,\qquad
L=\operatorname{pr}_B^*A\otimes\operatorname{pr}_{\mathbb P^1}^*\mathcal O(3).
\]
Choose a smooth divisor in \( |2L|\) and form its double cover \(\pi:X\to Y\). Both \(2L\) and \(K_Y+L\) are very ample, so Bertini supplies a smooth branch divisor and \(K_X=\pi^*(K_Y+L)\) is ample. The double-cover identities of Part~1, Section~\ref{cp1-v-7}, give
\[
\chi(\mathcal O_X)=2\chi(\mathcal O_Y)+\tfrac12 L(L+K_Y)=2a+1,
\qquad K_X^2=2(K_Y+L)^2=4a+8.
\]
Indeed \(\chi(\mathcal O_Y)=-1\), \(L^2=6a\), and \(K_YL=6-2a\). By K\"unneth, \(H^1(Y,L^{-1})=0\), since the two summands each contain the space of sections of a negative-degree line bundle. Thus \(q(X)=2\) and \(p_g(X)=2a+2\). As \(p_g(Y)=0\), the entire canonical system is the pullback of \( |K_Y+L|\), and its map has degree two onto the embedded surface \(Y\). It is not composed with a pencil. For \(a=4\),
\[
(\chi,p_g,q,K_X^2)=(9,10,2,24),\qquad24<3\cdot9.
\]
Equivalently \(p^{(1)}=25<3p_a+4=28\), contradicting the printed \(q=2\) bound.
\end{example}

\begin{example}[A double cover with \(q=3\)]
\label{cp2-xi-false-bound-q3}
Let \(Y=B\times E\), where \(g(B)=2\) and \(E\) is elliptic. Take \(L=A\boxtimes M\), with \(\deg A=a\ge4\) and \(\deg M=d\ge3\), and again take a smooth double cover branched in \( |2L|\). Here
\[
\chi(\mathcal O_Y)=0,\quad L^2=2ad,\quad K_YL=2d,
\]
and therefore
\[
\chi(\mathcal O_X)=d(a+1),\qquad
K_X^2=4d(a+2),\qquad q(X)=3.
\]
The irregularity follows from \(H^1(Y,L^{-1})=0\), by the same K\"unneth calculation. The canonical system contains the pullback of the very ample system \( |K_Y+L|\), so its image has dimension two. For \(a=4,d=3\),
\[
(\chi,p_g,q,K_X^2)=(15,17,3,72),\qquad72<6\cdot15-3=87.
\]
Thus \(p^{(1)}=73<6p_a+4=88\). This contradicts the printed \(q\ge3\) bound.
\end{example}

The calculations use the standard double-cover formulas, also written explicitly by Debarre (1982, \S1.4) and Xiao (1985, \S1.3). A valid general replacement is Debarre's inequality \(K_X^2\ge2p_g\) for minimal irregular surfaces of general type (1982, Theorem~6.1, with the 1983 Addendum completing the canonical-pencil case). Stronger inequalities depend on whether the canonical map is birational, has degree two, has another degree, or is composed with a pencil. The source's weaker phrase ``not composed with a pencil'' does not license a bound that needs more restrictive canonical-map hypotheses. The attribution to an unpublished communication of Castelnuovo in November 1942 cannot be used to recover such missing hypotheses.

\section{Superficie di genere numerico negativo}
\label{cp2-xi-8}

\subsection{The Castelnuovo--de Franchis inequality}

The central theorem on p.~459 is correct, with a useful strengthening of its wording: if \(X\) has no irrational pencil of base genus at least two, then
\[
p_g\ge2q-3,
\]
or, equivalently for integral invariants,
\[
p_g>2(q-2),\qquad p_g<2p_a+4.
\]
The numerical contrapositive is
\[
p_g\ge2p_a+4\quad\Longleftrightarrow\quad p_g\le2q-4,
\]
and it forces a fibration over a curve of genus at least two. The book says only ``irrational'', which also permits an elliptic base; the stronger base-genus conclusion matters in its applications.

For completeness, the proof explains exactly what the historical appeal to integrals supplies. Put \(V=H^0(X,\Omega_X^1)\), so \(\dim V=q\), and consider the wedge map
\[
\mu:\bigwedge^2V\longrightarrow H^0(X,K_X).
\]
If independent holomorphic one-forms \(\alpha,\beta\) satisfy \(\alpha\wedge\beta=0\), then \(\alpha=h\beta\) in the function field for a nonconstant meromorphic function \(h\). Holomorphic one-forms on a projective complex variety are closed; differentiating gives \(dh\wedge\beta=0\). Resolve the map defined by \(h\) and take its Stein factorization. The two forms are pulled back from meromorphic differentials on its smooth base curve \(B\). These differentials have no poles: locally, a differential with a pole remains polar after a finite map \(t=z^e\), so it could not pull back to a holomorphic form. They are independent, hence \(g(B)\ge2\). Any exceptional rational curve maps constantly to \(B\), so the fibration descends to \(X\). This is the Castelnuovo--de Franchis theorem, in the formulation of Mendes Lopes and Pardini (2011, Theorem~2.1.1).

If no such fibration exists, the projectivized kernel of \(\mu\) is disjoint from the Pl\"ucker Grassmannian \(\operatorname{Gr}(2,V)\subset\mathbb P(\bigwedge^2V)\), of dimension \(2q-4\). The projective dimension theorem then gives
\[
\operatorname{rank}\mu\ge2q-3,
\]
and \(p_g\ge\operatorname{rank}\mu\) proves the inequality. The cases \(q\le1\) are immediate. Thus the numerical argument is algebraic linear algebra after the closedness of one-forms and the fibration theorem. Enriques's account on pp.~461--462 correctly identifies dependent integrals and their Jacobians as the mechanism, but it omits this dimension calculation. His statement that an algebraic-geometric proof was then unavailable is a historical report, not a continuing limitation of the theorem. The references to De Franchis, Castelnuovo, and Enriques in 1905 record the historical attribution of the theorem.

\subsection{Negative Euler characteristic and the boundary case}

If \(p_a<-1\), then \(\chi(\mathcal O_X)<0\). The inequality forces an irrational pencil \(f:X\to B\) with \(b=g(B)\ge2\). Suppose that \(X\) is not ruled, and replace it by its minimal model. Then \(K_X\) is nef, so \(K_X^2\ge0\), and the general fibre has genus \(g\ge1\). The Euler-characteristic formula is
\[
e(X)=4(g-1)(b-1)+\sum_v\bigl(e(F_v)-e(F)\bigr)
=4(g-1)(b-1)+\Delta\ge0.
\]
Nonnegativity of each fibre contribution is the result treated in Chapter X, with the multiplicities and singular support properly accounted for. Noether's formula now yields
\[
12\chi(\mathcal O_X)=K_X^2+e(X)\ge0,
\]
a contradiction. Therefore \(X\) is ruled, and its base genus is \(b=q=-p_a\). This proves the assertion on pp.~459--460 without assuming \(p_g=0\) beforehand.

The displayed equation \(\Delta-4=12p_a-p^{(1)}+9\) on these pages omits \(4(g-1)(b-1)\) unless that term vanishes or has been absorbed into a differently defined \(\Delta\). The preceding chapter defines \(\Delta\) by fibre defects, so it cannot silently absorb this term. The correct inequality still proves the desired ruledness result. This is an error in a displayed intermediate formula, not a counterexample to the theorem \(p_a<-1\Rightarrow\) ruled.

For \(p_a=-1\), the same reasoning, or the minimal classification when the small-\(p_g\) cases do not meet the inequality, gives \(K_X^2=0\) on a nonruled minimal surface. The consequences asserted after this calculation need correction. The full list is:
\begin{enumerate}
\item ruled surfaces over an elliptic curve, with all \(P_m=0\);
\item abelian and bielliptic surfaces, with \(\kappa=0\);
\item properly elliptic surfaces with \(\chi=0\), namely smooth free quotients \( (C\times E)/G\) with \(g(C)\ge2\), including both trivial and nontrivial action on the elliptic differential.
\end{enumerate}
In the third case the fibres are smooth elliptic curves or multiples of smooth elliptic curves. The trivial-differential case has a positive-dimensional elliptic automorphism group and \(g(B)=p_g\); the nontrivial-differential case has \(\operatorname{Aut}^0(X)=0\) and \(g(B)=p_g+1\). The latter surfaces invalidate the identification on pp.~460--461 of all negative-arithmetic-genus surfaces with nonrational surfaces possessing a continuous birational group. Catanese, Frapporti, Glei\ss ner, Liu, and Sch\"utt (2025, Lemma~1.1) explicitly distinguish these two cases. The counterexample in Section~\ref{cp2-xi-monodromy-example} has \(p_g=1\), and its version with a higher-genus base gives \(p_g>1\). Chapter X supplies the \(p_g=0,q=1\) example. Thus none of the three geometric-genus subcases can be protected by the book's asserted exclusion.

There is a correct implication in the direction from groups to invariants. A nonrational ruled surface has negative \(p_a\); a nonruled surface with a positive-dimensional connected birational group is abelian or belongs to the translation case of the elliptic construction, and has \(p_a=-1\). The converse is the failed implication. Likewise \(p_g=P_4=1\) still characterizes an abelian surface within this nonruled \(\chi=0\) setting, whereas \(p_g=1,P_4>1\) does not imply an elliptic transformation group. On p.~460 the zero-plurigenus equalities describe the ruled cases to be excluded. If read as the conditions to be conjoined with \(p_a=-1\), the signs must instead be \(P_4+P_6>0\) or \(P_{12}>0\). The surrounding wording permits the first reading, so the displayed zero signs are compatible with that reading. Only the reference following ``In questo caso'', concerning the \(p_g=1\) analysis, should point to the present Chapter XI, \S4. The later sentence completing a result of Chapter X, \S4 refers correctly to that earlier chapter.

\subsection{The Rosenblatt--Comessatti range}

The historical note on p.~462 attributes to Rosenblatt and Comessatti a refinement under \(p_g\ge2p_a+4\), equivalently \(p_g\le2q-4\). Its ruled and Kodaira-dimension-one cases require the same monodromy correction as above. For general type the modern result is especially precise: Beauville's appendix to Debarre (1982) proves
\[
p_g\ge2q-4,
\]
with equality precisely for a product of a genus-two curve and a curve of genus at least two. Therefore a general-type surface satisfying the source's reverse inequality lies on the equality line and is such a product. If \(X=C_2\times C_g\), \(g\ge2\), then
\[
p_g=2g,\qquad q=g+2,\qquad
\chi=g-1,\qquad p_a=g-2,\qquad K_X^2=8(g-1).
\]
This verifies the source's third numerical type and supplies its missing modern proof. When \(g=1\), the product has Kodaira dimension one and belongs instead to the elliptic part of the list. The fourth type, \(K_X^2=0\) and \(p_a\ge0\), is a further genus-one-fibration class, not general type.

The note's broader claim about irregular double planes having an irrational pencil is quoted from De Franchis's earlier work; no proof or branch hypotheses are given here. Its use requires the normal model, resolution, and branch assumptions of the double-plane theorems in Part~1 and Chapter VIII. It is not a new proof for arbitrary singular coverings. Beauville's appendix (1982, pp.~343--346) explicitly reports an exceptional case left by Rosenblatt, an incorrect proof by Comessatti, and the subsequent absence of a complete proof. His repair combines Castelnuovo--de Franchis with positivity for a fibration of fibre genus \(g\ge2\) over a base of genus \(b\ge2\):
\[
\chi\ge(b-1)(g-1)=(b-2)(g-2)+b+g-3\ge q-3.
\]
Here \(q\le b+g\); equality in that inequality, for this fibre-genus-at-least-two situation, makes the fibration birationally a product. The hypothesis \(p_g\le2q-4\) says \(\chi\le q-3\), so every inequality is an equality. Hence one factor has genus two and minimality gives the asserted product. This is a documented later repair, not a presumption that the original proofs were complete. Mendes Lopes and Pardini (2011, \S3) present the same later inequality and equality characterization.

\section{Riassunto della classificazione precisata}
\label{cp2-xi-9}

\subsection{The bicanonical conclusion and its remaining hypothesis}

On a minimal surface of general type, pluricanonical vanishing and Riemann--Roch give
\[
P_m=\chi(\mathcal O_X)+\frac{m(m-1)}2K_X^2\quad(m\ge2).
\]
Consequently
\[
P_2-1=p_a+p^{(1)}-1.
\]
Since \(\chi\ge1\), \(p^{(1)}>3\), or \(K_X^2>2\), implies \(P_2-1\ge3\). The arithmetic on p.~463 is correct. A linear system of dimension at least three can nevertheless be composed with a pencil. Enriques explicitly marks this doubt and does not claim to have disposed of it on this page. His footnote's ``cap.~IV, \S\S18--20'' should be Chapter VIII, the relevant chapter and section range.

Xiao's Theorem~1 (1985, p.~24) supplies an exact modern repair: the bicanonical map of a minimal complex surface of general type is generically finite unless \(p_g=q=0\) and \(K_X^2=1\). In that exceptional numerical Godeaux case \(P_2=2\), so its image is a curve. Thus the conclusion about a two-dimensional bicanonical image is valid throughout the range \(p^{(1)}>3\), indeed in a larger range. ``A bicanonical model, possibly multiple'' means a generically finite map onto its image; it does not establish a birational embedding. The degree-two and other multiple maps discussed in Chapter VIII must remain visible.

The suggestion that the \(K_X^2=1,2\) surfaces might all be regular is not established by the source's constructions. For \(K_X^2=1\), Debarre's irregularity bound (1982, Theorem~6.1, completed by the 1983 Addendum) does imply regularity: if \(q>0\), then \(\chi\ge1\) implies \(p_g\ge q\ge1\), while \(K_X^2\ge2p_g\ge2\), a contradiction. For \(K_X^2=2\), the same argument permits \(p_g=q=1\), and such surfaces do exist. A precise construction uses \(Y=E^{(2)}\), the symmetric square of an elliptic curve, with section \(D_0\) and Albanese fibre \(G_0\). A general branch divisor in \( |6D_0-2G_0|\) is smooth, and its double cover has \(p_g=q=1,K_X^2=2\). Lewis and Lyons (2018, Theorem~2.1, reproducing the Bombieri--Catanese and Horikawa classification) verifies both smoothness and these invariants. Indeed \(D_0^2=D_0G_0=1\), \(G_0^2=0\), \(K_Y\equiv-2D_0+G_0\), and \(L=3D_0-G_0\), so \(K_X\equiv\pi^*D_0\) and
\[
K_X^2=2,\qquad\chi(\mathcal O_X)=\tfrac12L(L+K_Y)=1.
\]
Thus the regularity suggestion is false for \(p^{(1)}=3\). One must not promote the assumption \(q=0\) made in Chapter VIII to a theorem here. The two false bounds in the footnote on p.~459 cannot close this gap.

\subsection{Reading the final tables}

The tables on pp.~463--464 restate the four classes, so their entries inherit the hypotheses and repairs above rather than supplying additional proofs. In the ruled row, \(n>2\pi-2\) is an existence criterion for an appropriate moving system, not a requirement on every curve. In the \(P_{12}=1\) row, the integer polarization distinctions organize K3 and abelian models, whereas Enriques surfaces and the seven bielliptic deformation types have the stated finite type list. The second occurrence of \(P_4=0,1\) in the bielliptic subrow is a printed repetition; the intended fourth entry is \(P_6=0,1\). The preceding classifications determine the actual plurigenera from the canonical torsion order, rather than allowing these four zero-or-one entries independently.

For \(P_{12}>1,K_X^2=0,\chi=0\), the table must include both base genera \(p_g\) and \(p_g+1\), according to the Hodge line. Its printed \(p_g\) alone is incomplete. When \(\chi>0\), \(g(B)=q=p_g-p_a\) is correct. The slogan \(n<2\pi-2\) in this row and in general type again concerns a sufficiently positive system: a genus-one fibre always has equality. The general-type assertion about a canonical model for \(p_g>3\) needs a separate hypothesis on the canonical map. The dimension \(p_g-1\) by itself proves neither two-dimensional image nor birationality, and the canonical system may be composed with a pencil. The bicanonical assertion in its stated range is repaired by Xiao's theorem; a uniformly birational model is supplied by \(5K_X\), as explained above.

Finally ``a finite number of types for fixed \(p^{(1)}\)'' has the bounded-family interpretation proved in Section~\ref{cp2-xi-7}. It is not finiteness of individual surfaces. With these qualifications, the fourfold organization by \(P_{12}\) and the minimal canonical square is a durable achievement of the chapter. The closing comparison with a mountain ascent and the quotation concerning mathematical harmony are literary reflections; they make no additional mathematical assertion.

\section*{References}
\markboth{References}{References}

The following list identifies the modern results and historical works cited in Chapters~VII--XI. Citations explicitly marked as reported by Enriques preserve the information available in his text; omitted bibliographical fields and abbreviated titles are not silently completed. Where a published date differs from his reported date, both are distinguished. Author-manuscript and draft locators identify the stated version and need not coincide with the numbering of the final publication.

\begin{sloppypar}
\begin{itemize}
\item Bagnera, G., \& De Franchis, M. (1907). [Notes on hyperelliptic surfaces; titles not specified in these citations]. \emph{Rendiconti dell'Accademia dei Lincei}, ser.~5, Vol.~16, first half-year. Cited by Enriques, pp.~414, 425. Bibliographical details beyond that citation are unverified.

\item Bagnera, G., \& De Franchis, M. (1908). Le superficie algebriche le quali ammettono una rappresentazione parametrica mediante funzioni iperellittiche di due argomenti. \emph{Memorie della Società Italiana delle Scienze detta dei XL}. Cited by Enriques, p.~455. Bibliographical details beyond that citation are unverified.

\item Barth, W. P., Hulek, K., Peters, C. A. M., \& Van de Ven, A. (2004). \emph{Compact complex surfaces} (2nd enlarged ed.; Ergebnisse der Mathematik und ihrer Grenzgebiete, 3rd ser., Vol.~4). Springer. \href{https://doi.org/10.1007/978-3-642-57739-0}{doi:10.1007/978-3-642-57739-0}. For the canonical bundle formula over an arbitrary base curve, see the explicit statement in Bombieri and Mumford (1977), Theorem~2, below.

\item Bauer, I. C., Catanese, F., \& Pignatelli, R. (2006). The moduli space of surfaces with $K^2=6$ and $p_g=4$. \emph{Mathematische Annalen, 336}(2), 421--438. \href{https://doi.org/10.1007/s00208-006-0008-8}{doi:10.1007/s00208-006-0008-8}. \href{https://pignatelli.maths.unitn.it/papers/k26pg4.pdf}{Author manuscript}, Introduction and Theorem~0.1. Their notation $\mathrm{II}\to\mathrm{III}_b$ means a deformation of type $\mathrm{III}_b$ to type $\mathrm{II}$.

\item Beauville, A. (1978). \emph{Surfaces algébriques complexes}. Astérisque, No.~54. Société mathématique de France. \href{https://numdam.org/item/AST_1978__54__R1_0/}{NUMDAM edition}. General reference for minimal surfaces and the classification framework.

\item Beauville, A. (1982). L'inégalité $p_g\geq 2q-4$ pour les surfaces de type général. Appendix to O.~Debarre, \emph{Inégalités numériques pour les surfaces de type général}, \emph{Bulletin de la Société Mathématique de France, 110}, 343--346. \href{https://www.numdam.org/article/BSMF_1982__110__319_0.pdf}{Original article and appendix}.

\item Bianchi, L. (1900). \emph{Lezioni sulla teoria dei gruppi di sostituzioni e delle equazioni algebriche secondo Galois}. Pisa [publisher unspecified in the citation]. Cited by Enriques, pp.~404, 416, 419. Bibliographical details beyond that citation are unverified. Enriques cites Chapter~III, \S~30 and Chapter~IV, \S~76.

\item Bombieri, E. (1973). Canonical models of surfaces of general type. \emph{Publications Mathématiques de l'IHÉS, 42}, 171--219. \href{https://doi.org/10.1007/BF02685880}{doi:10.1007/BF02685880}. \href{https://www.numdam.org/article/PMIHES_1973__42__171_0.pdf}{Original text}, Main Theorem, p.~173, and \S~13, p.~218.

\item Bombieri, E., \& Catanese, F. (1978). The tricanonical map of a surface with $K^2=2$, $p_g=0$. In \emph{C.~P. Ramanujam---A tribute} (Tata Institute of Fundamental Research Studies in Mathematics, Vol.~8, pp.~279--290). Tata Institute of Fundamental Research. \href{https://www.mathe8.uni-bayreuth.de/en/team/prof-fabrizio-catanese/pdf/index.html}{Author's publication record}. The detailed argument used here is available in Catanese's 1981 exposition.

\item Bombieri, E., \& Mumford, D. (1977). Enriques' classification of surfaces in char. $p$, II. In W.~L. Baily, Jr., \& T.~Shioda (Eds.), \emph{Complex analysis and algebraic geometry: A collection of papers dedicated to K.~Kodaira} (pp.~23--42). Iwanami Shoten and Cambridge University Press. \href{https://doi.org/10.1017/CBO9780511569197.004}{doi:10.1017/CBO9780511569197.004}. \href{https://www.dam.brown.edu/people/mumford/alg_geom/papers/1976d--EnrClass-II-NC.pdf}{Author-hosted original}, \S~1, Theorem~2, p.~27, and the characteristic-zero specialization on p.~28. The theorem allows an arbitrary nonsingular base curve and includes the multiple-fibre contributions.

\item Bourniat, P. [name as printed by Enriques] (1936). Note sur quelques surfaces canoniques et sur des surfaces de genre un. \emph{Bulletin de l'Académie Royale de Belgique}, ser.~5, Vol.~22. Cited by Enriques, p.~294. Bibliographical details beyond that citation are unverified.

\item Bourniat, P. [name as printed by Enriques] (1936). Note sur quelques surfaces canoniques. \emph{Académie Royale de Belgique}, cited as ser.~3, Vol.~22. Cited by Enriques, p.~294. Bibliographical details beyond that citation are unverified.

\item Bourniat, P. [name as printed by Enriques] (1937). Sur des hypersurfaces et variétés canoniques. \emph{Académie Royale de Belgique}, ser.~5, Vol.~23. Cited by Enriques, p.~294. Bibliographical details beyond that citation are unverified.

\item Brion, M. (2017). Some structure theorems for algebraic groups. \emph{Proceedings of Symposia in Pure Mathematics, 94}, 53--125. \href{https://arxiv.org/abs/1509.03059}{arXiv:1509.03059}. Theorem~7.1.1 and Proposition~7.1.4 in author version~3, pp.~63--64, give the automorphism-scheme statements used here.

\item Campedelli, L. (1932). Sopra alcuni piani doppi notevoli con curva di diramazione del decimo ordine. \emph{Rendiconti della Reale Accademia Nazionale dei Lincei}, ser.~6, \emph{15}, 536--542. The page range distinguishes this paper from \emph{Sui piani doppi}, listed separately; some later references conflate the two titles and page ranges. \href{https://www.sciencedirect.com/science/article/pii/1385725881900615}{Primary bibliographic cross-reference in Oort and Peters}.

\item Campedelli, L. (1932). Sui piani doppi con curva di diramazione del decimo ordine. \emph{Rendiconti della Reale Accademia Nazionale dei Lincei}, ser.~6, \emph{15}, 358--362. This is distinct from \emph{Sopra alcuni piani doppi notevoli}, listed separately, and is relevant to the constructions invoked by Enriques, pp.~307--308 and~321. \href{https://www.persee.fr/doc/barb_0001-4141_1959_num_45_1_67638}{Godeaux's 1959 original article distinguishes the two references}.

\item Campedelli, L. (1933). Intorno alle superficie algebriche su cui esistono curve di genere $\pi$ e di grado $n>2\pi-2$. \emph{Rendiconti dell'Accademia dei Lincei}, second half-year. Cited by Enriques, p.~376. Bibliographical details beyond that citation are unverified.

\item Campedelli, L. (1934). Ancora sul computo dell'invariante di Zeuthen--Segre per una superficie algebrica. \emph{Rendiconti dell'Accademia dei Lincei}. Cited by Enriques, p.~377. Bibliographical details beyond that citation are unverified.

\item Campedelli, L. (1934). Sul computo dell'invariante di Zeuthen--Segre per una superficie algebrica. \emph{Rendiconti dell'Accademia dei Lincei}. Cited by Enriques, p.~377. Bibliographical details beyond that citation are unverified.

\item Campedelli, L. (1935). Intorno alle superficie ellittiche con un fascio di curve di genere due. \emph{Rendiconti del Seminario Matematico della Università di Padova}. Cited by Enriques, p.~435. Bibliographical details beyond that citation are unverified.

\item Castelnuovo, G. (1891). Osservazioni intorno alla geometria sopra una superficie. Nota II. \emph{Rendiconti dell'Istituto Lombardo}, ser.~2, Vol.~24; reprinted in \emph{Memorie scelte}, no.~XVIII. Cited by Enriques, p.~297. Bibliographical details beyond that citation are unverified.

\item Castelnuovo, G. (1905, May 14). Sulle superficie aventi il genere aritmetico negativo. \emph{Rendiconti del Circolo Matematico di Palermo}; \emph{Memorie scelte}, no.~XXVII, p.~501. Cited by Enriques, p.~461. Bibliographical details beyond that citation are unverified.

\item Castelnuovo, G. (1905, May--June). Sugli integrali semplici appartenenti ad una superficie irregolare. \emph{Rendiconti dell'Accademia dei Lincei}; reprinted in \emph{Memorie scelte}, pp.~473 ff.. Cited by Enriques, pp.~345, 369. Bibliographical details beyond that citation are unverified. The individual dates given are May~21, June~4, and June~18; no.~10 is located at p.~490 of the collected edition.

\item Castelnuovo, G. (1928). La geometria algebrica e la scuola italiana. Lecture to the International Congress of Mathematicians, Bologna, September. Cited by Enriques, pp.~427--428. Bibliographical details beyond that citation are unverified. The congress lecture is the source of the passage quoted by Enriques; no publication-page locator is inferred from the quotation.

\item Castelnuovo, G. (1942, November). [Unpublished communication on inequalities for irregular surfaces]. Reported by Enriques. Cited by Enriques, p.~459. Bibliographical details beyond that citation are unverified. This is a reported communication, not an identified published theorem. The separate communication at p.~445 concerns the Zeuthen--Segre invariant and is likewise cited through Enriques.

\item Castelnuovo, G., \& Enriques, F. (1895, July 29). Sur les surfaces algébriques admettant un groupe continu de transformations birationelles en elles-mêmes [title as cited]. \emph{Comptes rendus de l'Académie des Sciences}. Cited by Enriques, p.~432. Bibliographical details beyond that citation are unverified.

\item Castelnuovo, G., \& Enriques, F. (1901). Sopra alcune questioni fondamentali nella teoria delle superficie algebriche. \emph{Annali di Matematica}. Cited by Enriques, pp.~374, 377, 433. Bibliographical details beyond that citation are unverified. The p.~433 citation specifies no.~19.

\item Castelnuovo, G., \& Enriques, F. (1906). Sur les intégrales simples de première espèce d'une surface ou d'une variété algébrique à plusieurs dimensions. \emph{Annales scientifiques de l'École Normale Supérieure}, ser.~3, Vol.~23. Cited by Enriques, p.~358. Bibliographical details beyond that citation are unverified.

\item Catanese, F. (1981). Pluricanonical mappings of surfaces with $K^2=1,2$, $q=p_g=0$. In G.~Tomassini (Ed.), \emph{Algebraic surfaces: C.I.M.E., Cortona 1977} (pp.~247--266). Liguori. \href{https://www.mathe8.uni-bayreuth.de/de/team/prof-fabrizio-catanese-old/pdf/006.pdf}{Author-hosted scan}; \href{https://link.springer.com/book/10.1007/978-3-642-11087-0}{Springer reprint record}. Relevant passages: \S~III, Theorem~1, pp.~254--255; \S~IV, Theorem~2, pp.~258--259; and \S~V, pp.~260--264. The later Springer volume is a reprint of the 1981 publication.

\item Catanese, F., \& Debarre, O. (1989). Surfaces with $K^2=2$, $p_g=1$, $q=0$. \emph{Journal für die reine und angewandte Mathematik, 395}, 1--55. \href{https://www.mathe8.uni-bayreuth.de/de/team/prof-fabrizio-catanese-old/pdf/032.pdf}{Author-hosted original}. The introduction, pp.~1--3, distinguishes the five bicanonical cases and the stated moduli results.

\item Catanese, F., Frapporti, D., Glei\ss ner, C., Liu, W., \& Schütt, M. (2025). On the cohomologically trivial automorphisms of elliptic surfaces I: $\chi(S)=0$. \emph{Taiwanese Journal of Mathematics, 29}(6), 1209--1260. \href{https://doi.org/10.11650/tjm/241106}{doi:10.11650/tjm/241106}. \href{https://arxiv.org/abs/2408.16936v2}{Author version~2, November 22, 2024}, Lemma~1.1, \S~3.3, and Appendix~B. The translation/nontranslation alternatives and monodromy example are cited in that version.

\item Chisini, O. (1921). [Memoir on branch curves and their crossings; title not given in the citation]. \emph{Atti del Reale Istituto Veneto di Scienze, Lettere ed Arti}, cited as Vol.~XXX, part~2. Cited by Enriques, p.~397. Bibliographical details beyond that citation are unverified. The reported volume number and the precise identity of this abbreviated citation remain unverified.

\item Chisini, O. (1921). Le superficie ellittiche il cui determinante è un numero composto. \emph{Rendiconti dell'Accademia dei Lincei}, ser.~5, Vol.~30, second half-year. Cited by Enriques, pp.~414, 419. Bibliographical details beyond that citation are unverified.

\item Comessatti, A. (1919). [Note on irregular surfaces; title not specified]. \emph{Rendiconti dell'Accademia dei Lincei}. Cited by Enriques, p.~462. Bibliographical details beyond that citation are unverified.

\item Comessatti, A. (1922). Intorno alle superficie algebriche irregolari\ldots{} [title abbreviated by Enriques]. \emph{Rendiconti del Circolo Matematico di Palermo, 46}. Cited by Enriques, p.~462. Bibliographical details beyond that citation are unverified.

\item Cossec, F. R., Dolgachev, I. V., \& Liedtke, C. (2025). \emph{Enriques surfaces I} (2nd ed.). Springer. \href{https://doi.org/10.1007/978-981-96-1214-7}{doi:10.1007/978-981-96-1214-7}. Relevant chapters: Generalities, pp.~199--252; Linear Systems, pp.~253--297; Projective Models, pp.~299--358. Numbering from the earlier edition is not used for this edition.

\item d'Orgeval, B. (1945). \emph{Sur les surfaces algébriques dont tous les genres sont 1}. Gauthier-Villars. Cited by Enriques, p.~257. Bibliographical details beyond that citation are unverified.

\item De Franchis, M. (1903). Sulla varietà $\infty^2$ delle coppie di punti di due curve o di una curva algebrica. \emph{Rendiconti del Circolo Matematico di Palermo}. Cited by Enriques, p.~340. Bibliographical details beyond that citation are unverified.

\item De Franchis, M. (1905, April 25). Sulle superficie algebriche le quali contengono un fascio irrazionale di curve. \emph{Rendiconti del Circolo Matematico di Palermo}. Cited by Enriques, pp.~461--462. Bibliographical details beyond that citation are unverified. Enriques also reports the irregular-double-plane consequence and an earlier 1904 result, without identifying a separate 1904 paper.

\item Debarre, O. (1982). Inégalités numériques pour les surfaces de type général. \emph{Bulletin de la Société Mathématique de France, 110}, 319--346, with an appendix by A.~Beauville. \href{https://www.numdam.org/article/BSMF_1982__110__319_0.pdf}{Original text}. See \S~1.4 for double covers and Theorem~6.1 for $K^2\geq2p_g$; the 1983 addendum below is part of the reference for that theorem.

\item Debarre, O. (1983). Addendum à «Inégalités numériques pour les surfaces de type général». \emph{Bulletin de la Société Mathématique de France, 111}, 301--302. \href{https://doi.org/10.24033/bsmf.1990}{doi:10.24033/bsmf.1990}. \href{https://www.numdam.org/article/BSMF_1983__111__301_0.pdf}{Original text}. It repairs the canonical-pencil case in the proof of the preceding inequality, using an argument of Xiao.

\item Del Padrone, A., \& Mistretta, E. (2006). Families of curves and variation in moduli. \emph{Le Matematiche, 61}(1), 163--177. \href{https://lematematiche.dmi.unict.it/index.php/lematematiche/article/view/113}{Journal record}; \href{https://www.math.unipd.it/~ernesto/pub/pragmaticum-2.pdf}{Author-hosted text}. Theorem~1.3, pp.~165--166, proves isotriviality for smooth families of curves of genus at least two over an elliptic curve or $\mathbf P^1$, $\mathbf C$, or $\mathbf C^*$.

\item Dolgachev, I. V. (1981). Algebraic surfaces with $q=p_g=0$. In G.~Tomassini (Ed.), \emph{Algebraic surfaces: C.I.M.E., Cortona 1977}. Liguori. \href{https://sites.lsa.umich.edu/idolga-new/wp-content/uploads/sites/1467/2024/08/cortona81.pdf}{Author-hosted scan}; \href{https://link.springer.com/book/10.1007/978-3-642-11087-0}{reprint record}. In the author's PDF, Chapter~II, \S~1, Propositions~1--2 occupy pp.~24--25; the canonical-divisor theorem and first corollary for regular surfaces over $\mathbf P^1$ are on pp.~32--35; the Euler comparison is on pp.~36--37. For the Godeaux torsion obstruction, see Chapter~III, \S~2, Proposition~3, PDF pp.~62--63, and the remark on pp.~75--76. The author bibliography's original pagination, 97--217, differs from the reprint contents, 97--215; the PDF locators avoid that unresolved discrepancy.

\item Dolgachev, I. V. (2016). A brief introduction to Enriques surfaces. In \emph{Development of moduli theory---Kyoto 2013} (Advanced Studies in Pure Mathematics, Vol.~69, pp.~1--32). Mathematical Society of Japan. \href{https://arxiv.org/abs/1412.7744}{arXiv:1412.7744}, version~2. See \S~2 for normalization of the tetrahedral sextic, \S~4 for moduli, and \S~5 for nef isotropic sequences and projective models.

\item Dolgachev, I. V., \& Kond\=o, S. (2024, April 19). \emph{Enriques surfaces II} [Author draft]. \href{https://sites.lsa.umich.edu/idolga/wp-content/uploads/sites/1334/2024/08/EnriquesTwo.pdf}{Dated draft}. Definition~6.3.2 and Theorem~6.3.3, printed p.~29, concern special genus-one fibrations on nodal Enriques surfaces.

\item Dolgachev, I. V., \& Kond\=o, S. (2025). \emph{Enriques surfaces II}. Springer. \href{https://doi.org/10.1007/978-981-96-1513-1}{doi:10.1007/978-981-96-1513-1}. Reference for the broader nodal-surface and automorphism theory; the preceding draft's theorem numbers are not assigned to the final edition.

\item du Val, P. (1932). Osservazioni sulle superficie di genere uno che non sono base per un sistema di quadriche. \emph{Rendiconti della Reale Accademia dei Lincei}, February--March. Cited by Enriques, p.~257. Bibliographical details beyond that citation are unverified.

\item du Val, P. (1932). Superficie di genere uno che non sono base per un sistema di quadriche. \emph{Rendiconti della Reale Accademia dei Lincei}, February--March. Cited by Enriques, p.~257. Bibliographical details beyond that citation are unverified.

\item du Val, P. (1933). On rational surfaces whose prime sections are canonical curves. \emph{Proceedings of the London Mathematical Society}, ser.~2, \emph{35}(1), 1--13. Cited by Enriques, p.~250. Enriques gives 1932. \href{https://doi.org/10.1112/plms/s2-35.1.1}{doi:10.1112/plms/s2-35.1.1} identifies the published 1933 article.

\item Enriques, F. (1893). \emph{Ricerche di geometria sulle superficie algebriche}. Clausen. Cited by Enriques, p.~254. Enriques cites Part~III, no.~6 for Castelnuovo's examples; no separate Castelnuovo paper is identified by that citation.

\item Enriques, F. (1893). Sui gruppi continui di trasformazioni cremoniane nel piano. \emph{Rendiconti dell'Accademia dei Lincei}, ser.~5, Vol.~2, 468--473. Cited by Enriques, p.~429. The collected-works record identifies this paper behind the abbreviated reference.

\item Enriques, F. (1893). Sui sistemi lineari di superficie algebriche ad intersezioni variabili iperellittiche [title as cited]. \emph{Rendiconti dell'Accademia dei Lincei}; further reference to \emph{Mathematische Annalen, 46}. Cited by Enriques, p.~298. Bibliographical details beyond that citation are unverified.

\item Enriques, F. (1894). Sui sistemi lineari di superficie algebriche le cui intersezioni variabili sono curve ellittiche. \emph{Rendiconti dell'Accademia dei Lincei}, May--June; further reference to \emph{Mathematische Annalen, 46}. Cited by Enriques, p.~293. Bibliographical details beyond that citation are unverified.

\item Enriques, F. (1894). Sulla massima dimensione dei sistemi lineari di curve di dato genere appartenenti ad una superficie algebrica. \emph{Atti della Reale Accademia delle Scienze di Torino}. Cited by Enriques, p.~375. Bibliographical details beyond that citation are unverified.

\item Enriques, F. (1895, as cited). Un'osservazione relativa alla rappresentazione parametrica delle curve algebriche. \emph{Rendiconti del Circolo Matematico di Palermo}. Cited by Enriques, p.~339. The collected-works record gives publication in Vol.~10 (1896), pp.~30--35.

\item Enriques, F. (1896). Sopra le superficie algebriche di cui le curve canoniche sono iperellittiche. \emph{Rendiconti dell'Accademia dei Lincei}, ser.~5, Vol.~5, p.~191. Cited by Enriques, p.~296. Bibliographical details beyond that citation are unverified.

\item Enriques, F. (1896). Sui piani doppi di genere uno. \emph{Memorie della Società Italiana delle Scienze detta dei XL}. Cited by Enriques, p.~254. Bibliographical details beyond that citation are unverified.

\item Enriques, F. (1897, May). Sulle superficie algebriche di genere lineare $p^{(1)}=3$. \emph{Rendiconti dell'Accademia dei Lincei}. Cited by Enriques, p.~311. Bibliographical details beyond that citation are unverified.

\item Enriques, F. (1898, as cited). Sui piani doppi di genere lineare $p^{(1)}=1$. \emph{Rendiconti dell'Accademia dei Lincei}. Cited by Enriques, p.~265. Bibliographical details beyond that citation are unverified. The two notes and their exact publication dates are not reconstructed from the brief retrospective citation.

\item Enriques, F. (1899). Una proprietà delle serie continue di curve appartenenti ad una superficie algebrica regolare. \emph{Rendiconti del Circolo Matematico di Palermo}. Cited by Enriques, p.~339. Bibliographical details beyond that citation are unverified.

\item Enriques, F. (1904, December, as cited). Sulla proprietà caratteristica delle superficie algebriche irregolari. \emph{Rendiconti dell'Accademia delle Scienze dell'Istituto di Bologna}. Cited by Enriques, pp.~342, 368. Bibliographical details beyond that citation are unverified.

\item Enriques, F. (1905). Sulle superficie algebriche che ammettono un gruppo continuo di trasformazioni birazionali in se stesse. \emph{Rendiconti del Circolo Matematico di Palermo, 20}, 61--72. Cited by Enriques, p.~446. Enriques's repeated shortened citations refer to the same memoir.

\item Enriques, F. (1905). Sulle superficie algebriche di genere geometrico zero. \emph{Rendiconti del Circolo Matematico di Palermo, 20}, 1--33. \href{https://www.federigoenriques.org/indice-cronologico-delle-opere/}{Collected-works record, 1905\_2}. Enriques cites this memoir in abbreviated form on pp.~412 and~446.

\item Enriques, F. (1906, as cited). Sopra le superficie algebriche di bigenere uno. \emph{Memorie della Società Italiana delle Scienze detta dei XL}. Cited by Enriques, p.~246. Bibliographical details beyond that citation are unverified.

\item Enriques, F. (1906, December 9). Intorno alle superficie algebriche di genere lineare $p^{(1)}=1$. \emph{Rendiconti dell'Accademia delle Scienze dell'Istituto di Bologna}. Cited by Enriques, pp.~266, 357. Bibliographical details beyond that citation are unverified. The date is the date given by Enriques, not a separately established publication date.

\item Enriques, F. (1908, as cited). Le superficie di genere uno. \emph{Rendiconti dell'Accademia delle Scienze dell'Istituto di Bologna}. Cited by Enriques, p.~255. The collected-works record dates the publication 1909, Vol.~13, pp.~25--28; Enriques's 1914 note identifies the presentation date as December~13, 1908.

\item Enriques, F. (1912). Sulle superficie algebriche con un fascio di curve ellittiche. \emph{Rendiconti dell'Accademia dei Lincei}. Cited by Enriques, p.~259. Bibliographical details beyond that citation are unverified. Enriques explicitly acknowledges a calculation error in no.~3.

\item Enriques, F. (1914). Sulla classificazione delle superficie algebriche, in particolare di genere lineare $p^{(1)}=1$ [title as cited]. \emph{Rendiconti dell'Accademia dei Lincei}, February. Cited by Enriques, pp.~265, 411, 455. Bibliographical details beyond that citation are unverified.

\item Enriques, F. (1930, July). Sopra le superficie algebriche trasformabili in rigate. [Venue not specified in the citation]. Cited by Enriques, p.~413. Bibliographical details beyond that citation are unverified.

\item Enriques, F. (1934, February 18). Sulle superficie ellittiche di genere zero. \emph{Rendiconti dell'Accademia dei Lincei}, ser.~6, Vol.~19, 195--199. Cited by Enriques, p.~426. The title and page range identify the note abbreviated to its venue and date by Enriques.

\item Enriques, F. (1936). Curve infinitamente vicine sopra una superficie algebrica [notes as cited]. \emph{Rendiconti dell'Accademia dei Lincei}, April; and \emph{Rendiconti del Seminario Matematico della Università di Roma}. Cited by Enriques, p.~346. Bibliographical details beyond that citation are unverified.

\item Enriques, F. (1937, October). La proprietà caratteristica delle superficie algebriche irregolari e le curve infinitamente vicine. \emph{Rendiconti dell'Accademia dei Lincei}. Cited by Enriques, p.~346. Bibliographical details beyond that citation are unverified.

\item Enriques, F. (1938, June). Sulla proprietà caratteristica delle superficie algebriche irregolari. \emph{Rendiconti dell'Accademia dei Lincei}. Cited by Enriques, p.~348. Bibliographical details beyond that citation are unverified.

\item Enriques, F. (1939, January 3). Sur la propriété caractéristique des surfaces algébriques irrégulières. \emph{Comptes rendus de l'Académie des Sciences}. Cited by Enriques, p.~349. Bibliographical details beyond that citation are unverified.

\item Enriques, F. (1939, October 19). Des courbes paracanoniques appartenant à une surface algébrique irrégulière. \emph{Société Royale des Sciences de Liège}. Cited by Enriques, p.~354. Bibliographical details beyond that citation are unverified.

\item Enriques, F. (1942, March 20, received). Sui sistemi continui di curve appartenenti ad una superficie algebrica. \emph{Commentarii Mathematici Helvetici}. Cited by Enriques, pp.~349--350, 352. Bibliographical details beyond that citation are unverified. The date is explicitly a manuscript-receipt date in Enriques's citation, not a verified publication year.

\item Enriques, F. (1949). \emph{Le superficie algebriche}. Zanichelli. Chapters~VII--XI, pp.~237--464, are the primary text discussed here.

\item Enriques, F. (lectures collected by L.~Campedelli) (1932). \emph{Lezioni sulla teoria delle superficie algebriche} (Part~I). CEDAM. Cited by Enriques, pp.~255, 376. The first volume of the lectures precedes the 1934 classification installment listed separately. The abbreviated p.~311 locator reported on p.~255 does not identify an edition; that interior locator remains unverified.

\item Enriques, F. (lectures collected by L.~Campedelli) (1934). \emph{Sulla classificazione delle superficie algebriche particolarmente di genere zero}. \emph{Rendiconti del Seminario Matematico della Università di Roma}, ser.~3, Vol.~1, 7--190. Cited by Enriques, pp.~255, 376, 411--414. This is the classification installment of the work cited as Enriques--Campedelli, \emph{Lezioni}. The p.~311 locator on p.~255 has not been matched to a volume or edition; bibliographical details beyond that abbreviated citation are unverified.

\item Enriques, F., \& Chisini, O. (1915--1934). \emph{Lezioni sulla teoria geometrica delle equazioni e delle funzioni algebriche} (Vols.~I--IV). Zanichelli. Cited by Enriques, pp.~259, 272, 291--297, 329--369, 373--454. Books and physical volumes are distinct divisions. Volume~IV (1934) contains Book~VI; its Chapter~III, \S~40 begins on p.~246, and the Kummer quotient theorem is on p.~250. The citation on p.~451 of \emph{Le superficie algebriche}, Book~V and p.~24, should accordingly read Book~VI and p.~246 for the section beginning.

\item Enriques, F., \& Conforto, F. (1939). \emph{Le superficie razionali}. Zanichelli. Cited by Enriques, pp.~250, 314, 317, 372, 430. The title was published under Conforto's name; the collected-works record identifies Enriques's collaboration. Cited locators include Book~I, \S~41, p.~184 and \S~46, pp.~220, 223; and Book~II, Chapters~I--III.

\item Enriques, F., \& Severi, F. (1907). Intorno alle superficie iperellittiche. \emph{Rendiconti dell'Accademia dei Lincei}, ser.~5, Vol.~16, first half-year. Cited by Enriques, p.~425. Bibliographical details beyond that citation are unverified. The abbreviated source cross-reference to 1907 is distinguished from the 1908 note and the Acta memoir.

\item Enriques, F., \& Severi, F. (1908). Intorno alle superficie iperellittiche irregolari. \emph{Rendiconti dell'Accademia dei Lincei}, ser.~5, Vol.~17, first half-year. Cited by Enriques, p.~425. Bibliographical details beyond that citation are unverified.

\item Enriques, F., \& Severi, F. (1909--1910). Mémoire sur les surfaces hyperelliptiques. \emph{Acta Mathematica, 32}, 283--392; \emph{33}, 321--403. \href{https://archive.ymsc.tsinghua.edu.cn/pacm_paperurl/20170108203023008561205}{Institutional archive record for the second installment}. Cited by Enriques, pp.~256, 425, and~455. The published installment dates differ from the 1908 date on p.~455. The source's ``no.~27'' on p.~256 is retained as an internal locator, not a journal issue or a page number; pp.~425--426 cite \S\S~56--57.

\item Fano, G. (1909). Sulle varietà algebriche che sono intersezioni complete di più forme. \emph{Atti della Reale Accademia delle Scienze di Torino, 44}. Cited by Enriques, p.~256. Bibliographical details beyond that citation are unverified.

\item Fano, G. (1936). Su alcune varietà algebriche a 3 dimensioni aventi curve sezioni canoniche. In \emph{Scritti matematici offerti a L.~Berzolari}, p.~329. Cited by Enriques, pp.~293--294. Bibliographical details beyond that citation are unverified.

\item Fano, G. (1936). Superficie algebriche e varietà a 3 dimensioni a curve sezioni canoniche. \emph{Atti dell'Accademia dei Lincei}. Cited by Enriques, p.~293. Bibliographical details beyond that citation are unverified.

\item Fano, G. (1937). Sulle varietà algebriche a tre dimensioni a curve sezioni canoniche. \emph{Memorie dell'Accademia d'Italia}. Cited by Enriques, p.~294. Bibliographical details beyond that citation are unverified.

\item Felisetti, C., \& Fontanari, C. (2022). On the splitting principle of Beniamino Segre. \emph{Annali di Matematica Pura ed Applicata, 201}(4), 1577--1581. \href{https://doi.org/10.1007/s10231-021-01171-w}{doi:10.1007/s10231-021-01171-w}. \href{https://arxiv.org/abs/2105.00892}{Author manuscript}, Theorem~1.1. The flatness, normal-base, nodality, and generality hypotheses are essential.

\item Franchetta, A. (1939). Su alcuni esempi di superficie canoniche. \emph{Seminario Matematico di Roma}. Cited by Enriques, p.~275. Bibliographical details beyond that citation are unverified.

\item Franchetta, A. (1942). Sulle superficie le cui curve canoniche posseggono una $g^1_3$. \emph{Bollettino dell'Unione Matematica Italiana}. Cited by Enriques, p.~303. Bibliographical details beyond that citation are unverified.

\item Franchetta, A. (1942). Sulle superficie regolari di genere uno con curva canonica d'ordine zero. \emph{Istituto Lombardo}, June~18. Cited by Enriques, pp.~256--257. Bibliographical details beyond that citation are unverified.

\item Geppert, H. (1931, as cited). Die Klassification der algebraischen Flaechen [title as printed by Enriques]. \emph{Jahresbericht der Deutschen Mathematiker-Vereinigung}. Cited by Enriques, p.~411. Bibliographical details beyond that citation are unverified. The citation concerns the formal multiplicity triple $(2,5,5)$; it does not by itself assert that such a genus-one fibration exists.

\item Godeaux, L. (1920). Sur le calcul de l'invariant de Zeuthen--Segre. \emph{Mémoires de la Société des Sciences du Hainaut}. Cited by Enriques, p.~377. Bibliographical details beyond that citation are unverified. Enriques's citation misspells Hainaut as ``Hainart.''

\item Godeaux, L. (1936). Sur la construction d'une surface canonique. \emph{Bulletin de la Société Mathématique de Liège, 5}. Cited by Enriques, p.~294. Bibliographical details beyond that citation are unverified.

\item Godeaux, L. (1936). Sur une surface canonique appartenant à la variété de Segre représentant les couples de points de deux plans. \emph{Bulletin de l'Académie Royale de Belgique, 22}. Cited by Enriques, p.~294. Bibliographical details beyond that citation are unverified.

\item Godeaux, L. (1937). Sur les surfaces algébriques dont les sections hyperplanes sont des courbes canoniques (I--III). \emph{Bulletin de la Société Mathématique de Liège, 6}. Cited by Enriques, p.~294. Bibliographical details beyond that citation are unverified.

\item Godeaux, L. (1944). Construction d'une surface canonique du huitième ordre. \emph{Bulletin des Sciences Mathématiques}. Cited by Enriques, p.~294. Bibliographical details beyond that citation are unverified.

\item Godeaux, L. (1944). Construction d'une surface canonique du neuvième ordre. \emph{Bulletin de l'Académie Royale de Belgique}. Cited by Enriques, p.~294. Bibliographical details beyond that citation are unverified.

\item Godeaux, L. (1944). Construction d'une surface canonique du septième ordre. \emph{Bulletin de la Société des Sciences de Liège}. Cited by Enriques, p.~294. Bibliographical details beyond that citation are unverified.

\item Green, M., \& Lazarsfeld, R. (1987). Deformation theory, generic vanishing theorems, and some conjectures of Enriques, Catanese and Beauville. \emph{Inventiones Mathematicae, 90}(2), 389--407. \href{https://doi.org/10.1007/BF01388711}{doi:10.1007/BF01388711}. \href{https://www.math.stonybrook.edu/robert.lazarsfeld/Reprints/Green.Laz.GVT.pdf}{Author-hosted original}. The generic-vanishing application here retains the maximal-Albanese-dimension hypothesis.

\item Horikawa, E. (1973). On deformations of quintic surfaces. \emph{Proceedings of the Japan Academy, 49}(6), 377--379. \href{https://doi.org/10.3792/pja/1195519286}{doi:10.3792/pja/1195519286}. \href{https://www.jstage.jst.go.jp/article/pjab1945/49/6/49_6_377/_pdf}{Original text}, Theorems~1--3.

\item Humbert, G. (1893--1894). Sur une propriété d'une classe de surfaces algébriques [1893 note, as cited]. \emph{Comptes rendus de l'Académie des Sciences} (1893); related memoir in \emph{Journal de Mathématiques Pures et Appliquées}, ser.~4, Vol.~10 (1894). Cited by Enriques, pp.~339, 368. Bibliographical details beyond that citation are unverified. The later memoir's title and page range are not specified by the abbreviated citation.

\item Huybrechts, D. (2016). \emph{Lectures on K3 surfaces} (Cambridge Studies in Advanced Mathematics, Vol.~158). Cambridge University Press. \href{https://doi.org/10.1017/CBO9781316594193}{doi:10.1017/CBO9781316594193}. \href{https://www.math.uni-bonn.de/people/huybrech/K3Global.pdf}{Author's prepublication text}. The locators used here belong to that text: Chapter~2, Corollary~2.5, Proposition~3.1, and Remark~3.2; Chapter~6, \S~2 and Corollary~4.4.

\item Kleiman, S. L. (2005). The Picard scheme. In \emph{Fundamental algebraic geometry: Grothendieck's FGA explained} (Mathematical Surveys and Monographs, Vol.~123, pp.~235--321). American Mathematical Society. \href{https://arxiv.org/abs/math/0504020}{Author manuscript, arXiv:math/0504020}. Theorem~5.4, Theorem~5.11, Corollaries~5.13--5.14, Remarks~5.15 and~5.18, and Proposition~5.19 refer to the manuscript numbering, not the renumbered Chapter~9 of the book.

\item Kleiman, S. L. (2014). The Picard scheme. In L.~Schneps (Ed.), \emph{Alexandre Grothendieck: A mathematical portrait} (pp.~35--74). International Press. \href{https://dspace.mit.edu/entities/publication/b19e8bf7-8b44-404b-b54c-ddceddc1a5d9}{MIT record and author's final manuscript}. Historical account distinct from the 2005 technical exposition. The author's final manuscript in the MIT repository is dated September~2013; the published book bears the date 2014.

\item Kodaira, K., \& Spencer, D. C. (1959). A theorem of completeness of characteristic systems of complete continuous systems. \emph{American Journal of Mathematics, 81}(2), 477--500. \href{https://doi.org/10.2307/2372752}{doi:10.2307/2372752}. For the semiregularity formulation used here, see Kleiman (2005), Remark~5.18.

\item Lefschetz, S. (1921). On certain numerical invariants of algebraic varieties with application to Abelian varieties. \emph{Transactions of the American Mathematical Society, 22}, 327--406, with continuation 407--482. Cited by Enriques, p.~257. Enriques abbreviates the title.

\item Lefschetz, S. (1923). Concerning Noether's theorem on curves traced on non singular surfaces. \emph{Bulletin of the American Mathematical Society, 29}. Cited by Enriques, p.~257. Bibliographical details beyond that citation are unverified.

\item Lewis, P. D., \& Lyons, C. (2018). Some results on surfaces with $p_g=q=1$ and $K^2=2$. \emph{International Mathematics Research Notices, 2018}(6), 1878--1919. \href{https://doi.org/10.1093/imrn/rnw294}{doi:10.1093/imrn/rnw294}. \href{https://arxiv.org/abs/1508.01863}{Author manuscript}, \S~2 and Theorem~2.1, for smooth double-cover constructions over $E^{(2)}$ and genus-two Albanese fibres.

\item Matsuki, K. (2002). \emph{Introduction to the Mori program}. Springer. \href{https://doi.org/10.1007/978-1-4757-5602-9}{doi:10.1007/978-1-4757-5602-9}. ``Birational Geometry of Surfaces,'' pp.~9--108, supplies the general minimal-model framework.

\item Maxwell, E. A. (1937). Regular canonical surfaces of genus three and four. \emph{Proceedings of the Cambridge Philosophical Society, 33}. Cited by Enriques, p.~294. Bibliographical details beyond that citation are unverified.

\item Mendes Lopes, M., \& Pardini, R. (2011). The geography of irregular surfaces. In \emph{Current developments in algebraic geometry} (Mathematical Sciences Research Institute Publications, Vol.~59, pp.~349--378). Cambridge University Press. \href{https://library.slmath.org/books/Book59/files/66pardini.pdf}{MSRI original text}. Theorem~2.1.1, p.~350, is the Castelnuovo--de Franchis theorem; \S~3 discusses the relevant irregularity inequalities.

\item Mendes Lopes, M., Pardini, R., \& Pirola, G. P. (2013). Continuous families of divisors, paracanonical systems and a new inequality for varieties of maximal Albanese dimension. \emph{Geometry \& Topology, 17}(2), 1205--1223. \href{https://doi.org/10.2140/gt.2013.17.1205}{doi:10.2140/gt.2013.17.1205}. \href{https://arxiv.org/abs/1207.4516}{Author manuscript}, Theorem~1.1 and the introductory distinction between main and exorbitant components.

\item Milne, J. S. (2008). \emph{Abelian varieties} (Version~2.00). \href{https://www.jmilne.org/math/CourseNotes/AV.pdf}{Author's course notes}. Part~I, \S\S~7--11, particularly Proposition~10.1 and Theorem~11.1, together with Proposition~13.8, Remark~13.9, and Corollary~13.10 (p.~60); Part~III, \S\S~3--5 and Theorem~10.1. Complements are taken up to isogeny, not as product isomorphisms.

\item Miyaoka, Y. (1976). Tricanonical maps of numerical Godeaux surfaces. \emph{Inventiones Mathematicae, 34}, 99--111. \href{https://doi.org/10.1007/BF01425477}{doi:10.1007/BF01425477}. Theorem~3 gives tricanonical birationality; Theorems~2 and~$2'$ on p.~106 concern base points. \href{https://www.digizeitschriften.de/download/pdf/356556735_0034/log12.pdf}{Digitized primary text}.

\item Mumford, D. (1966). \emph{Lectures on curves on an algebraic surface}, with a section by G.~M. Bergman (Annals of Mathematics Studies, No.~59). Princeton University Press. For the Severi--Zappa example, see the account in Kleiman (2005), Remark~5.18.

\item Noether, M. (1870, as cited). [Work on the independence of surface invariants; title not specified in this passage]. Publication details not specified. Cited by Enriques, p.~267. Bibliographical details beyond that citation are unverified.

\item Noether, M. (1883). Zur Grundlegung der algebraischen Raumcurven. \emph{Abhandlungen der Akademie der Wissenschaften}, Berlin. Cited by Enriques, p.~256. Bibliographical details beyond that citation are unverified.

\item Noether, M. (1887). Ueber die totalen algebraischen Differentialausdrücke. \emph{Mathematische Annalen, 29}. Cited by Enriques, p.~369. Bibliographical details beyond that citation are unverified. Enriques identifies Noether as author in the sentence introducing the footnote.

\item Noether, M. (1896). Ueber die reductiblen algebraischen Curven. \emph{Acta Mathematica}. Cited by Enriques, p.~342. Bibliographical details beyond that citation are unverified.

\item Nuer, H. (2025). Stable sheaves on bielliptic surfaces: From the classical to the modern. \emph{Mathematische Zeitschrift, 309}, Article~39. \href{https://doi.org/10.1007/s00209-024-03659-6}{doi:10.1007/s00209-024-03659-6}. \S~2.1, Table~2, Theorem~2.2, and Table~3 give the seven bielliptic types, their multiple fibres, and canonical orders.

\item Painlevé, P. (1897). \emph{Leçons sur la théorie analytique des équations différentielles}. Hermann. Cited by Enriques, pp.~399, 446. Bibliographical details beyond that citation are unverified. At p.~446 Enriques instead gives Stockholm, 1895, p.~270; that lecture date is retained as his citation, not substituted for the Paris publication date.

\item Peters, C. A. M. (n.d.). \emph{Compact complex surfaces} [Author-hosted conference text; cover identifies the Canadian Mathematical Society Summer Meeting, June 1990]. \href{https://cpeters1.win.tue.nl/ConfsAndSchools/Canada/surface.pdf}{Author's text}. Theorem~7.2, p.~18, gives the Reider criteria; \S~8, p.~19, treats pluricanonical systems. The event date is not asserted to be the publication year.

\item Picard, E. (1889). Mémoire sur la théorie des fonctions algébriques de deux variables indépendantes. \emph{Journal de Mathématiques Pures et Appliquées}. Cited by Enriques, p.~446. Bibliographical details beyond that citation are unverified. Enriques cites Chapter~III and the later exposition in Picard--Simart, Vol.~II.

\item Picard, E. (1905). [Note on the numbers of Picard integrals; title not given by Enriques]. \emph{Comptes rendus de l'Académie des Sciences}. Cited by Enriques, p.~344. Bibliographical details beyond that citation are unverified. Enriques also directs the reader to Picard--Simart, Vol.~II, p.~417.

\item Picard, E., \& Simart, G. (1897--1906). \emph{Théorie des fonctions algébriques de deux variables indépendantes} (Vols.~I--II). Gauthier-Villars. Cited by Enriques, pp.~344, 369, 446. Enriques's cited locators are Vol.~I, Chapter~V, \S~15, and Vol.~II, p.~417; Vol.~II is explicitly dated 1906 in his p.~446 citation.

\item Poincaré, H. (1910). Sur les courbes tracées sur les surfaces algébriques. \emph{Annales scientifiques de l'École Normale Supérieure}, ser.~3, Vol.~27, starting at p.~3. Cited by Enriques, pp.~343, 345. Bibliographical details beyond that citation are unverified.

\item Pompilj, G. (1940). Sulle superficie algebriche le cui curve canoniche posseggono una $g^1_3$. \emph{Rendiconti dell'Istituto Lombardo}. Cited by Enriques, p.~273. Bibliographical details beyond that citation are unverified.

\item Pompilj, G. (1941, February 27). Osservazioni sui piani tripli. \emph{Rendiconti dell'Istituto Lombardo}. Cited by Enriques, p.~303. Bibliographical details beyond that citation are unverified.

\item Pompilj, G. (1941, March 27). [Note on the classification of surfaces with trigonal canonical curves; title not specified in the citation]. \emph{Rendiconti dell'Istituto Lombardo}. Cited by Enriques, p.~302. Bibliographical details beyond that citation are unverified.

\item Pompilj, G. (n.d.). Sulla rappresentazione algebrica dei piani tripli. \emph{Rendiconti del Seminario Matematico della Università di Roma}, ser.~4, Vol.~3. Cited by Enriques, p.~303. Bibliographical details beyond that citation are unverified.

\item Reid, M. (1978). Surfaces with $p_g=0$, $K^2=1$. \emph{Journal of the Faculty of Science, University of Tokyo, Section IA, 25}(1), 75--92. \href{https://mreid.warwick.ac.uk/Reprint/pg0K21.pdf}{Author's reprint with later updates}. The canonical-ring and torsion discussion is distinguished from the explicit examples calculated in the commentary.

\item Reider, I. (1988). Vector bundles of rank 2 and linear systems on algebraic surfaces. \emph{Annals of Mathematics, 127}(2), 309--316. \href{https://doi.org/10.2307/2007055}{doi:10.2307/2007055}. The precise adjoint criterion used in Chapter~VIII is the formulation in Peters's text cited above.

\item Rohn, K. (1897). Die Raumcurven auf der Flaechen vierter Ordnung [title as printed by Enriques]. \emph{Berichte der königlich Sächsischen Gesellschaft der Wissenschaften}, Leipzig. Cited by Enriques, p.~256. Bibliographical details beyond that citation are unverified. The grammatical form in the quoted title is not presented as an independently established original title.

\item Rosenblatt, A. (1912). Sur quelques inégalités\ldots{} [title abbreviated by Enriques]. \emph{Comptes rendus de l'Académie des Sciences, 154}. Cited by Enriques, p.~299. Bibliographical details beyond that citation are unverified.

\item Rosenblatt, A. (1913). Sur les surfaces irrégulières\ldots{} [title abbreviated by Enriques]. \emph{Rendiconti del Circolo Matematico di Palermo}. Cited by Enriques, p.~462. Bibliographical details beyond that citation are unverified.

\item Rosenblatt, A. (1928). Varietà algebriche a tre e più dimensioni. \emph{Atti del Congresso Internazionale dei Matematici}. Cited by Enriques, p.~289. Bibliographical details beyond that citation are unverified.

\item Saint-Donat, B. (1974). Projective models of K-3 surfaces. \emph{American Journal of Mathematics, 96}(4), 602--639. \href{https://doi.org/10.2307/2373709}{doi:10.2307/2373709}. For projective normality and linear systems on K3 surfaces, see also Huybrechts's account cited above.

\item Schütt, M., \& Shioda, T. (2010). Elliptic surfaces. In \emph{Algebraic geometry in East Asia---Seoul 2008} (Advanced Studies in Pure Mathematics, Vol.~60, pp.~51--160). Mathematical Society of Japan. \href{https://arxiv.org/abs/0907.0298}{Author version~3}. See \S\S~3--4 and Theorems~6.8 and~6.12. The survey's fibration-with-section hypothesis is retained; its canonical formula is not by itself a reference for multiple-fibre corrections on a non-Jacobian fibration.

\item Segre, B. (1932). Determinazione di certi gruppi covarianti di due o più serie lineari. \emph{Rendiconti del Circolo Matematico di Palermo}. Cited by Enriques, p.~364. Bibliographical details beyond that citation are unverified.

\item Segre, B. (1932). Sulle superficie algebriche aventi il sistema canonico composto con un'involuzione. \emph{Rendiconti dell'Accademia dei Lincei}, ser.~6, Vol.~16, p.~316. Cited by Enriques, p.~301. Bibliographical details beyond that citation are unverified.

\item Segre, B. (1936). Quelques résultats nouveaux dans la géométrie sur une $V_3$ algébrique. \emph{Mémoires de l'Académie Royale de Belgique}, ser.~2, Vol.~14. Cited by Enriques, p.~289. Bibliographical details beyond that citation are unverified.

\item Segre, B. (1938). Un teorema fondamentale della geometria sulle superficie algebriche e il principio di spezzamento. \emph{Annali di Matematica}. Cited by Enriques, p.~347. Bibliographical details beyond that citation are unverified.

\item Segre, B. (1939, January 9). Sur un théorème fondamental de géométrie sur les surfaces algébriques. \emph{Comptes rendus de l'Académie des Sciences}. Cited by Enriques, p.~349. Bibliographical details beyond that citation are unverified.

\item Severi, F. (1903). Su alcune questioni di postulazione. \emph{Rendiconti del Circolo Matematico di Palermo, 17}. Cited by Enriques, p.~291. Bibliographical details beyond that citation are unverified.

\item Severi, F. (1903, January 25). Sulle superficie che rappresentano le coppie di punti di una curva algebrica. \emph{Atti della Reale Accademia delle Scienze di Torino}. Cited by Enriques, pp.~340, 355. Bibliographical details beyond that citation are unverified.

\item Severi, F. (1904, February). Osservazioni sui sistemi continui di curve appartenenti ad una superficie algebrica. \emph{Atti della Reale Accademia delle Scienze di Torino}. Cited by Enriques, p.~341. Bibliographical details beyond that citation are unverified.

\item Severi, F. (1904--1905, as cited). Sulle superficie algebriche che posseggono integrali di Picard della seconda specie. \emph{Rendiconti dell'Accademia dei Lincei} (1904); further reference to \emph{Mathematische Annalen} (1905). Cited by Enriques, p.~344. Bibliographical details beyond that citation are unverified. The author is Severi in Enriques's surrounding sentence.

\item Severi, F. (1905). Intorno alla costruzione dei sistemi completi non lineari che appartengono ad una superficie irregolare. \emph{Rendiconti del Circolo Matematico di Palermo}. Cited by Enriques, pp.~342--343. Bibliographical details beyond that citation are unverified.

\item Severi, F. (1905). Sul teorema di Riemann--Roch e sulle serie continue di curve appartenenti ad una superficie algebrica. \emph{Atti della Reale Accademia delle Scienze di Torino}. Cited by Enriques, p.~345. Bibliographical details beyond that citation are unverified.

\item Severi, F. (1905). Sulla differenza fra i numeri degli integrali di Picard della prima e della seconda specie appartenenti ad una superficie irregolare. \emph{Atti della Reale Accademia delle Scienze di Torino}. Cited by Enriques, p.~344. Bibliographical details beyond that citation are unverified.

\item Severi, F. (1905, August, as cited). Il teorema d'Abel sulle superficie algebriche. \emph{Annali di Matematica}. Cited by Enriques, p.~345. The published article is dated 1906, ser.~3, Vol.~12, pp.~55--79, \href{https://doi.org/10.1007/BF02419496}{doi:10.1007/BF02419496}.

\item Severi, F. (1908). Sulle superficie algebriche che ammettono un gruppo continuo permutabile a due parametri di trasformazioni birazionali. \emph{Atti dell'Istituto Veneto}. Cited by Enriques, p.~446. Bibliographical details beyond that citation are unverified.

\item Severi, F. (1909). Fondamenti per la geometria sulle varietà algebriche. \emph{Rendiconti del Circolo Matematico di Palermo, 28}. Cited by Enriques, pp.~289, 291. Bibliographical details beyond that citation are unverified.

\item Severi, F. (1909, as cited). Le superficie algebriche con curva canonica d'ordine zero. \emph{Atti dell'Istituto Veneto}, cited as Vol.~LXV, part~2. Cited by Enriques, pp.~255--256. Bibliographical details beyond that citation are unverified. The source's volume/date combination is retained as a reported locator; conflicting catalog details are not silently resolved.

\item Severi, F. (1911). Sulle superficie e varietà algebriche irregolari di genere geometrico nullo. \emph{Rendiconti dell'Accademia dei Lincei}, ser.~5, Vol.~20. Cited by Enriques, pp.~364, 369. Bibliographical details beyond that citation are unverified.

\item Severi, F. (1916, April 2 and 16). Nuovi contributi alla teoria dei sistemi continui di curve appartenenti ad una superficie algebrica. \emph{Rendiconti dell'Accademia dei Lincei}. Cited by Enriques, pp.~343, 353. Bibliographical details beyond that citation are unverified.

\item Severi, F. (1921). Sulla teoria degli integrali semplici di prima specie appartenenti ad una superficie algebrica [series of notes]. \emph{Rendiconti dell'Accademia dei Lincei}. Cited by Enriques, p.~345. Bibliographical details beyond that citation are unverified.

\item Severi, F. (1941). La teoria generale dei sistemi continui di curve sopra una superficie algebrica. \emph{Memorie dell'Accademia d'Italia, 19}. Cited by Enriques, pp.~349--354. Bibliographical details beyond that citation are unverified.

\item Severi, F. (1942, January 16). [Note correcting a lemma on continuous systems; title not specified]. Note presented to the Reale Accademia d'Italia. Cited by Enriques, p.~350. Bibliographical details beyond that citation are unverified.

\item Severi, F. (1942, November 30, as cited). Intorno ai sistemi continui di curve sopra una superficie algebrica. \emph{Commentarii Mathematici Helvetici}. Cited by Enriques, p.~350. Bibliographical details beyond that citation are unverified. The citation's date is not converted into an independently established publication year.

\item Sheridan, J. (2022). Divisor varieties of symmetric products. \emph{International Mathematics Research Notices, 2022}(13), 9830--9863. \href{https://doi.org/10.1093/imrn/rnaa313}{doi:10.1093/imrn/rnaa313}. \href{https://arxiv.org/abs/1906.05465}{Author manuscript}, \S~1 and Example~1.3. General-curve assumptions in the broader component-count results are not suppressed.

\item The Stacks Project Authors. (2026). \emph{The Stacks Project}. \href{https://stacks.math.columbia.edu/tag/03RD}{Tag~03RD} (Tsen's theorem) and \href{https://stacks.math.columbia.edu/tag/03RF}{Tag~03RF} (vanishing of the Brauer group of a function field of a curve over an algebraically closed field). Accessed November~7, 2026. Tags, rather than changeable displayed theorem numbers, identify the cited results.

\item Torelli, R. (1914). Sulle serie algebriche semplicemente infinite di gruppi di punti appartenenti a una curva algebrica. \emph{Rendiconti del Circolo Matematico di Palermo, 37}. Cited by Enriques, p.~358. Bibliographical details beyond that citation are unverified. Enriques cites \S~II, no.~11.

\item Voight, J. (2021). QM abelian surfaces. In \emph{Quaternion algebras} (Graduate Texts in Mathematics, Vol.~288, pp.~799--830). Springer. \href{https://doi.org/10.1007/978-3-030-56694-4_43}{doi:10.1007/978-3-030-56694-4\_43}. Theorem~43.5.1 gives the genus-two Jacobian or product-of-elliptic-curves alternatives for a principally polarized complex abelian surface.

\item Werner, C. (1994). A surface of general type with $p_g=q=0$, $K^2=1$. \emph{Manuscripta Mathematica, 84}, 327--341. \href{https://doi.org/10.1007/BF02567460}{doi:10.1007/BF02567460}. \href{https://arxiv.org/abs/alg-geom/9408007}{Author manuscript}, \S~2, Proposition~1, and \S\S~3--4. The explicit construction is distinct from the historical unsuccessful configuration.

\item Xiao, G. (1985). Finitude de l'application bicanonique des surfaces de type général. \emph{Bulletin de la Société Mathématique de France, 113}, 23--51. \href{https://doi.org/10.24033/bsmf.2019}{doi:10.24033/bsmf.2019}. \href{https://www.numdam.org/article/BSMF_1985__113__23_0.pdf}{Original text}, Theorem~1, p.~24, and \S~1.3.

\item Zariski, O. (1935). \emph{Algebraic surfaces}. Springer. Cited by Enriques, p.~345. Bibliographical details beyond that citation are unverified. Enriques cites Chapter~VII, \S\S~4--5, pp.~128 and~133; these locators refer to the original edition, not a later enlarged edition.
\end{itemize}
\end{sloppypar}
